% AUTO-GENERATED FROM MATH/Math.md BY build_math_from_md.py.
% Edit Math.md or the builder, then regenerate.
% Options for packages loaded elsewhere
\PassOptionsToPackage{unicode}{hyperref}
\PassOptionsToPackage{hyphens}{url}
%
\documentclass[
  10pt,
]{article}
\usepackage{amsmath,amssymb}
\usepackage{lmodern}
\usepackage{iftex}
\ifPDFTeX
  \usepackage[T1]{fontenc}
  \usepackage[utf8]{inputenc}
  \usepackage{textcomp} % provide euro and other symbols
\else % if luatex or xetex
  \usepackage{unicode-math}
  \defaultfontfeatures{Scale=MatchLowercase}
  \defaultfontfeatures[\rmfamily]{Ligatures=TeX,Scale=1}
\fi
% Use upquote if available, for straight quotes in verbatim environments
\IfFileExists{upquote.sty}{\usepackage{upquote}}{}
\IfFileExists{microtype.sty}{% use microtype if available
  \usepackage[]{microtype}
  \UseMicrotypeSet[protrusion]{basicmath} % disable protrusion for tt fonts
}{}
\makeatletter
\@ifundefined{KOMAClassName}{% if non-KOMA class
  \IfFileExists{parskip.sty}{%
    \usepackage{parskip}
  }{% else
    \setlength{\parindent}{0pt}
    \setlength{\parskip}{6pt plus 2pt minus 1pt}}
}{% if KOMA class
  \KOMAoptions{parskip=half}}
\makeatother
\usepackage{xcolor}
\IfFileExists{xurl.sty}{\usepackage{xurl}}{} % add URL line breaks if available
\IfFileExists{bookmark.sty}{\usepackage{bookmark}}{\usepackage{hyperref}}
\hypersetup{
  pdftitle={Hubbard-Holstein Mathematical Notes (Symbolic-Only Edition)},
  pdfauthor={Jake Skyler Strobel},
  hidelinks,
  pdfcreator={LaTeX via pandoc}}
\urlstyle{same} % disable monospaced font for URLs
\usepackage[margin=0.8in]{geometry}
\usepackage{color}
\usepackage{fancyvrb}
\newcommand{\VerbBar}{|}
\newcommand{\VERB}{\Verb[commandchars=\\\{\}]}
\DefineVerbatimEnvironment{Highlighting}{Verbatim}{commandchars=\\\{\}}
% Add ',fontsize=\small' for more characters per line
\newenvironment{Shaded}{}{}
\newcommand{\AlertTok}[1]{\textcolor[rgb]{1.00,0.00,0.00}{\textbf{#1}}}
\newcommand{\AnnotationTok}[1]{\textcolor[rgb]{0.38,0.63,0.69}{\textbf{\textit{#1}}}}
\newcommand{\AttributeTok}[1]{\textcolor[rgb]{0.49,0.56,0.16}{#1}}
\newcommand{\BaseNTok}[1]{\textcolor[rgb]{0.25,0.63,0.44}{#1}}
\newcommand{\BuiltInTok}[1]{#1}
\newcommand{\CharTok}[1]{\textcolor[rgb]{0.25,0.44,0.63}{#1}}
\newcommand{\CommentTok}[1]{\textcolor[rgb]{0.38,0.63,0.69}{\textit{#1}}}
\newcommand{\CommentVarTok}[1]{\textcolor[rgb]{0.38,0.63,0.69}{\textbf{\textit{#1}}}}
\newcommand{\ConstantTok}[1]{\textcolor[rgb]{0.53,0.00,0.00}{#1}}
\newcommand{\ControlFlowTok}[1]{\textcolor[rgb]{0.00,0.44,0.13}{\textbf{#1}}}
\newcommand{\DataTypeTok}[1]{\textcolor[rgb]{0.56,0.13,0.00}{#1}}
\newcommand{\DecValTok}[1]{\textcolor[rgb]{0.25,0.63,0.44}{#1}}
\newcommand{\DocumentationTok}[1]{\textcolor[rgb]{0.73,0.13,0.13}{\textit{#1}}}
\newcommand{\ErrorTok}[1]{\textcolor[rgb]{1.00,0.00,0.00}{\textbf{#1}}}
\newcommand{\ExtensionTok}[1]{#1}
\newcommand{\FloatTok}[1]{\textcolor[rgb]{0.25,0.63,0.44}{#1}}
\newcommand{\FunctionTok}[1]{\textcolor[rgb]{0.02,0.16,0.49}{#1}}
\newcommand{\ImportTok}[1]{#1}
\newcommand{\InformationTok}[1]{\textcolor[rgb]{0.38,0.63,0.69}{\textbf{\textit{#1}}}}
\newcommand{\KeywordTok}[1]{\textcolor[rgb]{0.00,0.44,0.13}{\textbf{#1}}}
\newcommand{\NormalTok}[1]{#1}
\newcommand{\OperatorTok}[1]{\textcolor[rgb]{0.40,0.40,0.40}{#1}}
\newcommand{\OtherTok}[1]{\textcolor[rgb]{0.00,0.44,0.13}{#1}}
\newcommand{\PreprocessorTok}[1]{\textcolor[rgb]{0.74,0.48,0.00}{#1}}
\newcommand{\RegionMarkerTok}[1]{#1}
\newcommand{\SpecialCharTok}[1]{\textcolor[rgb]{0.25,0.44,0.63}{#1}}
\newcommand{\SpecialStringTok}[1]{\textcolor[rgb]{0.73,0.40,0.53}{#1}}
\newcommand{\StringTok}[1]{\textcolor[rgb]{0.25,0.44,0.63}{#1}}
\newcommand{\VariableTok}[1]{\textcolor[rgb]{0.10,0.09,0.49}{#1}}
\newcommand{\VerbatimStringTok}[1]{\textcolor[rgb]{0.25,0.44,0.63}{#1}}
\newcommand{\WarningTok}[1]{\textcolor[rgb]{0.38,0.63,0.69}{\textbf{\textit{#1}}}}
\usepackage{longtable,booktabs,array}
\usepackage{calc} % for calculating minipage widths
% Correct order of tables after \paragraph or \subparagraph
\usepackage{etoolbox}
\makeatletter
\patchcmd\longtable{\par}{\if@noskipsec\mbox{}\fi\par}{}{}
\makeatother
% Allow footnotes in longtable head/foot
\IfFileExists{footnotehyper.sty}{\usepackage{footnotehyper}}{\usepackage{footnote}}
\makesavenoteenv{longtable}
\setlength{\emergencystretch}{3em} % prevent overfull lines
\providecommand{\tightlist}{%
  \setlength{\itemsep}{0pt}\setlength{\parskip}{0pt}}
\setcounter{secnumdepth}{-\maxdimen} % remove section numbering
\ifLuaTeX
  \usepackage{selnolig}  % disable illegal ligatures
\fi

\title{Hubbard-Holstein Mathematical Notes (Symbolic-Only Edition)}
\author{Jake Skyler Strobel}
\date{March 20, 2026}

\begin{document}
\maketitle

{
\setcounter{tocdepth}{3}
\tableofcontents
}
\hypertarget{l2-noisy-surface-run-handoff}{%
\section{L=2 Noisy Surface Run
Handoff}\label{l2-noisy-surface-run-handoff}}

For the \textbf{local/offline L=2 HH noisy surface} (\texttt{phase3\_v1}
+ \texttt{pareto\_lean\_l2} + \texttt{FakeNighthawk}), see
\href{l2_noisy_surface_run.md}{MATH/l2\_noisy\_surface\_run.md}.

This is the terminal-agent handoff note for the three mitigation lanes:

\begin{itemize}
\tightlist
\item
  \texttt{none}
\item
  \texttt{readout}
\item
  \texttt{full}
\end{itemize}

It is intentionally the front-page pointer so the noisy-surface run
recipe is visible before the rest of the manuscript.

\begin{center}\rule{0.5\linewidth}{0.5pt}\end{center}

\hypertarget{qpu-l2-winner-framing-and-real-qpu-settings-april-8-2026}{%
\section{{[}QPU{]} L=2 Winner Framing and Real QPU Settings (April 8,
2026)}\label{qpu-l2-winner-framing-and-real-qpu-settings-april-8-2026}}

\begin{quote}
Front-page contract: distinguish the best raw HF-start ADAPT result from
the separately validated real-QPU deployment anchor. Do not conflate
either one with later prune/readapt JSON post-processing or with
fixed-scaffold runtime submissions.
\end{quote}

\hypertarget{two-front-page-objects}{%
\subsection{Two Front-Page Objects}\label{two-front-page-objects}}

\begin{itemize}
\tightlist
\item
  \textbf{Scientific raw HF-start winner.} This means direct ADAPT from
  the Hartree-Fock reference with \texttt{adapt\_ref\_json\ =\ null},
  with no later JSON prune/recovery pass substituted in. Under the
  canonical \texttt{L=2,\ U=4.0,\ g=0.5} gate
  \texttt{\textbar{}ΔE\textbar{}\ \textbackslash{}sim\ 10\^{}\{-4\}},
  the best substantiated raw HF-start line currently in hand is the
  direct \texttt{full\_meta} artifact
  \texttt{hh\_backend\_adapt\_fullhorse\_powell\_20260322T194155Z\_fakenighthawk.json},
  which reaches \[
  |\Delta E|=5.6178234645\times 10^{-5}
  \] at \texttt{81} transpiled two-qubit gates on
  \texttt{FakeMarrakesh}.
\item
  \textbf{Validated real-QPU deployment anchor.} This is the separately
  validated April 5 \texttt{pareto\_lean\ +\ phase3\_v1\ +\ SPSA} raw
  HF-start command that is currently trusted for rerunability on the
  live public CLI surface. Its scaffold burden is heavier, but it
  remains the validated deployment recipe: \[
  |\Delta E|=1.0822209459\times 10^{-4},
  \] with \texttt{218} transpiled two-qubit gates on the same
  \texttt{FakeMarrakesh} proxy.
\end{itemize}

\hypertarget{raw-hf-start-direct-adapt-table-delta-esim-10-4}{%
\subsection{\texorpdfstring{Raw HF-Start Direct ADAPT Table
(\texttt{\textbar{}\textbackslash{}Delta\ E\textbar{}\textbackslash{}sim\ 10\^{}\{-4\}})}{Raw HF-Start Direct ADAPT Table (\textbar\textbackslash Delta E\textbar\textbackslash sim 10\^{}\{-4\})}}\label{raw-hf-start-direct-adapt-table-delta-esim-10-4}}

\begin{longtable}[]{@{}
  >{\raggedright\arraybackslash}p{(\columnwidth - 14\tabcolsep) * \real{0.10}}
  >{\raggedright\arraybackslash}p{(\columnwidth - 14\tabcolsep) * \real{0.10}}
  >{\raggedright\arraybackslash}p{(\columnwidth - 14\tabcolsep) * \real{0.10}}
  >{\raggedleft\arraybackslash}p{(\columnwidth - 14\tabcolsep) * \real{0.14}}
  >{\raggedleft\arraybackslash}p{(\columnwidth - 14\tabcolsep) * \real{0.14}}
  >{\raggedleft\arraybackslash}p{(\columnwidth - 14\tabcolsep) * \real{0.14}}
  >{\raggedleft\arraybackslash}p{(\columnwidth - 14\tabcolsep) * \real{0.14}}
  >{\raggedleft\arraybackslash}p{(\columnwidth - 14\tabcolsep) * \real{0.14}}@{}}
\toprule
Role & Artifact / surface & Pool &
\texttt{\textbar{}\textbackslash{}Delta\ E\textbar{}} & Logical ops &
Runtime params & Transpiled 2Q & Depth \\
\midrule
\endhead
Best raw HF-start ADAPT winner &
\texttt{hh\_backend\_adapt\_fullhorse\_powell\_20260322T194155Z\_fakenighthawk}
& \texttt{full\_meta} & \texttt{5.6178234645e-05} & 16 & 27 & 81 &
151 \\
Best reduced-pool direct line &
\texttt{adapt\_hh\_L2\_ecut1\_pareto\_lean\_l2\_phase3\_powell\_rerun\_currentcode\_20260321T171615Z}
& \texttt{pareto\_lean\_l2} & \texttt{5.6178234644e-05} & 14 & 25 & 97 &
208 \\
Wide-shortlist direct comparison &
\texttt{campaign\_A7c\_L2\_shortlist\_wide\_phase3\_v1} &
\texttt{full\_meta} & \texttt{5.6178234644e-05} & 14 & 30 & 116 & 310 \\
Current-code reproduction of March 22 low-energy regime &
\texttt{20260409\_hh\_l2\_hist81\_repro\_powell\_histbare\_d16\_v1} &
\texttt{full\_meta} historicalized POWELL surface &
\texttt{5.7217492583e-05} & 17 & 54 & 265 & 786 \\
Best current native recovery branch on today's diagnostic surface &
\texttt{20260409\_hh\_l2\_children\_repeat\_bridge\_diag\_v1}
(\texttt{children\_off\ +\ repeats\_off\ +\ hist}) & \texttt{full\_meta}
diagnostic bridge surface & \texttt{5.6178241861e-05} & 13 & 36 & 160 &
408 \\
Validated April 5 deployment anchor &
\texttt{pareto\_lean\ +\ phase3\_v1\ +\ SPSA} & \texttt{pareto\_lean} &
\texttt{1.0822209459e-04} & 20 & 52 & 218 & 633 \\
\bottomrule
\end{longtable}

The final \texttt{218}-2Q row is the same raw HF-start scaffold family
as the confirmed April 5 authority command; it is heavier than the
scientific raw winner, but it is the currently validated deployment
anchor on the public CLI surface.

The post-processed \texttt{31/37/57}-2Q prune/readapt artifacts are
\textbf{not} the object a fresh real-QPU ADAPT run from HF would
execute. They remain useful comparison context only.

\hypertarget{current-native-recovery-note-april-9-2026}{%
\subsubsection{Current native recovery note (April 9,
2026)}\label{current-native-recovery-note-april-9-2026}}

The new \texttt{children\_off\ +\ repeats\_off\ +\ hist} branch from
\texttt{20260409\_hh\_l2\_children\_repeat\_bridge\_diag\_v1} deserves
front-page visibility because it is the best \textbf{current-code native
recovery branch} on today's diagnostic/public surface:

\[
|\Delta E|=5.6178241861\times 10^{-5},\qquad
F=0.9999719078,
\]

with \texttt{13} logical operators, \texttt{36} runtime parameters, and
\texttt{160} transpiled two-qubit gates at depth \texttt{408} on
\texttt{FakeMarrakesh}.

This is \textbf{not} the absolute historical raw HF-start winner; the
March 22 \texttt{81}-2Q \texttt{full\_meta} artifact remains lighter
overall. What makes the April 9 branch important is different: it is the
first current diagnostic line that re-enters the good native basin while
simultaneously showing that disabling children and literal repeats helps
native recovery.

It still does \textbf{not} recover the historical lean family exactly.
Relative to the historical targets, it reaches \texttt{9/10} bridge
overlap and \texttt{5/6} lean overlap, with the remaining extras
concentrated in

\begin{itemize}
\tightlist
\item
  \texttt{paop\_full:paop\_disp(site=1)}
\item
  \texttt{paop\_lf\_full:paop\_dbl\_p(site=0-\textgreater{}phonon=1)}
\item
  \texttt{paop\_lf\_full:paop\_dbl\_p(site=1-\textgreater{}phonon=0)}
\item
  \texttt{paop\_full:paop\_hopdrag(0,1)}
\end{itemize}

This diagnostic result matters because it isolates the current native
failure mode more sharply than the older heavy broad-pool lines: the
good basin is real, but some quadrature and inside-family extras are
still acting as \textbf{path operators} rather than merely removable
terminal clutter.

\hypertarget{current-code-reproduction-of-the-march-22-low-energy-regime}{%
\subsubsection{Current-code reproduction of the March 22 low-energy
regime}\label{current-code-reproduction-of-the-march-22-low-energy-regime}}

The March 22 \texttt{81}-2Q \texttt{full\_meta} line is still a real
historical raw HF-start winner, but current code does \textbf{not}
replay that exact burden from its historical settings block. What
current code \textbf{does} reproduce is the same low-energy regime on a
historicalized POWELL surface with a hard cap at the first in-run depth
that re-enters the target band.

The saved current-code replay is

\[
|\Delta E|=5.7217492583\times 10^{-5},
\]

from
\texttt{20260409\_hh\_l2\_hist81\_repro\_powell\_histbare\_d16\_v1},
with \texttt{17} logical blocks, \texttt{54} runtime parameters, and
\texttt{265} transpiled two-qubit gates at depth \texttt{786} on
\texttt{FakeMarrakesh}.

So this run should be interpreted as a \textbf{current-code low-energy
reproduction}, not as a faithful current reproduction of the historical
\texttt{81}-2Q scaffold burden.

\textbf{Exact current-code replay CLI}

\begin{Shaded}
\begin{Highlighting}[]
\ExtensionTok{python}\NormalTok{ pipelines/hardcoded/adapt\_pipeline.py }\DataTypeTok{\textbackslash{}}
  \AttributeTok{{-}{-}L}\NormalTok{ 2 }\AttributeTok{{-}{-}problem}\NormalTok{ hh }\AttributeTok{{-}{-}t}\NormalTok{ 1.0 }\AttributeTok{{-}{-}u}\NormalTok{ 4.0 }\AttributeTok{{-}{-}dv}\NormalTok{ 0.0 }\AttributeTok{{-}{-}omega0}\NormalTok{ 1.0 }\AttributeTok{{-}{-}g{-}ep}\NormalTok{ 0.5 }\AttributeTok{{-}{-}n{-}ph{-}max}\NormalTok{ 1 }\DataTypeTok{\textbackslash{}}
  \AttributeTok{{-}{-}boson{-}encoding}\NormalTok{ binary }\AttributeTok{{-}{-}ordering}\NormalTok{ blocked }\AttributeTok{{-}{-}boundary}\NormalTok{ open }\AttributeTok{{-}{-}term{-}order}\NormalTok{ sorted }\DataTypeTok{\textbackslash{}}
  \AttributeTok{{-}{-}adapt{-}pool}\NormalTok{ full\_meta }\AttributeTok{{-}{-}adapt{-}continuation{-}mode}\NormalTok{ phase3\_v1 }\DataTypeTok{\textbackslash{}}
  \AttributeTok{{-}{-}adapt{-}max{-}depth}\NormalTok{ 16 }\AttributeTok{{-}{-}adapt{-}eps{-}grad}\NormalTok{ 5e{-}7 }\AttributeTok{{-}{-}adapt{-}eps{-}energy}\NormalTok{ 1e{-}9 }\DataTypeTok{\textbackslash{}}
  \AttributeTok{{-}{-}adapt{-}inner{-}optimizer}\NormalTok{ POWELL }\AttributeTok{{-}{-}adapt{-}state{-}backend}\NormalTok{ compiled }\DataTypeTok{\textbackslash{}}
  \AttributeTok{{-}{-}adapt{-}reopt{-}policy}\NormalTok{ windowed }\AttributeTok{{-}{-}adapt{-}window{-}size}\NormalTok{ 999999 }\AttributeTok{{-}{-}adapt{-}window{-}topk}\NormalTok{ 999999 }\DataTypeTok{\textbackslash{}}
  \AttributeTok{{-}{-}adapt{-}full{-}refit{-}every}\NormalTok{ 8 }\AttributeTok{{-}{-}adapt{-}final{-}full{-}refit}\NormalTok{ true }\DataTypeTok{\textbackslash{}}
  \AttributeTok{{-}{-}adapt{-}beam{-}live{-}branches}\NormalTok{ 1 }\AttributeTok{{-}{-}adapt{-}beam{-}children{-}per{-}parent}\NormalTok{ 1 }\AttributeTok{{-}{-}adapt{-}beam{-}terminated{-}keep}\NormalTok{ 1 }\DataTypeTok{\textbackslash{}}
  \AttributeTok{{-}{-}phase1{-}lambda{-}F}\NormalTok{ 1.0 }\AttributeTok{{-}{-}phase1{-}lambda{-}compile}\NormalTok{ 0.05 }\AttributeTok{{-}{-}phase1{-}lambda{-}measure}\NormalTok{ 0.02 }\AttributeTok{{-}{-}phase1{-}lambda{-}leak}\NormalTok{ 0.0 }\DataTypeTok{\textbackslash{}}
  \AttributeTok{{-}{-}phase1{-}score{-}z{-}alpha}\NormalTok{ 0.0 }\AttributeTok{{-}{-}phase1{-}shortlist{-}size}\NormalTok{ 256 }\AttributeTok{{-}{-}phase1{-}probe{-}max{-}positions}\NormalTok{ 999999 }\DataTypeTok{\textbackslash{}}
  \AttributeTok{{-}{-}phase1{-}plateau{-}patience}\NormalTok{ 2 }\AttributeTok{{-}{-}phase1{-}trough{-}margin{-}ratio}\NormalTok{ 1.0 }\DataTypeTok{\textbackslash{}}
  \AttributeTok{{-}{-}phase1{-}prune{-}enabled} \AttributeTok{{-}{-}phase1{-}prune{-}mode}\NormalTok{ live }\AttributeTok{{-}{-}phase1{-}prune{-}fraction}\NormalTok{ 0.25 }\AttributeTok{{-}{-}phase1{-}prune{-}max{-}candidates}\NormalTok{ 6 }\AttributeTok{{-}{-}phase1{-}prune{-}max{-}regression}\NormalTok{ 1e{-}8 }\DataTypeTok{\textbackslash{}}
  \AttributeTok{{-}{-}phase2{-}shortlist{-}fraction}\NormalTok{ 1.0 }\AttributeTok{{-}{-}phase2{-}shortlist{-}size}\NormalTok{ 128 }\DataTypeTok{\textbackslash{}}
  \AttributeTok{{-}{-}phase2{-}lambda{-}H}\NormalTok{ 1e{-}6 }\AttributeTok{{-}{-}phase2{-}rho}\NormalTok{ 0.25 }\AttributeTok{{-}{-}phase2{-}gamma{-}N}\NormalTok{ 1.0 }\DataTypeTok{\textbackslash{}}
  \AttributeTok{{-}{-}phase2{-}enable{-}batching} \AttributeTok{{-}{-}phase2{-}batch{-}target{-}size}\NormalTok{ 8 }\AttributeTok{{-}{-}phase2{-}batch{-}size{-}cap}\NormalTok{ 16 }\AttributeTok{{-}{-}phase2{-}batch{-}near{-}degenerate{-}ratio}\NormalTok{ 0.98 }\DataTypeTok{\textbackslash{}}
  \AttributeTok{{-}{-}phase2{-}w{-}depth}\NormalTok{ 0.0 }\AttributeTok{{-}{-}phase2{-}w{-}group}\NormalTok{ 0.0 }\AttributeTok{{-}{-}phase2{-}w{-}shot}\NormalTok{ 0.0 }\AttributeTok{{-}{-}phase2{-}w{-}optdim}\NormalTok{ 0.0 }\AttributeTok{{-}{-}phase2{-}w{-}reuse}\NormalTok{ 0.0 }\AttributeTok{{-}{-}phase2{-}w{-}lifetime}\NormalTok{ 0.0 }\DataTypeTok{\textbackslash{}}
  \AttributeTok{{-}{-}phase2{-}eta{-}L}\NormalTok{ 0.0 }\AttributeTok{{-}{-}phase2{-}motif{-}bonus{-}weight}\NormalTok{ 0.0 }\AttributeTok{{-}{-}phase2{-}duplicate{-}penalty{-}weight}\NormalTok{ 0.0 }\DataTypeTok{\textbackslash{}}
  \AttributeTok{{-}{-}phase2{-}compat{-}overlap{-}weight}\NormalTok{ 0.0 }\AttributeTok{{-}{-}phase2{-}compat{-}comm{-}weight}\NormalTok{ 0.0 }\AttributeTok{{-}{-}phase2{-}compat{-}curv{-}weight}\NormalTok{ 0.0 }\AttributeTok{{-}{-}phase2{-}compat{-}sched{-}weight}\NormalTok{ 0.0 }\AttributeTok{{-}{-}phase2{-}compat{-}measure{-}weight}\NormalTok{ 0.0 }\DataTypeTok{\textbackslash{}}
  \AttributeTok{{-}{-}phase2{-}frontier{-}ratio}\NormalTok{ 1.0 }\AttributeTok{{-}{-}phase3{-}frontier{-}ratio}\NormalTok{ 1.0 }\AttributeTok{{-}{-}phase2{-}remaining{-}evaluations{-}proxy{-}mode}\NormalTok{ none }\DataTypeTok{\textbackslash{}}
  \AttributeTok{{-}{-}phase3{-}tie{-}beam{-}max{-}branches}\NormalTok{ 1 }\DataTypeTok{\textbackslash{}}
  \AttributeTok{{-}{-}phase3{-}symmetry{-}mitigation{-}mode}\NormalTok{ verify\_only }\AttributeTok{{-}{-}phase3{-}enable{-}rescue} \DataTypeTok{\textbackslash{}}
  \AttributeTok{{-}{-}phase3{-}lifetime{-}cost{-}mode}\NormalTok{ phase3\_v1 }\AttributeTok{{-}{-}phase3{-}runtime{-}split{-}mode}\NormalTok{ shortlist\_pauli\_children\_v1 }\DataTypeTok{\textbackslash{}}
  \AttributeTok{{-}{-}phase3{-}backend{-}cost{-}mode}\NormalTok{ transpile\_single\_v1 }\AttributeTok{{-}{-}phase3{-}backend{-}name}\NormalTok{ FakeNighthawk }\DataTypeTok{\textbackslash{}}
  \AttributeTok{{-}{-}phase3{-}backend{-}transpile{-}seed}\NormalTok{ 7 }\AttributeTok{{-}{-}phase3{-}backend{-}optimization{-}level}\NormalTok{ 1 }\DataTypeTok{\textbackslash{}}
  \AttributeTok{{-}{-}adapt{-}seed}\NormalTok{ 7 }\AttributeTok{{-}{-}adapt{-}finite{-}angle{-}fallback} \AttributeTok{{-}{-}adapt{-}finite{-}angle}\NormalTok{ 0.1 }\AttributeTok{{-}{-}adapt{-}finite{-}angle{-}min{-}improvement}\NormalTok{ 1e{-}12 }\DataTypeTok{\textbackslash{}}
  \AttributeTok{{-}{-}adapt{-}drop{-}floor}\NormalTok{ 1e{-}11 }\AttributeTok{{-}{-}adapt{-}drop{-}patience}\NormalTok{ 5 }\AttributeTok{{-}{-}adapt{-}drop{-}min{-}depth}\NormalTok{ 12 }\AttributeTok{{-}{-}adapt{-}grad{-}floor} \AttributeTok{{-}1.0} \DataTypeTok{\textbackslash{}}
  \AttributeTok{{-}{-}paop{-}r}\NormalTok{ 1 }\AttributeTok{{-}{-}paop{-}normalization}\NormalTok{ none }\DataTypeTok{\textbackslash{}}
  \AttributeTok{{-}{-}initial{-}state{-}source}\NormalTok{ adapt\_vqe }\DataTypeTok{\textbackslash{}}
  \AttributeTok{{-}{-}skip{-}pdf}
\end{Highlighting}
\end{Shaded}

\hypertarget{real-qpu-deployment-anchor-april-5-2026-validated}{%
\subsection{Real QPU Deployment Anchor (April 5, 2026
Validated)}\label{real-qpu-deployment-anchor-april-5-2026-validated}}

\begin{quote}
\textbf{This is the recommended configuration for real QPU deployment.}
Validated via narrow 36-case local sweep confirming rerunnable
deployment optimality in the \texttt{pareto\_lean} + \texttt{phase3\_v1}
regime. This is the deployment anchor, not the absolute lowest raw
HF-start two-qubit line.
\end{quote}

\hypertarget{anchor-configuration}{%
\subsection{Anchor Configuration}\label{anchor-configuration}}

\textbf{Use this exact command for real QPU L=2 ADAPT runs:}

\begin{Shaded}
\begin{Highlighting}[]
\ExtensionTok{python} \AttributeTok{{-}u} \AttributeTok{{-}m}\NormalTok{ pipelines.hardcoded.adapt\_pipeline }\DataTypeTok{\textbackslash{}}
  \AttributeTok{{-}{-}L}\NormalTok{ 2 }\AttributeTok{{-}{-}problem}\NormalTok{ hh }\AttributeTok{{-}{-}t}\NormalTok{ 1.0 }\AttributeTok{{-}{-}u}\NormalTok{ 4.0 }\AttributeTok{{-}{-}dv}\NormalTok{ 0.0 }\DataTypeTok{\textbackslash{}}
  \AttributeTok{{-}{-}omega0}\NormalTok{ 1.0 }\AttributeTok{{-}{-}g{-}ep}\NormalTok{ 0.5 }\AttributeTok{{-}{-}n{-}ph{-}max}\NormalTok{ 1 }\DataTypeTok{\textbackslash{}}
  \AttributeTok{{-}{-}boson{-}encoding}\NormalTok{ binary }\AttributeTok{{-}{-}ordering}\NormalTok{ blocked }\AttributeTok{{-}{-}boundary}\NormalTok{ open }\AttributeTok{{-}{-}term{-}order}\NormalTok{ sorted }\DataTypeTok{\textbackslash{}}
  \AttributeTok{{-}{-}adapt{-}pool}\NormalTok{ pareto\_lean }\AttributeTok{{-}{-}adapt{-}continuation{-}mode}\NormalTok{ phase3\_v1 }\DataTypeTok{\textbackslash{}}
  \AttributeTok{{-}{-}adapt{-}max{-}depth}\NormalTok{ 160 }\AttributeTok{{-}{-}adapt{-}eps{-}grad}\NormalTok{ 5e{-}7 }\AttributeTok{{-}{-}adapt{-}eps{-}energy}\NormalTok{ 1e{-}9 }\AttributeTok{{-}{-}adapt{-}seed}\NormalTok{ 5 }\DataTypeTok{\textbackslash{}}
  \AttributeTok{{-}{-}adapt{-}inner{-}optimizer}\NormalTok{ SPSA }\DataTypeTok{\textbackslash{}}
  \AttributeTok{{-}{-}adapt{-}spsa{-}a}\NormalTok{ 0.1 }\AttributeTok{{-}{-}adapt{-}spsa{-}c}\NormalTok{ 0.02 }\AttributeTok{{-}{-}adapt{-}spsa{-}A}\NormalTok{ 5.0 }\DataTypeTok{\textbackslash{}}
  \AttributeTok{{-}{-}adapt{-}spsa{-}callback{-}every}\NormalTok{ 5 }\AttributeTok{{-}{-}adapt{-}spsa{-}progress{-}every{-}s}\NormalTok{ 30 }\DataTypeTok{\textbackslash{}}
  \AttributeTok{{-}{-}adapt{-}maxiter}\NormalTok{ 3200 }\DataTypeTok{\textbackslash{}}
  \AttributeTok{{-}{-}adapt{-}state{-}backend}\NormalTok{ compiled }\DataTypeTok{\textbackslash{}}
  \AttributeTok{{-}{-}adapt{-}reopt{-}policy}\NormalTok{ windowed }\AttributeTok{{-}{-}adapt{-}window{-}size}\NormalTok{ 999999 }\AttributeTok{{-}{-}adapt{-}window{-}topk}\NormalTok{ 999999 }\DataTypeTok{\textbackslash{}}
  \AttributeTok{{-}{-}adapt{-}full{-}refit{-}every}\NormalTok{ 8 }\AttributeTok{{-}{-}adapt{-}final{-}full{-}refit}\NormalTok{ true }\DataTypeTok{\textbackslash{}}
  \AttributeTok{{-}{-}adapt{-}beam{-}live{-}branches}\NormalTok{ 2 }\AttributeTok{{-}{-}adapt{-}beam{-}children{-}per{-}parent}\NormalTok{ 2 }\AttributeTok{{-}{-}adapt{-}beam{-}terminated{-}keep}\NormalTok{ 2 }\DataTypeTok{\textbackslash{}}
  \AttributeTok{{-}{-}phase1{-}prune{-}enabled} \AttributeTok{{-}{-}phase1{-}prune{-}fraction}\NormalTok{ 0.25 }\AttributeTok{{-}{-}phase1{-}prune{-}max{-}candidates}\NormalTok{ 6 }\AttributeTok{{-}{-}phase1{-}prune{-}max{-}regression}\NormalTok{ 1e{-}8 }\DataTypeTok{\textbackslash{}}
  \AttributeTok{{-}{-}phase1{-}probe{-}max{-}positions}\NormalTok{ 999999 }\AttributeTok{{-}{-}phase1{-}trough{-}margin{-}ratio}\NormalTok{ 1.0 }\AttributeTok{{-}{-}phase1{-}shortlist{-}size}\NormalTok{ 64 }\DataTypeTok{\textbackslash{}}
  \AttributeTok{{-}{-}phase2{-}shortlist{-}fraction}\NormalTok{ 1.0 }\AttributeTok{{-}{-}phase2{-}shortlist{-}size}\NormalTok{ 64 }\DataTypeTok{\textbackslash{}}
  \AttributeTok{{-}{-}phase2{-}frontier{-}ratio}\NormalTok{ 0.85 }\DataTypeTok{\textbackslash{}}
  \AttributeTok{{-}{-}phase2{-}lambda{-}H}\NormalTok{ 1e{-}06 }\AttributeTok{{-}{-}phase2{-}rho}\NormalTok{ 0.25 }\AttributeTok{{-}{-}phase2{-}gamma{-}N}\NormalTok{ 1.0 }\DataTypeTok{\textbackslash{}}
  \AttributeTok{{-}{-}phase2{-}w{-}depth}\NormalTok{ 0.1 }\AttributeTok{{-}{-}phase2{-}w{-}group}\NormalTok{ 0.075 }\AttributeTok{{-}{-}phase2{-}w{-}shot}\NormalTok{ 0.075 }\AttributeTok{{-}{-}phase2{-}w{-}optdim}\NormalTok{ 0.05 }\AttributeTok{{-}{-}phase2{-}w{-}reuse}\NormalTok{ 0.05 }\AttributeTok{{-}{-}phase2{-}w{-}lifetime}\NormalTok{ 0.05 }\DataTypeTok{\textbackslash{}}
  \AttributeTok{{-}{-}phase2{-}eta{-}L}\NormalTok{ 0.0 }\AttributeTok{{-}{-}phase2{-}motif{-}bonus{-}weight}\NormalTok{ 0.05 }\AttributeTok{{-}{-}phase2{-}duplicate{-}penalty{-}weight}\NormalTok{ 0.0 }\DataTypeTok{\textbackslash{}}
  \AttributeTok{{-}{-}phase2{-}compat{-}overlap{-}weight}\NormalTok{ 0.4 }\AttributeTok{{-}{-}phase2{-}compat{-}comm{-}weight}\NormalTok{ 0.2 }\AttributeTok{{-}{-}phase2{-}compat{-}curv{-}weight}\NormalTok{ 0.2 }\DataTypeTok{\textbackslash{}}
  \AttributeTok{{-}{-}phase2{-}compat{-}sched{-}weight}\NormalTok{ 0.2 }\AttributeTok{{-}{-}phase2{-}compat{-}measure{-}weight}\NormalTok{ 0.2 }\DataTypeTok{\textbackslash{}}
  \AttributeTok{{-}{-}phase2{-}remaining{-}evaluations{-}proxy{-}mode}\NormalTok{ auto }\DataTypeTok{\textbackslash{}}
  \AttributeTok{{-}{-}phase3{-}frontier{-}ratio}\NormalTok{ 0.85 }\DataTypeTok{\textbackslash{}}
  \AttributeTok{{-}{-}phase3{-}symmetry{-}mitigation{-}mode}\NormalTok{ verify\_only }\AttributeTok{{-}{-}phase3{-}enable{-}rescue} \DataTypeTok{\textbackslash{}}
  \AttributeTok{{-}{-}phase3{-}backend{-}cost{-}mode}\NormalTok{ proxy }\AttributeTok{{-}{-}phase3{-}runtime{-}split{-}mode}\NormalTok{ off }\AttributeTok{{-}{-}phase3{-}lifetime{-}cost{-}mode}\NormalTok{ off }\DataTypeTok{\textbackslash{}}
  \AttributeTok{{-}{-}adapt{-}drop{-}floor}\NormalTok{ 0.0005 }\AttributeTok{{-}{-}adapt{-}drop{-}patience}\NormalTok{ 3 }\AttributeTok{{-}{-}adapt{-}drop{-}min{-}depth}\NormalTok{ 12 }\DataTypeTok{\textbackslash{}}
  \AttributeTok{{-}{-}initial{-}state{-}source}\NormalTok{ adapt\_vqe }\AttributeTok{{-}{-}skip{-}pdf} \AttributeTok{{-}{-}phase2{-}no{-}batching}
\end{Highlighting}
\end{Shaded}

\hypertarget{expected-performance}{%
\subsection{Expected Performance}\label{expected-performance}}

\begin{longtable}[]{@{}
  >{\raggedright\arraybackslash}p{(\columnwidth - 4\tabcolsep) * \real{0.36}}
  >{\raggedright\arraybackslash}p{(\columnwidth - 4\tabcolsep) * \real{0.32}}
  >{\raggedright\arraybackslash}p{(\columnwidth - 4\tabcolsep) * \real{0.32}}@{}}
\toprule
Metric & Value & Notes \\
\midrule
\endhead
Energy & 0.15877612622031317 & Exact ground state reference \\
\textbar ΔE\textbar{} & \textbf{1.082e-4} & Excellent accuracy for
NISQ \\
Logical scaffold size & 20 & Validated raw HF-start deployment
scaffold \\
Runtime parameters & 52 & Efficient parameterization \\
Marrakesh proxy burden & \texttt{218} 2Q, depth \texttt{633}, size
\texttt{939} & Same raw HF-start scaffold as the confirmed strong
template \\
Wall-clock & \textasciitilde43s & Fast convergence \\
\bottomrule
\end{longtable}

\hypertarget{validation-results-april-5-sweep}{%
\subsection{Validation Results (April 5
Sweep)}\label{validation-results-april-5-sweep}}

A 36-case narrow local sweep confirmed this configuration is optimal:

\begin{itemize}
\tightlist
\item
  {[}OK{]} \textbf{SPSA maxiter=3200}: Goldilocks zone (2400
  undershoots, 4000 overshoots)
\item
  {[}OK{]} \textbf{Frontier ratio 0.85}: All values (0.80--0.90)
  converge to same energy; 0.85 balances quality
\item
  {[}OK{]} \textbf{Beam (2,2,2)}: Sufficient; wider beam (3,2,3) costs
  50\% more time, zero quality gain
\item
  {[}OK{]} \textbf{Drop policy (0.0005, 3, 12)}: Optimal shallowness;
  tighter drop \(\to\) shallower circuits for QPU
\end{itemize}

\hypertarget{critical-finding}{%
\subsubsection{Critical Finding}\label{critical-finding}}

Increasing SPSA \texttt{maxiter} beyond 3200 \textbf{degrades}
performance (1.27--1.54e-4 vs 1.082e-4): - Suggests pool saturation:
\texttt{pareto\_lean} operators are exhausted at this budget - SPSA
noise amplification: longer optimization in finite ansatz space doesn't
help

\textbf{Do not deviate from maxiter=3200 in this regime.}

\hypertarget{what-not-to-change}{%
\subsection{What NOT to Change}\label{what-not-to-change}}

Based on sweep validation: - {[}X{]} Pool: Do not switch from
\texttt{pareto\_lean} - {[}X{]} Continuation: Do not move away from
\texttt{phase3\_v1} - {[}X{]} Runtime split: Keep \texttt{off} for L=2 -
{[}X{]} SPSA budget: Do not increase \texttt{maxiter} beyond 3200 -
{[}X{]} Frontier: Do not vary frontier ratio (insensitive in 0.80--0.90
range) - {[}X{]} Beam: Do not use (3,2,3); keep (2,2,2)

\hypertarget{for-real-qpu-deployment}{%
\subsection{For Real QPU Deployment}\label{for-real-qpu-deployment}}

\begin{enumerate}
\def\labelenumi{\arabic{enumi}.}
\tightlist
\item
  Use the command above exactly as written
\item
  Expect \textbar ΔE\textbar{} \(\approx\) 1.08e-4 (this is the
  best-known trustworthy result)
\item
  Ansatz depth = 20 is excellent for NISQ mitigation
\item
  No further local tuning will improve this configuration (sweep
  confirms)
\end{enumerate}

\hypertarget{companion-driven-candidate}{%
\subsection{Companion Driven
Candidate}\label{companion-driven-candidate}}

The leading \textbf{driven realtime-controller} candidate for the same
L=2 HH scaffold family is now the heavier-target
\texttt{secant\_lead100} setting from
\texttt{artifacts/agent\_runs/20260405\_controller\_drive\_heavier\_secant\_lead\_scan\_v1/json/secant\_lead100.json}.

This should be tracked alongside the winning static/noiseless L=2 ADAPT
anchor above when choosing real-QPU-facing settings, because it is
currently the best controller law we have for recovering visible driven
dynamics on the stronger \texttt{A=0.55}, \texttt{tbar=4.0},
\texttt{t\_final=8.0} benchmark.

\hypertarget{current-driven-controller-headline-metrics}{%
\subsubsection{Current driven-controller headline
metrics}\label{current-driven-controller-headline-metrics}}

\begin{longtable}[]{@{}
  >{\raggedright\arraybackslash}p{(\columnwidth - 4\tabcolsep) * \real{0.36}}
  >{\raggedright\arraybackslash}p{(\columnwidth - 4\tabcolsep) * \real{0.32}}
  >{\raggedright\arraybackslash}p{(\columnwidth - 4\tabcolsep) * \real{0.32}}@{}}
\toprule
Metric & Value & Notes \\
\midrule
\endhead
Controller law & \texttt{secant\_lead100} & Signed-energy-lead-tapered
secant + signed blend \\
Drive target & \texttt{A=0.55}, \texttt{omega=2.0}, \texttt{tbar=4.0},
\texttt{t0=0.0} & Heavier than the earlier fine-resolution regime \\
Mean \texttt{\textbar{}ΔE\textbar{}} & \textbf{1.633e-3} & Best
whole-window value among current heavier-target candidates \\
Max \texttt{\textbar{}ΔE\textbar{}} & \textbf{6.501e-3} & Large
improvement over prior signed-blend and raw secant anchors \\
Early-window MAE (\texttt{t≤3.35}) & \textbf{1.176e-3} & Recovers the
startup dip far better than \texttt{signed\_mild\_front\_current} \\
Late-window MAE (\texttt{t≥3.35}) & \textbf{1.964e-3} & Still slightly
better than the previous late-time winner \\
Energy span ratio & \texttt{1.032} & Slight overshoot, but close to
exact \\
Min fidelity & \texttt{0.9626} & Remains bounded with no append
events \\
\bottomrule
\end{longtable}

\hypertarget{leading-qpu-facing-pair}{%
\subsubsection{Leading QPU-facing pair}\label{leading-qpu-facing-pair}}

\begin{itemize}
\tightlist
\item
  \textbf{Static / ground-state anchor:} the validated
  \texttt{pareto\_lean\ +\ phase3\_v1} L=2 ADAPT command above.
\item
  \textbf{Driven / dynamics anchor:} the heavier-target
  \texttt{secant\_lead100} controller result above.
\end{itemize}

These are the two front-page candidates we should currently treat as the
leading settings when optimizing toward a real QPU run.

\begin{center}\rule{0.5\linewidth}{0.5pt}\end{center}

\begin{quote}
Canonical naming note: \texttt{MATH/Math.md} is the markdown authoring
source for manuscript sync, and \texttt{MATH/Math.tex} is the generated
typesetting twin used to build \texttt{MATH/Math.pdf}. Regenerate the
synchronized TeX/PDF pair with
\texttt{python\ MATH/build\_math\_from\_md.py} or
\texttt{MATH/build\_math.sh}. The archival repo-oriented manuscript is
\texttt{MATH/archaic\_repo\_math.md}.
\end{quote}

\hypertarget{parameter-manifest-and-reader-contract}{%
\section{1. Parameter Manifest and Reader
Contract}\label{parameter-manifest-and-reader-contract}}

This duplicate is intentionally symbolic-first. It keeps the same broad
manuscript shape as the existing mathematical notes,

\begin{enumerate}
\def\labelenumi{\arabic{enumi}.}
\tightlist
\item
  primitives first,
\item
  composite operators second,
\item
  explicit substitutions third,
\item
  reduced master forms last,
\end{enumerate}

but it suppresses almost all implementation labels, file names, mode
strings, and executable-surface commentary. The goal is to make the
document readable as mathematics rather than as a repository audit.

\hypertarget{required-parameter-manifest}{%
\subsection{1.1 Required parameter
manifest}\label{required-parameter-manifest}}

A complete symbolic run or derivation should specify at least the
following data.

\begin{itemize}
\tightlist
\item
  Model family: Hubbard or Hubbard-Holstein.
\item
  Lattice size: \(L\).
\item
  Ordering map: interleaved or blocked, i.e.~a choice of
  \(p(i,\sigma)\).
\item
  Boundary condition: open or periodic.
\item
  Fermionic parameters: hopping \(t\), onsite interaction \(U\), and
  static site potentials \(v_i\).
\item
  Bosonic parameters: phonon frequency \(\omega_0\), electron-phonon
  coupling \(g\), and cutoff \(n_{\mathrm{ph,max}}\).
\item
  Boson encoding choice: binary or unary.
\item
  Variational family, when relevant: layerwise, physical-termwise,
  excitation-based, or adaptive.
\item
  Control parameters, when relevant: optimizer family, propagation rule,
  drive waveform parameters, and benchmark grid.
\end{itemize}

\hypertarget{reader-contract}{%
\subsection{1.2 Reader contract}\label{reader-contract}}

This manuscript favors explicit substitution over compressed
abstraction.

\begin{itemize}
\tightlist
\item
  If a primitive operator exists, it is written first.
\item
  If a composite object is built from primitives, the primitive form is
  shown.
\item
  If a later equation can be simplified by inserting an earlier
  primitive explicitly, that insertion is made.
\item
  If multiple mathematical surfaces exist, such as different index
  orderings or boson encodings, the surface is named locally rather than
  hidden behind an unnamed default.
\end{itemize}

\hypertarget{quick-notation-glossary-for-adaptive-and-continuation-sections}{%
\subsubsection{1.2.1 Quick notation glossary for adaptive and
continuation
sections}\label{quick-notation-glossary-for-adaptive-and-continuation-sections}}

This subsection is a compact map for the symbols that recur in the
adaptive-selection, split-aware, replay, and beam-style sections.

\begin{itemize}
\tightlist
\item
  \(g,m\): generic candidate generators. The manuscript uses both
  letters in different sections; they play the same role.
\item
  \(p\): candidate insertion position in the ordered scaffold.
\item
  \(r=(m,p)\): a candidate-position record, i.e.~a generator together
  with a chosen insertion position.
\item
  \(\mathcal G\): the full master generator pool.
\item
  \(\mathcal G^{(1)},\mathcal G^{(2)},\mathcal G^{(3)}\): progressively
  more restricted candidate universes, with \[
  \mathcal G^{(3)}\subseteq\mathcal G^{(2)}\subseteq\mathcal G^{(1)}\subseteq\mathcal G.
  \]
\item
  \(\mathcal P_g\) or \(\mathcal P_m\): admissible insertion positions
  for candidate \(g\) or \(m\).
\item
  \(W(p)\): inherited refit window used if a candidate is inserted at
  position \(p\).
\item
  \(W_{\mathrm{new}}(p)\): the newest-parameter portion of the refit
  window.
\item
  \(W_{|\theta|}(p)\): the older large-amplitude carry subset of the
  refit window.
\item
  \(\mathcal P_{\mathrm{avail}}^{(3)}\): the full available set of
  Phase-3 candidate-position records.
\item
  \(\mathcal C_{\mathrm{cheap}}\): the horizon-sized coarse shortlist
  obtained by cheap-score ranking over the full admissible Phase-3
  candidate-position pool.
\item
  \(\mathcal S_2,\mathcal S_3\): the Phase-2 and Phase-3 shortlists that
  survive cheap screening before full reranking.
\item
  \(\mathcal C_{\mathrm{split}}(m)\): the family of split children
  generated from macro generator \(m\).
\item
  \(\Gamma_{\mathrm{prune}}^{\mathrm{perm}}\): post-admission
  prune-permissibility gate deciding whether local scaffold
  simplification is attempted.
\item
  \(\mathcal J_{\mathrm{exp}}\): the post-admission set of locally
  expendable scaffold coordinates or generators considered by the prune
  pass.
\item
  \(\mathcal O_*\): the finally selected adaptive scaffold, i.e.~the
  ordered generator sequence admitted by the adaptive stage.
\item
  \(\mathcal M_{\mathrm{warm}},\mathcal M_{\mathrm{refine}},\mathcal M_{\mathrm{replay}}\):
  archival appendix-only warm-start, optional refine, and replay
  manifolds.
\item
  \(\mathcal B\): a portable handoff/state bundle carrying manifest,
  amplitudes, energies, and continuation metadata.
\item
  \(\mathcal C\): continuation payload / provenance payload attached to
  a handoff bundle.
\item
  \(\Gamma_{\mathrm{stage}}\): stage-admissibility gate.
\item
  \(\Gamma_{\mathrm{sym}}\): symmetry-admissibility gate.
\item
  \(\mathfrak b\): a beam branch state. In §11 it indexes alternative
  ADAPT scaffolds; in §17 it indexes projected-dynamics branches over
  checkpoint segments.
\item
  \(\mathcal H_b\): admitted-record history carried by scaffold branch
  \(b\).
\item
  \(\mathcal A_b\): branch-local admissible admission set retained for
  scaffold-beam expansion.
\item
  \(\mathcal Q(b)\): scaffold-beam proposal family, usually
  \(\{\mathrm{stop}\}\cup\mathcal A_b\).
\item
  \(\eta_b\): branch-local stage label for the scaffold beam,
  e.g.~core/seed versus residual.
\item
  \(W_{\mathrm{act}}(\mathfrak b)\): the currently active position-probe
  window on branch \(b\).
\item
  \(\mathfrak F_c\): live frontier at beam round \(c\).
\item
  \(\mathfrak T_c\): terminal-branch pool at beam round \(c\).
\item
  \(\mathfrak W_{\mathrm{final}}\): final union of surviving frontier
  and terminal branches.
\item
  \(J_b\): cumulative branch objective tuple used to rank or prune
  branch \(b\); its coordinates are section-specific.
\end{itemize}

\hypertarget{quick-parameter-glossary-for-adaptive-and-continuation-sections}{%
\subsubsection{1.2.2 Quick parameter glossary for adaptive and
continuation
sections}\label{quick-parameter-glossary-for-adaptive-and-continuation-sections}}

The symbols below are the main scalar controls that appear in the
adaptive-selection, split-aware, and beam-style formulas.

\begin{itemize}
\tightlist
\item
  \(\lambda_{\mathrm{2q}}\): weight on two-qubit proxy burden.
\item
  \(\lambda_d\): weight on depth-style proxy burden.
\item
  \(\lambda_\theta\): weight on marginal parameter-count burden.
\item
  \(\lambda_F\): cheap-stage normalization weight multiplying the
  tangent metric term.
\item
  \(\lambda_H\): regularization strength for inherited-window Hessian
  blocks.
\item
  \(\rho\): trust-region radius scale used in one-dimensional local
  improvement models.
\item
  \(N_{\mathrm{cheap}}^{\max}\): absolute cap on the horizon-sized
  Phase-3 coarse shortlist.
\item
  \(N_{\max}\): maximum shortlist size.
\item
  \(f_{\mathrm{short}}\): shortlist fraction used in Phase 3.
\item
  \(w_R\): weight on useful-horizon lifetime burden.
\item
  \(w_D\): weight on lifetime depth/compile burden.
\item
  \(w_G\): weight on newly introduced grouped-measurement burden.
\item
  \(w_C\): weight on newly introduced compile or circuit burden.
\item
  \(w_c\): weight on any extra cheap-path correction burden.
\item
  \(\gamma_N\): exponent applied to novelty factors in full rerank
  scores.
\item
  \(\varepsilon\): small positive stabilizer used in denominators and
  tie-safe quotients.
\item
  \(\varepsilon_{\mathrm{nov}}\): novelty regularizer used when
  inverting or stabilizing tangent-overlap matrices.
\item
  \(z_\alpha\): confidence multiplier used to convert a raw gradient
  estimate into a lower-confidence gradient.
\item
  \(\sigma(g,p)\) or \(\sigma_r\): uncertainty estimate attached to the
  candidate-position gradient signal.
\item
  \(\tau_1\): Phase-2 shortlist threshold.
\item
  \(\tau_{\mathrm{split}}\): split-replacement margin required before a
  split child or child-set beats its parent.
\item
  \(\tau_{\mathrm{drop}}\): minimum per-depth absolute-error improvement
  regarded as meaningful by the stage controller.
\item
  \(M\): patience window used when counting repeated small-drop depths.
\item
  \(D_{\mathrm{left}}(t)\): remaining controller depth budget at
  selector step \(t\).
\item
  \(S_t(k)\): survival factor for having a still-useful admission
  opportunity \(k\) selector steps ahead.
\item
  \(Q_t(k)\): conditional usefulness factor for the \(k\)-ahead
  admission opportunity.
\item
  \(\widehat N_{\mathrm{rem}}(t)\): heuristic forecast of useful
  remaining admissions at selector step \(t\).
\item
  \(C_t\): confidence attached to the current useful-runway forecast.
\item
  \(K_{\mathrm{coarse}}(t)\): dynamic coarse-cap induced by the current
  useful-runway forecast.
\item
  \(H_t\): effective useful-depth horizon, \[
  H_t=\min\bigl(D_{\mathrm{left}}(t),\widehat N_{\mathrm{rem}}(t)\bigr).
  \]
\item
  \(W_{\mathrm{sat}}(t)\): saturation pressure induced by the current
  useful-runway forecast.
\item
  \(\pi_t\): selector mode tag, e.g.~fast / mixed / high-fidelity.
\item
  \(\mathrm{regime}_{\mathrm{sat}}(t)\): hungry / watch / plateau
  saturation regime.
\item
  \(\eta_{\mathrm{ret}}\): minimum retained-gain fraction required by
  the post-admission prune accept rule.
\item
  \(\beta_{\mathrm{keep}}\): repo-facing retained-gain floor used in
  prune acceptance; in this manuscript it plays the same role as
  \(\eta_{\mathrm{ret}}\).
\item
  \(\tau_{\mathrm{miss}}\): projected-miss threshold used in
  manifold-growth logic.
\item
  \(B_{\mathrm{child}}\): beam child cap per parent branch-expansion
  round.
\item
  \(B_{\mathrm{target}}(t)\): active batch target size at selector step
  \(t\).
\item
  \(B_{\mathrm{live}}\): live-frontier beam width retained after
  pruning.
\item
  \(B_{\mathrm{term}}\): cap on stored terminal branches.
\item
  \(M_{\mathrm{probe}}\): cap on distinct probed insertion positions
  retained in order.
\item
  \(a_j(t)\): admission age of scaffold coordinate or generator \(j\) at
  selector step \(t\).
\item
  \(\varrho_j(t)\): normalized admission-rank score of coordinate or
  generator \(j\) at selector step \(t\).
\item
  \(\mathcal P_{\mathrm{protect}}(t)\): hard-protected coordinate set
  excluded from the live prune pass.
\item
  \(c_j^{\mathrm{cool}}(t)\): local prune cooldown on coordinate \(j\).
\item
  \(S_\theta(j,t)\): optional small-angle prescreen score used only to
  save local prune work.
\item
  \(\tau_{\mathrm{stale}}\): minimum local expendability score required
  before a prune trial is even considered.
\item
  \(\Xi_t^{\mathrm{trial}}\): exact rollback snapshot taken before a
  local prune/remove-refit trial.
\item
  \(\Delta k_{\mathrm{cp}}\): checkpoint advance used in the
  projected-dynamics beam of §17, not in the scaffold beam of §11.
\item
  \(\Delta k_{\mathrm{probe}}\): probe advance used in the
  projected-dynamics beam of §17, not in the scaffold beam of §11.
\item
  \(C\): final beam round index / number of retained beam-update rounds.
\end{itemize}

In §11, the index \(t\) denotes a selector/controller step rather than
the real-time variable used later in the dynamics sections.

\hypertarget{non-negotiable-representational-conventions}{%
\subsection{1.3 Non-negotiable representational
conventions}\label{non-negotiable-representational-conventions}}

\hypertarget{internal-pauli-alphabet}{%
\subsubsection{1.3.1 Internal Pauli
alphabet}\label{internal-pauli-alphabet}}

In formulas below we use the standard alphabet \[
\{I,X,Y,Z\}.
\] If one prefers the lower-case alphabet \(\{e,x,y,z\}\), the
identification is simply \[
e \leftrightarrow I,
\qquad
x \leftrightarrow X,
\qquad
y \leftrightarrow Y,
\qquad
z \leftrightarrow Z.
\]

\hypertarget{pauli-word-and-qubit-ordering}{%
\subsubsection{1.3.2 Pauli-word and qubit
ordering}\label{pauli-word-and-qubit-ordering}}

Pauli words and computational-basis labels are written left-to-right as
\[
q_{N_q-1}\cdots q_1 q_0,
\] with qubit \(q_0\) the rightmost symbol and also the
least-significant bit in basis-index arithmetic.

\hypertarget{canonical-algebra-sources}{%
\subsubsection{1.3.3 Canonical algebra
sources}\label{canonical-algebra-sources}}

The canonical mathematical primitives are:

\begin{itemize}
\tightlist
\item
  Jordan-Wigner ladder operators \(\hat c_p^\dagger\) and \(\hat c_p\),
\item
  fermion number operators \(\hat n_p\),
\item
  boson ladder operators \(\hat b_i^\dagger\) and \(\hat b_i\),
\item
  Hermitian Pauli words and Pauli polynomials,
\item
  ordered exponentials generated from those primitives.
\end{itemize}

\hypertarget{surface-specific-defaults}{%
\subsubsection{1.3.4 Surface-specific
defaults}\label{surface-specific-defaults}}

This symbolic version does not pretend there is a single universal
surface. The mathematics admits several equally valid choices:

\begin{itemize}
\tightlist
\item
  interleaved or blocked fermion ordering,
\item
  open or periodic boundary conditions,
\item
  binary or unary boson encoding,
\item
  static or driven Hamiltonians,
\item
  fixed-structure or adaptive variational manifolds.
\end{itemize}

Whenever a formula depends on the choice of surface, that dependence is
shown locally.

\hypertarget{canonical-mathematical-anchors}{%
\subsection{1.4 Canonical mathematical
anchors}\label{canonical-mathematical-anchors}}

The main objects carried through the manuscript are:

\begin{itemize}
\tightlist
\item
  the ordering map \(p(i,\sigma)\),
\item
  the ladder primitives \(\hat c_p^\dagger\), \(\hat c_p\),
  \(\hat b_i^\dagger\), \(\hat b_i\),
\item
  the density operators \(\hat n_{i\sigma}\), \(\hat n_i\), and
  \(\hat D_i\),
\item
  the Hubbard and Hubbard-Holstein Hamiltonians,
\item
  the reference state \(|\psi_{\mathrm{ref}}\rangle\),
\item
  the variational manifold \(\mathcal M\) and its ansatz map
  \(\theta \mapsto |\psi(\theta)\rangle\),
\item
  the propagation rule \(U(t)\),
\item
  the continuation or handoff bundle \(\mathcal B\).
\end{itemize}

\hypertarget{ordering-indexing-and-register-layout}{%
\section{2. Ordering, Indexing, and Register
Layout}\label{ordering-indexing-and-register-layout}}

\hypertarget{site-spin-and-mode-indices}{%
\subsection{2.1 Site, spin, and mode
indices}\label{site-spin-and-mode-indices}}

The site index is \[
i \in \{0,1,\dots,L-1\},
\] and the spin label is stored as \[
\sigma \in \{\uparrow,\downarrow\} \equiv \{0,1\}.
\] The fermion mode index is the Jordan-Wigner qubit index.

\hypertarget{interleaved-ordering}{%
\subsubsection{2.1.1 Interleaved ordering}\label{interleaved-ordering}}

The interleaved map is \[
p(i,\sigma)=2i+\sigma,
\] so that \[
p(i,\uparrow)=2i,
\qquad
p(i,\downarrow)=2i+1.
\]

\hypertarget{blocked-ordering}{%
\subsubsection{2.1.2 Blocked ordering}\label{blocked-ordering}}

The blocked map is \[
p(i,\uparrow)=i,
\qquad
p(i,\downarrow)=L+i.
\]

\hypertarget{pauli-word-placement-and-basis-index-extraction}{%
\subsection{2.2 Pauli-word placement and basis-index
extraction}\label{pauli-word-placement-and-basis-index-extraction}}

If a Pauli letter acts on qubit \(q\), then in a printed word of length
\(N_q\) it sits at string position \[
\operatorname{pos}(q)=N_q-1-q.
\]

If the computational-basis index is \(k\), then the occupation bit on
qubit \(q\) is \[
b_q(k)=\left\lfloor\frac{k}{2^q}\right\rfloor \bmod 2.
\] Thus the printed bitstring and integer basis index obey the same
rightmost-\(q_0\) convention.

\hypertarget{full-hh-register-layout}{%
\subsection{2.3 Full HH register layout}\label{full-hh-register-layout}}

The fermion register uses \[
N_{\mathrm{ferm}}=2L
\] qubits.

If the local phonon cutoff is \(n_{\mathrm{ph,max}}\), then the local
phonon Hilbert dimension is \[
d=n_{\mathrm{ph,max}}+1.
\] The number of phonon qubits per site is \[
q_{\mathrm{pb}}=
\begin{cases}
\max\{1,\lceil \log_2 d\rceil\}, & \text{binary},\\
d, & \text{unary}.
\end{cases}
\] Therefore the total HH qubit count is \[
N_q=2L+Lq_{\mathrm{pb}}.
\]

In qubit-index order, the register is naturally partitioned as \[
[\text{fermions}\;|\;\text{site-0 phonons}\;|\;\cdots\;|\;\text{site-(L-1) phonons}].
\] In printed order \(q_{N_q-1}\cdots q_0\), the high-index phonon
blocks appear on the left.

\hypertarget{fermionic-primitives-and-direct-substitution}{%
\section{3. Fermionic Primitives and Direct
Substitution}\label{fermionic-primitives-and-direct-substitution}}

\hypertarget{jordan-wigner-ladder-primitives}{%
\subsection{3.1 Jordan-Wigner ladder
primitives}\label{jordan-wigner-ladder-primitives}}

For mode \(p\), the creation operator is \[
\hat c_p^\dagger
=
\frac{1}{2}(X_p-iY_p)\prod_{r=0}^{p-1}Z_r,
\] and the annihilation operator is \[
\hat c_p
=
\frac{1}{2}(X_p+iY_p)\prod_{r=0}^{p-1}Z_r.
\]

In fully expanded printed-word form, the operator acts on the word
ordered as \(q_{N_q-1}\cdots q_0\) with the Jordan-Wigner \(Z\)-string
extending through all modes below \(p\).

\hypertarget{number-primitive}{%
\subsection{3.2 Number primitive}\label{number-primitive}}

The fermion number operator is \[
\hat n_p=\hat c_p^\dagger\hat c_p=\frac{I-Z_p}{2}.
\]

\hypertarget{site-densities-and-doublon-operator}{%
\subsection{3.3 Site densities and doublon
operator}\label{site-densities-and-doublon-operator}}

At site \(i\), define \[
\hat n_i = \hat n_{i\uparrow}+\hat n_{i\downarrow},
\] and the doublon operator \[
\hat D_i = \hat n_{i\uparrow}\hat n_{i\downarrow}.
\] By direct substitution, \[
\hat D_i = \frac{1}{4}(I-Z_{p(i,\uparrow)})(I-Z_{p(i,\downarrow)}).
\]

\hypertarget{worked-ordering-substitutions}{%
\subsection{3.4 Worked ordering
substitutions}\label{worked-ordering-substitutions}}

\hypertarget{interleaved}{%
\subsubsection{3.4.1 Interleaved}\label{interleaved}}

Under interleaved ordering, \[
\hat n_i = I - \frac{1}{2}(Z_{2i}+Z_{2i+1}),
\] and \[
\hat D_i = \frac{1}{4}(I-Z_{2i})(I-Z_{2i+1}).
\] For nearest-neighbor spin-up hopping in a one-dimensional chain, \[
\hat c_{i\uparrow}^\dagger \hat c_{i+1,\uparrow} + \hat c_{i+1,\uparrow}^\dagger \hat c_{i\uparrow}
=
\frac{1}{2}\Bigl(X_{2i} Z_{2i+1} X_{2i+2} + Y_{2i} Z_{2i+1} Y_{2i+2}\Bigr).
\]

\hypertarget{blocked}{%
\subsubsection{3.4.2 Blocked}\label{blocked}}

Under blocked ordering, \[
\hat n_i = I - \frac{1}{2}(Z_i + Z_{L+i}),
\] and \[
\hat D_i = \frac{1}{4}(I-Z_i)(I-Z_{L+i}).
\] For spin-up nearest-neighbor hopping, the Jordan-Wigner string is
empty because the two blocked spin-up modes are adjacent, so \[
\hat c_{i\uparrow}^\dagger \hat c_{i+1,\uparrow} + \hat c_{i+1,\uparrow}^\dagger \hat c_{i\uparrow}
=
\frac{1}{2}\Bigl(X_i X_{i+1} + Y_i Y_{i+1}\Bigr).
\] The general blocked-ordering formula uses the same Jordan-Wigner
interval rule as before, with an empty product for adjacent modes.

\hypertarget{boson-primitives-and-encodings}{%
\section{4. Boson Primitives and
Encodings}\label{boson-primitives-and-encodings}}

\hypertarget{local-boson-hilbert-space}{%
\subsection{4.1 Local boson Hilbert
space}\label{local-boson-hilbert-space}}

For each site \(i\), the truncated boson Hilbert space is \[
\mathcal H_{b,i}=\operatorname{span}\{|n_i\rangle : 0\le n_i\le n_{\mathrm{ph,max}}\}.
\] The ladder operators satisfy \[
\hat b_i^\dagger |n\rangle = \sqrt{n+1}\,|n+1\rangle,
\qquad
\hat b_i |n\rangle = \sqrt{n}\,|n-1\rangle,
\] with \[
\hat n_{b,i}=\hat b_i^\dagger \hat b_i,
\qquad
\hat x_i=\hat b_i+\hat b_i^\dagger,
\qquad
\hat P_i=i(\hat b_i^\dagger-\hat b_i).
\]

\hypertarget{binary-encoding}{%
\subsection{4.2 Binary encoding}\label{binary-encoding}}

In binary encoding, each local occupation number \(n\in\{0,\dots,d-1\}\)
is mapped to a binary word on \(q_{\mathrm{pb}}=\lceil \log_2 d\rceil\)
qubits: \[
|n\rangle \longmapsto |\operatorname{bin}(n)\rangle.
\] The truncated number operator is then \[
\hat n_{b,i} = \sum_{n=0}^{d-1} n\,|n\rangle\langle n|,
\] understood as an operator on the binary-encoded subspace.

\hypertarget{unary-encoding}{%
\subsection{4.3 Unary encoding}\label{unary-encoding}}

In unary encoding, each local occupation \(n\) is represented by a
one-hot basis state on \(d\) qubits: \[
|n\rangle \longmapsto |0\cdots 010\cdots 0\rangle,
\] with the single \(1\) marking the occupied unary slot. The number
operator remains \[
\hat n_{b,i} = \sum_{n=0}^{d-1} n\,|n\rangle\langle n|,
\] but now the basis states are embedded in a one-hot code space.

\hypertarget{phonon-vacuum}{%
\subsection{4.4 Phonon vacuum}\label{phonon-vacuum}}

\hypertarget{binary-vacuum}{%
\subsubsection{4.4.1 Binary vacuum}\label{binary-vacuum}}

The local binary vacuum is \[
|0\rangle_{b,i} \longmapsto |00\cdots 0\rangle.
\]

\hypertarget{unary-vacuum}{%
\subsubsection{4.4.2 Unary vacuum}\label{unary-vacuum}}

The local unary vacuum is the one-hot state carrying occupancy zero, \[
|0\rangle_{b,i} \longmapsto |10\cdots 0\rangle,
\] or whichever one-hot convention is chosen consistently for the lowest
occupation.

\hypertarget{hubbard-hamiltonian-by-explicit-substitution}{%
\section{5. Hubbard Hamiltonian by Explicit
Substitution}\label{hubbard-hamiltonian-by-explicit-substitution}}

Write the static Hubbard Hamiltonian as \[
\hat H_{\mathrm{Hub}} = \hat H_t + \hat H_U + \hat H_v.
\]

\hypertarget{kinetic-term}{%
\subsection{5.1 Kinetic term}\label{kinetic-term}}

The kinetic term is \[
\hat H_t = -t\sum_{\langle i,j\rangle,\sigma}
\left(
\hat c_{i\sigma}^\dagger \hat c_{j\sigma}
+
\hat c_{j\sigma}^\dagger \hat c_{i\sigma}
\right).
\] For two fermion modes \(p<q\), direct Jordan-Wigner substitution
gives \[
\hat c_p^\dagger \hat c_q + \hat c_q^\dagger \hat c_p
=
\frac{1}{2}
\left(
X_p Z_{p+1}\cdots Z_{q-1} X_q
+
Y_p Z_{p+1}\cdots Z_{q-1} Y_q
\right).
\]

\hypertarget{onsite-interaction}{%
\subsection{5.2 Onsite interaction}\label{onsite-interaction}}

The onsite interaction term is \[
\hat H_U = U\sum_i \hat n_{i\uparrow}\hat n_{i\downarrow}
= U\sum_i \hat D_i.
\] Substituting the number primitive, \[
\hat H_U = \frac{U}{4}\sum_i (I-Z_{p(i,\uparrow)})(I-Z_{p(i,\downarrow)}).
\]

\hypertarget{static-potential-term}{%
\subsection{5.3 Static potential term}\label{static-potential-term}}

The static potential term is \[
\hat H_v = \sum_{i,\sigma} v_i \hat n_{i\sigma}.
\] After substitution, \[
\hat H_v = \sum_i v_i\left(I-\frac{1}{2}(Z_{p(i,\uparrow)}+Z_{p(i,\downarrow)})\right).
\]

\hypertarget{fully-substituted-hubbard-hamiltonian}{%
\subsection{5.4 Fully substituted Hubbard
Hamiltonian}\label{fully-substituted-hubbard-hamiltonian}}

Combining the pieces, \[
\hat H_{\mathrm{Hub}}
=
-t\sum_{\langle i,j\rangle,\sigma}
\left(
\hat c_{i\sigma}^\dagger \hat c_{j\sigma} + \hat c_{j\sigma}^\dagger \hat c_{i\sigma}
\right)
+
\frac{U}{4}\sum_i (I-Z_{p(i,\uparrow)})(I-Z_{p(i,\downarrow)})
+
\sum_i v_i\left(I-\frac{1}{2}(Z_{p(i,\uparrow)}+Z_{p(i,\downarrow)})\right),
\] with the kinetic term expanded further into Pauli strings as needed.

\hypertarget{hubbard-holstein-hamiltonian-by-explicit-substitution}{%
\section{6. Hubbard-Holstein Hamiltonian by Explicit
Substitution}\label{hubbard-holstein-hamiltonian-by-explicit-substitution}}

The driven Hubbard-Holstein Hamiltonian is \[
\hat H_{\mathrm{HH}}(t)
=
\hat H_t + \hat H_U + \hat H_v + \hat H_{\mathrm{ph}} + \hat H_g + \hat H_{\mathrm{drive}}(t).
\]

The driven onsite-density contribution is written in fixed-label Pauli
form as \[
\hat H_{\mathrm{drive}}(t)=\sum_{\lambda\in\mathcal L_{\mathrm{drive}}} c_\lambda(t) P_\lambda,
\] where the operator labels \(P_\lambda\) are static and only the
coefficients \(c_\lambda(t)\) vary in time. For the site-density drive
used in the active HH dynamics stack, \[
\Delta c[Z_{p(i,\uparrow)}](t)=\Delta c[Z_{p(i,\downarrow)}](t)=-\frac{1}{2}v_i(t),
\] and, if one retains the scalar offset explicitly, \[
\Delta c[I](t)=\sum_i v_i(t).
\]

\hypertarget{phonon-energy}{%
\subsection{6.1 Phonon energy}\label{phonon-energy}}

The phonon energy is \[
\hat H_{\mathrm{ph}} = \omega_0\sum_i \left(\hat n_{b,i}+\frac{1}{2}I\right).
\]

\hypertarget{unary-explicit-form}{%
\subsubsection{6.1.1 Unary explicit form}\label{unary-explicit-form}}

In unary form, \(\hat n_{b,i}\) is diagonal on the one-hot occupation
basis: \[
\hat n_{b,i} = \sum_{n=0}^{d-1} n\,|n\rangle_i\langle n|.
\] For the full qubit-level Pauli reduction in the spin-block JW + unary
setting, including the explicit \(I/Z\) form of
\(\hat H_{\mathrm{ph}}\), see Section 6.2.3.4(a).

\hypertarget{binary-explicit-form}{%
\subsubsection{6.1.2 Binary explicit form}\label{binary-explicit-form}}

In binary form, \[
\hat n_{b,i} = \sum_{n=0}^{d-1} n\,|\operatorname{bin}(n)\rangle_i\langle \operatorname{bin}(n)|,
\] understood as an operator on the truncated binary code space.

\hypertarget{electron-phonon-coupling}{%
\subsection{6.2 Electron-phonon
coupling}\label{electron-phonon-coupling}}

A natural density-shifted coupling is \[
\hat H_g = g\sum_i \hat x_i\bigl(\hat n_i-\bar n I\bigr),
\] where \(\bar n = N_e/L\) is the reference density, often \(\bar n=1\)
at half filling.

\hypertarget{unary-explicit-form-1}{%
\subsubsection{6.2.1 Unary explicit form}\label{unary-explicit-form-1}}

In unary encoding, \(\hat x_i\) acts as nearest-neighbor hopping on the
truncated local occupation chain: \[
\hat x_i = \sum_{n=0}^{d-2} \sqrt{n+1}\Bigl(|n+1\rangle_i\langle n| + |n\rangle_i\langle n+1|\Bigr).
\] Thus \[
\hat H_g = g\sum_i \hat x_i\bigl(\hat n_i-\bar n I\bigr)
\] with \(\hat x_i\) understood in that unary basis. For the detailed
spin-block JW + unary Pauli reduction of the Holstein sector, including
the explicit \(ZXX/ZYY\) expansion, see Sections 6.2.3.2--6.2.3.4.

\hypertarget{binary-explicit-form-1}{%
\subsubsection{6.2.2 Binary explicit
form}\label{binary-explicit-form-1}}

In binary encoding, the same operator is written on the binary code
space: \[
\hat x_i = \sum_{n=0}^{d-2} \sqrt{n+1}\Bigl(|\operatorname{bin}(n+1)\rangle_i\langle \operatorname{bin}(n)| + |\operatorname{bin}(n)\rangle_i\langle \operatorname{bin}(n+1)|\Bigr),
\] with the same density factor \((\hat n_i-\bar n I)\).

\hypertarget{worked-holstein-only-reduction-after-jordan-wigner-fermions-and-unary-phonons}{%
\subsection{6.2.3 Worked Holstein-only reduction after Jordan-Wigner
fermions and unary
phonons}\label{worked-holstein-only-reduction-after-jordan-wigner-fermions-and-unary-phonons}}

Set \[
N_b = n_{\mathrm{ph,max}}.
\] In this worked subsection, isolate the Holstein part of the
Hubbard-Holstein Hamiltonian: \[
\hat H_{\mathrm{Hol}}
=
\omega_0\sum_{i=0}^{L-1}\hat b_i^\dagger \hat b_i
\;+
\;g\sum_{i=0}^{L-1}(\hat b_i^\dagger+\hat b_i)(\hat n_i-1),
\qquad
\hat n_i=\hat n_{i\uparrow}+\hat n_{i\downarrow}.
\]

\hypertarget{qubit-indexing-used}{%
\subsubsection{6.2.3.1 Qubit indexing used}\label{qubit-indexing-used}}

For fermions, use spin-block Jordan-Wigner ordering: \[
q(i,\sigma)=i+\sigma L,
\qquad
\sigma=0\equiv\uparrow,
\quad
\sigma=1\equiv\downarrow.
\] Thus \[
q_{i\uparrow}=i,
\qquad
q_{i\downarrow}=i+L,
\qquad
i=0,\dots,L-1.
\]

For phonons, use a unary register with cutoff \(N_b\). Per site \(i\),
allocate \(N_b+1\) qubits labelled by \(n=0,\dots,N_b\), and place all
fermions first and phonons second: \[
q_{i,n}^{(b)} = 2L + i(N_b+1)+n,
\qquad
n=0,\dots,N_b.
\]

Let \(X_q\), \(Y_q\), and \(Z_q\) denote the Pauli operators on qubit
\(q\).

\hypertarget{jordan-wigner-fermionic-number-operators}{%
\subsubsection{6.2.3.2 Jordan-Wigner fermionic number
operators}\label{jordan-wigner-fermionic-number-operators}}

Jordan-Wigner gives \[
\hat n_{i\sigma} = \frac{I-Z_{q_{i\sigma}}}{2},
\] so \[
\hat n_i-1
=
\left(\frac{I-Z_{q_{i\uparrow}}}{2}+\frac{I-Z_{q_{i\downarrow}}}{2}\right)-1
=
-\frac{Z_{q_{i\uparrow}}+Z_{q_{i\downarrow}}}{2}.
\]

\hypertarget{unary-phonon-ladder-operators-in-pauli-form}{%
\subsubsection{6.2.3.3 Unary phonon ladder operators in Pauli
form}\label{unary-phonon-ladder-operators-in-pauli-form}}

Define the single-qubit combinations \[
(+)_{q}=\frac{X_q+iY_q}{2},
\qquad
(-)_{q}=\frac{X_q-iY_q}{2}.
\]

Then the unary truncated ladder operators are \[
\hat b_i^\dagger
=
\sum_{n=0}^{N_b-1}\sqrt{n+1}\,(+)_{q_{i,n}^{(b)}}\,(-)_{q_{i,n+1}^{(b)}},
\] \[
\hat b_i
=
\sum_{n=1}^{N_b}\sqrt{n}\,(-)_{q_{i,n-1}^{(b)}}\,(+)_{q_{i,n}^{(b)}}.
\]

A convenient real Pauli expansion is \[
\hat b_i^\dagger+\hat b_i
=
\frac12\sum_{n=0}^{N_b-1}\sqrt{n+1}
\Big(
X_{q_{i,n}^{(b)}}X_{q_{i,n+1}^{(b)}}
+
Y_{q_{i,n}^{(b)}}Y_{q_{i,n+1}^{(b)}}
\Big).
\]

\hypertarget{explicit-holstein-hamiltonian-in-the-qubit-basis}{%
\subsubsection{6.2.3.4 Explicit Holstein Hamiltonian in the qubit
basis}\label{explicit-holstein-hamiltonian-in-the-qubit-basis}}

\hypertarget{a-phonon-energy-term}{%
\paragraph{(a) Phonon energy term}\label{a-phonon-energy-term}}

On the unary one-hot code space, the phonon number operator is diagonal
and may be written as \[
\hat b_i^\dagger\hat b_i \equiv \hat n_i^{(b)}
=
\sum_{n=0}^{N_b} n\,\frac{I-Z_{q_{i,n}^{(b)}}}{2}
\qquad
\text{(exact on the one-hot subspace)}.
\] Therefore \[
\boxed{
\hat H_{\mathrm{ph}}
=
\omega_0\sum_{i=0}^{L-1}\sum_{n=0}^{N_b} n\,\frac{I-Z_{q_{i,n}^{(b)}}}{2}
}.
\]

Equivalently, as constant plus Pauli-\(Z\), \[
\hat H_{\mathrm{ph}}
=
\underbrace{
\omega_0\frac12\sum_i\sum_{n=0}^{N_b} n
}_{\displaystyle \omega_0 L\frac{N_b(N_b+1)}{4}}
I
-
\frac{\omega_0}{2}\sum_{i=0}^{L-1}\sum_{n=0}^{N_b} n\,Z_{q_{i,n}^{(b)}}.
\] The global constant shift may be dropped if desired.

\hypertarget{b-electron-phonon-coupling-term}{%
\paragraph{(b) Electron-phonon coupling
term}\label{b-electron-phonon-coupling-term}}

Using \[
\hat n_i-1 = -\frac{Z_{q_{i\uparrow}}+Z_{q_{i\downarrow}}}{2}
\] and the unary expansion of \(\hat b_i^\dagger+\hat b_i\), one obtains
\[
\boxed{
\hat H_{e\text{-}ph}
=
-\frac{g}{4}\sum_{i=0}^{L-1}\sum_{n=0}^{N_b-1}\sqrt{n+1}
\bigl(Z_{q_{i\uparrow}}+Z_{q_{i\downarrow}}\bigr)
\Big(
X_{q_{i,n}^{(b)}}X_{q_{i,n+1}^{(b)}}
+
Y_{q_{i,n}^{(b)}}Y_{q_{i,n+1}^{(b)}}
\Big)
}.
\]

The fully expanded Pauli-string sum is \[
\hat H_{e\text{-}ph}
=
-\frac{g}{4}\sum_{i=0}^{L-1}\sum_{n=0}^{N_b-1}\sqrt{n+1}
\Big[
Z_{q_{i\uparrow}}X_{q_{i,n}^{(b)}}X_{q_{i,n+1}^{(b)}}
+
Z_{q_{i\uparrow}}Y_{q_{i,n}^{(b)}}Y_{q_{i,n+1}^{(b)}}
+
Z_{q_{i\downarrow}}X_{q_{i,n}^{(b)}}X_{q_{i,n+1}^{(b)}}
+
Z_{q_{i\downarrow}}Y_{q_{i,n}^{(b)}}Y_{q_{i,n+1}^{(b)}}
\Big].
\] Each summand is three-local.

\hypertarget{final-holstein-hamiltonian-in-jw-unary-form}{%
\paragraph{Final Holstein Hamiltonian in JW + unary
form}\label{final-holstein-hamiltonian-in-jw-unary-form}}

Collecting the phonon-energy and electron-phonon pieces, \[
\boxed{
\hat H_{\mathrm{Hol}}^{\mathrm{(JW+unary)}}
=
\omega_0\sum_{i=0}^{L-1}\sum_{n=0}^{N_b} n\,\frac{I-Z_{q_{i,n}^{(b)}}}{2}
-
\frac{g}{4}\sum_{i=0}^{L-1}\sum_{n=0}^{N_b-1}\sqrt{n+1}
\bigl(Z_{q_{i\uparrow}}+Z_{q_{i\downarrow}}\bigr)
\Big(
X_{q_{i,n}^{(b)}}X_{q_{i,n+1}^{(b)}}
+
Y_{q_{i,n}^{(b)}}Y_{q_{i,n+1}^{(b)}}
\Big)
}.
\]

\hypertarget{time-dependent-density-drive}{%
\subsection{6.3 Time-dependent density
drive}\label{time-dependent-density-drive}}

A general density drive may be written as \[
\hat H_{\mathrm{drive}}(t) = \sum_i v_i(t)\hat n_i.
\] A useful waveform class is \[
v_i(t)=A s_i \sin(\omega t + \phi)\exp\!\left[-\frac{(t-t_0)^2}{2\bar t^2}\right],
\] where \(A\) is the amplitude, \(s_i\) is a static spatial pattern,
\(\omega\) is the drive frequency, \(\phi\) is the phase, and \(\bar t\)
is the envelope width.

\hypertarget{generic-density-drive-waveform-surface}{%
\subsubsection{6.3.1 Generic density-drive waveform
surface}\label{generic-density-drive-waveform-surface}}

After inserting
\(\hat n_i=I-\tfrac12(Z_{p(i,\uparrow)}+Z_{p(i,\downarrow)})\), the
drive separates into an identity offset and a sum of time-dependent
\(Z\)-coefficients. Thus the time dependence resides in the coefficients
rather than in the operator labels.

\hypertarget{two-hh-assembly-surfaces}{%
\subsection{6.4 Two HH assembly
surfaces}\label{two-hh-assembly-surfaces}}

\hypertarget{exact-hamiltonian-surface}{%
\subsubsection{6.4.1 Exact Hamiltonian
surface}\label{exact-hamiltonian-surface}}

The exact surface treats \(\hat H_{\mathrm{HH}}(t)\) as an operator to
be diagonalized, exponentiated, projected, or applied directly to
statevectors.

\hypertarget{variational-hh-ansatz-surface}{%
\subsubsection{6.4.2 Variational HH ansatz
surface}\label{variational-hh-ansatz-surface}}

The variational surface treats the same physical pieces as generator
families from which one builds a structured manifold \[
|\psi(\theta)\rangle = U(\theta)|\psi_{\mathrm{ref}}\rangle.
\]

\hypertarget{fully-substituted-hh-master-form}{%
\subsection{6.5 Fully substituted HH master
form}\label{fully-substituted-hh-master-form}}

The full symbolic master form is therefore \[
\hat H_{\mathrm{HH}}(t)
=
\hat H_{\mathrm{Hub}}
+
\omega_0\sum_i \left(\hat n_{b,i}+\frac{1}{2}I\right)
+
 g\sum_i \hat x_i(\hat n_i-\bar n I)
+
\sum_i v_i(t)\hat n_i,
\] with \(\hat n_i\), \(\hat n_{b,i}\), and \(\hat x_i\) expanded
according to the chosen fermion ordering and boson encoding.

\hypertarget{reference-states-and-exact-sector-ed}{%
\section{7. Reference States and Exact Sector
ED}\label{reference-states-and-exact-sector-ed}}

\hypertarget{fermionic-hartree-fock-determinant}{%
\subsection{7.1 Fermionic Hartree-Fock
determinant}\label{fermionic-hartree-fock-determinant}}

Let \(\mathcal O\) denote the occupied spin-orbital set in a Slater
determinant reference. Then \[
|\Phi_{\mathrm{HF}}\rangle = \prod_{p\in\mathcal O} \hat c_p^\dagger |\mathrm{vac}_f\rangle.
\]

\hypertarget{hubbard-holstein-reference-state}{%
\subsection{7.2 Hubbard-Holstein reference
state}\label{hubbard-holstein-reference-state}}

The natural HH reference is the tensor product of the phonon vacuum and
fermionic Hartree-Fock state: \[
|\psi_{\mathrm{ref}}^{\mathrm{HH}}\rangle = |\mathrm{vac}_{\mathrm{ph}}\rangle \otimes |\Phi_{\mathrm{HF}}\rangle.
\]

\hypertarget{exact-hh-sector-basis}{%
\subsection{7.3 Exact HH sector basis}\label{exact-hh-sector-basis}}

Fixing the fermion particle numbers \((N_\uparrow,N_\downarrow)\), the
exact sector basis is \[
\mathcal H_{\mathrm{sector}} = \operatorname{span}\{|f\rangle\otimes|\mathbf n_b\rangle : |f\rangle\in\mathcal H_{N_\uparrow,N_\downarrow},\; 0\le n_{b,i}\le n_{\mathrm{ph,max}}\}.
\]

\hypertarget{binary-index-map}{%
\subsubsection{7.3.1 Binary index map}\label{binary-index-map}}

For binary boson encoding, one may order basis states by the pair \[
(f,\mathbf n_b),
\] where \(f\) is the fermionic bit pattern and
\(\mathbf n_b=(n_{b,0},\dots,n_{b,L-1})\) is the vector of local phonon
occupations, each encoded in binary.

\hypertarget{unary-index-map}{%
\subsubsection{7.3.2 Unary index map}\label{unary-index-map}}

For unary boson encoding, the same sector is indexed by \[
(f,\mathbf e_{n_{b,0}},\dots,\mathbf e_{n_{b,L-1}}),
\] where \(\mathbf e_n\) denotes the unary one-hot representation of
occupation \(n\).

\hypertarget{exact-hh-matrix-elements}{%
\subsection{7.4 Exact HH matrix
elements}\label{exact-hh-matrix-elements}}

\hypertarget{diagonal-elements}{%
\subsubsection{7.4.1 Diagonal elements}\label{diagonal-elements}}

For basis state \(|f\rangle\otimes|\mathbf n_b\rangle\), the diagonal
contribution is \[
U\sum_i n_{i\uparrow}(f)n_{i\downarrow}(f)
+
\sum_i v_i\,n_i(f)
+
\omega_0\sum_i \left(n_{b,i}+\frac{1}{2}\right).
\]

\hypertarget{hopping-matrix-elements}{%
\subsubsection{7.4.2 Hopping matrix
elements}\label{hopping-matrix-elements}}

If \(|f'\rangle\) differs from \(|f\rangle\) by a single allowed
nearest-neighbor hop with the same phonon occupations, then \[
\langle f',\mathbf n_b|\hat H_t|f,\mathbf n_b\rangle
=
-t\,(-1)^{\chi(f;p,q)},
\] where \(\chi(f;p,q)\) is the Jordan-Wigner parity count between the
two fermion modes \(p\) and \(q\).

\hypertarget{electron-phonon-matrix-elements}{%
\subsubsection{7.4.3 Electron-phonon matrix
elements}\label{electron-phonon-matrix-elements}}

At fixed fermion configuration, the electron-phonon term changes the
local phonon number by one. For site \(i\), \[
\langle n_{b,i}+1|\hat x_i|n_{b,i}\rangle = \sqrt{n_{b,i}+1},
\qquad
\langle n_{b,i}-1|\hat x_i|n_{b,i}\rangle = \sqrt{n_{b,i}}.
\] Hence the off-diagonal electron-phonon matrix elements are
proportional to \[
g\,(n_i(f)-\bar n)\sqrt{n_{b,i}+1}
\quad\text{or}\quad
g\,(n_i(f)-\bar n)\sqrt{n_{b,i}},
\] according to whether the phonon occupation increases or decreases.

\hypertarget{statevector-action-expectation-and-exponential-primitives}{%
\section{8. Statevector Action, Expectation, and Exponential
Primitives}\label{statevector-action-expectation-and-exponential-primitives}}

\hypertarget{basis-state-primitive}{%
\subsection{8.1 Basis-state primitive}\label{basis-state-primitive}}

A statevector on \(N_q\) qubits is written as \[
|\psi\rangle = \sum_{k=0}^{2^{N_q}-1} a_k |k\rangle
= \sum_{b\in\{0,1\}^{N_q}} a_b |b\rangle.
\]

\hypertarget{explicit-pauli-word-action}{%
\subsection{8.2 Explicit Pauli-word
action}\label{explicit-pauli-word-action}}

Let \(P\) be a Pauli word. Then its action on a computational basis
state can be written in the form \[
P|b\rangle = \omega_P(b)\,|b\oplus \chi_P\rangle,
\] where \(\chi_P\) records which qubits are flipped by \(X\) or \(Y\),
and \(\omega_P(b)\) is the phase determined by the \(Y\) and \(Z\)
letters acting on the bitstring \(b\).

\hypertarget{expectation-values}{%
\subsection{8.3 Expectation values}\label{expectation-values}}

If \[
\hat H = \sum_\alpha h_\alpha P_\alpha,
\] then the expectation value is \[
E(\psi)=\langle \psi|\hat H|\psi\rangle = \sum_\alpha h_\alpha \langle \psi|P_\alpha|\psi\rangle.
\]

\hypertarget{pauli-rotations}{%
\subsection{8.4 Pauli rotations}\label{pauli-rotations}}

For any Hermitian Pauli word with \(P^2=I\), \[
e^{-i\theta P} = \cos\theta\,I - i\sin\theta\,P.
\] This identity is the primitive behind fast statevector application of
Pauli exponentials.

\hypertarget{exponential-of-a-pauli-polynomial}{%
\subsection{8.5 Exponential of a Pauli
polynomial}\label{exponential-of-a-pauli-polynomial}}

For a Pauli polynomial \[
A = \sum_\alpha a_\alpha P_\alpha,
\] one has \[
U_A(\theta)=e^{-i\theta A}.
\] If the \(P_\alpha\) commute pairwise, then \[
U_A(\theta)=\prod_\alpha e^{-i\theta a_\alpha P_\alpha}.
\] If they do not commute, one may use either exact exponentiation on
the working subspace or an ordered product formula.

\hypertarget{exact-sector-energy-target}{%
\subsection{8.6 Exact sector energy
target}\label{exact-sector-energy-target}}

The natural exact comparison target in a fixed fermion sector is \[
E_{\mathrm{exact,sector}} = \min_{\psi\in\mathcal H_{\mathrm{sector}}} \langle \psi|\hat H|\psi\rangle.
\] This is distinct from the unrestricted full-Hilbert minimum when a
sector constraint is imposed.

\hypertarget{fixed-vqe-ansatz-families}{%
\section{9. Fixed VQE Ansatz Families}\label{fixed-vqe-ansatz-families}}

This section records the main symbolic manifold families without tying
them to any implementation names.

\hypertarget{hubbard-layerwise-ansatz}{%
\subsection{9.1 Hubbard layerwise
ansatz}\label{hubbard-layerwise-ansatz}}

For the static Hubbard model, write \[
\hat H_{\mathrm{Hub}} = \hat H_t + \hat H_U + \hat H_v.
\] A layerwise ansatz uses one parameter per physical group per layer:
\[
U_{\mathrm{Hub}}^{(\ell)}
=
\exp(-i\theta_t^{(\ell)}\hat H_t)
\exp(-i\theta_U^{(\ell)}\hat H_U)
\exp(-i\theta_v^{(\ell)}\hat H_v),
\] with the understanding that each grouped generator may itself be
split into Pauli exponentials internally while sharing the same scalar
parameter within the group.

\hypertarget{hubbard-excitation-based-surface}{%
\subsection{9.2 Hubbard excitation-based
surface}\label{hubbard-excitation-based-surface}}

Relative to a Hartree-Fock occupied/virtual partition, single- and
double-excitation generators take the form \[
G_{i\to a}^{(\sigma)} = i\left(\hat c_{a\sigma}^\dagger \hat c_{i\sigma} - \hat c_{i\sigma}^\dagger \hat c_{a\sigma}\right),
\] and \[
G_{ij\to ab}^{(\sigma,\sigma')} = i\left(\hat c_{a\sigma}^\dagger \hat c_{b\sigma'}^\dagger \hat c_{j\sigma'}\hat c_{i\sigma} - \text{h.c.}\right).
\] A layered excitation manifold is then \[
U_{\mathrm{exc}}(\theta)=\prod_r e^{-i\theta_r G_r}.
\]

\hypertarget{hh-layerwise-ansatz}{%
\subsection{9.3 HH layerwise ansatz}\label{hh-layerwise-ansatz}}

For the HH model, define grouped generators \[
G_t,\;G_U,\;G_v,\;G_{\mathrm{ph}},\;G_g,\;G_{\mathrm{drive}}.
\] One layer is \[
U_{\mathrm{HH,lw}}^{(\ell)}
=
\exp(-i\theta_{t,\ell}G_t)
\exp(-i\theta_{U,\ell}G_U)
\exp(-i\theta_{v,\ell}G_v)
\exp(-i\theta_{\mathrm{ph},\ell}G_{\mathrm{ph}})
\exp(-i\theta_{g,\ell}G_g)
\exp(-i\theta_{\mathrm{drive},\ell}G_{\mathrm{drive}}),
\] with inactive groups omitted.

\hypertarget{hh-physical-termwise-ansatz}{%
\subsection{9.4 HH physical-termwise
ansatz}\label{hh-physical-termwise-ansatz}}

A physical-termwise layer sits between the grouped and fully split
extremes. Let \[
\mathcal A_{\mathrm{phys}} = \{A_m\}
\] be the collection of unsplit physical HH generators such as bond
hopping terms, onsite interactions, local phonon energies, local
couplings, and local drive terms. Then \[
U_{\mathrm{HH,phys}}^{(\ell)} = \prod_{A_m\in\mathcal A_{\mathrm{phys}}} e^{-i\phi_{m,\ell}A_m}.
\] Because the generators remain grouped at the physical-operator level,
this fixed-family surface stays aligned with the same conserved-sector
story as the HH layerwise surface. In this symbolic HH manuscript, the
canonical fixed-family surface is intentionally restricted to the
sector-preserving HH layerwise and HH physical-termwise families.

\hypertarget{reference-state-and-ansatz-pairing-principle}{%
\subsection{9.5 Reference-state and ansatz pairing
principle}\label{reference-state-and-ansatz-pairing-principle}}

The reference state and variational family should be paired by three
criteria:

\begin{enumerate}
\def\labelenumi{\arabic{enumi}.}
\tightlist
\item
  the conserved sector they preserve,
\item
  the expressive scale needed by the target Hamiltonian,
\item
  the conditioning and optimization burden induced by the chosen
  parameterization.
\end{enumerate}

\hypertarget{polaronic-operator-families}{%
\section{10. Polaronic Operator
Families}\label{polaronic-operator-families}}

\hypertarget{primitive-polaronic-ingredients}{%
\subsection{10.1 Primitive polaronic
ingredients}\label{primitive-polaronic-ingredients}}

\hypertarget{shifted-density}{%
\subsubsection{10.1.1 Shifted density}\label{shifted-density}}

If the total electron number is \(N_e\), define the mean density \[
\bar n = \frac{N_e}{L},
\] and the shifted density \[
\tilde n_i = \hat n_i - \bar n I.
\] At half filling \(N_e=L\), this reduces to \[
\tilde n_i = -\frac{1}{2}(Z_{p(i,\uparrow)}+Z_{p(i,\downarrow)}).
\]

\hypertarget{phonon-primitives-p_i-and-x_i}{%
\subsubsection{\texorpdfstring{10.1.2 Phonon primitives \(P_i\) and
\(x_i\)}{10.1.2 Phonon primitives P\_i and x\_i}}\label{phonon-primitives-p_i-and-x_i}}

Define \[
P_i=i(\hat b_i^\dagger-\hat b_i),
\qquad
x_i=\hat b_i+\hat b_i^\dagger.
\]

\hypertarget{local-squeeze-primitive}{%
\subsubsection{10.1.3 Local squeeze
primitive}\label{local-squeeze-primitive}}

A natural local squeeze generator is \[
S_i = i\left((\hat b_i^\dagger)^2-\hat b_i^2\right).
\]

\hypertarget{doublon-primitive}{%
\subsubsection{10.1.4 Doublon primitive}\label{doublon-primitive}}

The doublon projector is \[
D_i = \hat n_{i\uparrow}\hat n_{i\downarrow}.
\]

\hypertarget{even-hopping-channel}{%
\subsubsection{10.1.5 Even hopping channel}\label{even-hopping-channel}}

A bond-even hopping operator is \[
T_{ij}^{(+)} = \sum_\sigma\left(\hat c_{i\sigma}^\dagger\hat c_{j\sigma}+\hat c_{j\sigma}^\dagger\hat c_{i\sigma}\right).
\]

\hypertarget{odd-current-channel}{%
\subsubsection{10.1.6 Odd current channel}\label{odd-current-channel}}

A bond-odd current operator is \[
J_{ij}^{(-)} = i\sum_\sigma\left(\hat c_{i\sigma}^\dagger\hat c_{j\sigma}-\hat c_{j\sigma}^\dagger\hat c_{i\sigma}\right).
\]

\hypertarget{representative-dressed-channels}{%
\subsection{10.2 Representative dressed
channels}\label{representative-dressed-channels}}

The exact family boundaries can vary, but the symbolic equivalents of
the common channels can be written as follows.

\hypertarget{conditional-displacement}{%
\subsubsection{10.2.1 Conditional
displacement}\label{conditional-displacement}}

\[
G_i^{(\mathrm{cd})} = \tilde n_i P_i.
\]

\hypertarget{doublon-dressing}{%
\subsubsection{10.2.2 Doublon dressing}\label{doublon-dressing}}

\[
G_i^{(\mathrm{dd})} = D_i P_i.
\]

\hypertarget{dressed-hopping-drag}{%
\subsubsection{10.2.3 Dressed hopping drag}\label{dressed-hopping-drag}}

\[
G_{ij}^{(\mathrm{hd})} = T_{ij}^{(+)}(P_i-P_j).
\]

\hypertarget{odd-channel-drag}{%
\subsubsection{10.2.4 Odd-channel drag}\label{odd-channel-drag}}

\[
G_{ij}^{(\mathrm{od})} = J_{ij}^{(-)}(x_i-x_j).
\]

\hypertarget{second-order-even-channel}{%
\subsubsection{10.2.5 Second-order even
channel}\label{second-order-even-channel}}

\[
G_{ij}^{(2)} = T_{ij}^{(+)}(x_i-x_j)^2.
\]

\hypertarget{third-order-odd-current-channel}{%
\subsubsection{10.2.6 Third-order odd current
channel}\label{third-order-odd-current-channel}}

\[
G_{ij}^{(3)} = J_{ij}^{(-)}(x_i-x_j)^3.
\]

\hypertarget{fourth-order-even-hopping-channel}{%
\subsubsection{10.2.7 Fourth-order even hopping
channel}\label{fourth-order-even-hopping-channel}}

\[
G_{ij}^{(4)} = T_{ij}^{(+)}(x_i-x_j)^4.
\]

\hypertarget{bond-conditioned-displacement-and-squeeze-channels}{%
\subsubsection{10.2.8 Bond-conditioned displacement and squeeze
channels}\label{bond-conditioned-displacement-and-squeeze-channels}}

With a bond-sensitive scalar observable \(B_{ij}\), one may write \[
G_{ij}^{(\mathrm{bd})}=B_{ij}P_i,
\qquad
G_{ij}^{(\mathrm{bs})}=B_{ij}S_i.
\] A simple choice is \(B_{ij}=\hat n_i-\hat n_j\).

\hypertarget{local-squeeze-channels}{%
\subsubsection{10.2.9 Local squeeze
channels}\label{local-squeeze-channels}}

\[
G_i^{(\mathrm{sq})}=S_i,
\qquad
\tilde G_i^{(\mathrm{sq})}=\tilde n_i S_i.
\]

\hypertarget{extended-cloud-channels}{%
\subsubsection{10.2.10 Extended cloud
channels}\label{extended-cloud-channels}}

A symbolic extended cloud family is \[
G_i^{(\mathrm{cloud})}=\tilde n_i\sum_{r\in\mathcal N(i)} \alpha_{ir} P_r,
\] with coefficients \(\alpha_{ir}\) controlling cloud shape.

\hypertarget{doublon-translation-and-doublon-squeeze-channels}{%
\subsubsection{10.2.11 Doublon-translation and doublon-squeeze
channels}\label{doublon-translation-and-doublon-squeeze-channels}}

\[
G_{ij}^{(\mathrm{DT})}=D_i T_{ij}^{(+)},
\qquad
G_i^{(\mathrm{DS})}=D_i S_i.
\]

\hypertarget{pool-family-map}{%
\subsection{10.3 Pool-family map}\label{pool-family-map}}

A useful abstract partition is \[
\mathcal G = \mathcal G_{\mathrm{local}} \cup \mathcal G_{\mathrm{bond}} \cup \mathcal G_{\mathrm{squeeze}} \cup \mathcal G_{\mathrm{cloud}} \cup \mathcal G_{\mathrm{doublon}}.
\] These families can be opened gradually rather than all at once.

\hypertarget{compact-displacementsqueeze-macro-families}{%
\subsubsection{10.3.1 Compact displacement/squeeze macro
families}\label{compact-displacementsqueeze-macro-families}}

Two especially compact subfamilies are \[
\mathcal G_{\mathrm{VLF}} = \{\tilde n_i P_i\}_i,
\qquad
\mathcal G_{\mathrm{SQ}} = \{S_i,\tilde n_i S_i\}_i.
\] They act as coarse polaronic and squeeze envelopes.

\hypertarget{adaptive-selection-and-staged-continuation}{%
\section{11. Adaptive Selection and Staged
Continuation}\label{adaptive-selection-and-staged-continuation}}

\hypertarget{cumulative-phase-construction-for-adaptive-selection}{%
\subsection{11.1 Cumulative phase construction for adaptive
selection}\label{cumulative-phase-construction-for-adaptive-selection}}

If a shortlisted macro generator \(m\) admits a split family
\(\mathcal C_{\mathrm{split}}(m)=\{m_1,\dots,m_{K_m}\}\), then each
split child inherits the same position \(p\) and therefore defines a
child record \(r_j=(m_j,p)\). For simplicity we drop the \(j\), and
assume each \(m\) below is a unique generator.

Record-first view: the primary object is a candidate-position record
\(r:=(m,p)\), with projections \(\pi_m(m,p):=m\) and \(\pi_p(m,p):=p\).
For each phase \(k \in \{1,2,3\}\), the one-way selector flow is \[
\mathcal R_k\longrightarrow\mathcal S_k\longrightarrow r_k^\star,
\] with \[
\mathcal R_k:=\{(m,p):m\in\mathcal G_k,\ p\in\mathcal P_m^{k}\},
\qquad
\mathcal S_k:=
\{(m,p)\in\mathcal R_k:S_k(m,p)\ge\tau_k\}
\quad\text{or}\quad
\operatorname{Top}_{N_k}\!\left(S_k;\mathcal R_k\right),
\qquad
r_k^\star\in\mathcal S_k.
\] where \(S_k\) denotes the phase-\(k\) score used for shortlist
construction. Generator image and position image:
\(\mathcal G^{(k)}:=\pi_m(\mathcal S_k)\) and
\(\mathcal P_m^{(k)}:=\{p:(m,p)\in\mathcal S_k\}\). Record inheritance:
when later phases inherit the same candidate-position records, the clean
handoff is \(\mathcal R_{k+1}:=\mathcal S_k\). Hence the cumulative
record chain is
\(\mathcal R_3 \subseteq \mathcal R_2 \subseteq \mathcal R_1\), the
cumulative generator chain is
\(\mathcal G^{(3)} \subseteq \mathcal G^{(2)} \subseteq \mathcal G^{(1)} \subseteq \mathcal G\),
the first chain follows from
\(\mathcal R_{k+1}:=\mathcal S_k\subseteq\mathcal R_k\), and the
retained position images inherit analogously via the corresponding
record shortlists \(\mathcal S_1\to\mathcal S_2\to\mathcal S_3\).
Post-shortlist admission, batching, pruning, and beam continuation are
then defined in terms of the final Phase-3 surface.

Seed broadening: in addition, we use a smooth broadening of the initial
candidate record \(\mathcal R\) as expected generator utility saturates,
anti-comensurate to scaffold depth, where the initial seeds are
low-order generators and late-stage seeds are higher-order excitation
and coupling terms.

Chronology: thus the presentation chronologically accords with the
algorithmic implementation: Phase 1 \(\to\) Phase 2 \(\to\) Phase 3,
with each to-arrow filtering its preceding phase.

\hypertarget{phase-1}{%
\subsection{11.2 Phase 1}\label{phase-1}}

\hypertarget{core-signal-append-position-and-one-coordinate-local-model}{%
\subsubsection{11.2.1 Core signal, append position, and one-coordinate
local
model}\label{core-signal-append-position-and-one-coordinate-local-model}}

Core record, scaffold, and append slot: \(r=(m,p)\),
\(\mathcal O=(a_1,\dots,a_n)\), \(m\in\mathcal G\), and
\(p\in\mathcal P_m^{1}=\{0,\dots,n\}\). Here \(m\) is a candidate
generator already admitted by the generator gate (so
\(m\in\mathcal G_1\)), and \(p\) is the index of a slot in the
pre-insertion, ordered scaffold \(\mathcal O\). Window bias: note, the
insertion window near the terminal scaffold point is biased to be small
proportional to expected scaffold utility remaining.

Scaffold parameter space and energy map: WLOG consider \(\mathcal O\),
s.t. each \(a\) is assigned one scalar \(\theta (\mod 2 \pi)\), then
define \(\Theta_{\mathcal O}:=\mathbb R^{|\mathcal O|}\),
\(E(\mathcal O,\theta):=\langle\psi(\mathcal O,\theta)|H|\psi(\mathcal O,\theta)\rangle\),
and \(\theta\in\Theta_{\mathcal O}\). Scaffold insertion and parameter
insertion: define
\(\mathcal O\oplus_p m:=(a_1,\dots,a_p,m,a_{p+1},\dots,a_n)\) and
\(\theta\oplus_p\vartheta:=(\theta_1,\dots,\theta_p,\vartheta,\theta_{p+1},\dots,\theta_n)\in\Theta_{\mathcal O\oplus_p m}\).

Write \(|\psi\rangle:=U_{>p}U_{\le p}|\psi_0\rangle\),
\(\widetilde A_r:=U_{>p}A_mU_{>p}^\dagger\),
\(U_r(\alpha):=e^{-i\alpha \widetilde A_r}\), and
\(|\psi_r(\alpha)\rangle:=U_r(\alpha)|\psi\rangle\). Let
\(Q_\psi:=I-|\psi\rangle\langle\psi|\),
\(|d_r\rangle:=\left.\partial_\alpha|\psi_r(\alpha)\rangle\right|_{\alpha=0}=-i\widetilde A_r|\psi\rangle\),
\(\tau_r:=Q_\psi|d_r\rangle\), \(g_r:=2\Re\langle \psi|H|d_r\rangle\),
\(\sigma_r:=\operatorname{SE}\!\left[\widehat g(r)\right]\ge 0\),
\(g_{\mathrm{lcb}}(r):=\max\{|g_r|-z_\alpha\sigma_r,0\}\), and
\(F_{\mathrm{raw}}(r):=\langle \tau_r,\tau_r\rangle+\varepsilon_F\). The
live broad Phase-1 one-coordinate trust-region proxy is \[
\Delta E_{1,\mathrm{TR}}(r)
:=
\max_{|\alpha|\le \rho/\sqrt{F_{\mathrm{raw}}(r)}}
\left[
g_{\mathrm{lcb}}(r)|\alpha|
-\frac12 \lambda_F F_{\mathrm{raw}}(r)\alpha^2
\right].
\] The strict one-coordinate insertion drop is retained as the exact
auxiliary comparator \[
\delta E_1(m,p\mid\mathcal O,\theta)
=
E(\mathcal O,\theta)
-\inf_{\vartheta\in\mathbb R}E\bigl(\mathcal O\oplus_p m,\theta\oplus_p\vartheta\bigr).
\] The infimum notation may be replaced by \(\min\) or
\(\operatorname{Min}\) whenever the minimizer exists. Thus
\(\Delta E_{1,\mathrm{TR}}(r)\) is the live broad-pool second-order
local insertion proxy, whereas \(\delta E_1(m,p\mid\mathcal O,\theta)\)
is the exact one-coordinate insertion drop.
\footnote{User interest: How does window parameterization width scale with cost?}

\hypertarget{gates-burdens-and-feature-ingredients}{%
\subsubsection{11.2.2 Gates, burdens, and feature
ingredients}\label{gates-burdens-and-feature-ingredients}}

Base hardware ingredients: each candidate-position pair carries the
proxy triple \(C_{\mathrm{2q}}(m,p)\), \(D(m,p)\), and
\(\Delta K(m,p)\), with fixed nonnegative weights
\(\lambda_{\mathrm{2q}},\lambda_d,\lambda_\theta\).\footnote{Admissibility gate: For each selector stage $k$, let $\Gamma_k:\mathcal G_k^{\mathrm{in}}\to\{0,1\}$ be the generator-level admissibility gate and define $\mathcal G_k^{\mathrm{adm}}:=\{m\in\mathcal G_k^{\mathrm{in}}:\Gamma_k(m)=1\}$. Record domain: the working candidate-record domain is built only from this gated generator pool, $\mathcal R_k:=\{(m,p):m\in\mathcal G_k^{\mathrm{adm}},\ p\in\mathcal P_m^{\mathrm{in},k}\}$. Selector action: all selectors below act on $\mathcal R_k$ itself, so admissibility is enforced by pre-score domain construction rather than post hoc score masking. Tie order: fix once and for all a total order $\prec$ on records (for example, lexicographic order on $(m,p)$). Tie-broken maximizer: for any score $S:\mathcal R_k\to\mathbb R$, define $\arg\max^{\prec}_{r\in\mathcal R_k} S(r):=\min_{\prec}\{r\in\mathcal R_k:S(r)=\max_{u\in\mathcal R_k}S(u)\}$, and similarly for $\arg\min^{\prec}$. Shortlist surface: define $\operatorname{Top}_K^{\prec}(S;\mathcal R_k)$ as the ordered list of the first $K$ records obtained by sorting $\mathcal R_k$ in descending score, with ties broken by $\prec$. Unadorned $\arg\max$, $\arg\min$, and $\operatorname{Top}_K$ below follow this convention.}
Here \(C_{\mathrm{2q}}\) is a two-qubit proxy cost, \(D\) is an expected
depth-induced by \(m\) (at \(p\)) proxy, and \(\Delta K\) is the
marginal parameter count. Positivity assumptions: assume throughout that
\(\lambda_{\mathrm{2q}}\ge 0\), \(C_{\mathrm{2q}}(m,p)\ge 0\), and
\(\lambda_\theta,\lambda_d,\Delta K(m,p),D(m,p)>0\) on admissible
records.
\footnote{Depth and parameter count are constants for our L=2 pareto lean runs.}

\hypertarget{anticipated-utility-left-and-controller-schedule}{%
\subsubsection{11.2.3 Anticipated utility left and controller
schedule}\label{anticipated-utility-left-and-controller-schedule}}

Let \(t\) denote the current selector step, let \(D_{\mathrm{left}}(t)\)
be the remaining controller depth budget, and let
\(\widehat N_{\mathrm{rem}}(t)\) be a heuristic forecast (expected
count) of useful remaining admissions. For gain-process statistics,
define \(\Delta_{\mathrm{adm}}(s):=E(\theta^{(s-1)})-E(\theta^{(s)})\),
\(\Delta_{\mathrm{adm}}^{+}(s):=\operatorname{asinh}\!\left(\frac{\Delta_{\mathrm{adm}}(s)}{\tau_\Delta}\right)\),
and let \(\mathcal I_t^+\) be the index set of past selector steps with
valid admissions used for EWMA/MAD statistics.

Let \(D_{\mathrm{allow}}\) denote the total allowed depth budget, let
\(\ell_t^{\mathrm{plat}}:=\max\{1,\lceil P\,D_{\mathrm{left}}(t)/D_{\mathrm{allow}}\rceil\}\)
and
\(\ell_t^{\mathrm{low}}:=\max\{1,\lceil L\,D_{\mathrm{left}}(t)/D_{\mathrm{allow}}\rceil\}\),
and define the relative-flattening and absolute weak-gain thresholds by
\(\tau_{\mathrm{plat}}^{+}(t):=\bigl(\eta_0+\eta_1(1-D_{\mathrm{left}}(t)/D_{\mathrm{allow}})\bigr)m_t\)
and
\(\tau_{\mathrm{low}}^{+}(t):=\operatorname{asinh}(\tau_{\mathrm{use},t}/\tau_\Delta)\).
The corresponding depth-adaptive streak statistics are
\(\operatorname{plateau\_streak}(t):=\frac{P}{\ell_t^{\mathrm{plat}}}\max\{q\in\{0,\dots,\ell_t^{\mathrm{plat}}\}:\Delta_{\mathrm{adm}}^{+}(t-j+1)\le \tau_{\mathrm{plat}}^{+}(t)\ \forall\,1\le j\le q\}\)
and
\(\operatorname{low\_streak}(t):=\frac{L}{\ell_t^{\mathrm{low}}}\max\{q\in\{0,\dots,\ell_t^{\mathrm{low}}\}:\Delta_{\mathrm{adm}}^{+}(t-j+1)\le \tau_{\mathrm{low}}^{+}(t)\ \forall\,1\le j\le q\}\).

For stagnation, gain, decay, and scale summaries, define
\(u_{\mathrm{stag}}(t):=\operatorname{clip}\!\left(\max\!\left(\frac{\mathrm{plateau\_streak}(t)}{P},\frac{\mathrm{low\_streak}(t)}{L}\right),0,1\right)\),
\(m_t:=\operatorname{EWMA}_{s\in\mathcal I_t^+}\!\bigl(\Delta_{\mathrm{adm}}^{+}(s)\bigr)\),
\(\rho_t:=\operatorname{clip}\!\left(\frac{\operatorname{EWMA}_{\mathrm{recent}}(\Delta_{\mathrm{adm}})}{\operatorname{EWMA}_{\mathrm{older}}(\Delta_{\mathrm{adm}})+\varepsilon},\rho_{\min},1\right)\),
\(\gamma_t:=[-\log\rho_t]_+\), and
\(s_t:=\max\!\left(s_{\min},\operatorname{MAD}_{s\in\mathcal I_t^+}\!\bigl(\Delta_{\mathrm{adm}}^{+}(s)\bigr)\right)\).

For the frontier statistic, let
\(s_{\mathrm{raw},1}^{(1)}(t)\ge s_{\mathrm{raw},2}^{(1)}(t)\ge 0\)
denote the first two order statistics of the current broad Phase-1 score
multiset \(\{S_1(r;t):r\in\mathcal R_1\}\) before any
\(\tau_1(t)\)-thresholding or \(N_1(t)\)-truncation, with the
conventions
\(s_{\mathrm{raw},1}^{(1)}(t)=s_{\mathrm{raw},2}^{(1)}(t)=0\) if
\(\mathcal R_1=\varnothing\) and \(s_{\mathrm{raw},2}^{(1)}(t)=0\) if
\(|\mathcal R_1|=1\). Define \[
u_{\mathrm{front}}(t)
:=
\frac{s_{\mathrm{raw},2}^{(1)}(t)+\varepsilon}{s_{\mathrm{raw},1}^{(1)}(t)+\varepsilon}.
\] Thus \(u_{\mathrm{front}}(t)\in[0,1]\), with larger values
corresponding to a flatter, more crowded current frontier and smaller
values corresponding to a clearly separated leading
record.\footnote{With the empty-pool convention $s_{\mathrm{raw},1}^{(1)}(t)=s_{\mathrm{raw},2}^{(1)}(t)=0$, one gets $u_{\mathrm{front}}(t)=1$. This treats the absence of a usable current frontier as maximal compression rather than as a missing-data null signal. An explicit signal-to-noise controller statistic is a possible implementation here, but it may not be necessary if worsening signal-to-noise already manifests as frontier flattening, which the present expected-utility model scores directly through $u_{\mathrm{front}}(t)$. Here $h_t(j)$ is the conditional probability that usefulness terminates at future offset $j$, given survival through offset $j-1$; accordingly, $S_t(k)$ is the survival probability to the start of offset $k$.}

Then the survival hazard is
\(h_t(j):=\sigma\!\Bigl(\beta_0+\beta_{\mathrm{stag}}u_{\mathrm{stag}}(t)+\beta_{\mathrm{front}}u_{\mathrm{front}}(t)+\eta_h(j-1)\Bigr)\).

\[
S_t(k):=\prod_{j=1}^{k-1}\bigl(1-h_t(j)\bigr),
\qquad
Q_t(k):=\Phi\!\left(
\frac{m_t-\gamma_t(k-1)-\operatorname{asinh}\!\left(\frac{\tau_{\mathrm{use},t}}{\tau_\Delta}\right)}{s_t}
\right).
\] \[
\widehat N_{\mathrm{rem}}(t)
:=
\sum_{k=1}^{D_{\mathrm{left}}(t)}
S_t(k)\,Q_t(k).
\] For pessimistic/optimistic runway envelopes, set
\(h_t^{\mathrm{pess/opt}}(j):=\sigma\!\bigl(\operatorname{logit}h_t(j)\pm\Delta_h\bigr)\),
\(m_t^{\mathrm{pess/opt}}:=m_t\mp\Delta_m\),
\(s_t^{\mathrm{pess/opt}}:=s_t+\Delta_s\) (pess) or
\(s_t^{\mathrm{pess/opt}}:=\max\!\left(s_{\min},s_t-\Delta_s\right)\)
(opt),
\(S_t^{\mathrm{pess/opt}}(k):=\prod_{j=1}^{k-1}\bigl(1-h_t^{\mathrm{pess/opt}}(j)\bigr)\),
and
\(Q_t^{\mathrm{pess/opt}}(k):=\Phi\!\left( \frac{m_t^{\mathrm{pess/opt}}-\gamma_t(k-1)-\operatorname{asinh}\!\left(\frac{\tau_{\mathrm{use},t}}{\tau_\Delta}\right)}{s_t^{\mathrm{pess/opt}}} \right)\).
Then
\(\widehat N_{\mathrm{rem}}^{\mathrm{low/high}}(t):=\sum_{k=1}^{D_{\mathrm{left}}(t)} S_t^{\mathrm{pess/opt}}(k)\,Q_t^{\mathrm{pess/opt}}(k)\),
with
\((\mathrm{low},\mathrm{high})\leftrightarrow(\mathrm{pess},\mathrm{opt})\),
and the confidence ratio is
\(C_t:=\frac{\widehat N_{\mathrm{rem}}^{\mathrm{low}}(t)}{\widehat N_{\mathrm{rem}}^{\mathrm{high}}(t)+\varepsilon}\).

For the controller coordinate and directional transforms, define \[
R_t:=\frac{\widehat N_{\mathrm{rem}}(t)}{D_{\mathrm{left}}(t)+\varepsilon}\in[0,1],
\qquad
e_t:=R_t^{a},
\qquad
g_t:=(1-R_t)^{b},
\qquad
a,b>0.
\] Here \(R_t\) is remaining useful runway as a fraction of the
depth-budget permitted by hardware/run-time constraints, and \(g_t\) is
the late-stage activation coordinate that scales controller variables
depending on age, where age is a normalized quantity relative
to/determined by expected utility of more generators.

Then \(R_t\approx 1\) early and \(R_t\approx 0\) late. Use \(e_t\) for
early-dominant quantities and \(g_t\) for late-dominant quantities. In
particular, define \[
\tau_{\mathrm{pos}}(t):=\tau_{\mathrm{pos}}^{\min}+(\tau_{\mathrm{pos}}^{\max}-\tau_{\mathrm{pos}}^{\min})e_t,
\qquad
\lambda_F(t):=\lambda_F^{\min}+(\lambda_F^{\max}-\lambda_F^{\min})g_t,
\qquad
u_{\mathrm{sat}}(t):=1-R_t.
\] The quantity \(\widehat N_{\mathrm{rem}}(t)\) is controller telemetry
rather than a physical observable: it summarizes recent selector
progress and remaining useful runway, but it is not itself a hard
admissibility
certificate.\footnote{Interpretation: the model has three layers --- (i) gain process ($m_t,s_t$), (ii) horizon-decay process ($\gamma_t$), and (iii) survival process ($h_t(j)$); together these drive the runway forecast $\widehat N_{\mathrm{rem}}(t)$.}

For inherited refit windows, use the maturity-shrinking retention law \[
\rho_W(t):=\rho_W^{\min}+(\rho_W^{\max}-\rho_W^{\min})e_t,
\qquad
0\le \rho_W^{\min}\le \rho_W^{\max}\le 1,
\] so that the retained old-coordinate fraction is largest early and
contracts as useful runway is exhausted.

For measurement effort, use the Phase-1 baseline schedule \[
N_{\mathrm{shot},1}(t):=
\left\lceil
N_{\mathrm{shot}}^{\min}
+
\bigl(N_{\mathrm{shot}}^{\max}-N_{\mathrm{shot}}^{\min}\bigr)
\Bigl(
\alpha_{\mathrm{front}}\,u_{\mathrm{front}}(t)
+
\alpha_{\mathrm{util}}\,(1-R_t)
\Bigr)
\right\rceil,
\] with \(N_{\mathrm{shot}}^{\min}\le N_{\mathrm{shot}}^{\max}\),
\(\alpha_{\mathrm{front}},\alpha_{\mathrm{util}}\ge 0\), and
\(\alpha_{\mathrm{front}}+\alpha_{\mathrm{util}}=1\). Thus Phase 1
increases shot count when the frontier flattens and when useful runway
is being exhausted.

For \(k\in\{2,3\}\), once the phase-\(k\) uncertainty summary
\(u_{\sigma,k}(t)\) is available, use \[
N_{\mathrm{shot},k}(t):=
\left\lceil
N_{\mathrm{shot}}^{\min}
+
\bigl(N_{\mathrm{shot}}^{\max}-N_{\mathrm{shot}}^{\min}\bigr)
\Bigl(
\alpha_{\mathrm{front}}\,u_{\mathrm{front}}(t)
+
\alpha_{\mathrm{util}}\,(1-R_t)
+
\alpha_{\sigma}\,u_{\sigma,k}(t)
\Bigr)
\right\rceil,
\qquad
k\in\{2,3\},
\] with
\(\alpha_{\mathrm{front}},\alpha_{\mathrm{util}},\alpha_{\sigma}\ge 0\)
and
\(\alpha_{\mathrm{front}}+\alpha_{\mathrm{util}}+\alpha_{\sigma}=1\).\footnote{A one-weight simplification is the special case $\alpha_{\mathrm{front}}=\alpha_{\mathrm{util}}=\frac{1-\alpha_{\sigma}}{2}$.}
Thus Phase 2 and Phase 3 inherit the same controller baseline and then
increase shot count further when geometric uncertainty rises.

For phase thresholds, use the live monotone law \[
\tau_k(t):=\tau_k^{\min}+(\tau_k^{\max}-\tau_k^{\min})e_t,
\qquad
\tau_k^{\min}\le \tau_k^{\max},
\qquad
k\in\{1,2,3\},
\] so that shortlist admission is stricter early and more permissive
later.

For phase shortlist caps, use the live bounded law \[
N_k(t):=
\left\lceil
N_k^{\min}+(\,N_k^{\max}-N_k^{\min}\,)g_t
\right\rceil,
\qquad
N_k^{\min}\le N_k^{\max},
\qquad
k\in\{1,2,3\}.
\] Thus \(N_k(t)\) and \(\tau_k(t)\) are both live shortlist-control
quantities: \(N_k(t)\) ensures budget constraints are met, whereas
\(\tau_k(t)\) helps ensure logical admissibility constraints proceed
with time.

\hypertarget{phase-local-cost-and-structural-burden-definitions}{%
\subsubsection{11.2.4 Phase-local cost and structural burden
definitions}\label{phase-local-cost-and-structural-burden-definitions}}

For any current phase-input record family \(\mathcal R_k(t)\), with
\(\mathcal R_1(t):=\mathcal R_1\) and \(k\in\{1,2,3\}\), and any scalar
record feature \(x:\mathcal R_k(t)\to\mathbb R_{\ge 0}\), define the
empirical quantiles
\(q_\alpha^x(t):=\operatorname{Quantile}_\alpha\bigl(\{x(u):u\in\mathcal R_k(t)\}\bigr)\)
and \(m_x(t):=q_{0.5}^x(t)\), together with the robust scales
\(s_x^{\mathrm{MAD}}(t):=1.4826\,\operatorname{median}_{u\in\mathcal R_k(t)}\bigl|x(u)-m_x(t)\bigr|\),
\(s_x^{\mathrm{pct}}(t):=\frac{q_{0.9}^x(t)-q_{0.1}^x(t)}{2.563}\), and
\(s_x(t):=\max\!\bigl(\varepsilon_x,\;s_x^{\mathrm{MAD}}(t),\;s_x^{\mathrm{pct}}(t)\bigr)\).
Thus \(s_x(t)\) is a percentile/MAD hybrid reference scale, with
percentile fallback when the MAD degenerates on sparse discrete
features.

Define the one-sided robust normalized excess by
\(z_x(r;t):=\operatorname{asinh}\!\left(\frac{[x(r)-m_x(t)]_+}{s_x(t)}\right)\),
with \([y]_+:=\max(y,0)\). When a direct ratio form is preferred, use
the reference level \(x_{\mathrm{ref}}(t):=m_x(t)+s_x(t)\).

For a candidate-position record \(r=(m,p)\), let
\(\mathcal L_r:=\operatorname{supp}(m)=\{j:m_j\neq I\}\),
\(\mathrm{wt}(r):=|\mathcal L_r|\), and
\(\#_{X,Y}(r):=\left|\{j\in\mathcal L_r:m_j\in\{X,Y\}\}\right|\).

Define feature-specific robust excesses by
\(z_{\mathrm{wt}}(r;t):=z_x(r;t)\ \text{with }x=\mathrm{wt}\) and
\(z_{XY}(r;t):=z_x(r;t)\ \text{with }x=\#_{X,Y}\). The structural burden
term is \[
D_{\mathrm{struct}}(r;t):=
\beta_{\mathrm{wt}}\,z_{\mathrm{wt}}(r;t)
+\beta_{XY}\,z_{XY}(r;t),
\] with nonnegative coefficients
\(\beta_{\mathrm{wt}},\beta_{XY}\ge 0\).\footnote{Provenance note: generator novelty $G_{\mathrm{new}}(r)$, family novelty $\mathrm{groups}_{\mathrm{new}}(r)$ with single-valued $\mathfrak g:\mathcal G\to\mathcal C$, lifetime/staleness $\mathbf 1_{\mathrm{lifetime}}(r;t)$ with $t_{\mathrm{first}}(m)$ meaning first time ever seen in the run, the auxiliary raw burden $D_{\mathrm{raw}}(r;t):=D_{\mathrm{struct}}(r;t)+d_{\mathrm{pos}}(r)$ with $d_{\mathrm{pos}}(r):=|p-p_{\mathrm{app}}|$, and the normalized descendants $\bar D_{\mathrm{life}}^{\mathrm{struct}},\bar G_{\mathrm{new}},\bar C_{\mathrm{new}},\bar c,\bar P_{\mathrm{opt}}$ are omitted from the live hardware cost in the present draft and retained only as controller-side provenance/policy bookkeeping.}

For candidate-specific measurement burden, let
\(O_r:=i[H,\widetilde A_r]=\sum_{b\in\mathcal B_r} O_{r,b}\), with
\(O_{r,b}:=\sum_{\ell\in B_{r,b}} c_{r,b\ell} P_{r,b\ell}\), where
\(\mathcal B_r\) is the qubit-wise commuting partition of the
candidate's gradient observable and
\(\mathcal B_r^{\mathrm{new}}\subseteq\mathcal B_r\) is the subset not
already covered by measurement reuse. Here \(\widetilde A_r\) is the
dressed candidate generator at record \(r=(m,p)\), \(\mathcal B_r\) is
the partition of \(O_r=i[H,\widetilde A_r]\) into qubit-wise commuting
Pauli-string measurement groups, and \(\mathcal B_r^{\mathrm{new}}\) is
the subset that still requires new measurement effort. When exact
state-dependent variances are unavailable, use the commutator-based
proxy \[
\widehat v_{r,b}:=\sum_{\ell\in B_{r,b}} |c_{r,b\ell}|^2,
\qquad
\widehat{\bar C}_{\mathrm{shot}}(r):=
\frac{1}{\mathrm{shot}_{\mathrm{ref}}}
\cdot
\frac{
\left(
\sum_{b\in\mathcal B_r^{\mathrm{new}}}\sqrt{\widehat v_{r,b}}
\right)^2
}{
\sigma_{\star}^{2}
}.
\] Here \(\sigma_\star\) is the target standard error for the candidate
gradient estimate. The positive constant
\(\mathrm{shot}_{\mathrm{ref}}\) is a baseline measurement budget
against which candidate-specific shot cost is compared.

For the shared phase-\(k\) cost, define
\(z_{\mathrm{2q}}(r;t):=z_x(r;t)\ \text{with }x=C_{\mathrm{2q}}\),
\(z_d(r;t):=z_x(r;t)\ \text{with }x=D\),
\(z_{\theta}(r;t):=z_x(r;t)\ \text{with }x=\Delta K\),
\(z_{\mathrm{shot}}(r;t):=z_x(r;t)\ \text{with }x=\widehat{\bar C}_{\mathrm{shot}}\),
and \[
K_k(r;t):=
\lambda_{\mathrm{2q}}\,z_{\mathrm{2q}}(r;t)
+\lambda_d\,z_d(r;t)
+\lambda_\theta\,z_{\theta}(r;t)
+\lambda_{\mathrm{shot}}\,z_{\mathrm{shot}}(r;t)
+D_{\mathrm{struct}}(r;t).
\] Explicitly, \[
\begin{aligned}
K_k(r;t)
&=
\lambda_{\mathrm{2q}}\,z_{\mathrm{2q}}(r;t)
+\lambda_d\,z_d(r;t)
+\lambda_\theta\,z_{\theta}(r;t)
+\lambda_{\mathrm{shot}}\,z_{\mathrm{shot}}(r;t)\\
&\qquad
+\beta_{\mathrm{wt}}\,z_{\mathrm{wt}}(r;t)
+\beta_{XY}\,z_{XY}(r;t),
\qquad
k\in\{1,2,3\}.
\end{aligned}
\] \#\#\# 11.2.5 Final Phase I Scoring For Phase 1, the set
\(\mathcal R_1:=\{(m,p):m\in\mathcal G_1,\ p\in\mathcal P_m^{1}\}\) is
scored by \[
S_1(r;t)
:=
\frac{\Delta E_{1,\mathrm{TR}}(r)}{1+K_1(r;t)},
\qquad r\in\mathcal R_1.
\] Here \(\Delta E_{1,\mathrm{TR}}(r)\) is the Phase-1 one-coordinate
trust-region proxy defined in \S 11.2.1, and the phase-1 cost is
evaluated on the current Phase-1 input family. This Phase-1 score is
intentionally broad as a screening surface: it privileges cheap local
second-order insertion estimates while the controller still has long
useful runway. The corresponding broad Phase-1 probing/evaluation budget
is set by the baseline controller shot law \(N_{\mathrm{shot},1}(t)\).

In the present draft, the live Phase-1 numerator is the one-coordinate
trust-region proxy \(\Delta E_{1,\mathrm{TR}}(r)\), so that the broad
screen remains in the same gradient-plus-second-order family later
refined in Phases 2 and 3. The strict one-coordinate insertion drop
\(\delta E_1(r\mid\mathcal O,\theta)\) is retained only as an auxiliary
exact comparator for selective frontier confirmation, ablation studies,
or future confirm-stage refinement. A broader early-stage local-refit
numerator remains only a non-live hook. To make such a refinement live,
one would first have to specify (i) the admissible refit-window family,
(ii) the corresponding constrained local-refit drop, and (iii) a
controller activation law independent of the current raw Phase-1 score
multiset, so that the live Phase-1 law is not self-referential.

\hypertarget{position-probing-trough-detection-and-effective-selector}{%
\subsubsection{11.2.6 Position probing, trough detection, and effective
selector}\label{position-probing-trough-detection-and-effective-selector}}

The Phase-1 scoring mechanism produces the retained record shortlist for
further refined scoring. The retained model set is
\(\mathcal G^{(1)}:=\pi_m(\mathcal S_1(t))\), the retained position
fibers are
\(\mathcal P_m^{(1)}:=\{p:(m,p)\in\mathcal S_1(t)\}\ \forall m\in\mathcal G^{(1)}\),
the Phase-2 record pool is \(\mathcal R_2:=\mathcal S_1(t)\), and the
acceptance set is
\(\mathcal A_1(t):=\{r=(m,p)\in\mathcal R_1:S_1(r;t)\ge \tau_1(t)\}=\{r_{(1)},\dots,r_{(M)}\}\),
with \(M:=|\mathcal A_1(t)|\). The ordered Phase-1 scores are
\(s_k^{(1)}(t):=S_1(r_{(k)};t)\), with
\(s^{(1)}_1(t)\ge s^{(1)}_2(t)\ge\cdots\ge s^{(1)}_M(t)\), and the
candidate score distinction is
\(u_{\mathrm{front},k}^{(1)}(t):=\frac{s^{(1)}_{k+1}(t)+\varepsilon}{s^{(1)}_{k}(t)+\varepsilon}\),
with \(k=1,\dots,M-1\) and \(0<\eta_1<1\). \[
\begin{aligned}
k_1^\star(t)&:=
\begin{cases}
0, & M=0,\\[4pt]
\min\!\left\{k\in\{1,\dots,\min(M-1,N_1(t)-1)\}:u_{\mathrm{front},k}^{(1)}(t)\le \eta_1\right\},
& \text{if this set is nonempty},\\[6pt]
\min\{M,N_1(t)\}, & \text{otherwise},
\end{cases}\\
\mathcal D_1(t)&:=\{r_{(j)}(t):1\le j\le k_1^\star(t)\},
\qquad
\mathcal S_1(t):=\mathcal A_1(t)\cap \mathcal D_1(t)=\mathcal D_1(t).
\end{aligned}
\] Here \(\eta_1\) is a controller parameter chosen either from policy
intuition or from empirical calibration.

\hypertarget{phase-ii}{%
\subsection{11.3 Phase II}\label{phase-ii}}

Phase II is the raw geometric selector stage. It acts on the Phase-1
survivors, restores the coarse geometric signal that distinguishes
directional usefulness beyond insertion-drop alone, and produces the
filtered shortlist \(\mathcal S_2\).

\hypertarget{raw-geometry-restored-novelty-and-auxiliary-overlap-window}{%
\subsubsection{11.3.1 Raw geometry, restored novelty, and auxiliary
overlap
window}\label{raw-geometry-restored-novelty-and-auxiliary-overlap-window}}

The Phase-2 selector works over candidate-position records inherited
from Phase 1. In the record-inheritance mode used here, the inherited
and admissible domain is \(\mathcal R_2:=\mathcal S_1(t)\). For
\(r=(m,p)\), let \(W(r;t)\subseteq\{1,\dots,n+1\}\) denote the inherited
overlap window attached to \(r\). In the present three-phase
architecture this window enters the Phase-2 selector authoritatively
only through overlap against the current horizontal tangent family; the
same inherited window is then carried forward to Phase 3, where
reduced-window elimination and final reranking are
performed.\footnote{Here $p$ indexes the cut after $U_p$. Thus $U_{>p}:=U_K\cdots U_{p+1}$ and $U_{\le p}:=U_p\cdots U_1$, and insertion at cut $p$ means $ |\psi_r(\alpha)\rangle = U_{>p}e^{-i\alpha A_m}U_{\le p}|\psi_0\rangle = U_K\cdots U_{p+1}\,e^{-i\alpha A_m}\,U_p\cdots U_1|\psi_0\rangle $. So the new generator is inserted between $U_p$ and $U_{p+1}$.}

Reusing the record-local insertion objects already introduced in Phase
1, let \(|d_r^{(2)}\rangle:=-\widetilde A_r^2|\psi\rangle\) and
\(\chi_r:=\operatorname{clip}\!\left(1-\frac{z_\alpha\,\sigma_r}{|g_r|+\varepsilon_g},0,1\right)\).\footnote{Let $\widehat g(r)$ denote the empirical gradient estimator for record $r=(m,p)$ obtained from $N_r$ independent shot-level repetitions. If $X^{(r)}_1,\dots,X^{(r)}_{N_r}$ are the repetition-level gradient samples, then $\widehat g(r):=\frac{1}{N_r}\sum_{i=1}^{N_r} X^{(r)}_i$ and $\sigma_r:=\operatorname{SE}\!\big[\widehat g(r)\big]=\sqrt{\operatorname{Var}\!\big(\widehat g(r)\big)}$. Under independent repetitions, $\operatorname{Var}\!\big(\widehat g(r)\big)=\frac{\operatorname{Var}\!\big(X^{(r)}_1\big)}{N_r}$, so in practice one estimates $\sigma_r\approx \frac{s_r}{\sqrt{N_r}}$, where $s_r$ is the sample standard deviation of the repetition-level gradients. Thus $\sigma_r$ is the standard error of the mean gradient estimate, not the sample standard deviation itself.}
For phase-dependent shot allocation in later phases, define \[
u_{\sigma,k}(t):=
\operatorname{clip}\!\left(
\operatorname{median}_{r\in\mathcal R_k(t)}
\frac{z_\alpha\,\sigma_r}{|g_r|+\varepsilon_g},
\,0,\,1
\right),
\qquad
k\in\{2,3\},
\] where in Phase 2 the record-local gradient and uncertainty estimates
are evaluated using \(N_{\mathrm{shot},2}(t)\), and in Phase 3 the pair
\((g_r,\sigma_r)\) is inherited from the Phase-II raw geometry carried
by the surviving records.

The raw geometric quantities are \[
\begin{aligned}
F_{\mathrm{raw}}(r)&:=\langle \tau_r,\tau_r\rangle+\varepsilon_F,
\qquad \varepsilon_F>0,\\
\mathcal N_2(r;t)
&:=1-\max_{j\in W(r;t)}
\frac{\left|\langle d_r|Q_\psi|\partial_{\theta_j}\psi\rangle\right|}
{\sqrt{F_{\mathrm{raw}}(r)}}
\frac{1}{\sqrt{\langle \partial_{\theta_j}\psi|Q_\psi|\partial_{\theta_j}\psi\rangle+\varepsilon_F}},\\
h_r&:=2\Re\!\left(\langle d_r|H|d_r\rangle+\langle \psi|H|d_r^{(2)}\rangle\right),\\
\Delta E_{\mathrm{TR}}^{\mathrm{raw}}(r)
&:=\max_{|\alpha|\le \rho/\sqrt{F_{\mathrm{raw}}(r)}}
\left[g_{\mathrm{lcb}}(r)|\alpha|-\frac12\max\bigl(h_r,0\bigr)\alpha^2\right].
\end{aligned}
\] Here \(F_{\mathrm{raw}}(r)\) is the raw tangent norm,
\(\mathcal N_2(r;t)\) is the restored Phase-2 novelty score, \(h_r\) is
the raw local curvature, and
\(\Delta E_{\mathrm{TR}}^{\mathrm{raw}}(r)\) is the raw trust-region
gain before reduced-window elimination.

The authoritative Phase-II score is \[
S_2(r;t)
:=
\frac{\Delta E_{\mathrm{TR}}^{\mathrm{raw}}(r)\,\chi_r\,\mathcal N_2(r;t)}{1+K_2(r;t)},
\qquad r\in\mathcal R_2.
\] Thus Phase II remains the raw geometric selector stage: it restores
directional discrimination beyond insertion-drop alone, damps that
geometry when the candidate direction is poorly resolved against its own
gradient noise, and filters the Phase-1 survivors before any
reduced-window refit quantity is made authoritative.

For continuity with the older local-refit notation, retain the auxiliary
windowed drop \[
\mathcal C_W(r;\theta):=\{\widetilde\theta\in\Theta_{\mathcal O\oplus_p m}:\widetilde\theta_j=(\theta\oplus_p 0)_j\ \forall j\notin W(r;t)\}
\] and \[
\delta E_2(r\mid\mathcal O,\theta):=
E(\mathcal O,\theta)-\inf_{\widetilde\theta\in\mathcal C_W(r;\theta)}E(\mathcal O\oplus_p m,\widetilde\theta).
\] This \(\delta E_2\) quantity is retained only as an auxiliary
comparator for the inherited window. It may be used for notation
continuity, diagnostics, or ablation studies, but it does not define a
separate selector phase, it does not enter \(S_2(r;t)\), and it does not
modify the Phase-II shortlist rule. Any authoritative reduced-window
rerank remains deferred to Phase 3.

\hypertarget{shortlist-rule}{%
\subsubsection{11.3.2 Shortlist rule}\label{shortlist-rule}}

The Phase-II scoring mechanism produces the retained geometric shortlist
for the later reduced-window stage. The acceptance set is
\(\mathcal A_2(t):=\{r=(m,p)\in\mathcal R_2:S_2(r;t)\ge \tau_2(t)\}=\{r_{(1)},\dots,r_{(M)}\}\),
with \(M:=|\mathcal A_2(t)|\). Here \(S_2\) is the raw geometric gain
times restored novelty score. Let \(s_k^{(2)}(t):=S_2(r_{(k)};t)\), such
that \(s^{(2)}_1(t)\ge s^{(2)}_2(t)\ge\cdots\ge s^{(2)}_M(t)\), with
superscript denoting the phase scoring regime. Rate the candidate's
score distinction by
\(u_{\mathrm{front},k}^{(2)}(t):=\frac{s^{(2)}_{k+1}(t)+\varepsilon}{s^{(2)}_{k}(t)+\varepsilon}\),
with \(k=1,\dots,M-1\) and \(0<\eta_2<1\). Then \[
\begin{aligned}
k_2^\star(t)&:=
\begin{cases}
0, & M=0,\\[4pt]
\min\!\left\{k\in\{1,\dots,\min(M-1,N_2(t)-1)\}:u_{\mathrm{front},k}^{(2)}(t)\le \eta_2\right\},
& \text{if this set is nonempty},\\[6pt]
\min\{M,N_2(t)\}, & \text{otherwise},
\end{cases}\\
\mathcal D_2(t)&:=\{r_{(j)}:1\le j\le k_2^\star(t)\},
\qquad
\mathcal S_2(t):=\mathcal A_2(t)\cap \mathcal D_2(t)=\mathcal D_2(t).
\end{aligned}
\] Of course, \(\mathcal G^{(2)}:=\pi_m(\mathcal S_2(t))\), and
\(\mathcal P_m^{(2)}:=\{p:(m,p)\in\mathcal S_2(t)\}\).

\hypertarget{phase-3-reduced-window-rerank-and-post-shortlist-admission}{%
\subsection{11.4 Phase 3 reduced-window rerank and post-shortlist
admission}\label{phase-3-reduced-window-rerank-and-post-shortlist-admission}}

This section refines the surviving Phase-2 records with the more
expensive reduced-window geometry, then uses that final rerank surface
for shortlist formation, batched admission, pruning, and continuation.

\hypertarget{reduced-window-geometry-curvature-and-trust-region-gain}{%
\subsubsection{11.4.1 Reduced-window geometry, curvature, and
trust-region
gain}\label{reduced-window-geometry-curvature-and-trust-region-gain}}

Phase 3 starts from the Phase-II raw quantities \(\tau_r\), \(g_r\),
\(h_r\), \(g_{\mathrm{lcb}}(r)\), \(F_{\mathrm{raw}}(r)\), and
\(\mathcal N_2(r;t)\), then refines them using the inherited refit
window. The coarse novelty score \(\mathcal N_2(r;t)\) is not recomputed
here; instead the candidate direction is reduced against the inherited
window before the expensive rerank surface is formed.

\hypertarget{refit-window-reduced-geometry-and-phase-3-record-family}{%
\subsubsection{11.4.2 Refit window, reduced geometry, and Phase-3 record
family}\label{refit-window-reduced-geometry-and-phase-3-record-family}}

In the record-inheritance mode used here, the Phase-3 input set is
\(\mathcal R_3(t):=\mathcal S_2(t)\). The inherited refit window
attached to position \(p\) at selector step \(t\) is denoted \(W(p;t)\),
with decomposition
\(W(p;t)=W_{\mathrm{new}}^{(p)}\cup W_{\mathrm{age}}^{(p)}\), where
\(W_{\mathrm{new}}^{(p)}\) keeps the newest coordinates and
\(W_{\mathrm{age}}^{(p)}\) preserves the maturity-scaled set of oldest
retained old coordinates defined later by the branch-local admission
order. For each shortlisted record \(r=(m,p)\), let \(W_r:=W(r;t)\) be
the inherited window that would actually be refit if \(m\) were admitted
at position \(p\), and let \(\{\tau_j\}_{j\in W_r}\) be the current
horizontal tangents in that inherited window. Phase 3 evaluates its
inherited reduced-window geometry under the later-phase controller shot
law \(N_{\mathrm{shot},3}(t)\).

Using the Phase-II raw tangent data, define
\((Q_r)_{jk}=\Re\langle \tau_j,\tau_k\rangle\),
\((q_r)_j=\Re\langle \tau_j,\tau_r\rangle\),
\(E_0:=E(\theta;\mathcal O)\),
\((b_r)_j:=\left.\partial_\alpha\partial_{\theta_j}E(\alpha,\delta\theta_{W_r};r)\right|_{\alpha=0,\,\delta\theta_{W_r}=0}\)
for \(j\in W_r\), and
\((H_{W_r})_{jk}:=\left.\partial_{\theta_j}\partial_{\theta_k}E(\alpha,\delta\theta_{W_r};r)\right|_{\alpha=0,\,\delta\theta_{W_r}=0}\)
for \(j,k\in W_r\). The local reduced-window energy model is \[
E(\alpha,\delta\theta_{W_r};r)
\approx
E_0+g_r\alpha+\tfrac12 h_r\alpha^2+\alpha b_r^\top\delta\theta_{W_r}+\tfrac12\delta\theta_{W_r}^\top H_{W_r}\delta\theta_{W_r}.
\] For reduction and reduced novelty,
\(\mathbf R_r=\tfrac12(H_{W_r}+H_{W_r}^{\top})+\lambda_H I\),
\(\widetilde h_r=h_r-b_r^\top \mathbf R_r^{-1}b_r\),
\(F_r^{\mathrm{red}}=F_{\mathrm{raw}}(r)-2q_r^\top \mathbf R_r^{-1}b_r+b_r^\top \mathbf R_r^{-1}Q_r\mathbf R_r^{-1}b_r\),
\(F_{r,\mathrm{safe}}^{\mathrm{red}}=\max\!\left(F_r^{\mathrm{red}},\varepsilon_{\mathrm{red}}\right)\)
with \(\varepsilon_{\mathrm{red}}>0\),
\(q_r^{\mathrm{red}}=q_r-Q_r\mathbf R_r^{-1}b_r\), and
\(z_r:=\frac{(q_r^{\mathrm{red}})^\top (Q_r+\varepsilon_{\mathrm{nov}}I)^{-1} q_r^{\mathrm{red}}}{F_{r,\mathrm{safe}}^{\mathrm{red}}}\),
\(\mathcal N_3(r;t):=\operatorname{clip}_{[0,1]}\!\left(1-z_r\right)\).
The reduced trust-region gain is \[
\Delta E_{\mathrm{TR}}(r)
:=
\max_{|\alpha|\le \rho/\sqrt{F_{r,\mathrm{safe}}^{\mathrm{red}}}}
\left[g_{\mathrm{lcb}}(r)|\alpha|-\tfrac12\max(\widetilde h_r,0)\alpha^2\right].
\] \#\#\# 11.4.3 Phase-3 rerank score and effective selector The
authoritative Phase-3 rerank score is \[
S_3(r;t)
:=
\frac{\Delta E_{\mathrm{TR}}(r)\,\chi_r\,\mathcal N_3(r;t)}{1+K_3(r;t)}
+\beta_{\mathrm{motif}}S_{\mathrm{motif}}(r)-\beta_{\mathrm{dup}}D_{\mathrm{dup}}(r),
\qquad r\in\mathcal R_3(t).
\] Here \(S_{\mathrm{motif}}\) measures motif compatibility or
transferable local pattern alignment, and \(D_{\mathrm{dup}}\) penalizes
near-duplicate continuation directions. These are controller-side
feature functionals normalized to comparable scale and layered on top of
the explicit reduced geometric product
\(\Delta E_{\mathrm{TR}}(r)\,\chi_r\,\mathcal N_3(r;t)\); they are not
part of the phase-invariant cost denominator.

The retained Phase-3 shortlist is obtained from
\(\mathcal A_3(t):=\{r=(m,p)\in\mathcal R_3(t):S_3(r;t)\ge \tau_3(t)\}=\{r_{(1)},\dots,r_{(M)}\}\),
with \(M:=|\mathcal A_3(t)|\). Let \(s_k^{(3)}(t):=S_3(r_{(k)};t)\),
such that \(s^{(3)}_1(t)\ge s^{(3)}_2(t)\ge\cdots\ge s^{(3)}_M(t)\).
Rate the candidate's score distinction by
\(u_{\mathrm{front},k}^{(3)}(t):=\frac{s^{(3)}_{k+1}(t)+\varepsilon}{s^{(3)}_{k}(t)+\varepsilon}\),
with \(k=1,\dots,M-1\) and \(0<\eta_3<1\). Then \[
\begin{aligned}
k_3^\star(t)&:=
\begin{cases}
0, & M=0,\\[4pt]
\min\!\left\{k\in\{1,\dots,\min(M-1,N_3(t)-1)\}:u_{\mathrm{front},k}^{(3)}(t)\le \eta_3\right\},
& \text{if this set is nonempty},\\[6pt]
\min\{M,N_3(t)\}, & \text{otherwise},
\end{cases}\\
\mathcal D_3(t)&:=\{r_{(j)}:1\le j\le k_3^\star(t)\},
\qquad
\mathcal S_3(t):=\mathcal A_3(t)\cap \mathcal D_3(t)=\mathcal D_3(t).
\end{aligned}
\] Of course, \(\mathcal G^{(3)}:=\pi_m(\mathcal S_3(t))\), and
\(\mathcal P_m^{(3)}:=\{p:(m,p)\in\mathcal S_3(t)\}\).

\hypertarget{reduced-plane-geometric-batch-selection-after-shortlisting}{%
\subsubsection{11.4.4 Reduced-plane geometric batch selection after
shortlisting}\label{reduced-plane-geometric-batch-selection-after-shortlisting}}

After the Phase-3 shortlisting has produced the record family
\(\mathcal S_3(t)\subseteq \mathcal R_3(t)\), batching is resolved by
reduced-state-space geometry on a common inherited tangent plane, rather
than by heuristic pairwise compatibility penalties. Support-overlap,
commutation, and related proxy features may still be used as optional
prescreens, but they are not part of the principled batching law below.
Assume throughout that \(\mathcal S_3(t)\neq\varnothing\); otherwise no
Phase-3 batch is formed.

Near-degenerate shell and structural feasibility: let
\(s_t(r):=S_3(r;t)\) and
\(r_{\max}(t)\in \arg\max_{r\in \mathcal S_3(t)} s_t(r)\). Fix a
near-degenerate shell parameter \(\eta_{\mathrm{nd}}\in(0,1]\) and a
batch cap \(B_{\max}(t)\in\mathbb N\). Define the near-degenerate record
shell
\(\mathcal N_3(t):=\left\{ r\in \mathcal S_3(t): s_t(r)\ge \eta_{\mathrm{nd}}\,s_t\!\bigl(r_{\max}(t)\bigr) \right\}\)
and the corresponding near-degenerate batch family
\(\mathfrak B_{\mathrm{nd}}(t):=\left\{ B\subseteq \mathcal N_3(t): r_{\max}(t)\in B,\quad 1\le |B|\le B_{\max}(t) \right\}\).
Define a hard batch-structural gate
\(\Gamma_{\mathrm{batch}}^{\mathrm{struct}}(B;t)\in\{0,1\}\) to mean
that the records in \(B\) are jointly insertable on the current
branch/state, satisfy all hard cross-record constraints, and determine a
well-defined batched scaffold update \(\mathcal O_t\oplus B\) in the
fixed canonical insertion order. The structurally feasible batch family
is
\(\mathfrak B_{\mathrm{struct}}(t):=\left\{ B\subseteq \mathcal S_3(t): \Gamma_{\mathrm{batch}}^{\mathrm{struct}}(B;t)=1 \right\}\).
A controller-side transient beam fallback may also be defined on the
same rerank
shell.\footnote{One convenient rule is to set $\mathcal T_{\mathrm{tie}}(t):=\{\,r\in\mathcal S_3(t): s_t(r)\ge \max(\eta_{\mathrm{tie}}\,s_t(r_{\max}(t)),\,s_t(r_{\max}(t))-\delta_{\mathrm{tie}})\,\}$. If $|\mathcal T_{\mathrm{tie}}(t)|>1$ and the controller elects not to resolve the shell by an immediate one-scaffold batch, it may retain a highest-scoring subset $\mathcal A_{\mathrm{tie}}(t)\subseteq \mathcal T_{\mathrm{tie}}(t)$ with $1\le |\mathcal A_{\mathrm{tie}}(t)|\le B_{\mathrm{tie}}(t)$ and pass one retained alternative per record into the scaffold-beam machinery of \S 11.5.2.}

Common-window reduced geometry: write the current state as
\(|\psi_t\rangle:=|\psi(\mathcal O_t,\theta_t)\rangle, \qquad Q_{\psi_t}:=I-|\psi_t\rangle\langle \psi_t|\),
where \(Q_{\psi_t}\) is the horizontal projector. For any batch
\(B\in \mathfrak B_{\mathrm{struct}}(t)\), define the common inherited
refit window \(W_B(t):=\bigcup_{r\in B} W_r(t)\) and the corresponding
inherited horizontal tangent matrix
\(T_B(t):=[\,t_j(t)\,]_{j\in W_B(t)}, \qquad t_j(t):=Q_{\psi_t}\,\partial_{\theta_j}|\psi_t\rangle\).
Let \(E_t(\alpha,\delta\theta_{W_B};B)\) denote the local energy
obtained by inserting the records in \(B\) with batch amplitudes
\(\alpha=(\alpha_r)_{r\in B}\) and varying the inherited coordinates on
\(W_B(t)\) by \(\delta\theta_{W_B}\). For each \(r\in B\), let
\(|\psi_{t,r}(\alpha_r)\rangle\) denote the state obtained by inserting
only record \(r\) at amplitude \(\alpha_r\) into the current scaffold
while keeping inherited coordinates fixed. Its zero-initialized
horizontal candidate tangent is
\(d_r(t):=Q_{\psi_t}\,\partial_{\alpha_r}|\psi_{t,r}(\alpha_r)\rangle\Big|_{\alpha_r=0}\).
Let \(\mathsf H_{W_B}(t)\) denote the inherited-window Hessian block of
\(E_t(\alpha,\delta\theta_{W_B};B)\) with respect to the coordinates on
\(W_B(t)\) at \((\alpha,\delta\theta_{W_B})=(0,0)\), and set
\(M_B(t):=\frac12\Bigl(\mathsf H_{W_B}(t)+\mathsf H_{W_B}(t)^\top\Bigr)+\lambda_H I, \qquad \lambda_H>0.\)
For each \(r\in B\), define the mixed candidate-window curvature vector
by
\(\bigl(b_r^{(B)}(t)\bigr)_j:=\left.\partial_{\alpha_r}\partial_{\theta_j} E_t(\alpha,\delta\theta_{W_B};B)\right|_{(\alpha,\delta\theta_{W_B})=(0,0)}, \qquad j\in W_B(t)\).
The reduced candidate tangent after linear relaxation of the common
inherited window is then
\(u_r^{(B)}(t):=d_r(t)-T_B(t)\,M_B(t)^{-1}b_r^{(B)}(t)\), hence the
effective batch plane is
\(\Pi_B(t):=\operatorname{span}\{u_r^{(B)}(t):r\in B\}\subseteq T_{\psi_t}\mathbb{P}(\mathcal H)\).
Its reduced Gram matrix is
\(\bigl(G_B(t)\bigr)_{rs}:=\Re\langle u_r^{(B)}(t),u_s^{(B)}(t)\rangle, \qquad r,s\in B\),
and in particular \(\bigl(G_B(t)\bigr)_{rr}=F_r^{\mathrm{red},(B)}(t)\)
is the reduced norm of record \(r\) in the common batch reduction. Next
define the symmetrized raw candidate-candidate curvature on the same
common reduction by
\(h_{rs}^{(B)}(t):=\frac12 \left. \bigl( \partial_{\alpha_r}\partial_{\alpha_s} + \partial_{\alpha_s}\partial_{\alpha_r} \bigr) E_t(\alpha,\delta\theta_{W_B};B) \right|_{(\alpha,\delta\theta_{W_B})=(0,0)}, \qquad r,s\in B,\)
and then the reduced curvature matrix
\(\bigl(\widetilde H_B(t)\bigr)_{rs}:=h_{rs}^{(B)}(t)-\bigl(b_r^{(B)}(t)\bigr)^\top M_B(t)^{-1} b_s^{(B)}(t)\).
For \(r=s\), the entry \(\bigl(\widetilde H_B(t)\bigr)_{rr}\) is the
usual single-record reduced curvature in the common batch window; for
\(r\neq s\), the entry \(\bigl(\widetilde H_B(t)\bigr)_{rs}\) is the
mixed second-order coupling after inherited-window relaxation. The
diagnostic coherences are
\(\mu_{\mathrm{tan}}(B;t):=\max_{r\neq s\in B} \frac{|(G_B(t))_{rs}|} {\sqrt{(G_B(t))_{rr}(G_B(t))_{ss}}+\varepsilon}\)
and
\(\mu_{\mathrm{curv}}(B;t):=\max_{r\neq s\in B} \frac{|(\widetilde H_B(t))_{rs}|} {\sqrt{\bigl((\widetilde H_B(t))_{rr}\bigr)^{(+)} \bigl((\widetilde H_B(t))_{ss}\bigr)^{(+)} }+\varepsilon}, \qquad x^{(+)}:=\max\{x,0\}\);
these may be used as optional early pruning diagnostics, but the actual
batch objective is the joint reduced trust-region gain below.

Joint reduced gain and additivity: let
\(g_{B,\mathrm{lcb}}(t):=\bigl(g_{\mathrm{lcb}}(r;t)\bigr)_{r\in B}\),
and let \(\bigl(\widetilde H_B(t)\bigr)^{(+)}\) denote the
positive-semidefinite part of \(\widetilde H_B(t)\). The batch
trust-region objective below is therefore a conservative convexified
local gain model on the reduced batch plane: \[
\Delta E_B^{\mathrm{TR}}(t)
:=
\max_{\alpha\in\mathbb R^{|B|},\ \alpha^\top G_B(t)\alpha\le \rho_B(t)^2}
\left[
g_{B,\mathrm{lcb}}(t)^\top \alpha
-
\frac12 \alpha^\top \bigl(\widetilde H_B(t)\bigr)^{(+)}\alpha
\right].
\] Here \(\rho_B(t)>0\) is the Phase-3 batch trust-region radius. If
record identity already includes the favorable orientation, one may
additionally impose \(\alpha\ge 0\) componentwise. For each \(r\in B\),
define the contextual single-record gain in the same common reduction by
\(\Delta E_{r\mid B}^{\mathrm{TR}}(t):=\max_{\beta\in\mathbb R,\ \beta^2 (G_B(t))_{rr}\le \rho_B(t)^2} \left[ g_{\mathrm{lcb}}(r;t)\,\beta - \frac12 \bigl((\widetilde H_B(t))_{rr}\bigr)^{(+)}\beta^2 \right].\)
Then the additivity defect is
\(\delta_{\mathrm{add}}(B;t):=\left[ 1- \frac{\Delta E_B^{\mathrm{TR}}(t)} {\sum_{r\in B}\Delta E_{r\mid B}^{\mathrm{TR}}(t)+\varepsilon} \right]_{+}.\)
Small \(\delta_{\mathrm{add}}(B;t)\) means that the joint reduced model
is approximately additive on the common batch plane.

Batch choice and greedy implementation: the geometric admissible batch
family is
\(\mathfrak B_{\mathrm{geom}}(t):=\left\{ B\in \mathfrak B_{\mathrm{nd}}(t)\cap \mathfrak B_{\mathrm{struct}}(t): \lambda_{\min}\!\bigl(G_B(t)\bigr) \ge \tau_{\mathrm{rank}}\, \frac{\operatorname{tr}G_B(t)}{|B|}, \quad \delta_{\mathrm{add}}(B;t)\le \tau_{\mathrm{add}} \right\}\).
When a more aggressive pruning rule is desired, one may additionally
require
\(\mu_{\mathrm{tan}}(B;t)\le \tau_{\mathrm{tan}}, \qquad \mu_{\mathrm{curv}}(B;t)\le \tau_{\mathrm{curv}}.\)
The principled post-shortlist batch choice is \[
B_t^\star
\in
\arg\max_{B\in \mathfrak B_{\mathrm{geom}}(t)}
\Delta E_B^{\mathrm{TR}}(t).
\] If \(\mathfrak B_{\mathrm{geom}}(t)=\varnothing\), no Phase-3 batch
is admitted from the current shortlist. A greedy implementation avoids
exhaustive search: initialize \(B_t^{(1)}:=\{r_{\max}(t)\}\). Given a
current batch \(B_t^{(k)}\), define the admissible add-on family
\(\mathcal A_k(t):=\left\{ r\in \mathcal N_3(t)\setminus B_t^{(k)}: B_t^{(k)}\cup\{r\}\in \mathfrak B_{\mathrm{nd}}(t)\cap \mathfrak B_{\mathrm{struct}}(t) \right\}\).
For each \(r\in \mathcal A_k(t)\), define the marginal joint gain
\(\Delta_{\mathrm{marg}}(r\mid B_t^{(k)};t):=\Delta E_{B_t^{(k)}\cup\{r\}}^{\mathrm{TR}}(t)-\Delta E_{B_t^{(k)}}^{\mathrm{TR}}(t).\)
Choose
\(r_k^{\mathrm{best}}(t)\in \arg\max_{r\in \mathcal A_k(t)} \Delta_{\mathrm{marg}}(r\mid B_t^{(k)};t)\)
subject to
\(\lambda_{\min}\!\bigl(G_{B_t^{(k)}\cup\{r\}}(t)\bigr) \ge \tau_{\mathrm{rank}}\, \frac{\operatorname{tr}G_{B_t^{(k)}\cup\{r\}}(t)}{|B_t^{(k)}|+1}, \qquad \delta_{\mathrm{add}}(B_t^{(k)}\cup\{r\};t)\le \tau_{\mathrm{add}},\)
and, when enabled,
\(\mu_{\mathrm{tan}}(B_t^{(k)}\cup\{r\};t)\le \tau_{\mathrm{tan}}, \qquad \mu_{\mathrm{curv}}(B_t^{(k)}\cup\{r\};t)\le \tau_{\mathrm{curv}}.\)
If no such \(r\) exists, or if the best feasible marginal gain is
nonpositive, stop and set \(B_t^\star:=B_t^{(k)}\). Otherwise append the
best add-on, \(B_t^{(k+1)}:=B_t^{(k)}\cup\{r_k^{\mathrm{best}}(t)\}\),
and continue until no further feasible positive-gain append exists or
until \(|B_t^{(k)}|=B_{\max}(t)\). The deployed Phase-3 action is the
final batch \(B_t^\star\). If \(|B_t^\star|=1\), this reduces to the
usual single-record admission. If \(|B_t^\star|>1\), the controller
inserts the batch on the current branch in the canonical batch order and
performs the common-window refit on \(W_{B_t^\star}(t)\).

\hypertarget{post-m_p-application-entails-windowed-refitting}{%
\subsection{\texorpdfstring{11.5 Post \(m_p\) Application Entails
Windowed-Refitting}{11.5 Post m\_p Application Entails Windowed-Refitting}}\label{post-m_p-application-entails-windowed-refitting}}

In ADAPT-VQE the variational problem is
\(\min_{\boldsymbol\theta} E_L(\boldsymbol\theta)\) with
\(E_L(\boldsymbol\theta)=\langle\phi_0|(\prod_{\ell=L}^{1}e^{i\theta_\ell G_\ell})H(\prod_{\ell=1}^{L}e^{-i\theta_\ell G_\ell})|\phi_0\rangle\),
and the measurement bottleneck arises because one iteration must
estimate not only \(E_L\) but also pool-selection scores
\(s_\mu(\boldsymbol\theta)=|\langle\phi_0|(\prod_{\ell=L}^{1}e^{i\theta_\ell G_\ell}),i[H,A_\mu],(\prod_{\ell=1}^{L}e^{-i\theta_\ell G_\ell})|\phi_0\rangle|\)
for \(\mu\in\mathcal P\), reoptimization gradients
\(\partial_{\theta_j}E_L(\boldsymbol\theta)\), and, for geometry-aware
updates, metric entries
\(M_{jk}(\boldsymbol\theta)=\Re\langle \partial_{\theta_j}[(\prod_{\ell=1}^{L}e^{-i\theta_\ell G_\ell})|\phi_0\rangle]|\partial_{\theta_k}[(\prod_{\ell=1}^{L}e^{-i\theta_\ell G_\ell})|\phi_0\rangle]\rangle\),
so at precisions \(\epsilon_s,\epsilon_g,\epsilon_M\) the shot cost
obeys
\(N_{\mathrm{iter}}\gtrsim \sum_{\mu\in\mathcal P}\mathrm{Var}(\hat s_\mu)/\epsilon_s^2+\sum_{j=1}^{L}\mathrm{Var}(\widehat{\partial_{\theta_j}E_L})/\epsilon_g^2+\sum_{j,k=1}^{L}\mathrm{Var}(\hat M_{jk})/\epsilon_M^2\);
the proposed strategy is to choose disjoint sets
\(F,A\subset{1,\dots,L}\) with \(F\cup A={1,\dots,L}\), impose
\(\dot\theta_j=0\) for \(j\in F\), and optimize only \(j\in A\), which
reduces the reparameterization sector to
\(\sum_{j\in A}\mathrm{Var}(\widehat{\partial_{\theta_j}E_L})/\epsilon_g^2+\sum_{j,k\in A}\mathrm{Var}(\hat M_{jk})/\epsilon_M^2\)
and therefore removes the growth coming from repeatedly updating old
coordinates, but it leaves the pool-screening term
\(\sum_{\mu\in\mathcal P}\mathrm{Var}(\hat s_\mu)/\epsilon_s^2\)
unchanged and it does not shorten the physical circuit, since every shot
still executes the full depth-\(L\) scaffold with entangling-gate count
\(N_{2q}(L)\) and runtime \(t_{\mathrm{circ}}(L)\), so hardware fidelity
still degrades roughly as
\(e^{-t_{\mathrm{circ}}(L)/T_1}e^{-t_{\mathrm{circ}}(L)/T_2}(1-\varepsilon_{2q})^{N_{2q}(L)}\);
consequently the freezing strategy cures only the reoptimization
component of the ADAPT bottleneck, and the limiting cost after that is
either the unchanged operator-pool screening overhead or, if screening
is separately controlled, the conventional VQE bottleneck of resolving
\(E_L\) on a deep noisy ansatz, with shots scaling as
\(N_E\gtrsim \mathrm{Var}(H_{\mathrm{meas}}\mid \prod_{\ell=1}^{L}e^{-i\theta_\ell G_\ell},\phi_0,\text{noise})/(\Delta E)^2\)
to distinguish an energy change \(\Delta E\).

Let the length-\((m)\) scaffold be
\(U^{(m)}(\boldsymbol{\theta})=\prod_{j=1}^{m} e^{-i\theta_j G_j}\),
where \((G_j)\) is the \((j)\)-th appended generator and \((\theta_j)\)
its coefficient. The reparameterization motif is to optimize only an
active subset \((\mathcal A_m\subseteq {1,\dots,m})\), chosen by an
activity functional
\(a_j^{(m)}:=\lambda_{\mathrm{rec}}\,r_j^{(m)}+\lambda_{\mathrm{mob}}\,v_j^{(m)}-\lambda_{\mathrm{age}}\,\phi(m-j)-\lambda_{\mathrm{stag}}\,s_j^{(m)}\),
and then set \((\mathcal A_m=\{j: a_j^{(m)}>\tau_m\})\). Here
\((r_j^{(m)})\) is a recency weight, large for recently appended
generators;
\((v_j^{(m)}:=\sum_{k=j}^{m-1}\lvert \Delta\theta_j^{(k)}\rvert)\) is a
lifetime mobility measure, large for parameters that have historically
moved substantially; \((\phi(m-j))\) is an age penalty, biasing against
early-added generators; and \((s_j^{(m)})\) is a stagnation indicator,
large when a parameter has remained nearly inert over recent refits.
Thus the inclusion bias is toward recent and historically labile
generators, while the exclusion bias is toward old and typically
stagnant generators.

This is the same scaffold language used above:
\(U^{(m)}(\boldsymbol\theta)\) is the scaffold \(\mathcal O\), the
active set \(\mathcal A_m\) plays the role of the inherited refit window
\(W(r;t)\), the frozen set is \(\{1,\ldots,n+1\}\setminus W(r;t)\), and
the activity functional \(a_j^{(m)}\) is the same maturity-facing
bookkeeping later summarized by the eligible set \(\mathcal M_t\).

Concomitantly, as scaffold maturity increases, the controller shot laws
\(N_{\mathrm{shot},1}(t)\) and \(N_{\mathrm{shot},k}(t)\) should also
increase so as to preserve the decision-level signal-to-noise ratio on
the active block. Formally, if \(\delta_t\) denotes a typical useful
signal scale among the unfrozen directions and \(\sigma_t\) the
corresponding estimator spread, then one chooses the phase-appropriate
shot law so that
\(\mathrm{SNR}_t\sim \frac{\delta_t\sqrt{N_{\mathrm{shot},1}(t)}}{\sigma_t}\ge \kappa\)
in Phase 1 and
\(\mathrm{SNR}_{t,k}\sim \frac{\delta_{t,k}\sqrt{N_{\mathrm{shot},k}(t)}}{\sigma_{t,k}}\ge \kappa\)
for \(k\in\{2,3\}\) in later phases. Hence the control law is: as
scaffold maturity grows, shrink the reoptimized subspace by freezing
archaic or stagnant directions while increasing shots on the surviving
active directions to maintain resolvable updates.

\hypertarget{branch-local-mature-scaffold-pruning-permissibility-frozen-ablation-ranking-and-exact-rollback}{%
\subsubsection{11.5.1 Branch-local mature-scaffold pruning:
permissibility, frozen-ablation ranking, and exact
rollback}\label{branch-local-mature-scaffold-pruning-permissibility-frozen-ablation-ranking-and-exact-rollback}}

All quantities in this subsection are branch-local. On a single active
branch we suppress the branch label. Let
\((\mathcal O_t^-,\theta_t^-,E_t^-)\) be the committed branch state
immediately before the accepted Phase-3 payload \(\mathcal A_t^\star\)
is applied, and let \((\mathcal O_t^+,\theta_t^+,E_t^+)\) be the
committed post-admission state after the ordinary local refit on that
same branch has completed. The admitted-step gain is
\(\Delta_{\mathrm{adm}}(t):=E_t^- - E_t^+\).

Each committed coordinate of \(\mathcal O_t^+\) carries the metadata
\((m_j,\iota_j,r_j,\kappa_j,c_j)\): generator label \(m_j\),
branch-local admission step \(\iota_j\), admission record \(r_j\),
stored selector burden \(\kappa_j:=K_3(r_j;\iota_j)\), and
prune-cooldown counter \(c_j(t)\in\{0,1,\dots,c_{\mathrm{cool}}\}\),
transported with that coordinate under later scaffold edits. Here
\(t_{\mathrm{first}}(m_j)\) denotes the first time the generator \(m_j\)
was ever seen in the run.

Use the maturity-increasing prune-pressure law
\(\lambda_{\mathrm{prune}}(t):=\lambda_{\mathrm{prune}}^{\min}+(\lambda_{\mathrm{prune}}^{\max}-\lambda_{\mathrm{prune}}^{\min})g_t\),
with
\(0\le \lambda_{\mathrm{prune}}^{\min}\le \lambda_{\mathrm{prune}}^{\max}\).
Let the baseline gain-resolution scale be
\(\Delta_{\mathrm{scr}}(t):=\max\!\bigl(\tau_{\mathrm{use},t},\operatorname{EWMA}_{\mathrm{recent}}(\Delta_{\mathrm{adm}}),\Delta_{\mathrm{num}}\bigr)\),
use the noise-aware screening scale
\(\Delta_{\mathrm{scr}}^{\mathrm{noise}}(t):=\bigl(1+u_{\sigma,3}(t)\bigr)\Delta_{\mathrm{scr}}(t)\),
and define the admitted-gain signal-to-noise ratio by
\(\mathrm{SNR}_{\mathrm{adm}}(t):=\Delta_{\mathrm{adm}}(t)\big/\bigl(\Delta_{\mathrm{scr}}^{\mathrm{noise}}(t)+\varepsilon\bigr)\).
A prune pass may open only after an accepted admission and its ordinary
local refit have already committed. Let
\(\mathcal T_{\mathrm{chk}}\subseteq\mathbb N\) be the scheduled
checkpoint set, define
\(C_t^{\mathrm{perm}}:=\bigl(u_{\mathrm{sat}}(t)\ge \tau_{\mathrm{med}}\bigr)\wedge\Bigl(\mathrm{SNR}_{\mathrm{adm}}(t)\le \kappa_{\mathrm{pr}}\ \vee\ t\in\mathcal T_{\mathrm{chk}}\Bigr)\),
and
\(\Gamma_{\mathrm{prune}}^{\mathrm{perm}}(t):=\mathbf 1[\text{accepted admission at }t]\mathbf 1[\text{post-admission local refit completed}]\mathbf 1[|\mathcal O_t^+|\ge 2]\mathbf 1[C_t^{\mathrm{perm}}]\).
Thus pruning is a maturity-controlled maintenance action that opens only
when the scaffold is already in the mature regime and the admitted gain
is no longer cleanly resolved above the noise-aware screening floor, or
when a maturity-gated checkpoint explicitly requests a simplification
test.

Newly admitted coordinates are protected for \(\ell_{\mathrm{protect}}\)
admitted steps via
\(\mathcal P_{\mathrm{protect}}(t):=\{\,j:\ t-\iota_j<\ell_{\mathrm{protect}}\,\}\).
Let \(s_j^{\mathrm{stag}}(t)\) denote the branch-local \(t\)-indexed
lift of the earlier coordinate stagnation indicator \(s_j^{(m)}\), with
larger values meaning less recent parameter motion and therefore greater
local stagnation. The mature prune-eligible set is \[
\mathcal M_t
:=
\Bigl\{
j:\ 
t-t_{\mathrm{first}}(m_j)\ge A_{\mathrm{stale}},
\ 
s_j^{\mathrm{stag}}(t)\ge s_{\mathrm{stale}},
\ 
j\notin\mathcal P_{\mathrm{protect}}(t),
\ 
c_j(t)=0
\Bigr\}.
\] Hence prune eligibility is branch-local and conservative: old by
first-seen age, stale by recent inactivity, not recently admitted, and
not under temporary cooldown.

Small angle is used only as a cheap nomination device. If
\(\mathcal M_t\neq\varnothing\), set
\(\theta_{\mathrm{small}}(t):=\max\!\bigl(\theta_{\mathrm{abs}},c_\theta\,\operatorname{median}_{j\in\mathcal M_t}|\operatorname{wrap}(\theta_j(t))|\bigr)\)
and
\(\mathcal J_{\angle}(t):=\{\,j\in\mathcal M_t:|\operatorname{wrap}(\theta_j(t))|\le \theta_{\mathrm{small}}(t)\,\}\);
let \(\mathcal P_{\mathrm{probe}}(t)\subseteq\mathcal J_{\angle}(t)\) be
the set of the \(\min\{K_p,|\mathcal J_{\angle}(t)|\}\) smallest values
of \(|\operatorname{wrap}(\theta_j(t))|\) over
\(j\in\mathcal J_{\angle}(t)\). If \(\mathcal M_t=\varnothing\), set
\(\mathcal P_{\mathrm{probe}}(t)=\varnothing\). Thus old/stale and
small-angle coordinates are nominated cheaply, but neither condition is
itself an admissible proof of redundancy.

For each \(j\in\mathcal P_{\mathrm{probe}}(t)\), let
\(\mathcal O_{t,-j}^+\) denote the scaffold obtained from
\(\mathcal O_t^+\) by deleting coordinate \(j\), and let
\((\theta_t^+)_{-j}\) be \(\theta_t^+\) with the same entry removed. The
frozen-ablation energy and immediate regression are
\(E_{-j}^{\mathrm{frz}}:=E(\mathcal O_{t,-j}^+,(\theta_t^+)_{-j})\) and
\(L_j^{\mathrm{frz}}(t):=E_{-j}^{\mathrm{frz}}-E_t^+\). Use the cheap
prune score \[
P_j^{\mathrm{pr}}(t)
:=
\frac{[L_j^{\mathrm{frz}}(t)]_+}{1+\kappa_j},
\qquad
j\in\mathcal P_{\mathrm{probe}}(t),
\] and, when \(\mathcal P_{\mathrm{probe}}(t)\neq\varnothing\), choose
\[
j_t^\star
:=
\arg\min_{j\in\mathcal P_{\mathrm{probe}}(t)}
P_j^{\mathrm{pr}}(t).
\] Thus the cheap ranking reuses the same \(K\)-style burden logic
already used for admission: low retained frozen-ablation utility
relative to stored selector burden means good prune-test priority, so
smaller \(P_j^{\mathrm{pr}}(t)\) is better.

Before the authoritative trial, store the exact rollback snapshot
\(\Xi_t^{\mathrm{pr}}:=(\mathcal O_t^+,\theta_t^+,E_t^+,\{c_i(t)\}_i,\text{current local optimizer state})\).
Only one generator is tested per prune pass. If
\(\Gamma_{\mathrm{prune}}^{\mathrm{perm}}(t)=0\) or
\(\mathcal P_{\mathrm{probe}}(t)=\varnothing\), no prune trial is
launched and the committed branch state remains
\((\mathcal O_t^+,\theta_t^+,E_t^+)\).

For the authoritative trial on \(j_t^\star\), let
\(n_t^+:=|\mathcal O_t^+|\),
\(\omega_{\mathrm{eff}}^{(-j)}:=\max\{1,\min(\omega_{\mathrm{pr}},n_t^+-1)\}\),
\(\ell_j:=\max\!\left\{1,\min\!\left(j-\left\lfloor\frac{\omega_{\mathrm{eff}}^{(-j)}-1}{2}\right\rfloor,\ n_t^+-\omega_{\mathrm{eff}}^{(-j)}\right)\right\}\),
\(W_{\mathrm{loc}}(j):=\{\ell_j,\dots,\ell_j+\omega_{\mathrm{eff}}^{(-j)}-1\}\),
\(C_{\mathrm{old}}^{(-j)}:=\{1,\dots,n_t^+-1\}\setminus W_{\mathrm{loc}}(j)\),
and
\(n_{\mathrm{old}}^{\mathrm{pr}}(t,j):=\left\lceil \rho_W(t)\,|C_{\mathrm{old}}^{(-j)}| \right\rceil\).
Let \(W_{\mathrm{age}}^{(-j)}(t)\) be the set of the
\(n_{\mathrm{old}}^{\mathrm{pr}}(t,j)\) oldest surviving coordinates in
\(C_{\mathrm{old}}^{(-j)}\) under the stored admission-order values
\(\iota_i\). The prune-refit window is \[
W_{\mathrm{pr}}(j;t)
:=
W_{\mathrm{loc}}(j)\cup W_{\mathrm{age}}^{(-j)}(t).
\] Hence the prune refit window contracts automatically as maturity
increases, because \(\rho_W(t)\) decreases as useful runway is
exhausted.

The authoritative post-ablation refit is \[
\theta_t^{(-j)}
:=
\arg\min_{\widetilde\theta\in\Theta_{\mathcal O_{t,-j}^+}}
\left\{
E(\mathcal O_{t,-j}^+,\widetilde\theta):
\widetilde\theta_i=((\theta_t^+)_{-j})_i
\ \forall i\notin W_{\mathrm{pr}}(j;t)
\right\},
\] with corresponding post-refit energy
\(E_t^{(-j)}:=E(\mathcal O_{t,-j}^+,\theta_t^{(-j)})\). This
ablation-plus-local-refit result is the authoritative prune test. No
extra geometric score is required at this stage.

Define the retained admitted-step gain after removing \(j\) by
\(\Delta_{\mathrm{keep}}^{(-j)}(t):=E_t^- - E_t^{(-j)}\). Accept the
prune of \(j_t^\star\) iff
\(E_t^{(-j_t^\star)}\le E_t^+ + \Delta_{\mathrm{safe}}\) and
\(\Delta_{\mathrm{keep}}^{(-j_t^\star)}(t)\ge \eta_{\mathrm{ret}}\Delta_{\mathrm{adm}}(t)\).
Thus the final accept rule is recoverability-based: the reduced scaffold
may be kept only when the ablation-plus-refit energy regression is
numerically small and the pruned scaffold still retains a controlled
fraction of the gain produced by the last admitted payload. The burden
term \(\kappa_j\) is used only for cheap ranking; the final accept rule
remains energy-based and conservative.

If the prune is accepted, commit
\((\mathcal O_{t,-j_t^\star}^+,\theta_t^{(-j_t^\star)},E_t^{(-j_t^\star)})\)
as the new branch state and delete the metadata carried by coordinate
\(j_t^\star\). If the prune is rejected, restore the branch state
exactly to \(\Xi_t^{\mathrm{pr}}\), set
\(c_{j_t^\star}(t)\leftarrow c_{\mathrm{cool}}\), and launch no second
prune trial at the same selector step. On later admitted steps the
cooldown counter decrements by one until it reaches zero. Hence prune
failure never invalidates the already committed admission, and rollback
is exact.

\hypertarget{beam-adapt-over-scaffold-alternatives-branch-state-pruning-and-effective-selector}{%
\subsubsection{11.5.2 Beam-adapt over scaffold alternatives: branch
state, pruning, and effective
selector}\label{beam-adapt-over-scaffold-alternatives-branch-state-pruning-and-effective-selector}}

In this section beam-adapt is \textbf{purely structural}: each branch
carries an alternative ADAPT scaffold together with its locally refit
amplitudes. There is no time checkpoint, probe rollout, or
projected-dynamics integral in this section; those belong to the
separate time-dynamics beam surface developed later.

The cleanest way to read a live scaffold branch is \[
\mathfrak b=
\begin{aligned}[t]
\bigl(&\text{identity/history};\ \text{current scaffold and amplitudes};\\
&\text{current energy/depth};\ \text{cumulative selector totals/controller telemetry/status}\bigr).
\end{aligned}
\] More explicitly, write \[
\mathfrak b=\bigl(\operatorname{id}_b,\pi_b,\mathcal H_b,\mathcal O_b,\theta_b,E_b,d_b,S_b^{\mathrm{cum}},K_b^{\mathrm{cum}},\mathcal T_b^{\mathrm{ctrl}},\mathrm{status}_b,\mathrm{term}_b\bigr),
\] with \(\mathcal H_b=(r_{b,1},\dots,r_{b,d_b})\) and
\(r_{b,j}=(m_{b,j}^{\mathrm{adm}},p_{b,j}^{\mathrm{adm}})\). Here
\(\operatorname{id}_b\) and \(\pi_b\) are the branch and parent
identifiers, \(\mathcal H_b\) is the ordered history of admitted
records, \(\mathcal O_b\) is the scaffold currently carried by branch
\(b\), \(\theta_b\) is the refit parameter vector on that scaffold,
\(E_b=E(\theta_b;\mathcal O_b)\) is the current variational energy,
\(d_b\) is the branch ADAPT depth,
\(S_b^{\mathrm{cum}}=\sum_{j=1}^{d_b}s_{b,j}\) is the cumulative
admitted selector value,
\(K_b^{\mathrm{cum}}=\sum_{j=1}^{d_b}K_{b,j}^{\mathrm{sel}}\) is the
cumulative admitted selector burden, \(\mathcal T_b^{\mathrm{ctrl}}\)
stores branch-local controller telemetry, and
\((\mathrm{status}_b,\mathrm{term}_b)\) record live/terminal status and
termination label. The increment \(K_{b,j}^{\mathrm{sel}}\) is the
selector-cost contribution attached to the \(j\)-th admission on branch
\(b\); for a newly materialized child this contribution is the
branch-local Phase-3 cost \(K_3(r;\mathfrak b)\). Thus
\(K_b^{\mathrm{cum}}\) is cumulative selector-side cost bookkeeping
rather than a raw hardware count.

If beam search is entered from an incumbent scaffold
\((\mathcal O^{(0)},\theta^{(0)})\), the root branch is
\(\mathfrak b_{\mathrm{root}}=\bigl(\operatorname{id}_0,\varnothing,\mathcal H^{(0)},\mathcal O^{(0)},\theta^{(0)},E(\theta^{(0)};\mathcal O^{(0)}),d_0,S_{\mathrm{root}}^{\mathrm{cum}},K_{\mathrm{root}}^{\mathrm{cum}},\mathcal T_{\mathrm{root}}^{\mathrm{ctrl}},\mathrm{frontier},\varnothing\bigr)\),
with \(d_0=|\mathcal H^{(0)}|\),
\(S_{\mathrm{root}}^{\mathrm{cum}}=\sum_{j=1}^{d_0}s_j^{(0)}\), and
\(K_{\mathrm{root}}^{\mathrm{cum}}=\sum_{j=1}^{d_0}K_j^{\mathrm{sel},(0)}\),
where \(K_j^{\mathrm{sel},(0)}\) denotes the selector-cost contribution
attached to the \(j\)-th historical admission on the incumbent scaffold.
If the beam is started from a fresh scaffold, then
\(\mathcal H^{(0)}=\varnothing\), \(d_0=0\), and both cumulative sums
vanish.

For clock and budget, set \(t_b:=d_b\) and
\(D_{\mathrm{left}}(\mathfrak b):=d_{\max}-d_b\). Now apply the branch
lift rule: for any live-state quantity \(X\) from Phase 3 whose
definition is written in terms of
\((\mathcal H,\mathcal O,\theta,t,\mathcal T^{\mathrm{ctrl}})\), define
\(X(\mathfrak b):=X\big|_{(\mathcal H,\mathcal O,\theta,t,\mathcal T^{\mathrm{ctrl}})\mapsto(\mathcal H_b,\mathcal O_b,\theta_b,t_b,\mathcal T_b^{\mathrm{ctrl}})}\);
for any record-local quantity \(X(r;t)\), define
\(X(r;\mathfrak b):=X(r;t)\big|_{(\mathcal H,\mathcal O,\theta,t,\mathcal T^{\mathrm{ctrl}})\mapsto(\mathcal H_b,\mathcal O_b,\theta_b,t_b,\mathcal T_b^{\mathrm{ctrl}})}\).
In particular, \(R(\mathfrak b)\), \(u_{\mathrm{sat}}(\mathfrak b)\),
\(\widehat N_{\mathrm{rem}}(\mathfrak b)\), \(K_3(r;\mathfrak b)\), and
\(g_{\mathrm{lcb}}(r;\mathfrak b)\) are branch-lifted evaluations of
their previously defined live-state counterparts.

The Phase-3 chain is branch-local by construction: all of its
dependencies factor through the state tuple
\((\mathcal O_b,\theta_b,t_b,\mathcal T_b^{\mathrm{ctrl}})\). Hence each
live branch induces its own raw candidate-position surface
\(\mathcal R_b^{\mathrm{raw}}\), its own branch-local shortlist
\(\mathcal S_3(\mathfrak b)\), and its own retained admission set
\(\mathcal A_b\).

For stage and enabled family, let
\(\eta_b\in\{\mathrm{core/seed},\mathrm{residual}\}\) and \[
\mathcal G_{\eta_b}^{(3)}(\mathfrak b)=
\begin{cases}
\mathcal G_{\mathrm{core/seed}}^{(3)}(\mathfrak b),&\eta_b=\mathrm{core/seed},\\
\mathcal G_{\mathrm{residual}}^{(3)}(\mathfrak b),&\eta_b=\mathrm{residual}.
\end{cases}
\] Let
\(\mathcal P_{\mathrm{residual}}(\mathfrak b)\subseteq\mathcal G^{(3)}\)
denote the residual-only generator family on branch \(b\). Then the
controller-stage gate and branch-local available Phase-3 generator
family are \[
\Gamma_{\mathrm{stage},b}(m)=
\begin{cases}
1,&\eta_b=\mathrm{residual},\\
1,&\eta_b=\mathrm{core/seed}\ \text{and}\ m\notin\mathcal P_{\mathrm{residual}}(\mathfrak b),\\
0,&\eta_b=\mathrm{core/seed}\ \text{and}\ m\in\mathcal P_{\mathrm{residual}}(\mathfrak b),
\end{cases}
\qquad
\mathcal G_{\mathrm{avail}}^{(3)}(\mathfrak b)=\{m\in\mathcal G^{(3)}:\Gamma_{\mathrm{stage},b}(m)=1\}.
\] For scaffold length, append slot, and active probe window, let
\(n_b=|\mathcal O_b|\), \(p_{\mathrm{app}}(\mathfrak b)=n_b\), and
\(W_{\mathrm{act}}(\mathfrak b)\subseteq\{0,\dots,n_b-1\}\). The ordered
probe list is
\(\mathcal L_{\mathrm{probe}}(\mathfrak b):=\bigl(p_{\mathrm{app}}(\mathfrak b),0,W_{\mathrm{act}}(\mathfrak b)\bigr)\),
and the distinct retained probe positions are
\(\mathcal P_{\mathrm{probe}}(\mathfrak b)=\operatorname{Head}_{M_{\mathrm{probe}}}\!\Bigl(\operatorname{Dedup}_{\mathrm{ord}}\bigl(\mathcal L_{\mathrm{probe}}(\mathfrak b)\bigr)\Bigr)\subseteq\{0,\dots,n_b\}\).
For symmetry and admissible positions, let \[
\Gamma_{\mathrm{sym},b}(m,p)=
\begin{cases}
0,&\text{the symmetry audit for }(m,p;\mathcal O_b)\text{ returns a hard violation},\\
1,&\text{otherwise},
\end{cases}
\] so the branch-local raw candidate-position surface is \[
\mathcal R_b^{\mathrm{raw}}
:=
\{\,r=(m,p):m\in\mathcal G_{\mathrm{avail}}^{(3)}(\mathfrak b),\ p\in\mathcal P_{\mathrm{probe}}(\mathfrak b),\ \Gamma_{\mathrm{sym},b}(m,p)=1\,\}.
\] The resulting local shortlist and retained admission set are \[
\mathcal S_3(\mathfrak b)
:=
\operatorname{Top}_{N_3(t_b)}\!\left(S_3(\cdot;\mathfrak b);\mathcal R_b^{\mathrm{raw}}\right),
\qquad
\mathcal A_b
:=
\operatorname{Top}_{K_b}\!\left(S_3(\cdot;\mathfrak b);\{r\in\mathcal S_3(\mathfrak b):S_3(r;\mathfrak b)>0\}\right).
\] When a branch-local rerank shell is judged tie-degenerate, the
retained admission set may be replaced by a transient tie-shell
selector.\footnote{For example, if $r_{b,\max}\in\arg\max_{r\in\mathcal S_3(\mathfrak b)}S_3(r;\mathfrak b)$, one may define $\mathcal T_{\mathrm{tie}}(\mathfrak b):=\{\,r\in\mathcal S_3(\mathfrak b): S_3(r;\mathfrak b)\ge \max(\eta_{\mathrm{tie}}\,S_3(r_{b,\max};\mathfrak b),\,S_3(r_{b,\max};\mathfrak b)-\delta_{\mathrm{tie}})\,\}$ and replace $\mathcal A_b$ by a retained subset $\mathcal A_b^{\mathrm{tie}}\subseteq \mathcal T_{\mathrm{tie}}(\mathfrak b)$ with $1\le |\mathcal A_b^{\mathrm{tie}}|\le B_{\mathrm{tie}}(t_b)$, chosen by descending $S_3(\cdot;\mathfrak b)$.}

\% {[}ARCHAIC/REVIEW: split-aware promotion machinery.{]} \% Origin:
§11.1 split family definition. \% Status: removed from the main beam
surface and retained only for provenance. \% If reinstated, it should
act on children of a single generator family at fixed position p, \% not
on unrelated branch children from different global generators.

\textbf{(ii) Materialize either stop or one new admission.} The proposal
family is \(\mathcal Q(b)=\{\mathrm{stop}\}\cup\mathcal A_b\). The
symbol \(\mathrm{stop}\) means
\texttt{terminate\ this\ scaffold\ as-is,\textquotesingle{}\textquotesingle{}\ while\ each\ \$r=(m,p)\textbackslash{}in\textbackslash{}mathcal\ A\_b\$\ means}insert
one new generator into this branch and refit.'\,' If
\(\mathcal O_b=(m_{b,1},\dots,m_{b,n_b})\) and \(0\le p\le n_b\), then
the scaffold insertion map is
\(\operatorname{Insert}(\mathcal O_b,m,p)=(m_{b,1},\dots,m_{b,p},m,m_{b,p+1},\dots,m_{b,n_b})\),
the insertion index transport is
\(I_p(j)=\begin{cases} j,&1\le j\le p,\\ j+1,&p<j\le n_b, \end{cases}\),
and the zero-lifted parameter vector is
\(\theta_b^{\uparrow(r)}=(\theta_{b,1},\dots,\theta_{b,p},0,\theta_{b,p+1},\dots,\theta_{b,n_b})\).

For contiguous inserted-window cap \(\omega_{\mathrm{new}}\), let
\(\omega_{\mathrm{eff}}^{(p)}:=\max\{1,\min(\omega_{\mathrm{new}},n_b+1)\}\)
and
\(\ell_p:=\max\!\left\{1,\min\!\left(p+1-\left\lfloor\frac{\omega_{\mathrm{eff}}^{(p)}-1}{2}\right\rfloor,\;n_b+2-\omega_{\mathrm{eff}}^{(p)}\right)\right\}\).
Then
\(W_{\mathrm{new}}^{(p)}=\{\ell_p,\;\ell_p+1,\;\ldots,\;\ell_p+\omega_{\mathrm{eff}}^{(p)}-1\}\),
so \(|W_{\mathrm{new}}^{(p)}|=\omega_{\mathrm{eff}}^{(p)}\) and
\(p+1\in W_{\mathrm{new}}^{(p)}\). Let the retained old-coordinate
fraction shrink with maturity according to
\(\rho_W(t):=\rho_W^{\min}+(\rho_W^{\max}-\rho_W^{\min})e_t\), with
\(0\le \rho_W^{\min}\le \rho_W^{\max}\le 1\), and define the
old-coordinate complement by
\(C_{\mathrm{old}}^{(p)}:=\{1,\dots,n_b+1\}\setminus W_{\mathrm{new}}^{(p)}\),
so \(|C_{\mathrm{old}}^{(p)}|=n_b+1-\omega_{\mathrm{eff}}^{(p)}\). Then
the retained old-coordinate count is
\(n_{\mathrm{old}}(t,p):=\left\lceil \rho_W(t)\,|C_{\mathrm{old}}^{(p)}|\right\rceil=\left\lceil \rho_W(t)\,\bigl(n_b+1-\omega_{\mathrm{eff}}^{(p)}\bigr)\right\rceil\).
For each current scaffold coordinate \(j\in\{1,\dots,n_b\}\), let
\(a_{b,j}^{\mathrm{adm}}\) denote the admission step at which that
coordinate entered the current scaffold on branch \(b\), with smaller
values meaning older coordinates, and let the insertion-lifted
admission-order vector be
\(a_b^{\uparrow(r)}:=(a_{b,1}^{\mathrm{adm}},\dots,a_{b,p}^{\mathrm{adm}},+\infty,a_{b,p+1}^{\mathrm{adm}},\dots,a_{b,n_b}^{\mathrm{adm}})\),
so the newly inserted slot is never selected as an old coordinate. The
retained old-coordinate window is then
\(W_{\mathrm{age}}^{(p)}:=\{\,j\in C_{\mathrm{old}}^{(p)}:\#\{\,i\in C_{\mathrm{old}}^{(p)}:(a_b^{\uparrow(r)})_i<(a_b^{\uparrow(r)})_j\,\}<n_{\mathrm{old}}(t,p)\,\}\),
namely the set of the \(n_{\mathrm{old}}(t,p)\) oldest old coordinates
outside \(W_{\mathrm{new}}^{(p)}\), and hence the inherited refit window
on the inserted scaffold is
\(W_r=W(r;t):=W_{\mathrm{new}}^{(p)}\cup W_{\mathrm{age}}^{(p)}\).

Then, for \(r=(m,p)\) and
\(\mathcal O_b\oplus_p m:=\operatorname{Insert}(\mathcal O_b,m,p)\), the
branch-local refit map is
\(\operatorname{Refit}_{W_r}(\theta_b;r):=\arg\min_{\widetilde\theta\in\Theta_{\mathcal O_b\oplus_p m}}\left\{E(\mathcal O_b\oplus_p m,\widetilde\theta):\ \widetilde\theta_j=(\theta_b^{\uparrow(r)})_j\ \forall j\notin W_r\right\}\).

The transfer map is
\(\mathcal T(\mathfrak b,\mathrm{stop})=\mathfrak b'\) with
\(\mathcal H_{b'}=\mathcal H_b\), \(\mathcal O_{b'}=\mathcal O_b\),
\(\theta_{b'}=\theta_b\), \(E_{b'}=E_b\), \(d_{b'}=d_b\),
\(S_{b'}^{\mathrm{cum}}=S_b^{\mathrm{cum}}\),
\(K_{b'}^{\mathrm{cum}}=K_b^{\mathrm{cum}}\),
\(\mathrm{status}_{b'}=\mathrm{terminal}\), and
\(\mathrm{term}_{b'}=\begin{cases} \mathrm{empty},&\mathcal A_b=\varnothing,\\ \mathrm{stop},&\mathcal A_b\neq\varnothing. \end{cases}\)
For \(r=(m,p)\), \(\mathcal T(\mathfrak b,r)=\mathfrak b'\) with
\(\mathcal O_{b'}=\operatorname{Insert}(\mathcal O_b,m,p)\),
\(\theta_{b'}=\operatorname{Refit}_{W_r}(\theta_b;r)\),
\(\mathcal H_{b'}=(\mathcal H_b,r)\),
\(E_{b'}=E(\theta_{b'};\mathcal O_{b'})\), \(d_{b'}=d_b+1\),
\(S_{b'}^{\mathrm{cum}}=S_b^{\mathrm{cum}}+S_3(r;\mathfrak b)\),
\(K_{b'}^{\mathrm{cum}}=K_b^{\mathrm{cum}}+K_3(r;\mathfrak b)\),
\(\mathrm{status}_{b'}=\begin{cases} \mathrm{terminal},&d_{b'}\ge d_{\max}\ \text{or}\ \mathcal A_{b'}=\varnothing,\\ \mathrm{frontier},&d_{b'}<d_{\max}\ \text{and}\ \mathcal A_{b'}\neq\varnothing, \end{cases}\),
and
\(\mathrm{term}_{b'}=\begin{cases} \mathrm{depth\_cap},&d_{b'}\ge d_{\max},\\ \mathrm{empty},&d_{b'}<d_{\max}\ \text{and}\ \mathcal A_{b'}=\varnothing,\\ \varnothing,&d_{b'}<d_{\max}\ \text{and}\ \mathcal A_{b'}\neq\varnothing. \end{cases}\)
where \(\mathcal A_{b'}\) is recomputed from step (i) on the child
branch itself. So a non-stop child is exactly one more ADAPT admission
applied to the parent scaffold, followed by the inherited-window refit
already defined above, and its cumulative \(K\)-coordinate is updated by
the admitted branch-local Phase-3 selector burden
\(K_3(r;\mathfrak b)\). If \(\mathcal A_b=\varnothing\), then
\(\mathcal Q(b)=\{\mathrm{stop}\}\) and the branch contributes only a
terminal child.

\textbf{(iii) Classify, deduplicate, and prune the scaffold children.}
The frontier and terminal descendants of branch \(b\) are
\(\mathfrak C_b^{\mathrm{front}}=\{\mathfrak b' : \mathfrak b'=\mathcal T(\mathfrak b,r),\ r\in\mathcal A_b,\ \mathrm{status}_{b'}=\mathrm{frontier}\}\)
and
\(\mathfrak C_b^{\mathrm{term}}=\{\mathfrak b' : \mathfrak b'=\mathcal T(\mathfrak b,q),\ q\in\mathcal Q(b),\ \mathrm{status}_{b'}=\mathrm{terminal}\}\).
Two branches are treated as duplicates when they represent the same
scaffold and essentially the same refit point, for example under the
fingerprint
\(\operatorname{rnd}_{10}(x):=10^{-10}\operatorname{round}(10^{10}x)\),
\(\operatorname{round}_{10}(\theta_b):=\bigl(\operatorname{rnd}_{10}(\theta_{b,1}),\dots,\operatorname{rnd}_{10}(\theta_{b,n_b})\bigr)\),
and
\(\mathrm{fp}(\mathfrak b)=\bigl(d_b,\operatorname{labels}(\mathcal O_b),\operatorname{round}_{10}(\theta_b)\bigr)\).
Among near-equivalent branches, keep the one with best lexicographic
prune key
\(\pi_{\mathrm{prune}}(\mathfrak b)=\bigl(E_b,-S_b^{\mathrm{cum}},K_b^{\mathrm{cum}},|\mathcal O_b|,\operatorname{labels}(\mathcal O_b),\operatorname{round}_{10}(\theta_b),\operatorname{id}_b\bigr)\).

The beam operators are therefore
\(\operatorname{Dedup}(\mathcal X)=\{\text{best branch per fingerprint under }\pi_{\mathrm{prune}}\}\)
and
\(\operatorname{Prune}(\mathcal X;B)=\operatorname{Top}^{\min}_B\!\left(\pi_{\mathrm{prune}};\operatorname{Dedup}(\mathcal X)\right)\).
Hence the live and terminal pools update by
\(\mathfrak F_{c+1}=\operatorname{Prune}\!\left(\bigcup_{b\in\mathfrak F_c}\mathfrak C_b^{\mathrm{front}};B_{\mathrm{live}}\right)\)
and
\(\mathfrak T_{c+1}=\operatorname{Prune}\!\left(\mathfrak T_c\cup\bigcup_{b\in\mathfrak F_c}\mathfrak C_b^{\mathrm{term}};B_{\mathrm{term}}\right)\).

\textbf{(iv) Choose the final scaffold branch.} After \(C\)
scaffold-beam rounds, the surviving candidate set is
\(\mathfrak W_{\mathrm{final}}=\mathfrak F_C\cup\mathfrak T_C\), and the
effective beam selector is
\(\mathfrak b_*=\arg\min_{\mathfrak b\in\mathfrak W_{\mathrm{final}}}\pi_{\mathrm{prune}}(\mathfrak b)\).

A compact scaffold-beam algorithm block is therefore \[
\begin{aligned}
\text{A: }&
\mathfrak F_0=\{\mathfrak b_{\mathrm{root}}\},
\qquad
\mathfrak T_0=\varnothing,\\
\text{B: }&
\mathcal R_b^{\mathrm{raw}}
\to
\mathcal S_3(\mathfrak b)
\to
\mathcal A_b
\qquad(\mathfrak b\in\mathfrak F_c),\\
\text{C: }&
\mathcal Q(b)=\{\mathrm{stop}\}\cup\mathcal A_b,
\qquad
\mathfrak C_b=\{\mathcal T(\mathfrak b,q):q\in\mathcal Q(b)\},\\
\text{D: }&
\mathfrak C_b^{\mathrm{front}}
=
\{\mathfrak b'\in\mathfrak C_b:\mathrm{status}_{b'}=\mathrm{frontier}\},
\qquad
\mathfrak C_b^{\mathrm{term}}
=
\{\mathfrak b'\in\mathfrak C_b:\mathrm{status}_{b'}=\mathrm{terminal}\},\\
\text{E: }&
\mathfrak F_{c+1}
=
\operatorname{Prune}\!\left(\bigcup_{b\in\mathfrak F_c}\mathfrak C_b^{\mathrm{front}};B_{\mathrm{live}}\right),
\qquad
\mathfrak T_{c+1}
=
\operatorname{Prune}\!\left(\mathfrak T_c\cup\bigcup_{b\in\mathfrak F_c}\mathfrak C_b^{\mathrm{term}};B_{\mathrm{term}}\right),\\
\text{F: }&
\mathfrak W_{\mathrm{final}}=\mathfrak F_C\cup\mathfrak T_C,
\qquad
\mathfrak b_*=\arg\min_{\mathfrak b\in\mathfrak W_{\mathrm{final}}}\pi_{\mathrm{prune}}(\mathfrak b).
\end{aligned}
\] \#\# 11.6 Active HH pool surface and selected scaffold Identify the
nested active HH pool with the retained generator images by
\(\Omega_{\mathrm{HH}}^{(k)}:=\mathcal G^{(k)}\) for \(k\in\{1,2,3\}\).
The active HH adaptive surface is therefore carried by
\(\Omega_{\mathrm{HH}}^{(3)}\subseteq\Omega_{\mathrm{HH}}^{(2)}\subseteq\Omega_{\mathrm{HH}}^{(1)}\),
with
\(\mathcal G_{\mathrm{adapt}}^{(k)}=\{\tau_m\}_{m\in\Omega_{\mathrm{HH}}^{(k)}}\).

If beam continuation is used, set
\(\mathcal O_*:=\mathcal O_{\mathfrak b_*}\) and
\(\theta_*^{\mathrm{adapt}}:=\theta_{\mathfrak b_*}\); otherwise
\((\mathcal O_*,\theta_*^{\mathrm{adapt}})\) denotes the final admitted
scaffold and its refit amplitudes on the main branch. The resulting
state is
\(|\psi_*\rangle=U(\theta_*^{\mathrm{adapt}};\mathcal O_*)|\phi_0\rangle\),
so the active HH scaffold manifold selected by ADAPT may be written as
\(\mathcal M_{\mathrm{scaf}}(\mathcal O_*)=\{\,U(\vartheta;\mathcal O_*)|\phi_0\rangle:\vartheta\in\mathbb R^{|\mathcal O_*|}\,\}\).

\hypertarget{spsa-and-optimizer-semantics}{%
\section{12. SPSA and Optimizer
Semantics}\label{spsa-and-optimizer-semantics}}

\hypertarget{spsa-schedules}{%
\subsection{12.1 SPSA schedules}\label{spsa-schedules}}

The standard SPSA schedules are \[
c_k=\frac{c}{(k+1)^\gamma},
\qquad
a_k=\frac{a}{(A+k+1)^\alpha}.
\]

\hypertarget{two-point-stochastic-gradient}{%
\subsection{12.2 Two-point stochastic
gradient}\label{two-point-stochastic-gradient}}

At iteration \(k\), with Rademacher direction \(\Delta_k\), evaluate \[
y_+=f(x_k+c_k\Delta_k),
\qquad
y_-=f(x_k-c_k\Delta_k),
\] and form the gradient estimate \[
\hat g_k = \frac{y_+-y_-}{2c_k}\,\Delta_k.
\]

\hypertarget{update-and-projection}{%
\subsection{12.3 Update and projection}\label{update-and-projection}}

The update is \[
x_{k+1}=x_k-a_k\hat g_k.
\] If projection or clipping is imposed, it is applied to both the
perturbed points and the updated iterate.

\hypertarget{repeat-aggregation}{%
\subsection{12.4 Repeat aggregation}\label{repeat-aggregation}}

If the same point is sampled multiple times, the repeated evaluations
may be aggregated by \[
\operatorname{mean}
\qquad\text{or}\qquad
\operatorname{median}.
\]

\hypertarget{return-policy}{%
\subsection{12.5 Return policy}\label{return-policy}}

Two standard return policies are:

\begin{itemize}
\tightlist
\item
  return an averaged late-trajectory iterate,
\item
  return the best sampled point observed during the run.
\end{itemize}

\hypertarget{surface-note}{%
\subsection{12.6 Surface note}\label{surface-note}}

Stochastic inner optimizers are especially natural for large, noisy, or
nondifferentiable effective surfaces, while deterministic optimizers are
often effective on smaller and smoother fixed-structure manifolds.

\hypertarget{a.-projective-mclachlan-real-time-dynamics-and-exact-forecast-checkpoint-adaptive-snake-control}{%
\section{17A. Projective McLachlan Real-Time Dynamics and Exact-Forecast
Checkpoint-Adaptive Snake
Control}\label{a.-projective-mclachlan-real-time-dynamics-and-exact-forecast-checkpoint-adaptive-snake-control}}

The canonical real-time law is piecewise smooth in physical time and
piecewise discrete in scaffold structure. On each open checkpoint
interval the scaffold is frozen and the state follows the projective
McLachlan projection of the sampled Schrödinger drift. At the checkpoint
itself one chooses exactly one structural action: keep the scaffold,
append one new ordered block, or delete one old block. In what follows,
each object is defined once, the main quantities are expanded linearly
down to the operator/expectation-value primitives, and the condensed
shorthand is given only after the fully expanded form has already
appeared.

The driven \(L=2\) controller printed at the front of this manuscript is
one calibrated instance of the law below, but the formulas are written
for a general scaffold length, a general driven Hamiltonian, a general
finite forecast horizon, and a general left/midpoint/right
drive-sampling rule.

\hypertarget{a.1-fixed-scaffold-projective-mclachlan-geometry}{%
\subsection{17A.1 Fixed-scaffold projective McLachlan
geometry}\label{a.1-fixed-scaffold-projective-mclachlan-geometry}}

Here \(n_k:=|\mathcal O_k|\) is the current scaffold length,
\(d:=\dim\mathcal H\) is the Hilbert-space dimension,
\(a\in\{1,\dots,d\}\) is the Hilbert-space row index, and
\(i,j\in\{1,\dots,n_k\}\) are scaffold-coordinate indices. The
quantities introduced in this subsection are the scaffold state, its
projected tangents, the projected Schrödinger drift, the tangent metric,
the force vector, the regularized normal matrix, the quadratic McLachlan
fit, and the projected velocity on the frozen scaffold.

\[
\begin{aligned}
|\psi_k\rangle
&:=|\psi(\theta_k;\mathcal O_k)\rangle
 =U(\theta_k;\mathcal O_k)|\psi_{\mathrm{ref}}\rangle,
\qquad
U(\theta_k;\mathcal O_k):=\prod_{j=n_k}^{1}e^{-i\theta_{k,j}A_{k,j}},\\[1mm]
U_{k,>j}
&:=\prod_{\ell=n_k}^{j+1}e^{-i\theta_{k,\ell}A_{k,\ell}},
\qquad
U_{k,<j}:=\prod_{\ell=j-1}^{1}e^{-i\theta_{k,\ell}A_{k,\ell}},
\qquad
\widetilde A_{k,j}:=U_{k,>j}A_{k,j}U_{k,>j}^{\dagger},\\[1mm]
\tau_{k,j}
&:=\partial_{\theta_{k,j}}|\psi_k\rangle
 =-iU_{k,>j}A_{k,j}U_{k,<j}|\psi_{\mathrm{ref}}\rangle
 =-i\widetilde A_{k,j}|\psi_k\rangle,\\[1mm]
\bar\tau_{k,j}
&:=\tau_{k,j}-\langle\psi_k,\tau_{k,j}\rangle|\psi_k\rangle
 =Q_{\psi_k}\tau_{k,j}
 =-i\bigl(\widetilde A_{k,j}-\langle\widetilde A_{k,j}\rangle_{\psi_k}\bigr)|\psi_k\rangle,
\qquad
Q_{\psi_k}:=I-|\psi_k\rangle\langle\psi_k|,\\[1mm]
\bar T_k
&:=\bigl[(\bar T_k)_{aj}\bigr]_{a=1,\dots,d}^{\,j=1,\dots,n_k}
 =\bigl[(\bar\tau_{k,j})_a\bigr]_{a=1,\dots,d}^{\,j=1,\dots,n_k}\\
&=\left[
\begin{array}{ccc}
-i\bigl((\widetilde A_{k,1}-\langle\widetilde A_{k,1}\rangle_{\psi_k})\psi_k\bigr)_1 & \cdots & -i\bigl((\widetilde A_{k,n_k}-\langle\widetilde A_{k,n_k}\rangle_{\psi_k})\psi_k\bigr)_1\\
\vdots & \ddots & \vdots\\
-i\bigl((\widetilde A_{k,1}-\langle\widetilde A_{k,1}\rangle_{\psi_k})\psi_k\bigr)_d & \cdots & -i\bigl((\widetilde A_{k,n_k}-\langle\widetilde A_{k,n_k}\rangle_{\psi_k})\psi_k\bigr)_d
\end{array}
\right],\\[1mm]
\bar b_k
&:=Q_{\psi_k}\bigl(-iH_k^{\sharp}|\psi_k\rangle\bigr)
 =-i\bigl(H_k^{\sharp}-E_{\psi_k,k}^{\sharp}\bigr)|\psi_k\rangle,
\qquad
E_{\psi_k,k}^{\sharp}:=\langle\psi_k|H_k^{\sharp}|\psi_k\rangle,\\[1mm]
&=\left[
\begin{array}{c}
\bigl(-i(H_k^{\sharp}-E_{\psi_k,k}^{\sharp})\psi_k\bigr)_1\\
\vdots\\
\bigl(-i(H_k^{\sharp}-E_{\psi_k,k}^{\sharp})\psi_k\bigr)_d
\end{array}
\right].
\end{aligned}
\]

The metric, force, regularized system, quadratic fit, and projected
velocity are then defined in one chain:

\[
\begin{aligned}
G_k
&:=\bigl[(G_k)_{ij}\bigr]_{i,j=1}^{n_k}
 =\Re(\bar T_k^{\dagger}\bar T_k)
 =\left[\Re\!\left(\sum_{a=1}^{d}(\bar T_k)_{ai}^{*}(\bar T_k)_{aj}\right)\right]_{i,j=1}^{n_k}\\
&=\left[\Re\!\Bigl(\langle\psi_k|\widetilde A_{k,i}\widetilde A_{k,j}|\psi_k\rangle-\langle\widetilde A_{k,i}\rangle_{\psi_k}\langle\widetilde A_{k,j}\rangle_{\psi_k}\Bigr)\right]_{i,j=1}^{n_k},\\[1mm]
f_k
&:=\bigl[(f_k)_i\bigr]_{i=1}^{n_k}
 =\Re(\bar T_k^{\dagger}\bar b_k)
 =\left[\Re\!\left(\sum_{a=1}^{d}(\bar T_k)_{ai}^{*}(\bar b_k)_a\right)\right]_{i=1}^{n_k}\\
&=\left[\Re\!\Bigl(\langle\psi_k|\widetilde A_{k,i}H_k^{\sharp}|\psi_k\rangle-\langle\widetilde A_{k,i}\rangle_{\psi_k}E_{\psi_k,k}^{\sharp}\Bigr)\right]_{i=1}^{n_k},\\[1mm]
K_k
&:=\bigl[(K_k)_{ij}\bigr]_{i,j=1}^{n_k}
 =G_k+\Lambda_k
 =\left[(G_k)_{ij}+(\Lambda_k)_{ij}\right]_{i,j=1}^{n_k}\\
&=\left[\Re\!\Bigl(\langle\psi_k|\widetilde A_{k,i}\widetilde A_{k,j}|\psi_k\rangle-\langle\widetilde A_{k,i}\rangle_{\psi_k}\langle\widetilde A_{k,j}\rangle_{\psi_k}\Bigr)+(\Lambda_k)_{ij}\right]_{i,j=1}^{n_k},\\[1mm]
\mathfrak m_k(v)
&:=\|\bar T_kv-\bar b_k\|_2^2+v^{\top}\Lambda_kv\\
&=\left\|\sum_{j=1}^{n_k}v_j\bigl[-i(\widetilde A_{k,j}-\langle\widetilde A_{k,j}\rangle_{\psi_k})|\psi_k\rangle\bigr]+i(H_k^{\sharp}-E_{\psi_k,k}^{\sharp})|\psi_k\rangle\right\|_2^2\\
&=\sum_{i,j=1}^{n_k}v_i(K_k)_{ij}v_j-2\sum_{i=1}^{n_k}v_i(f_k)_i+\langle\psi_k|(H_k^{\sharp}-E_{\psi_k,k}^{\sharp})^2|\psi_k\rangle,\\[1mm]
\mathcal Q_k(v)
&:=2v^{\top}f_k-v^{\top}G_kv\\
&=2\Re\!\left\langle\sum_{j=1}^{n_k}v_j\bigl[-i(\widetilde A_{k,j}-\langle\widetilde A_{k,j}\rangle_{\psi_k})|\psi_k\rangle\bigr],-i(H_k^{\sharp}-E_{\psi_k,k}^{\sharp})|\psi_k\rangle\right\rangle\\
&\qquad-\left\|\sum_{j=1}^{n_k}v_j\bigl[-i(\widetilde A_{k,j}-\langle\widetilde A_{k,j}\rangle_{\psi_k})|\psi_k\rangle\bigr]\right\|_2^2,\\[1mm]
\dot\theta_k
&:=K_k^{+}f_k\\
&=\left(\left[\Re\!\Bigl(\langle\psi_k|\widetilde A_{k,i}\widetilde A_{k,j}|\psi_k\rangle-\langle\widetilde A_{k,i}\rangle_{\psi_k}\langle\widetilde A_{k,j}\rangle_{\psi_k}\Bigr)+(\Lambda_k)_{ij}\right]_{i,j=1}^{n_k}\right)^{+}
\left[\Re\!\Bigl(\langle\psi_k|\widetilde A_{k,i}H_k^{\sharp}|\psi_k\rangle-\langle\widetilde A_{k,i}\rangle_{\psi_k}E_{\psi_k,k}^{\sharp}\Bigr)\right]_{i=1}^{n_k},\\[1mm]
\operatorname{Var}_{\psi_k}(H_k^{\sharp})
&:=\|\bar b_k\|_2^2
 =\langle\psi_k|(H_k^{\sharp}-E_{\psi_k,k}^{\sharp})^2|\psi_k\rangle.
\end{aligned}
\]

Equivalently, the regularized fixed-scaffold McLachlan solve may be
written in condensed form as \(K_k\dot\theta_k=f_k\), while the
intrinsic unregularized projective optimum is
\(\dot\theta_k^{\mathrm{proj}}:=G_k^{+}f_k\) with
\(\mathcal Q_k(\dot\theta_k^{\mathrm{proj}})=f_k^{\top}G_k^{+}f_k\).

\hypertarget{a.2-piecewise-constant-driven-hamiltonian-on-each-checkpoint-interval}{%
\subsection{17A.2 Piecewise-constant driven Hamiltonian on each
checkpoint
interval}\label{a.2-piecewise-constant-driven-hamiltonian-on-each-checkpoint-interval}}

Here \(H_{\mathrm{stat}}\) is the static Pauli polynomial,
\(H_{\mathrm{drv}}(t)\) is the driven Pauli polynomial,
\(I_k:=[\tau_k,\tau_{k+1}]\) is the \(k\)th checkpoint interval,
\(\Delta\tau_k:=\tau_{k+1}-\tau_k\), \(M_k\ge 1\) is the exact-reference
micro-step count, and \(\sigma_k\in\{0,\tfrac12,1\}\) selects left,
midpoint, or right time sampling. The driven Hamiltonian is kept simple
and inline:

\(H(t):=H_{\mathrm{stat}}+H_{\mathrm{drv}}(t)=\sum_{\alpha\in\mathcal I_{\mathrm{stat}}}h_{\alpha}P_{\alpha}+\sum_{\beta\in\mathcal I_{\mathrm{drv}}}c_{\beta}(t)P_{\beta}\),
\(t_k^{\sharp}:=t_0+\tau_k+\sigma_k\Delta\tau_k\), and
\(H_k^{\sharp}:=H(t_k^{\sharp})=\sum_{\alpha\in\mathcal I_{\mathrm{stat}}}h_{\alpha}P_{\alpha}+\sum_{\beta\in\mathcal I_{\mathrm{drv}}}c_{\beta}(t_k^{\sharp})P_{\beta}\).

With this sampled Hamiltonian, the force entries of §17A.1 become
\((f_k)_i=\sum_{\alpha\in\mathcal I_{\mathrm{stat}}}h_{\alpha}\Re\langle\bar\tau_{k,i},-iP_{\alpha}\psi_k\rangle+\sum_{\beta\in\mathcal I_{\mathrm{drv}}}c_{\beta}(t_k^{\sharp})\Re\langle\bar\tau_{k,i},-iP_{\beta}\psi_k\rangle\),
while \(G_k\) is unchanged because it depends only on tangent overlaps.
For the refined exact reference path one sets
\(\delta\tau_k:=\Delta\tau_k/M_k\),
\(t_{k,\mu}^{\sharp}:=t_0+\tau_k+(\mu+\sigma_k)\delta\tau_k\),
\(H_{k,\mu}^{\sharp}:=H(t_{k,\mu}^{\sharp})\), and
\(|\psi_{k,\mu+1}^{\mathrm{ex}}\rangle:=\exp\bigl(-i\delta\tau_k H_{k,\mu}^{\sharp}\bigr)|\psi_{k,\mu}^{\mathrm{ex}}\rangle\).
Thus midpoint sampling is simply the special case \(\sigma_k=\tfrac12\).

\hypertarget{a.3-residual-geometry-and-local-inadequacy-of-the-scaffold}{%
\subsection{17A.3 Residual geometry and local inadequacy of the
scaffold}\label{a.3-residual-geometry-and-local-inadequacy-of-the-scaffold}}

This subsection measures how much of the sampled horizontal Schrödinger
drift is not captured by the current scaffold tangent plane. The main
quantities are the regularized residual \(\bar r_k\), the realized step
defect \(\epsilon_{\mathrm{step},k}\), the intrinsic projective defect
\(\epsilon_{\mathrm{proj},k}\), and the normalized miss ratio
\(\rho_{\mathrm{miss},k}\).

\[
\begin{aligned}
\bar r_k
&:=\bar T_k\dot\theta_k-\bar b_k\\
&=\sum_{j=1}^{n_k}(\dot\theta_k)_j\bigl[-i(\widetilde A_{k,j}-\langle\widetilde A_{k,j}\rangle_{\psi_k})|\psi_k\rangle\bigr]+i(H_k^{\sharp}-E_{\psi_k,k}^{\sharp})|\psi_k\rangle,\\[1mm]
\epsilon_{\mathrm{step},k}^2
&:=\|\bar r_k\|_2^2\\
&=\left\|\sum_{j=1}^{n_k}(\dot\theta_k)_j\bigl[-i(\widetilde A_{k,j}-\langle\widetilde A_{k,j}\rangle_{\psi_k})|\psi_k\rangle\bigr]+i(H_k^{\sharp}-E_{\psi_k,k}^{\sharp})|\psi_k\rangle\right\|_2^2,\\[1mm]
\epsilon_{\mathrm{proj},k}^2
&:=\min_{v\in\mathbb R^{n_k}}\|\bar T_kv-\bar b_k\|_2^2\\
&=\operatorname{Var}_{\psi_k}(H_k^{\sharp})-\sum_{i,j=1}^{n_k}(f_k)_i(G_k^{+})_{ij}(f_k)_j\\
&=\langle\psi_k|(H_k^{\sharp}-E_{\psi_k,k}^{\sharp})^2|\psi_k\rangle
-\sum_{i,j=1}^{n_k}\Re\!\Bigl(\langle\psi_k|\widetilde A_{k,i}H_k^{\sharp}|\psi_k\rangle-\langle\widetilde A_{k,i}\rangle_{\psi_k}E_{\psi_k,k}^{\sharp}\Bigr)
(G_k^{+})_{ij}\\
&\qquad\qquad\times\Re\!\Bigl(\langle\psi_k|\widetilde A_{k,j}H_k^{\sharp}|\psi_k\rangle-\langle\widetilde A_{k,j}\rangle_{\psi_k}E_{\psi_k,k}^{\sharp}\Bigr),\\[1mm]
\rho_{\mathrm{miss},k}
&:=\frac{\epsilon_{\mathrm{proj},k}^2}{\operatorname{Var}_{\psi_k}(H_k^{\sharp})+\varepsilon}\\
&=\frac{\langle\psi_k|(H_k^{\sharp}-E_{\psi_k,k}^{\sharp})^2|\psi_k\rangle
-\sum_{i,j=1}^{n_k}\Re\!\Bigl(\langle\psi_k|\widetilde A_{k,i}H_k^{\sharp}|\psi_k\rangle-\langle\widetilde A_{k,i}\rangle_{\psi_k}E_{\psi_k,k}^{\sharp}\Bigr)(G_k^{+})_{ij}\Re\!\Bigl(\langle\psi_k|\widetilde A_{k,j}H_k^{\sharp}|\psi_k\rangle-\langle\widetilde A_{k,j}\rangle_{\psi_k}E_{\psi_k,k}^{\sharp}\Bigr)}{\langle\psi_k|(H_k^{\sharp}-E_{\psi_k,k}^{\sharp})^2|\psi_k\rangle+\varepsilon}.
\end{aligned}
\]

After this fully expanded form one may use the shorthand
\(\rho_{\mathrm{miss},k}\) alone.

\hypertarget{a.4-checkpoint-local-decision-surfaces-and-calmness-gate}{%
\subsection{17A.4 Checkpoint-local decision surfaces and calmness
gate}\label{a.4-checkpoint-local-decision-surfaces-and-calmness-gate}}

The projected differential law determines motion only while the scaffold
is frozen; it does not decide whether the scaffold itself should change.
The checkpoint law does. Here \(x\in\mathbb R^{m}\) and
\(y\in\mathbb R^{n}\) are coordinate vectors of possibly different
lengths, \(r:=\min(m,n)\) is their overlap length,
\(\dot\theta_{k-1}^{\mathrm{acc}}\) is the previously accepted
checkpoint velocity after zero padding to the current coordinate length,
and
\(\tau_{\mathrm{calm}}^{\cos},\tau_{\mathrm{calm}}^{\mathrm{rate}},\tau_{\mathrm{miss}},\tau_{\mathrm{pr}}>0\)
are the calmness and miss thresholds.

\[
\begin{aligned}
\|x-y\|_{\mathrm{ov}}^2
&:=\sum_{j=1}^{r}(x_j-y_j)^2+\sum_{j=r+1}^{m}x_j^2+\sum_{j=r+1}^{n}y_j^2,\\[1mm]
c_k^{\mathrm{dir}}
&:=\frac{\langle \dot\theta_k,\dot\theta_{k-1}^{\mathrm{acc}}\rangle}{\max\{\|\dot\theta_k\|_2\,\|\dot\theta_{k-1}^{\mathrm{acc}}\|_2,\varepsilon\}},
\qquad
\eta_k^{\mathrm{rate}}:=\frac{\|\dot\theta_k-\dot\theta_{k-1}^{\mathrm{acc}}\|_{\mathrm{ov}}}{\max\{\|\dot\theta_k\|_2,\|\dot\theta_{k-1}^{\mathrm{acc}}\|_2,\varepsilon\}},\\[1mm]
\mathfrak C_k
&:=\mathbf 1\bigl[c_k^{\mathrm{dir}}\ge \tau_{\mathrm{calm}}^{\cos}\bigr]\,\mathbf 1\bigl[\eta_k^{\mathrm{rate}}\le \tau_{\mathrm{calm}}^{\mathrm{rate}}\bigr],\\[1mm]
\mathcal U_k^{\mathrm{stay}}
&:=\{(\mathrm{stay},p,a):p\in\mathcal P_k^{\mathrm{stay}},\ a\in\mathcal A_k\},
\qquad
\mathcal U_k^{\mathrm{app}}:=\{(\mathrm{append},r,a):r\in\mathcal R_k^{\mathrm{conf}},\ a\in\mathcal A_k\},\\[1mm]
\nu_k^{\star}
&:=\begin{cases}
\displaystyle \arg\min_{u\in\mathcal U_k^{\mathrm{stay}}\cup\mathcal U_k^{\mathrm{app}}}\mathcal J_k^{\mathrm{forecast}}(u),&\rho_{\mathrm{miss},k}>\tau_{\mathrm{miss}},\\[2mm]
\text{best admissible prune of §17A.9 or stay},&\rho_{\mathrm{miss},k}\le \tau_{\mathrm{pr}}\text{ and }\mathfrak C_k=1,\\[1mm]
\mathrm{stay},&\text{otherwise.}
\end{cases}
\end{aligned}
\]

Thus the miss-dominant regime compares stay and append through the
forecast score, the low-miss calm regime opens prune, and all other
cases reduce to stay.

\hypertarget{a.5-position-aware-append-geometry-and-confirm-grade-gain}{%
\subsection{17A.5 Position-aware append geometry and confirm-grade
gain}\label{a.5-position-aware-append-geometry-and-confirm-grade-gain}}

An append move is a record \(r=(m,p)\) rather than a bare generator
because the new block is inserted at a definite ordered position. Let
\(m_r\) denote the number of new coordinates carried by the appended
block, let \(\eta=(\eta_1,\dots,\eta_{m_r})\) denote those new
coordinates, let \(\vartheta=(\theta,\eta)\) denote the augmented
coordinate vector, and let \(\eta=0\) at the append instant. The
physical state therefore does not jump at the checkpoint, but the
tangent plane acquires new columns. The append gain comes from solving
the augmented McLachlan normal system and eliminating the old
coordinates exactly.

\[
\begin{aligned}
\bar U_{r,k}
&:=\bigl[(\bar U_{r,k})_{sa}\bigr]_{s=1,\dots,d}^{\,a=1,\dots,m_r}
 =\bigl[(Q_{\psi_k}\partial_{\eta_a}|\psi(\theta_k,\eta;\mathcal O_k^{+}(r))\rangle\vert_{\eta=0})_s\bigr]_{s=1,\dots,d}^{\,a=1,\dots,m_r}\\
&=\left[
\begin{array}{ccc}
(\bar u_{r,1})_1 & \cdots & (\bar u_{r,m_r})_1\\
\vdots & \ddots & \vdots\\
(\bar u_{r,1})_d & \cdots & (\bar u_{r,m_r})_d
\end{array}
\right],\\[1mm]
B_{r,k}
&:=\bigl[(B_{r,k})_{ia}\bigr]_{i=1,\dots,n_k}^{\,a=1,\dots,m_r}
 =\Re(\bar T_k^{\dagger}\bar U_{r,k})
 =\left[\Re\!\left(\sum_{s=1}^{d}(\bar T_k)_{si}^{*}(\bar U_{r,k})_{sa}\right)\right]_{i=1,\dots,n_k}^{\,a=1,\dots,m_r},\\[1mm]
C_{r,k}
&:=\bigl[(C_{r,k})_{ab}\bigr]_{a,b=1}^{m_r}
 =\Re(\bar U_{r,k}^{\dagger}\bar U_{r,k})
 =\left[\Re\!\left(\sum_{s=1}^{d}(\bar U_{r,k})_{sa}^{*}(\bar U_{r,k})_{sb}\right)\right]_{a,b=1}^{m_r},\\[1mm]
q_{r,k}
&:=\bigl[(q_{r,k})_{a}\bigr]_{a=1}^{m_r}
 =\Re(\bar U_{r,k}^{\dagger}\bar b_k)
 =\left[\Re\!\left(\sum_{s=1}^{d}(\bar U_{r,k})_{sa}^{*}(\bar b_k)_s\right)\right]_{a=1}^{m_r},\\[1mm]
\mathfrak m_{r,k}(v,\xi)
&:=\|\bar T_kv+\bar U_{r,k}\xi-\bar b_k\|_2^2+v^{\top}\Lambda_kv+\xi^{\top}\Lambda_r\xi\\
&=\begin{bmatrix}v\\ \xi\end{bmatrix}^{\!\top}
\begin{bmatrix}K_k & B_{r,k}\\ B_{r,k}^{\top} & C_{r,k}+\Lambda_r\end{bmatrix}
\begin{bmatrix}v\\ \xi\end{bmatrix}
-2\begin{bmatrix}v\\ \xi\end{bmatrix}^{\!\top}\begin{bmatrix}f_k\\ q_{r,k}\end{bmatrix}
+\langle\psi_k|(H_k^{\sharp}-E_{\psi_k,k}^{\sharp})^2|\psi_k\rangle,\\[1mm]
\begin{bmatrix}\dot\theta_{r,k}^{\mathrm{old}}\\ \dot\eta_{r,k}\end{bmatrix}
&:=\begin{bmatrix}K_k & B_{r,k}\\ B_{r,k}^{\top} & C_{r,k}+\Lambda_r\end{bmatrix}^{+}\begin{bmatrix}f_k\\ q_{r,k}\end{bmatrix},
\qquad
v^{\star}(\xi):=K_k^{+}(f_k-B_{r,k}\xi),\\[1mm]
\mathfrak m_{r,k}(v^{\star}(\xi),\xi)
&=\langle\psi_k|(H_k^{\sharp}-E_{\psi_k,k}^{\sharp})^2|\psi_k\rangle-f_k^{\top}K_k^{+}f_k+\xi^{\top}S_{r,k}\xi-2\xi^{\top}w_{r,k},\\[1mm]
w_{r,k}
&:=q_{r,k}-B_{r,k}^{\top}K_k^{+}f_k\\
&=\left[\Re\!\left(\sum_{s=1}^{d}(\bar U_{r,k})_{sa}^{*}(\bar b_k)_s\right)-\sum_{i=1}^{n_k}(B_{r,k})_{ia}(K_k^{+}f_k)_i\right]_{a=1}^{m_r},\\[1mm]
S_{r,k}
&:=C_{r,k}+\Lambda_r-B_{r,k}^{\top}K_k^{+}B_{r,k}\\
&=\left[(C_{r,k})_{ab}+(\Lambda_r)_{ab}-\sum_{i,j=1}^{n_k}(B_{r,k})_{ia}(K_k^{+})_{ij}(B_{r,k})_{jb}\right]_{a,b=1}^{m_r},\\[1mm]
\Delta_{r,k}^{\mathrm{exact}}
&:=w_{r,k}^{\top}S_{r,k}^{+}w_{r,k}\\
&=(q_{r,k}-B_{r,k}^{\top}K_k^{+}f_k)^{\top}(C_{r,k}+\Lambda_r-B_{r,k}^{\top}K_k^{+}B_{r,k})^{+}(q_{r,k}-B_{r,k}^{\top}K_k^{+}f_k),\\[1mm]
\rho_{r,k}^{\mathrm{gain}}
&:=\frac{\Delta_{r,k}^{\mathrm{exact}}}{\langle\psi_k|(H_k^{\sharp}-E_{\psi_k,k}^{\sharp})^2|\psi_k\rangle+\varepsilon}.
\end{aligned}
\]

The confirm-grade compressed path keeps the same append object and
resolves only the old-coordinate contribution on a retained spectral
support. If \(K_k=V_k\Sigma_k^2V_k^{\top}\) is the supported
factorization, \(z_k:=\Sigma_k^{-1}V_k^{\top}f_k\), and
\(\Gamma_{r,k}:=\Sigma_k^{-1}V_k^{\top}B_{r,k}\), then
\(w_{r,k}=q_{r,k}-\Gamma_{r,k}^{\top}z_k\), the compressed Schur block
on a retained mode set \(J_r\) is
\(\widetilde S_{r,k}(J_r):=C_{r,k}+\Lambda_r-\Gamma_{r,k;J_r}^{\top}\Gamma_{r,k;J_r}\),
the compressed gain is
\(\widetilde\Delta_{r,k}(J_r):=(q_{r,k}-\Gamma_{r,k}^{\top}z_k)^{\top}\widetilde S_{r,k}(J_r)^{+}(q_{r,k}-\Gamma_{r,k}^{\top}z_k)\),
and with direction penalty
\(\delta_{r,k}^{\mathrm{dir}}:=\|\dot\vartheta_{r,k}^{\mathrm{aug}}-\dot\theta_{k-1}^{\mathrm{acc}}\|_{\mathrm{ov}}\)
the confirm score is \[
\begin{aligned}
S_{r,k}^{\mathrm{conf}}
:=\begin{cases}
\dfrac{(q_{r,k}-B_{r,k}^{\top}K_k^{+}f_k)^{\top}(C_{r,k}+\Lambda_r-B_{r,k}^{\top}K_k^{+}B_{r,k})^{+}(q_{r,k}-B_{r,k}^{\top}K_k^{+}f_k)}{\langle\psi_k|(H_k^{\sharp}-E_{\psi_k,k}^{\sharp})^2|\psi_k\rangle+\varepsilon}-\lambda_{\mathrm{dir}}\,\delta_{r,k}^{\mathrm{dir}}-\lambda_{\mathrm{meas}}g_{r,k}^{\mathrm{new}},&\text{exact confirm},\\[3mm]
\dfrac{(q_{r,k}-\Gamma_{r,k}^{\top}z_k)^{\top}\widetilde S_{r,k}(J_r)^{+}(q_{r,k}-\Gamma_{r,k}^{\top}z_k)}{\langle\psi_k|(H_k^{\sharp}-E_{\psi_k,k}^{\sharp})^2|\psi_k\rangle+\varepsilon}-\lambda_{\mathrm{dir}}\,\delta_{r,k}^{\mathrm{dir}}-\lambda_{\mathrm{meas}}g_{r,k}^{\mathrm{new}},&\text{compressed confirm}.
\end{cases}
\end{aligned}
\] Thus the confirmed append shortlist is
\(\mathcal R_k^{\mathrm{conf}}:=\{r:\rho_{r,k}^{\mathrm{gain}}\ge \tau_{\mathrm{gain}},\ S_{r,k}^{\mathrm{conf}}\ge \tau_{\mathrm{conf}}\}\).

\hypertarget{a.6-proposal-family-for-within-scaffold-motion}{%
\subsection{17A.6 Proposal family for within-scaffold
motion}\label{a.6-proposal-family-for-within-scaffold-motion}}

A stay action is a finite family of same-scaffold forecast proposals,
not a single velocity. Here \(M_k\) is the proposal metric,
\(I_k^{\mathrm{drv}}\subseteq\{1,\dots,n_k\}\) is the drive-aligned
coordinate block, \(\mathcal B_k\) is the finite blend-weight set, and
\(\mathcal A_k=\{a_{k,1},\dots,a_{k,m_k}\}\subset(0,\infty)\) is the
finite step-scale set.

\[
\begin{aligned}
M_k
&:=\begin{cases}G_k,&\text{when the tangent metric is used on its supported subspace},\\ I,&\text{otherwise},\end{cases}
\qquad
\|v\|_{M_k}:=\sqrt{v^{\top}M_kv},
\qquad
\langle x,y\rangle_{M_k}:=x^{\top}M_ky,\\[1mm]
v_k^{\mathrm{McL}}
&:=\dot\theta_k,\\[1mm]
v_{k,\mathrm{drv}}^{\mathrm{cur}}
&:=\text{the unique vector with }v_{k,\mathrm{drv}}^{\mathrm{cur}}[I_k^{\mathrm{drv}}]=K_k[I_k^{\mathrm{drv}},I_k^{\mathrm{drv}}]^{+}f_k[I_k^{\mathrm{drv}}]\text{ and }v_{k,\mathrm{drv}}^{\mathrm{cur}}[(I_k^{\mathrm{drv}})^{c}]=0,\\[1mm]
v_{k,\mathrm{drv}}^{\mathrm{next}}
&:=\text{the unique vector with }v_{k,\mathrm{drv}}^{\mathrm{next}}[I_k^{\mathrm{drv}}]=K_{k+1}^{\sharp}[I_k^{\mathrm{drv}},I_k^{\mathrm{drv}}]^{+}f_{k+1}^{\sharp}[I_k^{\mathrm{drv}}]\text{ and }v_{k,\mathrm{drv}}^{\mathrm{next}}[(I_k^{\mathrm{drv}})^{c}]=0,\\[1mm]
d_k^{\perp}
&:=d_k-\dfrac{\langle v_k^{\mathrm{McL}},d_k\rangle_{M_k}}{\|v_k^{\mathrm{McL}}\|_{M_k}^2}v_k^{\mathrm{McL}},
\qquad
\widetilde d_k^{\perp}:=\sqrt{\dfrac{\|v_k^{\mathrm{McL}}\|_{M_k}^2}{\|d_k^{\perp}\|_{M_k}^2}}\,d_k^{\perp}\ \text{when }d_k^{\perp}\neq 0,\\[1mm]
v_{k,\beta}^{\mathrm{blend}}
&:=\sqrt{\dfrac{\|v_k^{\mathrm{McL}}\|_{M_k}^2}{\|v_k^{\mathrm{McL}}+\beta\widetilde d_k^{\perp}\|_{M_k}^2}}\,(v_k^{\mathrm{McL}}+\beta\widetilde d_k^{\perp}),
\qquad \beta\in\mathcal B_k,\\[1mm]
\widehat v
&:=v/\|v\|_{M_k}\ \text{for }v\neq 0,
\qquad
\mathcal P_k^{\mathrm{stay}}\subseteq\{\widehat v_k^{\mathrm{McL}},\widehat v_{k,\mathrm{drv}}^{\mathrm{cur}},\widehat v_{k,\mathrm{drv}}^{\mathrm{next}},\widehat v_{k,\beta}^{\mathrm{blend}},\widehat v_k^{\mathrm{sec}}\},\\[1mm]
\mathcal U_k^{\mathrm{stay}}
&:=\{(\mathrm{stay},p,a):p\in\mathcal P_k^{\mathrm{stay}},\ a\in\mathcal A_k\}.
\end{aligned}
\]

Thus every stay action is already written as an explicit proposal
direction together with an explicit discrete scale.

\hypertarget{a.7-tangent-secant-forecast-proposal-and-signed-energy-lead-taper}{%
\subsection{17A.7 Tangent-secant forecast proposal and
signed-energy-lead
taper}\label{a.7-tangent-secant-forecast-proposal-and-signed-energy-lead-taper}}

The tangent-secant proposal is the tangent inverse image of the
phase-aligned exact one-step secant displacement, followed by a metric
trust clip and a signed-energy-lead taper. Here \(R_k>0\) is the trust
radius and \(\eta_E>0\) is the relative energy-lead cap.

\[
\begin{aligned}
|\widetilde\psi_{k+1}^{\mathrm{ex}}\rangle
&:=\exp\bigl(-i\arg\langle\psi_k,\psi_{k+1}^{\mathrm{ex}}\rangle\bigr)|\psi_{k+1}^{\mathrm{ex}}\rangle,\\[1mm]
\delta_k^{\mathrm{sec}}
&:=Q_{\psi_k}\bigl(|\widetilde\psi_{k+1}^{\mathrm{ex}}\rangle-|\psi_k\rangle\bigr)
 =|\widetilde\psi_{k+1}^{\mathrm{ex}}\rangle-|\psi_k\rangle-\langle\psi_k,\widetilde\psi_{k+1}^{\mathrm{ex}}-\psi_k\rangle|\psi_k\rangle,\\[1mm]
\xi_k^{\mathrm{sec}}
&:=\Re(\bar T_k^{\dagger}\delta_k^{\mathrm{sec}})
 =\left[\Re\!\left(\sum_{a=1}^{d}(\bar T_k)_{aj}^{*}(\delta_k^{\mathrm{sec}})_a\right)\right]_{j=1}^{n_k},\\[1mm]
v_k^{\mathrm{sec}}
&:=\frac{1}{\Delta\tau_k}K_k^{+}\xi_k^{\mathrm{sec}},
\qquad
v_k^{\mathrm{sec}}\leftarrow \min\!\left\{1,\frac{R_k}{\|v_k^{\mathrm{sec}}\|_{M_k}}\right\}v_k^{\mathrm{sec}},\\[1mm]
b_k^{E}
&:=E_k^{\mathrm{ctrl}}-E_k^{\mathrm{ex}},
\qquad
\Delta E_k^{\mathrm{ex}}:=E_{k+1}^{\mathrm{ex}}-E_k^{\mathrm{ex}},
\qquad
\ell_k^{E}:=\bigl[\operatorname{sgn}(\Delta E_k^{\mathrm{ex}})b_k^{E}\bigr]_{+},\\[1mm]
\chi_k^{E}
&:=\left[1-\frac{\ell_k^{E}}{\eta_E|\Delta E_k^{\mathrm{ex}}|+\varepsilon}\right]_{+},
\qquad
v_k^{\mathrm{sec}}\leftarrow \chi_k^{E}v_k^{\mathrm{sec}}.
\end{aligned}
\]

After the fully expanded chain above, one may use
\(\widehat v_k^{\mathrm{sec}}:=v_k^{\mathrm{sec}}/\|v_k^{\mathrm{sec}}\|_{M_k}\)
as shorthand when it is inserted into \(\mathcal P_k^{\mathrm{stay}}\).

\hypertarget{a.8-exact-forecast-horizon-score-and-stay-append-decision-law}{%
\subsection{17A.8 Exact-forecast horizon score and stay-append decision
law}\label{a.8-exact-forecast-horizon-score-and-stay-append-decision-law}}

This subsection defines the forecast functional used when the miss ratio
is too large to stay on the current scaffold blindly. The horizon index
is \(h\in\{1,\dots,H_k\}\), the horizon length is \(H_k\ge 1\), the
positive horizon weights are \(\omega_h\), and
\(W_k:=\sum_{h=1}^{H_k}\omega_h\). The action \(u\) is either a stay
action \(u=(\mathrm{stay},p,a)\) with
\(p\in\mathcal P_k^{\mathrm{stay}}\) and \(a\in\mathcal A_k\), or an
append action \(u=(\mathrm{append},r,a)\) with
\(r=(m,p_{\mathrm{ins}})\in\mathcal R_k^{\mathrm{conf}}\) and
\(a\in\mathcal A_k\). The nonnegative weights
\(\lambda_F,\lambda_s,\lambda_D,\lambda_n,\lambda_E,\lambda_{\mathrm{slope}},\lambda_{\mathrm{curv}},\lambda_{\mathrm{under}},\lambda_{\mathrm{over}}\)
determine the relative importance of fidelity, staggered density,
doublon, site occupations, total energy, and the optional energy-shape
and excursion penalties.

A stay forecast keeps the current scaffold and moves along a normalized
within-scaffold proposal, while an append forecast lifts to the
zero-initialized augmented scaffold and moves along the augmented append
velocity. Thus \[
\begin{aligned}
\theta_{k+h\mid k}^{(\mathrm{stay},p,a)}
&:=\theta_k+h\,a\,\Delta\tau_k\,p,
\qquad
|\psi_{k+h}^{(\mathrm{stay},p,a)}\rangle:=U\bigl(\theta_{k+h\mid k}^{(\mathrm{stay},p,a)};\mathcal O_k\bigr)|\psi_{\mathrm{ref}}\rangle,\\[1mm]
\vartheta_k^{(r)}&:=(\theta_k,0),
\qquad
\vartheta_{k+h\mid k}^{(\mathrm{append},r,a)}:=\vartheta_k^{(r)}+h\,a\,\Delta\tau_k\,\dot\vartheta_{r,k}^{\mathrm{aug}},
\qquad
|\psi_{k+h}^{(\mathrm{append},r,a)}\rangle:=U\bigl(\vartheta_{k+h\mid k}^{(\mathrm{append},r,a)};\mathcal O_k^{+}(r)\bigr)|\psi_{\mathrm{ref}}\rangle,\\[1mm]
|\psi_{k+h}^{\mathrm{ex}}\rangle&:=\text{exact driven reference state at the same future checkpoint time.}
\end{aligned}
\]

The tracked observables are the site occupations, the doublon, the
staggered density, and the total energy. If
\(|\phi\rangle=\sum_x\phi_x|x\rangle\), then
\(b_{j,\sigma}(x)\in\{0,1\}\) is the occupation bit of basis state
\(|x\rangle\) at site-spin pair \((j,\sigma)\), and \[
\begin{aligned}
\hat n_j&:=\hat n_{j,\uparrow}+\hat n_{j,\downarrow},
\qquad
n_{j,\sigma}[\phi]:=\langle\phi|\hat n_{j,\sigma}|\phi\rangle=\sum_x|\phi_x|^2b_{j,\sigma}(x),
\qquad
n_j[\phi]:=\langle\phi|\hat n_j|\phi\rangle=\sum_x|\phi_x|^2\bigl(b_{j,\uparrow}(x)+b_{j,\downarrow}(x)\bigr),\\[1mm]
\hat D&:=\sum_{j=1}^{L}\hat n_{j,\uparrow}\hat n_{j,\downarrow},
\qquad
D[\phi]:=\langle\phi|\hat D|\phi\rangle=\sum_x|\phi_x|^2\sum_{j=1}^{L}b_{j,\uparrow}(x)b_{j,\downarrow}(x),\\[1mm]
\hat S&:=\frac{1}{L}\sum_{j=1}^{L}(-1)^{j-1}\hat n_j,
\qquad
s[\phi]:=\langle\phi|\hat S|\phi\rangle=\frac{1}{L}\sum_{j=1}^{L}(-1)^{j-1}n_j[\phi],\\[1mm]
E_{k+h}^{\mathrm{ctrl}}(u)&:=\langle\psi_{k+h}^{(u)}|H_{k+h}^{\sharp}|\psi_{k+h}^{(u)}\rangle,
\qquad
E_{k+h}^{\mathrm{ex}}:=\langle\psi_{k+h}^{\mathrm{ex}}|H_{k+h}^{\sharp}|\psi_{k+h}^{\mathrm{ex}}\rangle,
\qquad
H_{k+h}^{\sharp}:=H(t_{k+h}^{\sharp}).
\end{aligned}
\]

From these primitives the first-order forecast discrepancies are defined
once by \[
\begin{aligned}
\varepsilon_{k+h}^{F}(u)
&:=1-\bigl|\langle\psi_{k+h}^{\mathrm{ex}},\psi_{k+h}^{(u)}\rangle\bigr|^2,\\[1mm]
\varepsilon_{k+h}^{s}(u)
&:=\left|\frac{1}{L}\sum_{j=1}^{L}(-1)^{j-1}\bigl(n_j[\psi_{k+h}^{(u)}]-n_j[\psi_{k+h}^{\mathrm{ex}}]\bigr)\right|,\\[1mm]
\varepsilon_{k+h}^{D}(u)
&:=\left|\sum_{j=1}^{L}\Bigl(\langle\psi_{k+h}^{(u)}|\hat n_{j,\uparrow}\hat n_{j,\downarrow}|\psi_{k+h}^{(u)}\rangle-\langle\psi_{k+h}^{\mathrm{ex}}|\hat n_{j,\uparrow}\hat n_{j,\downarrow}|\psi_{k+h}^{\mathrm{ex}}\rangle\Bigr)\right|,\\[1mm]
\varepsilon_{k+h}^{n}(u)
&:=\max_{1\le j\le L}\bigl|n_j[\psi_{k+h}^{(u)}]-n_j[\psi_{k+h}^{\mathrm{ex}}]\bigr|,\\[1mm]
\varepsilon_{k+h}^{E}(u)
&:=\bigl|\langle\psi_{k+h}^{(u)}|H_{k+h}^{\sharp}|\psi_{k+h}^{(u)}\rangle-\langle\psi_{k+h}^{\mathrm{ex}}|H_{k+h}^{\sharp}|\psi_{k+h}^{\mathrm{ex}}\rangle\bigr|.
\end{aligned}
\]

The energy-shape corrections compare first and second forward
differences of the predicted and exact energy paths, while the excursion
corrections compare the signed forecast excursion against a relative
exact band. Writing \(\varpi_h^{(1)}:=(\omega_h+\omega_{h+1})/2\),
\(\varpi_h^{(2)}:=(\omega_h+\omega_{h+1}+\omega_{h+2})/3\), and
\(\varrho_E\in[0,1)\), \[
\begin{aligned}
\Delta E_{k+h}^{\mathrm{ctrl}}(u)&:=E_{k+h+1}^{\mathrm{ctrl}}(u)-E_{k+h}^{\mathrm{ctrl}}(u),
\qquad
\Delta E_{k+h}^{\mathrm{ex}}:=E_{k+h+1}^{\mathrm{ex}}-E_{k+h}^{\mathrm{ex}},\\[1mm]
\Delta^2E_{k+h}^{\mathrm{ctrl}}(u)&:=E_{k+h+2}^{\mathrm{ctrl}}(u)-2E_{k+h+1}^{\mathrm{ctrl}}(u)+E_{k+h}^{\mathrm{ctrl}}(u),
\qquad
\Delta^2E_{k+h}^{\mathrm{ex}}:=E_{k+h+2}^{\mathrm{ex}}-2E_{k+h+1}^{\mathrm{ex}}+E_{k+h}^{\mathrm{ex}},\\[1mm]
\Xi_k^{\mathrm{shape,slope}}(u)&:=\frac{1}{\sum_h\varpi_h^{(1)}}\sum_h\varpi_h^{(1)}\bigl|\Delta E_{k+h}^{\mathrm{ctrl}}(u)-\Delta E_{k+h}^{\mathrm{ex}}\bigr|,\\[1mm]
\Xi_k^{\mathrm{shape,curv}}(u)&:=\frac{1}{\sum_h\varpi_h^{(2)}}\sum_h\varpi_h^{(2)}\bigl|\Delta^2E_{k+h}^{\mathrm{ctrl}}(u)-\Delta^2E_{k+h}^{\mathrm{ex}}\bigr|,\\[1mm]
x_{k+h}^{\mathrm{ctrl}}(u)&:=E_{k+h}^{\mathrm{ctrl}}(u)-E_k^{\mathrm{ctrl}},
\qquad
x_{k+h}^{\mathrm{ex}}:=E_{k+h}^{\mathrm{ex}}-E_k^{\mathrm{ex}},
\qquad
\Pi_{k+h}^{\mathrm{ctrl}}(u):=\operatorname{sgn}(x_{k+h}^{\mathrm{ex}})x_{k+h}^{\mathrm{ctrl}}(u),\\[1mm]
L_{k+h}^{\mathrm{ex}}&:=(1-\varrho_E)|x_{k+h}^{\mathrm{ex}}|,
\qquad
U_{k+h}^{\mathrm{ex}}:=(1+\varrho_E)|x_{k+h}^{\mathrm{ex}}|,\\[1mm]
\Xi_k^{\mathrm{exc,under}}(u)&:=\frac{1}{W_k}\sum_{h=1}^{H_k}\omega_h\bigl[L_{k+h}^{\mathrm{ex}}-\Pi_{k+h}^{\mathrm{ctrl}}(u)\bigr]_{+},
\qquad
\Xi_k^{\mathrm{exc,over}}(u):=\frac{1}{W_k}\sum_{h=1}^{H_k}\omega_h\bigl[\Pi_{k+h}^{\mathrm{ctrl}}(u)-U_{k+h}^{\mathrm{ex}}\bigr]_{+}.
\end{aligned}
\]

The final quantity of interest in the miss-dominant regime is the full
horizon score itself, so it is written first in fully expanded primitive
form: \[
\begin{aligned}
\mathcal J_k^{\mathrm{forecast}}(u)
&:=\frac{1}{W_k}\sum_{h=1}^{H_k}\omega_h\Biggl[
\lambda_F\Bigl(1-\bigl|\langle\psi_{k+h}^{\mathrm{ex}},\psi_{k+h}^{(u)}\rangle\bigr|^2\Bigr)
+\lambda_s\left|\frac{1}{L}\sum_{j=1}^{L}(-1)^{j-1}\bigl(n_j[\psi_{k+h}^{(u)}]-n_j[\psi_{k+h}^{\mathrm{ex}}]\bigr)\right|\\
&\qquad\qquad
+\lambda_D\left|\sum_{j=1}^{L}\Bigl(\langle\psi_{k+h}^{(u)}|\hat n_{j,\uparrow}\hat n_{j,\downarrow}|\psi_{k+h}^{(u)}\rangle-\langle\psi_{k+h}^{\mathrm{ex}}|\hat n_{j,\uparrow}\hat n_{j,\downarrow}|\psi_{k+h}^{\mathrm{ex}}\rangle\Bigr)\right|\\
&\qquad\qquad
+\lambda_n\max_{1\le j\le L}\bigl|n_j[\psi_{k+h}^{(u)}]-n_j[\psi_{k+h}^{\mathrm{ex}}]\bigr|
+\lambda_E\bigl|\langle\psi_{k+h}^{(u)}|H_{k+h}^{\sharp}|\psi_{k+h}^{(u)}\rangle-\langle\psi_{k+h}^{\mathrm{ex}}|H_{k+h}^{\sharp}|\psi_{k+h}^{\mathrm{ex}}\rangle\bigr|\\
&\qquad\qquad
+\lambda_{\mathrm{under}}\bigl[(1-\varrho_E)|E_{k+h}^{\mathrm{ex}}-E_k^{\mathrm{ex}}|-\operatorname{sgn}(E_{k+h}^{\mathrm{ex}}-E_k^{\mathrm{ex}})(E_{k+h}^{\mathrm{ctrl}}(u)-E_k^{\mathrm{ctrl}})\bigr]_{+}\\
&\qquad\qquad
+\lambda_{\mathrm{over}}\bigl[\operatorname{sgn}(E_{k+h}^{\mathrm{ex}}-E_k^{\mathrm{ex}})(E_{k+h}^{\mathrm{ctrl}}(u)-E_k^{\mathrm{ctrl}})-(1+\varrho_E)|E_{k+h}^{\mathrm{ex}}-E_k^{\mathrm{ex}}|\bigr]_{+}
\Biggr]\\[1mm]
&\qquad
+\lambda_{\mathrm{slope}}\frac{1}{\sum_h\varpi_h^{(1)}}\sum_h\varpi_h^{(1)}\bigl|\Delta E_{k+h}^{\mathrm{ctrl}}(u)-\Delta E_{k+h}^{\mathrm{ex}}\bigr|
+\lambda_{\mathrm{curv}}\frac{1}{\sum_h\varpi_h^{(2)}}\sum_h\varpi_h^{(2)}\bigl|\Delta^2E_{k+h}^{\mathrm{ctrl}}(u)-\Delta^2E_{k+h}^{\mathrm{ex}}\bigr|.
\end{aligned}
\] For shorthand only, one may now write
\(\mathcal J_k^{\mathrm{forecast}}=\mathcal J_k^{\mathrm{trk}}+\lambda_{\mathrm{slope}}\Xi_k^{\mathrm{shape,slope}}+\lambda_{\mathrm{curv}}\Xi_k^{\mathrm{shape,curv}}+\lambda_{\mathrm{under}}\Xi_k^{\mathrm{exc,under}}+\lambda_{\mathrm{over}}\Xi_k^{\mathrm{exc,over}}\).

The stay-versus-append comparison is then defined once by \[
\begin{aligned}
\nu_{k,\mathrm{stay}}^{\star}&\in\arg\min_{u\in\mathcal U_k^{\mathrm{stay}}}\mathcal J_k^{\mathrm{forecast}}(u),
\qquad
\nu_{k,\mathrm{app}}^{\star}\in\arg\min_{u\in\mathcal U_k^{\mathrm{app}}}\mathcal J_k^{\mathrm{forecast}}(u),\\[1mm]
\Delta F_k&:=\bigl|\langle\psi_{k+1}^{\mathrm{ex}},\psi_{k+1}^{(\nu_{k,\mathrm{app}}^{\star})}\rangle\bigr|^2-\bigl|\langle\psi_{k+1}^{\mathrm{ex}},\psi_{k+1}^{(\nu_{k,\mathrm{stay}}^{\star})}\rangle\bigr|^2,
\qquad
\Delta\varepsilon_k^{E}:=\varepsilon_{k+1}^{E}(\nu_{k,\mathrm{app}}^{\star})-\varepsilon_{k+1}^{E}(\nu_{k,\mathrm{stay}}^{\star}),\\[1mm]
\nu_k^{\star}
&:=\begin{cases}
\nu_{k,\mathrm{stay}}^{\star},&1-\bigl|\langle\psi_{k+1}^{\mathrm{ex}},\psi_{k+1}^{(\nu_{k,\mathrm{stay}}^{\star})}\rangle\bigr|^2\le 10^{-3},\ \varepsilon_{k+1}^{s}(\nu_{k,\mathrm{stay}}^{\star})\le 2\times 10^{-2},\ \varepsilon_{k+1}^{D}(\nu_{k,\mathrm{stay}}^{\star})\le 2\times 10^{-3},\ \varepsilon_{k+1}^{n}(\nu_{k,\mathrm{stay}}^{\star})\le 2\times 10^{-2},\ \varepsilon_{k+1}^{E}(\nu_{k,\mathrm{stay}}^{\star})\le 2\times 10^{-3},\\[2mm]
\nu_{k,\mathrm{stay}}^{\star},&\Delta F_k<-\tau_F\text{ and }\Delta\varepsilon_k^{E}>\tau_E,\\[1mm]
\nu_{k,\mathrm{app}}^{\star},&\mathcal J_k^{\mathrm{forecast}}(\nu_{k,\mathrm{app}}^{\star})<\mathcal J_k^{\mathrm{forecast}}(\nu_{k,\mathrm{stay}}^{\star}),\\[1mm]
\nu_{k,\mathrm{stay}}^{\star},&\text{otherwise.}
\end{cases}
\end{aligned}
\]

\hypertarget{a.9-exact-local-prune-law-in-the-calm-regime}{%
\subsection{17A.9 Exact local prune law in the calm
regime}\label{a.9-exact-local-prune-law-in-the-calm-regime}}

The prune law deletes one old block only in the calm low-miss regime.
Here \(j\) is the logical block index, \(I_j\) is the runtime coordinate
block occupied by block \(j\), \(I_j^{c}\) is its complement,
\(m_{j,k}^{\mathrm{mov}}\) and \(m_{j,k}^{\mathrm{fit}}\) are the recent
normalized motion and fit-activity indicators,
\(\alpha_{\mathrm{pr}}\in[0,1]\) is the stagnation mixing weight,
\(N_{\mathrm{prot}}\) is the protection age, and
\(\tau_{\mathrm{stale}},\tau_{\mathrm{loss}}^{\mathrm{pr}},\tau_{\theta}^{\mathrm{pr}},\tau_{\psi}^{\mathrm{pr}},\tau_{\rho}^{\mathrm{pr}}>0\)
are the prune thresholds.

\[
\begin{aligned}
\sigma_{j,k}^{\mathrm{stag}}
&:=\alpha_{\mathrm{pr}}\bigl(1-m_{j,k}^{\mathrm{mov}}\bigr)+(1-\alpha_{\mathrm{pr}})\bigl(1-m_{j,k}^{\mathrm{fit}}\bigr),\\[1mm]
j\in\mathcal D_k^{\mathrm{pr}}
&\Longleftrightarrow \operatorname{age}_{j,k}\ge N_{\mathrm{prot}},\ \operatorname{cooldown}_{j,k}=0,\ \sigma_{j,k}^{\mathrm{stag}}\ge \tau_{\mathrm{stale}},\ \mathfrak C_k=1,\\[1mm]
G_{k,-j}
&:=G_k[I_j^{c},I_j^{c}],
\qquad
f_{k,-j}:=f_k[I_j^{c}],
\qquad
\dot\theta_{k,-j}^{\mathrm{proj}}:=G_{k,-j}^{+}f_{k,-j},\\[1mm]
L_{j,k}^{\mathrm{cache}}
&:=\frac{\bigl[f_k^{\top}G_k^{+}f_k-f_{k,-j}^{\top}G_{k,-j}^{+}f_{k,-j}\bigr]_{+}}{\langle\psi_k|(H_k^{\sharp}-E_{\psi_k,k}^{\sharp})^2|\psi_k\rangle+\varepsilon},\\[1mm]
|\psi_k^{(-j)}\rangle
&:=U(\theta_k^{(-j)};\mathcal O_k^{(-j)})|\psi_{\mathrm{ref}}\rangle,
\qquad
\Theta_{j,k}^{\mathrm{blk}}:=\|\theta_{k,I_j}\|_2,
\qquad
\Delta\psi_{j,k}^{\mathrm{pr}}:=\|U(\theta_k^{(-j)};\mathcal O_k^{(-j)})|\psi_{\mathrm{ref}}\rangle-U(\theta_k;\mathcal O_k)|\psi_{\mathrm{ref}}\rangle\|_2,\\[1mm]
\rho_{\mathrm{miss},k}^{(-j)}
&:=\frac{\langle\psi_k^{(-j)}|(H_k^{\sharp}-E_{\psi_k^{(-j)},k}^{\sharp})^2|\psi_k^{(-j)}\rangle-f_{k,-j}^{\top}G_{k,-j}^{+}f_{k,-j}}{\langle\psi_k^{(-j)}|(H_k^{\sharp}-E_{\psi_k^{(-j)},k}^{\sharp})^2|\psi_k^{(-j)}\rangle+\varepsilon},
\qquad
\Delta\rho_{j,k}^{\mathrm{pr}}:=\rho_{\mathrm{miss},k}^{(-j)}-\rho_{\mathrm{miss},k},\\[1mm]
\operatorname{accept\_prune}(j)
&\Longleftrightarrow L_{j,k}^{\mathrm{cache}}\le \tau_{\mathrm{loss}}^{\mathrm{pr}},\ \Theta_{j,k}^{\mathrm{blk}}\le \tau_{\theta}^{\mathrm{pr}},\ \Delta\psi_{j,k}^{\mathrm{pr}}\le \tau_{\psi}^{\mathrm{pr}},\ \Delta\rho_{j,k}^{\mathrm{pr}}\le \tau_{\rho}^{\mathrm{pr}}.
\end{aligned}
\]

Thus a block is deleted only if the cached projective loss,
deleted-block norm, exact reduced-state jump, and reduced miss-ratio
increase are all simultaneously small.

\hypertarget{a.10-canonical-piecewise-smooth-law}{%
\subsection{17A.10 Canonical piecewise-smooth
law}\label{a.10-canonical-piecewise-smooth-law}}

We now assemble the previous pieces into the final law. Here
\(\mathcal O(\tau)\) is the active scaffold at physical time \(\tau\),
\(\theta(\tau)\) is its coordinate vector on the open checkpoint
interval, and the checkpoint times satisfy
\(\tau_0<\tau_1<\tau_2<\cdots\). The full law is written first in
expanded form and only then summarized.

\[
\begin{aligned}
|\psi(\tau)\rangle
&:=U(\theta(\tau);\mathcal O_k)|\psi_{\mathrm{ref}}\rangle,
\qquad \tau\in[\tau_k,\tau_{k+1}),\\[1mm]
\dot\theta(\tau)
&:=\left(\left[\Re\!\Bigl(\langle\psi_k|\widetilde A_{k,i}\widetilde A_{k,j}|\psi_k\rangle-\langle\widetilde A_{k,i}\rangle_{\psi_k}\langle\widetilde A_{k,j}\rangle_{\psi_k}\Bigr)+(\Lambda_k)_{ij}\right]_{i,j=1}^{n_k}\right)^{+}
\left[\Re\!\Bigl(\langle\psi_k|\widetilde A_{k,i}H_k^{\sharp}|\psi_k\rangle-\langle\widetilde A_{k,i}\rangle_{\psi_k}E_{\psi_k,k}^{\sharp}\Bigr)\right]_{i=1}^{n_k},\\[1mm]
\nu_k^{\star}
&:=\begin{cases}
\displaystyle \arg\min_{u\in\mathcal U_k^{\mathrm{stay}}\cup\mathcal U_k^{\mathrm{app}}}\mathcal J_k^{\mathrm{forecast}}(u),&\rho_{\mathrm{miss},k}>\tau_{\mathrm{miss}},\\[2mm]
\text{the best }j\in\mathcal D_k^{\mathrm{pr}}\text{ satisfying }\operatorname{accept\_prune}(j),&\rho_{\mathrm{miss},k}\le \tau_{\mathrm{pr}}\text{ and }\mathfrak C_k=1,\\[1mm]
\mathrm{stay},&\text{otherwise,}
\end{cases}\\[1mm]
(\mathcal O_{k+1},\theta_{k+1})
&:=\begin{cases}
(\mathcal O_k,\theta_{k+1}^{\mathrm{stay}}),&\nu_k^{\star}=\nu_{k,\mathrm{stay}}^{\star},\\[1mm]
(\mathcal O_k^{+}(r_k^{\star}),\vartheta_{k+1}^{\mathrm{app}}),&\nu_k^{\star}=\nu_{k,\mathrm{app}}^{\star}=(\mathrm{append},r_k^{\star},a_k^{\star}),\\[1mm]
(\mathcal O_k^{(-j_k^{\star})},\theta_{k+1}^{(-j_k^{\star})}),&\nu_k^{\star}=\mathrm{prune}(j_k^{\star}).
\end{cases}
\end{aligned}
\]

In condensed words only after the expanded law: project the sampled
Schrödinger drift onto the current projective tangent space on each
checkpoint interval; when the intrinsic miss ratio is large, choose
between stay and append by the fully expanded horizon score; when the
miss ratio is small and the motion is calm, choose deletion only by the
fully expanded exact local prune law; then restart the next interval on
the resulting scaffold.

\hypertarget{b.-projective-mclachlan-real-time-dynamics-and-checkpoint-adaptive-snake-control}{%
\section{17B. Projective McLachlan Real-Time Dynamics and
Checkpoint-Adaptive Snake
Control}\label{b.-projective-mclachlan-real-time-dynamics-and-checkpoint-adaptive-snake-control}}

This chapter is the real-time analogue of §11. The governing surface is
the same staged-controller picture: fixed local geometry on the current
scaffold, checkpoint-local structural search over candidate-position
records \(r=(m,p)\), scout/confirm/commit screening for append actions,
and checkpoint-local scaffold pruning to keep the tangent plane small.
The intended snake doctrine is: \[
\text{McLachlan}=\text{continuous within-scaffold motion law},
\qquad
\text{controller}=\text{discrete append/prune law at checkpoints}.
\] Thus the scaffold is piecewise constant in physical time, and may
change only at checkpoint events. Ordinary McLachlan dynamics is a local
projection law on a fixed scaffold: at the current state it asks for the
parameter-velocity vector whose tangent combination best matches the
physical Schr"odinger drift after removal of the physically irrelevant
global-phase direction. Equivalently, \[
\dot\theta_\lambda\in\arg\min_{v\in\mathbb R^{n_\theta}}
\Bigl\{
\|\bar T v-\bar b\|^2+v^\top\Lambda v
\Bigr\}.
\] Its practical outputs are the tangent metric \(\bar G\), the force
vector \(\bar f\), and the projected residual \(\bar r\); this residual
is the correct local inadequacy signal for the current scaffold, and its
natural normalization scale is the energy variance
\(\operatorname{Var}_\psi(H)\). Therefore McLachlan is treated here as
the continuous within-scaffold motion oracle, while the checkpoint
controller alone decides when the scaffold should expand or contract.
The live controller-defining expressions are \[
\rho_{\mathrm{miss},k},
\qquad
\Gamma_k^{\mathrm{live}},
\qquad
\Sigma_{\mathrm{scout}}(r;k),
\qquad
\Delta_{\mathrm{McL}}(r;\tau_k),
\qquad
\Sigma_{\mathrm{conf}}(r;k),
\qquad
P_j^{\mathrm{pr}}(k),
\] and the displays below isolate these quantities in terms of the
primitive projective objects \(Q_\psi\), \(\bar b\), \(\bar T\), \(K\),
and \(\bar r\). The older branch-level rollout and beam doctrine is
retained only as a future extension, not as the canonical live runtime
surface.

\hypertarget{b.1-fixed-scaffold-projective-geometry}{%
\subsection{17B.1 Fixed-scaffold projective
geometry}\label{b.1-fixed-scaffold-projective-geometry}}

Fix a scaffold \(\mathcal O\), state family
\(|\psi(\theta;\mathcal O)\rangle\) with
\(\theta\in\mathbb R^{n_\theta}\), and Hamiltonian \(H(\tau)\). At the
current state, let \(d:=\dim\mathcal H\), let \[
Q_\psi:=I-|\psi\rangle\langle\psi|,
\qquad
E_\psi:=\langle\psi|H|\psi\rangle,
\] and define the projected Schr"odinger drift vector by \[
\bar b
:=
Q_\psi(-iH\psi)
=
\left[
\begin{array}{c}
\bigl(Q_\psi(-iH\psi)\bigr)_1\\
\vdots\\
\bigl(Q_\psi(-iH\psi)\bigr)_d
\end{array}
\right]
=
\left[
\begin{array}{c}
\bigl(-i(H-E_\psi)\psi\bigr)_1\\
\vdots\\
\bigl(-i(H-E_\psi)\psi\bigr)_d
\end{array}
\right]
\in\mathbb C^d.
\] Define the projected tangent matrix by \[
\bar T
:=
\left[
\begin{array}{ccc}
\bigl(Q_\psi\,\partial_{\theta_1}|\psi(\theta;\mathcal O)\rangle\bigr)_1
&
\cdots
&
\bigl(Q_\psi\,\partial_{\theta_{n_\theta}}|\psi(\theta;\mathcal O)\rangle\bigr)_1
\\
\vdots & \ddots & \vdots\\
\bigl(Q_\psi\,\partial_{\theta_1}|\psi(\theta;\mathcal O)\rangle\bigr)_d
&
\cdots
&
\bigl(Q_\psi\,\partial_{\theta_{n_\theta}}|\psi(\theta;\mathcal O)\rangle\bigr)_d
\end{array}
\right]
\in\mathbb C^{d\times n_\theta}.
\] Define the regularizer by \[
\Lambda
:=
\left[
\begin{array}{ccc}
\Lambda_{11} & \cdots & \Lambda_{1n_\theta}\\
\vdots & \ddots & \vdots\\
\Lambda_{n_\theta 1} & \cdots & \Lambda_{n_\theta n_\theta}
\end{array}
\right]
\in\mathbb R^{n_\theta\times n_\theta},
\qquad
\Lambda^\top=\Lambda,
\qquad
\Lambda\succeq 0,
\] with simplest default \(\Lambda=\lambda_{\mathrm{ridge}}I\). Then the
real McLachlan metric, force, and regularized system are \[
\bar G:=\Re(\bar T^\dagger\bar T)
=
\left[
\begin{array}{ccc}
\Re\!\bigl(\sum_{a=1}^{d}(\bar T)_{a1}^*(\bar T)_{a1}\bigr)
&
\cdots
&
\Re\!\bigl(\sum_{a=1}^{d}(\bar T)_{a1}^*(\bar T)_{a n_\theta}\bigr)
\\
\vdots & \ddots & \vdots\\
\Re\!\bigl(\sum_{a=1}^{d}(\bar T)_{a n_\theta}^*(\bar T)_{a1}\bigr)
&
\cdots
&
\Re\!\bigl(\sum_{a=1}^{d}(\bar T)_{a n_\theta}^*(\bar T)_{a n_\theta}\bigr)
\end{array}
\right],
\] \[
\bar f
:=
\Re(\bar T^\dagger\bar b)
=
\left[
\begin{array}{c}
\Re\!\bigl(\sum_{a=1}^{d}(\bar T)_{a1}^*\,\bar b_a\bigr)\\
\vdots\\
\Re\!\bigl(\sum_{a=1}^{d}(\bar T)_{a n_\theta}^*\,\bar b_a\bigr)
\end{array}
\right],
\qquad
K:=\bar G+\Lambda,
\qquad
K\dot\theta_\lambda=\bar f.
\] Equivalently, on the fixed scaffold \(\mathcal O\), ordinary
McLachlan dynamics solves \[
\dot\theta_\lambda\in\arg\min_{v\in\mathbb R^{n_\theta}}
\Bigl\{
\|\bar T v-\bar b\|^2+v^\top\Lambda v
\Bigr\},
\] so the only role of the current scaffold is to define the available
tangent plane. The solve itself is therefore a within-scaffold motion
law, not a scaffold-construction law. On a fixed scaffold
\(\mathcal O_k\), the runtime first solves
\(K\dot\theta_\lambda=\bar f\) and propagates \(\theta(\tau)\) over the
next checkpoint segment; it then measures \(\rho_{\mathrm{miss},k}\),
and only if that miss is too large does the discrete checkpoint
controller open the structural action set
\(\{\mathrm{stay},\mathrm{append}(r),\mathrm{prune}(j)\}\). The baseline
miss telemetry is \[
\epsilon_{\mathrm{proj}}^2:=\|\bar b\|^2-\bar f^\top\bar G^+\bar f,
\qquad
\epsilon_{\mathrm{step}}^2:=\|\bar T\dot\theta_\lambda-\bar b\|^2,
\qquad
\rho_{\mathrm{miss}}:=\frac{\epsilon_{\mathrm{proj}}^2}{\|\bar b\|^2+\varepsilon}
=\frac{\epsilon_{\mathrm{proj}}^2}{\operatorname{Var}_\psi(H)+\varepsilon}.
\] So the miss statistic itself may be read as \[
\begin{aligned}
\epsilon_{\mathrm{proj}}^2
&=
\|\bar b\|^2-\bar f^\top\bar G^+\bar f\\
&=
\|Q_\psi(-iH\psi)\|^2
-
\bigl(\Re(\bar T^\dagger\bar b)\bigr)^\top
\bigl(\Re(\bar T^\dagger\bar T)\bigr)^+
\Re(\bar T^\dagger\bar b),\\
\rho_{\mathrm{miss}}
&=
\frac{\epsilon_{\mathrm{proj}}^2}{\operatorname{Var}_\psi(H)+\varepsilon}.
\end{aligned}
\] For any admissible checkpoint-local tangent baseline \(A\) with basis
matrix \(\bar T_{k,A}\) (for example a coordinate subset of the current
scaffold tangents, or any chosen real tangent subspace with an explicit
basis), define \[
\bar G_{k,A}:=\Re(\bar T_{k,A}^\dagger \bar T_{k,A}),
\qquad
\bar f_{k,A}:=\Re(\bar T_{k,A}^\dagger \bar b_k),
\] and the intrinsic reduced-space projective miss \[
\epsilon_{\mathrm{proj},k}^2(A)
:=
\|\bar b_k\|^2-\bar f_{k,A}^\top \bar G_{k,A}^+\bar f_{k,A},
\qquad
\rho_{\mathrm{miss},k}(A)
:=
\frac{\epsilon_{\mathrm{proj},k}^2(A)}{\|\bar b_k\|^2+\varepsilon}.
\] This quantity is exact for the reduced tangent space
\(\operatorname{span}_{\mathbb R}(\bar T_{k,A})\). If \(A\subseteq A'\)
as tangent subspaces, then \[
\epsilon_{\mathrm{proj},k}^2(A')\le \epsilon_{\mathrm{proj},k}^2(A),
\qquad
\rho_{\mathrm{miss},k}(A')\le \rho_{\mathrm{miss},k}(A),
\] because enlarging the tangent baseline can only decrease the
orthogonal projection distance of \(\bar b_k\). In particular, if
\(A_k^{\mathrm{sc}}\subseteq A_k^{\mathrm{cm}}\), then \[
\rho_{\mathrm{miss},k}(A_k^{\mathrm{cm}})
\le
\rho_{\mathrm{miss},k}(A_k^{\mathrm{sc}}),
\] so the scout-baseline miss is conservative for the larger shared
commit baseline. The projective and raw residuals are \[
\bar r:=\bar T\dot\theta_\lambda-\bar b,
\qquad
r_{\mathrm{raw}}:=T\dot\theta_\lambda+iH\psi,
\qquad
\bar r=Q_\psi r_{\mathrm{raw}}.
\] Thus \(\bar r\) is the local defect after the best available tangent
velocity has been chosen on the current scaffold, and
\(\rho_{\mathrm{miss}}\) is the corresponding local inadequacy ratio
relative to exact projective drift. So, for
\(\tau\in[\tau_k,\tau_{k+1})\), the live controller holds
\(\mathcal O(\tau)=\mathcal O_k\) fixed and uses McLachlan only to
evolve \(\theta(\tau)\) on that scaffold.

\hypertarget{b.2-checkpoint-miss-test-and-telemetry-driven-controller-outputs}{%
\subsection{17B.2 Checkpoint miss test and telemetry-driven controller
outputs}\label{b.2-checkpoint-miss-test-and-telemetry-driven-controller-outputs}}

At checkpoint \(k\), with physical time \(\tau_k\), the controller
baseline summary is \[
\Xi_k^{\mathrm{base}}
:=
\bigl(
E_k,\,
\operatorname{Var}_k,\,
\epsilon_{\mathrm{proj},k}^2,\,
\epsilon_{\mathrm{step},k}^2,\,
\rho_{\mathrm{miss},k},\,
\|\dot\theta_{\lambda,k}\|_2,\,
\operatorname{rank}(K_k),\,
\operatorname{cond}(K_k)
\bigr).
\] Let \[
A_k^{\mathrm{sc}}\subseteq A_k^{\mathrm{cm}}\subseteq\{1,\dots,n_{\theta,k}\}
\] denote the scout and shared commit baselines at checkpoint \(k\).
Unless decorated otherwise, the live checkpoint shorthand is \[
\epsilon_{\mathrm{proj},k}^2:=\epsilon_{\mathrm{proj},k}^2(A_k^{\mathrm{sc}}),
\qquad
\rho_{\mathrm{miss},k}:=\rho_{\mathrm{miss},k}(A_k^{\mathrm{sc}}).
\] The clean rule is \[
\text{stay on }\mathcal O_k\text{ if }\rho_{\mathrm{miss},k}\le \tau_{\mathrm{miss}},
\qquad
\text{open structural search otherwise.}
\] The \texttt{fullclone\_3} runtime enriches this baseline with
checkpoint cadence and telemetry-driven controller outputs. If
\(\dot\theta_k,\dot\theta_{k-1},\dot\theta_{k-2}\) denote the aligned
recent solved steps, define \[
\delta_k:=\dot\theta_k-\dot\theta_{k-1},
\qquad
\alpha_k:=\delta_k-\delta_{k-1},
\qquad
\rho_k^{\mathrm{rate}}
:=
\frac{\|\delta_k\|_2}{\max(\|\dot\theta_k\|_2,\|\dot\theta_{k-1}\|_2,\varepsilon)},
\] \[
c_k^{\mathrm{dir}}
:=
\cos(\dot\theta_k,\dot\theta_{k-1}),
\qquad
c_k^{\mathrm{curv}}
:=
\cos(\alpha_k,\alpha_{k-1}),
\qquad
d_k^{\mathrm{pred}}:=\text{predicted displacement at checkpoint }k.
\] Then the controller uses the indicators \[
\mathbf 1_k^{\mathrm{rev}}
:=
\mathbf 1\!\bigl[
c_k^{\mathrm{dir}}\le \tau_{\mathrm{rev}}
\ \wedge\
\rho_k^{\mathrm{rate}}\ge \tau_{\mathrm{calm}}^{\mathrm{rate}}
\ \wedge\
\|\delta_k\|_2\ge \tau_{\mathrm{acc}}
\bigr],
\] \[
\mathbf 1_k^{\mathrm{flip}}
:=
\mathbf 1\!\bigl[
c_k^{\mathrm{curv}}\le \tau_{\mathrm{flip}}
\ \wedge\
\|\alpha_k\|_2\ge \tau_{\mathrm{acc}}
\ \wedge\
\|\alpha_{k-1}\|_2\ge \tau_{\mathrm{acc}}
\bigr],
\] \[
\mathbf 1_k^{\mathrm{kink}}
:=
\mathbf 1\!\bigl[
\mathbf 1_k^{\mathrm{rev}}=1
\ \vee\
\mathbf 1_k^{\mathrm{flip}}=1
\ \vee\
\bigl(\|\delta_k\|_2\ge \tau_{\mathrm{acc}}
\ \wedge\
\rho_k^{\mathrm{rate}}\ge \tau_{\mathrm{kink}}^{\mathrm{rate}}\bigr)
\bigr],
\] \[
\mathbf 1_k^{\mathrm{calm}}
:=
\mathbf 1\!\bigl[
k>0
\ \wedge\
c_k^{\mathrm{dir}}\ge \tau_{\mathrm{calm}}^{\mathrm{dir}}
\ \wedge\
\rho_k^{\mathrm{rate}}\le \tau_{\mathrm{calm}}^{\mathrm{rate}}
\ \wedge\
\mathbf 1_k^{\mathrm{rev}}=0
\ \wedge\
\mathbf 1_k^{\mathrm{flip}}=0
\ \wedge\
d_k^{\mathrm{pred}}\le 0.05
\bigr].
\] If \(N_{\mathrm{short}}^0\) and \(f_{\mathrm{short}}^0\) are the base
shortlist size and fraction, then \[
N_{\mathrm{short}}(k)=
\begin{cases}
\max\!\bigl(1,N_{\mathrm{short}}^0+b_{\mathrm{kink}}^{\mathrm{short}}\bigr),&\mathbf 1_k^{\mathrm{kink}}=1,\\
\max\!\bigl(1,\lceil s_{\mathrm{calm}}^{\mathrm{short}}N_{\mathrm{short}}^0\rceil\bigr),&\mathbf 1_k^{\mathrm{calm}}=1,\\
N_{\mathrm{short}}^0,&\text{otherwise,}
\end{cases}
\] \[
f_{\mathrm{short}}(k)=
\begin{cases}
\min\!\bigl(1,1.5\,f_{\mathrm{short}}^0\bigr),&\mathbf 1_k^{\mathrm{kink}}=1,\\
\max\!\bigl(0.05,s_{\mathrm{calm}}^{\mathrm{short}}f_{\mathrm{short}}^0\bigr),&\mathbf 1_k^{\mathrm{calm}}=1,\\
f_{\mathrm{short}}^0,&\text{otherwise,}
\end{cases}
\] while the effective refresh pressure is \[
\pi_k:=\max_{\mathrm{ord}}\!\bigl(\pi_k^{\mathrm{base}},\pi_k^{\mathrm{tele}}\bigr),
\qquad
\pi_k^{\mathrm{tele}}
=
\begin{cases}
\mathrm{medium},&k=0,\\
\mathrm{high},&\mathbf 1_k^{\mathrm{kink}}=1,\\
\mathrm{low},&\text{otherwise.}
\end{cases}
\] The oracle-budget multiplier is then \[
\beta_{\mathrm{orc}}(k)=
\begin{cases}
\max\!\bigl(1,s_{\mathrm{kink}}^{\mathrm{orc}}\bigr),&\mathbf 1_k^{\mathrm{kink}}=1\ \vee\ \pi_k=\mathrm{high},\\
\max\!\bigl(0.25,s_{\mathrm{calm}}^{\mathrm{orc}}\bigr),&\mathbf 1_k^{\mathrm{calm}}=1\ \wedge\ \pi_k=\mathrm{low},\\
1,&\text{otherwise.}
\end{cases}
\] \#\# 17B.3 Live checkpoint controller: stay, append, and prune

At each checkpoint, the discrete structural controller applies \[
\Gamma_k^{\mathrm{live}}:
(\mathcal O_k,\theta_k^-,\Xi_k^{\mathrm{base}})
\longmapsto
(\mathcal O_{k+1},\theta_{k+1}^+),
\] with action set \[
\mathcal U_k^{\mathrm{live}}
:=
\{\mathrm{stay}\}
\cup
\{\mathrm{append}(r):r\in\mathcal R_k^{\mathrm{probe}}\}
\cup
\{\mathrm{prune}(j):j\in\mathcal D_k^{\mathrm{prune}}\}.
\] Thus the live controller separates the continuous within-scaffold
McLachlan motion from the discrete checkpoint-local scaffold edits. The
append surface is the candidate-position record set from the current
checkpoint, while the prune surface is the mature scaffold-coordinate
set selected by the local prune controller. Only one discrete action is
committed per checkpoint. The controller is lane-ordered rather than
symmetric: \[
\operatorname{lane}_k
:=
\begin{cases}
\mathrm{append},&\rho_{\mathrm{miss},k}>\tau_{\mathrm{miss}},\\[2mm]
\mathrm{prune},&\rho_{\mathrm{miss},k}\le \tau_{\mathrm{miss}}\ \wedge\ \Gamma_k^{\mathrm{prune,perm}}=1,\\[2mm]
\mathrm{stay},&\text{otherwise.}
\end{cases}
\] Hence a large miss opens the append lane, pruning is visible only in
the calm low-miss regime, and failed append confirmation falls back to
\(\mathrm{stay}\), not to prune at the same checkpoint.

\hypertarget{b.4-append-broad-screen-over-candidate-position-records}{%
\subsection{17B.4 Append broad screen over candidate-position
records}\label{b.4-append-broad-screen-over-candidate-position-records}}

Candidate append is position-aware: \[
r=(m,p),
\qquad
\mathcal O_r^+:=\mathcal O_k\oplus_p m,
\qquad
\theta_r^+:=\theta_k\oplus_p 0.
\] The probed positions are restricted to a capped active set built from
a fixed active suffix window, \[
W_k^{\mathrm{act}}
:=
\bigl\{\max(0,n_{\theta,k}-\omega_{\mathrm{act}}),\dots,n_{\theta,k}-1\bigr\},
\] \[
\mathcal P_k
:=
\operatorname{AllowedPositions}\!\bigl(
n_{\theta,k},\,
\mathrm{append}=n_{\theta,k},\,
W_k^{\mathrm{act}},\,
N_{\mathrm{probe}}^{\max}
\bigr)
\subseteq
\{0,\dots,n_{\theta,k}\},
\qquad
|\mathcal P_k|\le N_{\mathrm{probe}}^{\max}.
\] At zero initialization, \[
|\psi_r^+(\theta_k,0)\rangle=|\psi_k\rangle,
\qquad
Q_{\psi_r^+}=Q_{\psi_k},
\] and the old tangent columns transport exactly. If \(r=(m,p)\) inserts
a \(q_r\)-parameter candidate block at position \(p\), define the
old-index transport \[
I_{r,p}(j)
:=
\begin{cases}
j,& j\le p,\\
j+q_r,& j>p.
\end{cases}
\] Then \[
\bar\tau_{k,r,I_{r,p}(j)}^+=\bar\tau_{k,j},
\] hence for every old tangent baseline \(A\), \[
\bar G_{k,I_{r,p}(A),I_{r,p}(A)}^+
=
\bar G_{k,A},
\qquad
\bar f_{k,I_{r,p}(A)}^+
=
\bar f_{k,A},
\qquad
K_{k,I_{r,p}(A),I_{r,p}(A)}^+
=
K_{k,A}.
\] Thus zero-init append transports the full old-old geometry exactly
and requires no remeasurement of old-old blocks.

At checkpoint \(k\), let \(x_k\in\mathbb R^{M_k}\) collect the
measurement bank already available on the current scaffold. Then every
current old-old block is a fixed linear image of \(x_k\): \[
\operatorname{vec}(\bar G_{k,A})=L_{G,A}x_k,
\qquad
\bar f_{k,A}=L_{f,A}x_k,
\qquad
\|\bar b_k\|^2=\ell_{b,k}^\top x_k.
\] For a candidate \(r\), let \[
x_{k,r}^{\mathrm{aug}}
:=
\begin{bmatrix}
x_k\\
y_{k,r}
\end{bmatrix}
\] augment the bank only by those candidate-local observables not
already present. Then scout, confirm, and commit are nested linear
images of this common candidate-local augmentation bank: \[
w_{r,k}^{\mathrm{cm}}=L_{w,r}^{\mathrm{sc}}x_{k,r}^{\mathrm{aug}},
\qquad
\operatorname{vec}(B_{r,k}^{\mathrm{cm}})=L_{B,r}^{\mathrm{cm}}x_{k,r}^{\mathrm{aug}},
\qquad
\operatorname{vec}(C_{r,k})=L_{C,r}x_{k,r}^{\mathrm{aug}}.
\] Scout therefore needs only \(w_{r,k}^{\mathrm{cm}}\) and \(C_{r,k}\);
confirm requests only selected \(K_k^{\mathrm{cm}}\)-whitened rows of
\(B_{r,k}^{\mathrm{cm}}\); commit may materialize the full old-new
block.

Now fix the shared commit baseline \(A_k^{\mathrm{cm}}\). Write \[
\bar T_k^{\mathrm{cm}}:=\bar T_{k,A_k^{\mathrm{cm}}},
\qquad
\bar G_k^{\mathrm{cm}}:=\bar G_{k,A_k^{\mathrm{cm}}},
\qquad
\bar f_k^{\mathrm{cm}}:=\bar f_{k,A_k^{\mathrm{cm}}},
\qquad
K_k^{\mathrm{cm}}:=\bar G_k^{\mathrm{cm}}+\Lambda_k^{\mathrm{cm}},
\] and let \[
K_k^{\mathrm{cm}}\dot\theta_{k,\lambda}^{\mathrm{cm}}=\bar f_k^{\mathrm{cm}},
\qquad
\bar r_k^{\mathrm{cm}}:=\bar T_k^{\mathrm{cm}}\dot\theta_{k,\lambda}^{\mathrm{cm}}-\bar b_k.
\] The tangent space still gains a new column \[
u_r:=\left.\partial_{\vartheta_r}|\psi_r^+(\theta_k,\vartheta_r)\rangle\right|_{\vartheta_r=0},
\qquad
\bar u_r:=Q_\psi u_r.
\] For the candidate block \(\bar U_r\in\mathbb C^{d\times q_r}\),
define \[
B_{r,k}^{\mathrm{cm}}
:=
\Re\!\bigl((\bar T_k^{\mathrm{cm}})^\dagger \bar U_r\bigr),
\qquad
C_{r,k}
:=
\Re(\bar U_r^\dagger \bar U_r),
\qquad
q_{r,k}
:=
\Re(\bar U_r^\dagger \bar b_k),
\] and the residualized force on the shared commit baseline \[
w_{r,k}^{\mathrm{cm}}
:=
q_{r,k}-(B_{r,k}^{\mathrm{cm}})^\top (K_k^{\mathrm{cm}})^{-1}\bar f_k^{\mathrm{cm}}
=
-\Re(\bar U_r^\dagger \bar r_k^{\mathrm{cm}}).
\] The scout-safe lower-bound Schur gain is \[
\underline\Delta_{\mathrm{McL}}(r;\tau_k)
:=
(w_{r,k}^{\mathrm{cm}})^\top (C_{r,k}+\Lambda_r)^+ w_{r,k}^{\mathrm{cm}}
=
\max_{\eta\in\mathbb R^{q_r}}
\left\{
2\eta^\top w_{r,k}^{\mathrm{cm}}
-\eta^\top(C_{r,k}+\Lambda_r)\eta
\right\}.
\] This is a conservative lower bound to the exact commit-grade gain
because \[
S_{r,k}^{\mathrm{cm}}
:=
C_{r,k}+\Lambda_r-(B_{r,k}^{\mathrm{cm}})^\top (K_k^{\mathrm{cm}})^{-1} B_{r,k}^{\mathrm{cm}}
\preceq
C_{r,k}+\Lambda_r.
\] The normalized scout score is therefore \[
\underline\rho_{\mathrm{gain}}(r;k)
:=
\frac{\underline\Delta_{\mathrm{McL}}(r;\tau_k)}{\|\bar b_k\|^2+\varepsilon},
\] \[
\Sigma_{\mathrm{scout}}(r;k)
:=
\underline\rho_{\mathrm{gain}}(r;k)
+b_{\mathrm{temp}}(r;k)
-\beta_{\mathrm{comp}}\,C_{\mathrm{comp}}(r;k)
-\beta_{\mathrm{meas}}\,G_{\mathrm{new}}(r;k)
-\beta_{\mathrm{jump}}\,J_{\mathrm{jump}}(r;k).
\] Hence scout is already written on the same McLachlan geometry as
commit, but uses only the residualized force and the candidate
self-block rather than the full old-new Schur data.

\hypertarget{b.5-append-confirm-compressed-rerank-and-exact-commit}{%
\subsection{17B.5 Append confirm, compressed rerank, and exact
commit}\label{b.5-append-confirm-compressed-rerank-and-exact-commit}}

The runtime remains organized in three tiers, \[
\mathrm{scout},
\qquad
\mathrm{confirm},
\qquad
\mathrm{commit},
\] but confirm now enriches the same candidate on the shared commit
baseline by a monotone family of \(K_k^{\mathrm{cm}}\)-whitened old-mode
responses.

Let \[
K_k^{\mathrm{cm}}=V_k\Sigma_k^2V_k^\top
\] be the spectral factorization of the shared commit baseline on its
numerical support, and define the corresponding
\(K_k^{\mathrm{cm}}\)-whitened coordinates \[
z_k:=\Sigma_k^{-1}V_k^\top \bar f_k^{\mathrm{cm}},
\qquad
\Gamma_{r,k}:=\Sigma_k^{-1}V_k^\top B_{r,k}^{\mathrm{cm}}.
\] Then \[
w_{r,k}^{\mathrm{cm}}
=
q_{r,k}-\Gamma_{r,k}^\top z_k,
\qquad
S_{r,k}^{\mathrm{cm}}
=
C_{r,k}+\Lambda_r-\Gamma_{r,k}^\top \Gamma_{r,k}.
\] Let \[
J_r^{(0)}\subseteq J_r^{(1)}\subseteq \cdots \subseteq J_r^{(m_{\max})}
\subseteq
\{1,\dots,\operatorname{rank}(K_k^{\mathrm{cm}})\}
\] be any nested candidate-local family of whitened baseline mode sets,
and let \(\Pi_{J_r^{(m)}}\) denote the coordinate projector onto
\(J_r^{(m)}\). A deterministic implementation may, for example, order
the modes by decreasing row norm of \(\Gamma_{r,k}\), but the ladder
below does not depend on the chosen ordering. Define the compressed
Schur family \[
\widetilde S_{r,k}^{(m)}
:=
C_{r,k}+\Lambda_r-\Gamma_{r,k}^\top \Pi_{J_r^{(m)}}\Gamma_{r,k},
\] and the corresponding confirm-grade lower-bound gains \[
\widetilde\Delta_{\mathrm{McL}}^{(m)}(r;\tau_k,J_r)
:=
(w_{r,k}^{\mathrm{cm}})^\top
(\widetilde S_{r,k}^{(m)})^+
w_{r,k}^{\mathrm{cm}}.
\] Equivalently, \[
\widetilde\Delta_{\mathrm{McL}}^{(m)}(r;\tau_k,J_r)
=
\max_{\eta\in\mathbb R^{q_r},\;\zeta\in\operatorname{range}(\Pi_{J_r^{(m)}})}
\left\{
2\eta^\top w_{r,k}^{\mathrm{cm}}
-\eta^\top(C_{r,k}+\Lambda_r)\eta
-2\zeta^\top\Gamma_{r,k}\eta
-\|\zeta\|_2^2
\right\}.
\] Thus the confirm ladder is monotone: \[
\underline\Delta_{\mathrm{McL}}(r;\tau_k)
=
\widetilde\Delta_{\mathrm{McL}}^{(0)}(r;\tau_k,J_r)
\le
\widetilde\Delta_{\mathrm{McL}}^{(1)}(r;\tau_k,J_r)
\le
\cdots
\le
\widetilde\Delta_{\mathrm{McL}}^{(m_{\max})}(r;\tau_k,J_r)
\le
\Delta_{\mathrm{McL}}(r;\tau_k,A_k^{\mathrm{cm}}),
\] where the exact commit-grade gain on the shared baseline is \[
\Delta_{\mathrm{McL}}(r;\tau_k,A_k^{\mathrm{cm}})
:=
(w_{r,k}^{\mathrm{cm}})^\top (S_{r,k}^{\mathrm{cm}})^+ w_{r,k}^{\mathrm{cm}}.
\] The normalized confirm and commit gains are \[
\widetilde\rho_{\mathrm{gain}}^{(m)}(r;k)
:=
\frac{\widetilde\Delta_{\mathrm{McL}}^{(m)}(r;\tau_k,J_r)}{\|\bar b_k\|^2+\varepsilon},
\qquad
\rho_{\mathrm{gain}}^{\mathrm{cm}}(r;k)
:=
\frac{\Delta_{\mathrm{McL}}(r;\tau_k,A_k^{\mathrm{cm}})}{\|\bar b_k\|^2+\varepsilon}.
\] Scout keeps \[
\mathcal S_k^{\mathrm{scout}}
:=
\operatorname{Top}_{N_{\mathrm{short}}(k)}
\bigl(\Sigma_{\mathrm{scout}}(r;k);\,r\in\mathcal R_k^{\mathrm{probe}}\bigr).
\] Confirm evaluates one chosen compressed level \(m_r\) for each scout
survivor and ranks with \[
\Sigma_{\mathrm{conf}}(r;k)
:=
\widetilde\rho_{\mathrm{gain}}^{(m_r)}(r;k)
-\beta_{\mathrm{meas}}\,G_{\mathrm{new}}(r;k)
-\beta_{\mathrm{dir}}\,D_{\mathrm{dir}}(r;k).
\] The confirm shortlist is therefore \[
\mathcal S_k^{\mathrm{conf}}
:=
\operatorname{Top}_{\max\{1,\lceil f_{\mathrm{short}}(k)\,|\mathcal S_k^{\mathrm{scout}}|\rceil\}}
\bigl(\Sigma_{\mathrm{conf}}(r;k);\,r\in\mathcal S_k^{\mathrm{scout}}\bigr).
\] The selected append record is \[
r_k^\star
:=
\arg\max_{r\in\mathcal S_k^{\mathrm{conf}}}\Sigma_{\mathrm{conf}}(r;k).
\] The clean commit rule is \[
\text{append }r_k^\star
\iff
\rho_{\mathrm{miss},k}(A_k^{\mathrm{cm}})>\tau_{\mathrm{miss}},
\quad
\rho_{\mathrm{gain}}^{\mathrm{cm}}(r_k^\star;k)\ge \tau_{\mathrm{gain}},
\quad
\Delta_{\mathrm{McL}}(r_k^\star;\tau_k,A_k^{\mathrm{cm}})\ge \delta_{\mathrm{append}}.
\] Otherwise the controller stays on \(\mathcal O_k\). Thus scout,
confirm, and commit all remain inside the same McLachlan geometry, but
use a staged Schur family that narrows measurement effort from
residualized force, to selected old-mode responses, to the exact
shared-baseline Schur complement.

\hypertarget{b.6-local-scaffold-pruning-and-rollback}{%
\subsection{17B.6 Local scaffold pruning and
rollback}\label{b.6-local-scaffold-pruning-and-rollback}}

Pruning acts on scaffold coordinates, not on branch objects. Let the
current scaffold be \(\mathcal O_k=(m_1,\dots,m_{n_k})\) with
\(n_k:=|\mathcal O_k|\), and let each coordinate carry the metadata
\((\iota_j,\kappa_j,c_j)\), where \(\iota_j\) is the checkpoint at which
coordinate \(j\) was appended, \(\kappa_j\) is its stored append-time
burden, and \(c_j(k)\in\{0,1,\dots,c_{\mathrm{cool}}\}\) is its
prune-cooldown counter. The hard protection set is \[
\mathcal P_{\mathrm{protect}}(k)
:=
\{\,j:\ k-\iota_j<\ell_{\mathrm{protect}}\,\},
\] so coordinates inside \(\mathcal P_{\mathrm{protect}}(k)\) are never
considered for pruning. For the recent-history stagnation window, let \[
\mathcal I_j^{\mathrm{stag}}(k)
:=
\{\,\ell:\max(\iota_j+1,k-L_{\mathrm{stag}}+1)\le \ell\le k\,\},
\qquad
\Delta\theta_{j,\ell}
:=
\theta_j(\tau_\ell^-)-\theta_j(\tau_{\ell-1}^+),
\qquad
\bar f_{j,\ell}
:=
\bigl(\bar f_\ell\bigr)_j.
\] Then the normalized recent-motion and fit-contribution summaries are
\[
m_j^{\mathrm{mov}}(k)
:=
\frac{1}{|\mathcal I_j^{\mathrm{stag}}(k)|}
\sum_{\ell\in\mathcal I_j^{\mathrm{stag}}(k)}
\frac{|\Delta\theta_{j,\ell}|}{\|\Delta\theta_\ell\|_\infty+\varepsilon_\theta},
\] \[
m_j^{\mathrm{fit}}(k)
:=
\frac{1}{|\mathcal I_j^{\mathrm{stag}}(k)|}
\sum_{\ell\in\mathcal I_j^{\mathrm{stag}}(k)}
\frac{|\Delta\theta_{j,\ell}|\,|\bar f_{j,\ell}|}{\sum_i |\Delta\theta_{i,\ell}|\,|\bar f_{i,\ell}|+\varepsilon_f}.
\] The hybrid checkpoint-local stagnation score is \[
s_j^{\mathrm{stag}}(k)
:=
\alpha_{\mathrm{stag}}
\Bigl(1-\operatorname{clip}(m_j^{\mathrm{mov}}(k),0,1)\Bigr)
+
(1-\alpha_{\mathrm{stag}})
\Bigl(1-\operatorname{clip}(m_j^{\mathrm{fit}}(k),0,1)\Bigr),
\qquad
0\le \alpha_{\mathrm{stag}}\le 1,
\] so larger values mean less recent motion and less recent contribution
to local drift tracking. For coordinates outside the hard protection
window, the optional soft early-removal penalty is \[
\Pi_j^{\mathrm{age}}(k)
:=
1+\lambda_{\mathrm{age}}^{\mathrm{pr}}
\exp\!\left(
-\frac{k-\iota_j-\ell_{\mathrm{protect}}}{\ell_{\mathrm{age}}}
\right),
\qquad
j\notin\mathcal P_{\mathrm{protect}}(k).
\] The mature eligible set is \[
\mathcal E_k^{\mathrm{prune}}
:=
\{\,j:\ j\notin\mathcal P_{\mathrm{protect}}(k),\ c_j(k)=0,\ |\mathcal I_j^{\mathrm{stag}}(k)|\ge 1,\ s_j^{\mathrm{stag}}(k)\ge s_{\mathrm{stale}}\,\}.
\] Pruning is opened only when the current scaffold is already tracking
the drift adequately. The permissibility gate is \[
\Gamma_k^{\mathrm{prune,perm}}
:=
\mathbf 1[\rho_{\mathrm{miss},k}\le \tau_{\mathrm{prune}}^{\mathrm{miss}}]\,
\mathbf 1[c_k^{\mathrm{dir}}\ge \tau_{\mathrm{calm}}^{\mathrm{dir}}]\,
\mathbf 1[\rho_k^{\mathrm{rate}}\le \tau_{\mathrm{calm}}^{\mathrm{rate}}]\,
\mathbf 1[n_k\ge 2],
\qquad
\tau_{\mathrm{prune}}^{\mathrm{miss}}\le \tau_{\mathrm{miss}}.
\] Once the current baseline geometry
\((\bar G_k^{\mathrm{cm}},\bar f_k^{\mathrm{cm}})\) is cached, broad
prune nomination is shot-free. For each
\(j\in\mathcal E_k^{\mathrm{prune}}\cap A_k^{\mathrm{cm}}\), let \[
A_{k,-j}^{\mathrm{cm}}
:=
A_k^{\mathrm{cm}}\setminus\{j\},
\qquad
\bar G_{k,-j}
:=
\bar G_{k,A_{k,-j}^{\mathrm{cm}}},
\qquad
\bar f_{k,-j}
:=
\bar f_{k,A_{k,-j}^{\mathrm{cm}}}.
\] Define the cached frozen-ablation miss increase \[
L_{j,\mathrm{frz}}^{\mathrm{cache}}(k)
:=
\left[
\frac{
\epsilon_{\mathrm{proj},k}^2(A_{k,-j}^{\mathrm{cm}})
-
\epsilon_{\mathrm{proj},k}^2(A_k^{\mathrm{cm}})
}{\|\bar b_k\|^2+\varepsilon}
\right]_+
=
\left[
\frac{
(\bar f_k^{\mathrm{cm}})^\top (\bar G_k^{\mathrm{cm}})^+\bar f_k^{\mathrm{cm}}
-
\bar f_{k,-j}^\top \bar G_{k,-j}^+\bar f_{k,-j}
}{\|\bar b_k\|^2+\varepsilon}
\right]_+.
\] This quantity depends only on cached submatrices and subvectors of
the current geometry and therefore requires no new shots. The local
prune probe set is then \[
\mathcal D_k^{\mathrm{prune}}
:=
\begin{cases}
\operatorname{Top}_{K_{\mathrm{pr}}}\!\left(-L_{j,\mathrm{frz}}^{\mathrm{cache}}(k);\,j\in\mathcal E_k^{\mathrm{prune}}\cap A_k^{\mathrm{cm}}\right),&\Gamma_k^{\mathrm{prune,perm}}=1,\\[2mm]
\varnothing,&\Gamma_k^{\mathrm{prune,perm}}=0.
\end{cases}
\] When cached losses tie numerically, ties are broken in favor of
larger stored burden \(\kappa_j\). Inside the nominated set, use the
cheap shot-free prune score \[
P_j^{\mathrm{pr}}(k)
:=
L_{j,\mathrm{frz}}^{\mathrm{cache}}(k),
\qquad
j_k^\star
:=
\arg\min_{j\in\mathcal D_k^{\mathrm{prune}}}
P_j^{\mathrm{pr}}(k).
\] Thus broad prune ranking is geometry-exact on the cached baseline and
requires no new measurements; the later remove-and-reintegrate test
remains the authoritative prune accept/reject stage. The short local
prune horizon is fixed as \[
\Delta\tau_{\mathrm{pr}}
:=
\eta_{\mathrm{pr}}(\tau_{k+1}-\tau_k),
\qquad
0<\eta_{\mathrm{pr}}\le 1.
\] The authoritative prune test removes \(j_k^\star\), resolves
McLachlan on the reduced scaffold over the short local probe segment
\([\tau_k,\tau_k+\Delta\tau_{\mathrm{pr}}]\), and compares the reduced
run against the incumbent by \[
J_{\mathrm{seg},k}
:=
\int_{\tau_k}^{\tau_k+\Delta\tau_{\mathrm{pr}}}\|\bar r(\tau)\|^2\,d\tau,
\] \[
\Delta\rho_j^{\mathrm{pr}}(k)
:=
\rho_{\mathrm{miss},k}^{(-j)}-\rho_{\mathrm{miss},k},
\qquad
\Delta J_j^{\mathrm{pr}}(k)
:=
\int_{\tau_k}^{\tau_k+\Delta\tau_{\mathrm{pr}}}
\Bigl(\|\bar r^{(-j)}(\tau)\|^2-\|\bar r(\tau)\|^2\Bigr)\,d\tau.
\] Accept the prune iff \[
\Delta\rho_{j_k^\star}^{\mathrm{pr}}(k)
\le
\Delta_{\mathrm{safe}}^{\mathrm{miss}},
\qquad
\Delta J_{j_k^\star}^{\mathrm{pr}}(k)
\le
\Delta_{\mathrm{safe}}^{\mathrm{seg}}.
\] Otherwise rollback to the pre-prune scaffold and set
\(c_{j_k^\star}(k^+)\leftarrow c_{\mathrm{cool}}\). Thus the live
checkpoint action is \[
\Gamma_k^{\mathrm{live}}\in\{\mathrm{stay},\mathrm{append}(r_k^\star),\mathrm{prune}(j_k^\star)\},
\] where append is selected by the shortlist/confirm/commit controller
above, and prune is selected by the local remove-and-reintegrate
controller above. Between checkpoints, the scaffold is fixed; at
checkpoints, only one discrete append or prune action is committed.

\hypertarget{b.7-future-branch-level-rollout-and-beam-extension}{%
\subsection{17B.7 Future branch-level rollout and beam
extension}\label{b.7-future-branch-level-rollout-and-beam-extension}}

If one later lifts the local stay/append/prune controller to a branch
controller, a natural future extension is to compare short rollout
alternatives by segment residual cost. On a fixed scaffold over a
segment \([\tau_a,\tau_b]\), the raw and projective segment costs are \[
J_{\mathrm{seg}}^{\mathrm{raw}}
:=
\int_{\tau_a}^{\tau_b}\|r_{\mathrm{raw}}(\tau)\|^2\,d\tau,
\qquad
J_{\mathrm{seg}}^{\mathrm{proj}}
:=
\int_{\tau_a}^{\tau_b}\|\bar r(\tau)\|^2\,d\tau.
\] A clean scalar branch objective is \[
J_{\mathrm{dyn}}(\mathfrak b)
:=
\sum_{k=0}^{K-1}J_{\mathrm{seg},k}
+\beta\,C_{\mathrm{hw}}(\mathfrak b)
+\gamma\,N_{\mathrm{adapt}}(\mathfrak b).
\] The associated residual ledgers may be summarized by \[
\Delta J^{\epsilon}
:=
\int \epsilon_{\mathrm{step}}^2,
\qquad
\Delta J^{\mathrm{model}}
:=
\int \epsilon_{\mathrm{proj}}^2,
\qquad
\Delta J^{\mathrm{damp}}
:=
\int \max\!\bigl(\epsilon_{\mathrm{step}}^2-\epsilon_{\mathrm{proj}}^2,0\bigr),
\] then culls rollout alternatives lexicographically after deduplicating
by checkpoint index, scaffold labels, and rounded parameters. Thus the
branch objective may be read as the integrated residual ledger plus
explicit hardware and adaptation penalties: \[
\begin{aligned}
J_{\mathrm{dyn}}(\mathfrak b)
&=
\sum_{k=0}^{K-1}J_{\mathrm{seg},k}
+\beta\,C_{\mathrm{hw}}(\mathfrak b)
+\gamma\,N_{\mathrm{adapt}}(\mathfrak b),\\
J_{\mathrm{seg},k}
&\sim
\int_{\tau_k}^{\tau_{k+1}}\|\bar r(\tau)\|^2\,d\tau,\\
\Delta J^{\epsilon}
&=
\int \|\bar T\dot\theta_\lambda-\bar b\|^2,\\
\Delta J^{\mathrm{model}}
&=
\int \bigl(\|\bar b\|^2-\bar f^\top\bar G^+\bar f\bigr),\\
\Delta J^{\mathrm{damp}}
&=
\int \max\!\bigl(\epsilon_{\mathrm{step}}^2-\epsilon_{\mathrm{proj}}^2,0\bigr).
\end{aligned}
\] This branch-level rollout and beam surface is retained only as a
future extension. It is not the canonical live runtime surface of
\texttt{fullclone\_3}. \medskip

\noindent\textbf{Implementation status.} The repo surface in
\texttt{fullclone\_3} is already more structured than the original
long-form derivation. It already implements checkpoint-local geometry
caches, three measurement tiers, active-window and probe-position caps,
position-aware zero-initialized single-parameter growth, and
compile/measurement/directional/temporal-prior controller modifiers. The
condensed doctrine is therefore:

\begin{quote}
Propagate on the current scaffold by projective McLachlan. Measure miss at each checkpoint. If the miss is large, open a position-aware append search over records $r=(m,p)$. Scout by residual overlap and hardware penalties, confirm by local Schur-complement gain, and append only if the miss and gain thresholds are passed. If the scaffold is already calm and well fit, prune stale scaffold generators by a local remove-and-reintegrate test with rollback. Treat branch-level rollout and beam competition only as a future extension.
\end{quote}

\hypertarget{c.-projective-mclachlan-real-time-geometry-and-piecewise-beam-adaptive-control}{%
\section{17C. Projective McLachlan Real-Time Geometry and Piecewise
Beam-Adaptive
Control}\label{c.-projective-mclachlan-real-time-geometry-and-piecewise-beam-adaptive-control}}

Below is the cleanest mathematical story for what the projected
real-time implementation is doing. It is most naturally written as a
\textbf{piecewise-smooth control problem on a family of variational manifolds},
with \(\hbar=1\), normalized pure states, and real parameters. The two
geometric facts underneath everything are:

\begin{enumerate}
\item the McLachlan principle projects the exact Schrödinger velocity onto a \textbf{real-linear} tangent space using the real inner product $g(u,v)=\Re\langle u,v\rangle$, and
\item physical pure states are rays, so the natural state space is \textbf{projective Hilbert space}, which means global-phase directions should be removed from the local scoring geometry.

\end{enumerate}

\hypertarget{c.1-what-problem-is-being-solved}{%
\subsection{17C.1 What problem is being
solved?}\label{c.1-what-problem-is-being-solved}}

Fix a Hilbert space \(\mathcal H\cong\mathbb C^d\), a Hamiltonian
\(H(\tau)\), and a current ansatz scaffold \(\mathcal O\). That scaffold
defines a parameterized family of states \[
|\psi(\theta;\mathcal O)\rangle,
\qquad
\theta\in\mathbb R^{n_\theta}.
\] At physical time \(\tau\), the exact Schrödinger equation asks for
the state velocity \[
\frac{d}{d\tau}|\psi_{\mathrm{ex}}(\tau)\rangle
=
-iH(\tau)|\psi_{\mathrm{ex}}(\tau)\rangle.
\] The local projected-dynamics question is therefore:

\begin{quote}
\textbf{At the current state, how closely can the current ansatz manifold reproduce the exact instantaneous quantum velocity?}
\end{quote}

This is the local question answered by real-time McLachlan dynamics. It
is also the quantity that drives adaptive manifold growth: enlarge the
ansatz when the instantaneous projected defect becomes too large.

\hypertarget{c.2-why-projective-hilbert-space-is-the-right-setting}{%
\subsection{17C.2 Why projective Hilbert space is the right
setting}\label{c.2-why-projective-hilbert-space-is-the-right-setting}}

A normalized ket \(|\psi\rangle\) and \(e^{i\phi}|\psi\rangle\)
represent the same physical state. So the physical state space is
projective Hilbert space \(\mathbb P(\mathcal H)\) rather than the unit
sphere itself. At a normalized state \(|\psi\rangle\), define the
horizontal projector \[
Q_\psi:=I-|\psi\rangle\langle\psi|.
\] For each parameter \(\theta_j\), the raw tangent vector is \[
\tau_j:=\partial_{\theta_j}|\psi(\theta;\mathcal O)\rangle,
\] and the horizontal tangent is \[
\bar\tau_j:=Q_\psi\tau_j.
\] Collect the horizontal tangents as columns: \[
\bar T:=[\bar\tau_1,\dots,\bar\tau_{n_\theta}]\in\mathbb C^{d\times n_\theta}.
\] The exact Schrödinger drift is \[
b:=-iH\psi,
\] but the physically relevant local velocity is its horizontal part \[
\bar b:=Q_\psi b
=
-i(H-E_\psi)\psi,
\qquad
E_\psi:=\langle\psi|H|\psi\rangle.
\] Its squared norm is \[
\|\bar b\|^2
=
\langle\psi|(H-E_\psi)^2|\psi\rangle
=
\operatorname{Var}_\psi(H).
\] So the squared norm of the exact projective Schrödinger velocity is
the energy variance. This is why \(\operatorname{Var}_\psi(H)\) is the
natural normalization scale for projected McLachlan defects.

\hypertarget{c.3-the-current-manifold-and-its-allowed-velocities}{%
\subsection{17C.3 The current manifold and its allowed
velocities}\label{c.3-the-current-manifold-and-its-allowed-velocities}}

Because the parameters are real, the attainable instantaneous velocities
form the \textbf{real-linear} tangent cone \[
T_\psi\mathcal M_{\mathcal O}
=
\left\{
\bar T\dot\theta:\dot\theta\in\mathbb R^{n_\theta}
\right\}.
\] This is not a complex-linear subspace in general. So the correct
local geometry is built from the real inner product \[
g(u,v):=\Re\langle u,v\rangle.
\] The variational velocity associated with parameter rates
\(\dot\theta\) is \[
v(\dot\theta)=\bar T\dot\theta.
\] \#\# 17C.4 The fixed-structure McLachlan step

The fixed-structure local problem is \[
\dot\theta_0=\arg\min_{\dot\theta\in\mathbb R^{n_\theta}}\|\bar T\dot\theta-\bar b\|^2.
\] Its quadratic form, projective metric, force vector, and solved
stationarity system close to \[
\begin{aligned}
\|\bar T\dot\theta-\bar b\|^2
&=\|\bar b\|^2-2\Re\langle \bar T\dot\theta,\bar b\rangle+\dot\theta^\top\Re(\bar T^\dagger\bar T)\dot\theta,\\
(\bar G)_{ij}&=\Re\langle \bar\tau_i,\bar\tau_j\rangle=\Re\langle \partial_{\theta_i}\psi|Q_\psi|\partial_{\theta_j}\psi\rangle,\\
\bar f_i&=\Re\langle \bar\tau_i,\bar b\rangle=\Re\langle Q_\psi\partial_{\theta_i}\psi,-i(H-E_\psi)|\psi\rangle,\\
\mathcal L_0(\dot\theta)&=\|\bar b\|^2-2\bar f^\top\dot\theta+\dot\theta^\top\bar G\dot\theta,\\
\bar G\dot\theta_0&=\bar f,\\
\dot\theta_0&=\bar G^+\bar f\qquad\text{if }\bar G\text{ is singular or rank-deficient},\\
0&=\Re\langle \bar\tau_i,\bar T\dot\theta_0-\bar b\rangle\qquad \forall i.
\end{aligned}
\] \#\# 17C.5 What regularization means

A stabilized runtime step solves \[
\dot\theta_\lambda=\arg\min_{\dot\theta\in\mathbb R^{n_\theta}}\left(\|\bar T\dot\theta-\bar b\|^2+\dot\theta^\top\Lambda\dot\theta\right),
\qquad
\Lambda\succeq 0.
\] With \[
K:=\bar G+\Lambda,
\] the regularized geometry closes to \[
\begin{aligned}
K\dot\theta_\lambda&=\bar f,\\
\epsilon_{\mathrm{proj}}^2&:=\min_{\dot\theta}\|\bar T\dot\theta-\bar b\|^2=\|\bar b\|^2-\bar f^\top\bar G^+\bar f,\\
\epsilon_{\mathrm{step}}^2&:=\|\bar T\dot\theta_\lambda-\bar b\|^2,\\
\mathcal L_\lambda^\star&:=\min_{\dot\theta}\left(\|\bar T\dot\theta-\bar b\|^2+\dot\theta^\top\Lambda\dot\theta\right)=\|\bar b\|^2-\bar f^\top K^{-1}\bar f\qquad (K\text{ invertible}),\\
\mathcal L_\lambda^\star&=\epsilon_{\mathrm{step}}^2+\dot\theta_\lambda^\top\Lambda\dot\theta_\lambda.
\end{aligned}
\] So \(\epsilon_{\mathrm{proj}}^2\) measures intrinsic model
inadequacy, \(\epsilon_{\mathrm{step}}^2\) measures the actual mismatch
of the damped runtime step, and \(\mathcal L_\lambda^\star\) is the
stabilized solver objective.

\hypertarget{c.6-why-the-residual-is-the-right-local-error-quantity}{%
\subsection{17C.6 Why the residual is the right local error
quantity}\label{c.6-why-the-residual-is-the-right-local-error-quantity}}

Choose a local gauge in which the representatives are horizontal at the
current physical time \(\tau\). Then the exact and variational
projective velocities are \[
\bar b
\qquad\text{and}\qquad
\bar T\dot\theta.
\] Define the horizontal residual \[
\bar r:=\bar T\dot\theta-\bar b.
\] Then, for a short step \(h\), \[
|\psi_{\mathrm{ex}}(\tau+h)\rangle
=
|\psi\rangle+h\,\bar b+O(h^2),
\] \[
|\psi_{\mathrm{var}}(\tau+h)\rangle
=
|\psi\rangle+h\,\bar T\dot\theta+O(h^2),
\] so \[
|\psi_{\mathrm{var}}(\tau+h)\rangle-|\psi_{\mathrm{ex}}(\tau+h)\rangle
=
h\,\bar r+O(h^2).
\] So the correct local score is an
\textbf{instantaneous derivative-mismatch score}. Static energy scores
or novelty scores can help in secondary ranking, but they are not the
primary real-time control quantity.

\hypertarget{c.7-the-geometric-projector-behind-the-formulas}{%
\subsection{17C.7 The geometric projector behind the
formulas}\label{c.7-the-geometric-projector-behind-the-formulas}}

Because the tangent space is real-linear, the correct projector is \[
P_{\bar T}^{\mathbb R}v
=
\bar T\,\bar G^+\,\Re(\bar T^\dagger v).
\] The unregularized optimal residual is therefore \[
\bar r_0
=
\bar T\dot\theta_0-\bar b
=
\left(P_{\bar T}^{\mathbb R}-I\right)\bar b.
\] So the ansatz follows the component of the exact projective
Schrödinger vector field that lies in the current real tangent plane,
and the residual is the orthogonal complement.

\hypertarget{c.8-what-a-candidate-generator-really-is}{%
\subsection{17C.8 What a candidate generator really
is}\label{c.8-what-a-candidate-generator-really-is}}

Suppose the current scaffold is an ordered product ansatz \[
|\psi(\theta;\mathcal O)\rangle
=
e^{-i\theta_{n_\theta} a_{n_\theta}}\cdots e^{-i\theta_1 a_1}|\phi_0\rangle.
\] Now fix a candidate-position record \(r=(m,p)\), insert the candidate
generator \(A_m\) at position \(p\), and let the new parameter
\(\vartheta_r\) be initialized at zero: \[
|\psi_r^+(\theta,\vartheta_r)\rangle
=
e^{-i\theta_{n_\theta} a_{n_\theta}}\cdots e^{-i\theta_p a_p}
e^{-i\vartheta_r A_m}
e^{-i\theta_{p-1}a_{p-1}}\cdots e^{-i\theta_1 a_1}|\phi_0\rangle.
\] At \(\vartheta_r=0\), \[
|\psi_r^+(\theta,0)\rangle=|\psi(\theta)\rangle.
\] So the state is continuous, but the tangent space changes
immediately. The new tangent column is \[
u_r:=\left.\partial_{\vartheta_r}|\psi_r^+(\theta,\vartheta_r)\rangle\right|_{\vartheta_r=0},
\qquad
\bar u_r:=Q_\psi u_r.
\] So a candidate is not merely ``a promising operator.'' It is a
\textbf{new tangent direction that becomes available at the current state without moving the state}.
For a block candidate carrying \(q\) new real parameters, collect the
new columns into \[
\bar U_r\in\mathbb C^{d\times q}.
\] \#\# 17C.9 The exact local gain formula

Now augment the local problem: \[
\min_{\dot\theta,\eta}\Big(\|\bar T\dot\theta+\bar U_r\eta-\bar b\|^2+\dot\theta^\top\Lambda\dot\theta+\eta^\top\Lambda_r\eta\Big),
\qquad
\Lambda_r\succeq 0.
\] With the mixed overlaps \[
\begin{aligned}
B_{r,ia}&=\Re\langle \bar\tau_i,\bar u_{r,a}\rangle,&
B_r&=\Re(\bar T^\dagger\bar U_r)\in\mathbb R^{n_\theta\times q},\\
C_{r,ab}&=\Re\langle \bar u_{r,a},\bar u_{r,b}\rangle,&
C_r&=\Re(\bar U_r^\dagger\bar U_r)\in\mathbb R^{q\times q},\\
q_{r,a}&=\Re\langle \bar u_{r,a},\bar b\rangle,&
q_r&=\Re(\bar U_r^\dagger\bar b)\in\mathbb R^q,
\end{aligned}
\] the augmented normal equations, Schur complement, and gain reduce to
\[
\begin{aligned}
\begin{bmatrix}K&B_r\\B_r^\top&C_r+\Lambda_r\end{bmatrix}\begin{bmatrix}\dot\theta\\\eta\end{bmatrix}&=\begin{bmatrix}\bar f\\q_r\end{bmatrix},\\
\dot\theta(\eta)&=K^{-1}(\bar f-B_r\eta),\\
\mathcal L_{\mathrm{aug},\lambda}^\star(\eta)&=\mathcal L_\lambda^\star-2\eta^\top w_r+\eta^\top S_r\eta,\\
S_r&=C_r+\Lambda_r-B_r^\top K^{-1}B_r,\\
w_r&=q_r-B_r^\top K^{-1}\bar f,\\
\eta_\star&=S_r^+w_r,\\
\Delta_{\mathrm{McL}}(r;\tau)&:=\mathcal L_\lambda^\star-\mathcal L_{\mathrm{aug},\lambda}^\star=w_r^\top S_r^+w_r\\
&=\bigl(q_r-B_r^\top K^{-1}\bar f\bigr)^\top\bigl(C_r+\Lambda_r-B_r^\top K^{-1}B_r\bigr)^+\bigl(q_r-B_r^\top K^{-1}\bar f\bigr).
\end{aligned}
\] For a single real candidate direction (\(q=1\)), this reduces to \[
\Delta_{\mathrm{McL}}(r;\tau)=\frac{w_r^2}{S_r}.
\] \#\# 17C.10 What \(S\) and \(w\) mean

The Schur-complement pair and resulting score may be read in one chain:
\[
\begin{aligned}
S_r&=C_r+\Lambda_r-B_r^\top K^{-1}B_r,\\
w_r&=q_r-B_r^\top K^{-1}\bar f,\\
\Delta_{\mathrm{McL}}(r;\tau)&=w_r^\top S_r^+w_r\\
&=\bigl(q_r-B_r^\top K^{-1}\bar f\bigr)^\top\bigl(C_r+\Lambda_r-B_r^\top K^{-1}B_r\bigr)^+\bigl(q_r-B_r^\top K^{-1}\bar f\bigr).
\end{aligned}
\] So \(S_r\) is the metric of the new block after removing the part
already explainable by the old tangent space, \(w_r\) is the
residualized overlap with the exact projective drift, and
\(\Delta_{\mathrm{McL}}(r;\tau)\) is residualized force squared divided
by residualized norm.

\hypertarget{c.11-the-unregularized-geometric-meaning}{%
\subsection{17C.11 The unregularized geometric
meaning}\label{c.11-the-unregularized-geometric-meaning}}

Set \(\Lambda=\Lambda_c=0\). For a single real candidate direction
\(\bar u\), define \[
\bar u_\perp:=(I-P_{\bar T}^{\mathbb R})\bar u,
\qquad
\bar r_0:=(I-P_{\bar T}^{\mathbb R})\bar b.
\] Then the exact one-direction gain closes to \[
\begin{aligned}
\Delta_0
&=\frac{\bigl(\Re\langle \bar u_\perp,\bar r_0\rangle\bigr)^2}{\|\bar u_\perp\|^2}\\
&=\frac{\bigl(\Re\langle (I-P_{\bar T}^{\mathbb R})\bar u,(I-P_{\bar T}^{\mathbb R})\bar b\rangle\bigr)^2}{\|(I-P_{\bar T}^{\mathbb R})\bar u\|^2}.
\end{aligned}
\] The numerator uses \(\Re\langle \bar u_\perp,\bar r_0\rangle\) rather
than \(|\langle \bar u_\perp,\bar r_0\rangle|\) because one is adding
\textbf{one real parameter}, not an arbitrary complex tangent line.

\hypertarget{c.12-when-should-the-ansatz-adapt}{%
\subsection{17C.12 When should the ansatz
adapt?}\label{c.12-when-should-the-ansatz-adapt}}

The intrinsic normalized miss closes to \[
\begin{aligned}
\rho_{\mathrm{miss}}
&:=\frac{\epsilon_{\mathrm{proj}}^2}{\|\bar b\|^2+\varepsilon}
=\frac{\|\bar b\|^2-\bar f^\top\bar G^+\bar f}{\|\bar b\|^2+\varepsilon}\\
&=\frac{\operatorname{Var}_\psi(H)-\bar f^\top\bar G^+\bar f}{\operatorname{Var}_\psi(H)+\varepsilon}
=\frac{\langle \psi|(H-E_\psi)^2|\psi\rangle-\bar f^\top\bar G^+\bar f}{\langle \psi|(H-E_\psi)^2|\psi\rangle+\varepsilon}.
\end{aligned}
\] Interpretation:

\begin{itemize}
\item $\rho_{\mathrm{miss}}\approx 0$: the current manifold already captures the projective velocity well,
\item $\rho_{\mathrm{miss}}$ large: the manifold is missing important directions.

\end{itemize}

A clean rule is therefore \[
\text{adapt if }\rho_{\mathrm{miss}}>\tau_{\mathrm{miss}}.
\] One may also adapt for numerical unreliability or after a maximum
interval without any structural update.

\hypertarget{c.13-which-candidates-should-be-admissible}{%
\subsection{17C.13 Which candidates should be
admissible?}\label{c.13-which-candidates-should-be-admissible}}

A candidate may have nonzero gain and still be numerically poor because
it is nearly redundant. So before ranking by \(\Delta_\lambda\), one
should inspect the residualized geometry:

\begin{itemize}
\item compute singular values of $S$, or of the corresponding whitened residualized block,
\item reject candidates whose retained support is too ill-conditioned.

\end{itemize}

An invariant way to do this is to whiten the old tangent metric. If \[
\bar G\approx V\Sigma^2V^\top
\] on its numerical support, define whitened coordinates \[
z:=\Sigma V^\top\dot\theta.
\] Then penalizing \(\|z\|^2\) is a penalty on
\textbf{physical tangent speed} rather than coordinate speed. This is
the least chart-dependent way to stabilize an overparameterized
manifold.

\hypertarget{c.14-two-different-residuals-should-be-tracked}{%
\subsection{17C.14 Two different residuals should be
tracked}\label{c.14-two-different-residuals-should-be-tracked}}

The projective residual is \[
\bar r:=\bar T\dot\theta_\lambda-\bar b.
\] This is the correct gauge-invariant quantity for local candidate
scoring and local manifold-growth decisions.

The raw Hilbert-space residual is \[
r_{\mathrm{raw}}
:=
T\dot\theta_\lambda-(-iH\psi)
=
T\dot\theta_\lambda+iH\psi.
\] These are related by \[
\bar r=Q_\psi r_{\mathrm{raw}}.
\] So:

\begin{itemize}
\item use $\bar r$ for \textbf{local structure growth},
\item use $r_{\mathrm{raw}}$ when a rigorous Hilbert-space accumulated-error bound is desired.

\end{itemize}

\hypertarget{c.15-why-integrated-residual-is-the-right-branch-pruning-quantity}{%
\subsection{17C.15 Why integrated residual is the right branch-pruning
quantity}\label{c.15-why-integrated-residual-is-the-right-branch-pruning-quantity}}

Let a branch be propagated on a fixed scaffold over a segment
\([\tau_a,\tau_b]\). Let \(|\psi_{\mathrm{ex}}\rangle\) be the exact
state and \(|\psi_{\mathrm{var}}\rangle\) the variational state with the
same initial condition at \(\tau_a\), and define \[
e(\tau):=|\psi_{\mathrm{var}}(\tau)\rangle-|\psi_{\mathrm{ex}}(\tau)\rangle.
\] Then the error evolution, Duhamel representation, norm bound, and
segment cost close to \[
\begin{aligned}
\dot\psi_{\mathrm{ex}}&=-iH\psi_{\mathrm{ex}},&
\dot\psi_{\mathrm{var}}&=T\dot\theta_\lambda,&
r_{\mathrm{raw}}&=T\dot\theta_\lambda+iH\psi,\\
\dot e&=-iHe+r_{\mathrm{raw}},\\
e(\tau)&=\int_{\tau_a}^\tau U(\tau,s)\,r_{\mathrm{raw}}(s)\,ds,\\
\|e(\tau)\|&\le\int_{\tau_a}^\tau\|r_{\mathrm{raw}}(s)\|\,ds\le\sqrt{\tau-\tau_a}\left(\int_{\tau_a}^\tau\|r_{\mathrm{raw}}(s)\|^2\,ds\right)^{1/2},\\
J_{\mathrm{seg}}^{\mathrm{raw}}&:=\int_{\tau_a}^{\tau_b}\|r_{\mathrm{raw}}(\tau)\|^2\,d\tau=\int_{\tau_a}^{\tau_b}\|T\dot\theta_\lambda(\tau)+iH(\tau)\psi(\tau)\|^2\,d\tau.
\end{aligned}
\] So \(J_{\mathrm{seg}}^{\mathrm{raw}}\) is a natural upper-bound
surrogate for accumulated state-vector error. If one wants a fully
gauge-invariant pruning metric instead, replace
\(\|r_{\mathrm{raw}}(\tau)\|^2\) by
\(\|\bar r(\tau)\|^2=\|Q_\psi r_{\mathrm{raw}}(\tau)\|^2\) and interpret
the result as a ray/projective error surrogate.

\hypertarget{c.16-the-branch-objective}{%
\subsection{17C.16 The branch
objective}\label{c.16-the-branch-objective}}

A branch \(\mathfrak b\) is a sequence of scaffolds
\(\mathcal O_0,\mathcal O_1,\dots,\mathcal O_J\) used over consecutive
checkpoint segments, and the cumulative objective closes as \[
\begin{aligned}
J_{\mathrm{dyn}}(\mathfrak b)
&=
\sum_{j=0}^{J-1}J_{\mathrm{seg},j}
+\beta\,C_{\mathrm{hw}}
+\gamma\,N_{\mathrm{adapt}}
+\mu_E\sum_{j=0}^{J-1}\int_{\tau_j}^{\tau_{j+1}}
\left|
\frac{d}{d\tau}\langle H(\tau)\rangle_j-\langle\partial_\tau H(\tau)\rangle_j
\right|^2d\tau\\
&=
\sum_{j=0}^{J-1}\int_{\tau_j}^{\tau_{j+1}}\|r_{\mathrm{raw}}^{(j)}(\tau)\|^2\,d\tau
+\beta\,C_{\mathrm{hw}}
+\gamma\,N_{\mathrm{adapt}}
+\mu_E\sum_{j=0}^{J-1}\int_{\tau_j}^{\tau_{j+1}}
\left|
\frac{d}{d\tau}\langle H(\tau)\rangle_j-\langle\partial_\tau H(\tau)\rangle_j
\right|^2d\tau\\
&=
\sum_{j=0}^{J-1}\int_{\tau_j}^{\tau_{j+1}}
\bigl\|T^{(j)}(\tau)\dot\theta_\lambda^{(j)}(\tau)+iH(\tau)\psi^{(j)}(\tau)\bigr\|^2\,d\tau
+\beta\,C_{\mathrm{hw}}
+\gamma\,N_{\mathrm{adapt}}
+\mu_E\sum_{j=0}^{J-1}\int_{\tau_j}^{\tau_{j+1}}
\left|
\frac{d}{d\tau}\langle H(\tau)\rangle_j-\langle\partial_\tau H(\tau)\rangle_j
\right|^2d\tau.
\end{aligned}
\] If one wants a fully gauge-invariant branch surrogate instead,
replace \(\|r_{\mathrm{raw}}^{(j)}(\tau)\|^2\) by
\(\|\bar r^{(j)}(\tau)\|^2=\|Q_{\psi^{(j)}}r_{\mathrm{raw}}^{(j)}(\tau)\|^2\).

\hypertarget{c.17-the-full-process-in-one-coherent-sequence}{%
\subsection{17C.17 The full process, in one coherent
sequence}\label{c.17-the-full-process-in-one-coherent-sequence}}

At each physical time \(\tau\), with current scaffold \(\mathcal O\),
the full loop can be written compactly as \[
\begin{aligned}
\text{A: }&
|\psi\rangle=|\psi(\theta;\mathcal O)\rangle,
\qquad
Q_\psi=I-|\psi\rangle\langle\psi|,
\qquad
\bar T=Q_\psi T,
\qquad
\bar b=Q_\psi(-iH\psi)=-i(H-E_\psi)\psi,\\
\text{B: }&
\bar G=\Re(\bar T^\dagger\bar T),
\qquad
\bar f=\Re(\bar T^\dagger\bar b),
\qquad
K=\bar G+\Lambda,
\qquad
\dot\theta_\lambda=K^+\bar f,\\
&
\epsilon_{\mathrm{proj}}^2=\|\bar b\|^2-\bar f^\top\bar G^+\bar f,
\qquad
\epsilon_{\mathrm{step}}^2=\|\bar T\dot\theta_\lambda-\bar b\|^2,
\qquad
\bar r=\bar T\dot\theta_\lambda-\bar b,
\qquad
r_{\mathrm{raw}}=T\dot\theta_\lambda+iH\psi,\\
\text{C: }&
\rho_{\mathrm{miss}}
=
\frac{\epsilon_{\mathrm{proj}}^2}{\|\bar b\|^2+\varepsilon}
=
\frac{\epsilon_{\mathrm{proj}}^2}{\operatorname{Var}_\psi(H)+\varepsilon},
\qquad
\text{adapt iff }\rho_{\mathrm{miss}}>\tau_{\mathrm{miss}},\\
\text{D: }&
B_r=\Re(\bar T^\dagger\bar U_r),
\qquad
C_r=\Re(\bar U_r^\dagger\bar U_r),
\qquad
q_r=\Re(\bar U_r^\dagger\bar b),\\
&
S_r=C_r+\Lambda_r-B_r^\top K^+B_r,
\qquad
w_r=q_r-B_r^\top K^+\bar f,
\qquad
\Delta_{\mathrm{McL}}(r;\tau)=w_r^\top S_r^+w_r.
\end{aligned}
\] For each shortlisted admissible candidate-position record
\(r=(m,p)\), one then creates a zero-initialized child, propagates it
over the next probe segment, and prunes on the accumulated branch
objective: \[
\begin{aligned}
\text{E: }&
\mathcal O\mapsto\mathcal O_r^+,
\qquad
|\psi(\theta,0;\mathcal O_r^+)\rangle=|\psi(\theta;\mathcal O)\rangle,\\
\text{F: }&
J_{\mathrm{seg},r}^{\mathrm{raw}}
=
\int_{\tau_a}^{\tau_b}\|r_{\mathrm{raw},r}(\tau)\|^2\,d\tau
=
\int_{\tau_a}^{\tau_b}\bigl\|T_r(\tau)\dot\theta_{\lambda,r}(\tau)+iH(\tau)\psi_r(\tau)\bigr\|^2\,d\tau,\\
\text{G: }&
J_{\mathrm{dyn}}^{\mathrm{new}}(r;\mathfrak b)
=
J_{\mathrm{dyn}}^{\mathrm{old}}(\mathfrak b)
+
J_{\mathrm{seg},r}^{\mathrm{raw}}
+
\beta\,C_{\mathrm{hw},r}
+
\gamma\,N_{\mathrm{adapt},r},
\qquad
\mathfrak B_{\mathrm{next}}
=
\operatorname{Prune}\!\left(\{(\mathcal O_r^+,J_{\mathrm{dyn}}^{\mathrm{new}}(r;\mathfrak b))\}_r\right).
\end{aligned}
\] That is the whole piecewise-manifold control loop.

\hypertarget{c.18-one-concise-interpretation}{%
\subsection{17C.18 One concise
interpretation}\label{c.18-one-concise-interpretation}}

The method is doing this:

\begin{quote}
The exact quantum dynamics asks for a projective velocity $\bar b$. \\
The current ansatz only provides a real tangent plane. \\
McLachlan chooses the closest available tangent velocity. \\
If the miss is too large, the algorithm enlarges the tangent plane by adding zero-initialized candidate directions. \\
The local gain formula tells us which candidate most reduces the current miss. \\
Short probe rollouts then decide which structural choices remain best after actual propagation.
\end{quote}

\hypertarget{c.19-explicit-symbolicmathematical-differences-between-the-clean-story-and-the-current-implementation}{%
\subsection{17C.19 Explicit symbolic--mathematical differences between
the clean story and the current
implementation}\label{c.19-explicit-symbolicmathematical-differences-between-the-clean-story-and-the-current-implementation}}

The clean story above is the right mathematical idealization. The
current implementation differs from that idealization in a few explicit
ways:

\begin{enumerate}
\item \textbf{Sign convention.} The clean story uses
\end{enumerate}

\[
   \bar b=-iQ_\psi H\psi,
   \qquad
   \bar r=\bar T\dot\theta-\bar b.
\] The implementation instead uses \[
   r=\bar T\dot\theta+iQ_\psi H\psi,
   \qquad
   f_i=\Im\langle \bar\tau_i|H|\psi\rangle,
   \qquad
   \Re\langle \bar\tau_i,-iH\psi\rangle=\Im\langle \bar\tau_i,H\psi\rangle.
\]

\begin{enumerate}
\item \textbf{Regularization surface.} Rather than a general positive-semidefinite penalty $\Lambda$ or $\Lambda_c$, the implementation uses a scalar $\lambda\ge 0$ on either chart coordinates or whitened metric coordinates:
\end{enumerate}

\[
   \begin{aligned}
   K_{\mathrm{chart}}&=\bar G+\lambda I,\\
   \bar G&\approx V\Sigma^2V^\top,
   \qquad
   z=\Sigma V^\top\dot\theta,
   \qquad
   (I+\lambda I)z=\Sigma^{-1}V^\top\bar f,
   \qquad
   \dot\theta=V\Sigma^{-1}z.
   \end{aligned}
\]

\begin{enumerate}
\item \textbf{Propagation solve versus policy solve.} The clean story presents one stabilized local solve, whereas the implementation may separate the propagated step from the policy/telemetry step:
\end{enumerate}

\[
   (\dot\theta_{\mathrm{prop}},\epsilon_{\mathrm{step}}^{2,\mathrm{prop}})
   \neq
   (\dot\theta_{\mathrm{policy}},\epsilon_{\mathrm{step}}^{2,\mathrm{policy}})
   \quad\text{in general},
\] even when both are derived from the same local geometry.

\begin{enumerate}
\item \textbf{Variance normalization and guards.} The clean story writes
\end{enumerate}

\[
   \rho_{\mathrm{miss}}=\frac{\epsilon_{\mathrm{proj}}^2}{\operatorname{Var}_\psi(H)+\varepsilon}.
\] The implemented ratio surface is closer to \[
   \rho_{\mathrm{miss}}
   =
   \begin{cases}
   \epsilon_{\mathrm{proj}}^2/(\operatorname{Var}_\psi(H)+\varepsilon),&\operatorname{Var}_\psi(H)\ge \nu_{\min},\\
   \text{suppressed},&\operatorname{Var}_\psi(H)<\nu_{\min},
   \end{cases}
\] with an auxiliary guard flag recording the suppressed regime.

\begin{enumerate}
\item \textbf{Active-parameter windows.} The clean story is written on the full tangent space, but the implementation may restrict to an active suffix or oracle-supplied index set:
\end{enumerate}

\[
   T\mapsto T_{W_{\mathrm{act}}},
   \qquad
   \dot\theta\mapsto(\dot\theta_{W_{\mathrm{act}}},0),
   \qquad
   \bar G\mapsto \bar G_{W_{\mathrm{act}}W_{\mathrm{act}}},
\] so inactive coordinates are frozen.

\begin{enumerate}
\item \textbf{Candidate gain evaluation.} The clean story scores candidates by the Schur-complement gain
\end{enumerate}

\[
   \Delta_\lambda=w^\top S^+w.
\] The implementation instead zero-initializes one append candidate,
rebuilds the augmented local geometry, and scores by direct defect drops
\[
   \Delta_0=\epsilon_{\mathrm{proj,current}}^2-\epsilon_{\mathrm{proj,aug}}^2,
   \qquad
   \Delta_\lambda=\epsilon_{\mathrm{step,current}}^2-\epsilon_{\mathrm{step,aug}}^2.
\]

\begin{enumerate}
\item \textbf{Position-aware current branch growth.} The clean story allows arbitrary insertion positions and multi-parameter blocks. The helper layer now supports position-aware insertion and position-centered post-insert refit windows,
\end{enumerate}

\[
   \mathcal T(\mathfrak b,m,p)=\mathfrak b',
   \qquad
   \mathcal O_{b'}=\mathcal O_b\oplus_p m,
   \qquad
   \theta_{b'}=\theta_b\oplus_p 0,
\] and the active probe set is built from the append slot, the left
boundary, and the active reoptimization window. Hence, when the active
window is made maximally large and the probe-position cap is also made
large enough, the implementation can evaluate the full insertion set
\(p\in\{0,\dots,|\theta|\}\) rather than append only. The live surface
still remains zero-initialized \textbf{single-parameter} growth, but not
append-only growth.

\begin{enumerate}
\item \textbf{Checkpoint-trigger logic.} The clean story says “adapt when $\rho_{\mathrm{miss}}$ is too large.” The implemented trigger surface is richer:
\end{enumerate}

\[
   \mathbf 1_{\mathrm{checkpoint}}
   =
   \mathbf 1_{\mathrm{cadence}}
   \lor
   \mathbf 1\!\bigl(\epsilon_{\mathrm{proj}}^2>\tau_{\mathrm{proj}}\bigr)
   \lor
   \mathbf 1\!\bigl(\epsilon_{\mathrm{step}}^2>\tau_{\mathrm{step}}\bigr)
   \lor
   \mathbf 1\!\bigl(\rho_{\mathrm{miss}}>\tau_{\mathrm{miss}}\text{ when defined}\bigr),
\] together with a model-versus-damping failure classification.

\begin{enumerate}
\item \textbf{Branch objective surface.} The clean story uses a scalar objective such as
\end{enumerate}

\[
   J_{\mathrm{scalar}}=\sum J_{\mathrm{seg}}+\beta C_{\mathrm{hw}}+\gamma N_{\mathrm{adapt}}.
\] The implementation instead prunes by the lexicographic key \[
   \kappa_{\mathrm{prune}}
   =
   \bigl(J^{\mathrm{rt}},J^{\mathrm{res}},J^{\mathrm{model}},J^{\mathrm{damp}},J^{\mathrm{cons}},|\mathcal O|,k,\operatorname{labels}(\mathcal O),\operatorname{round}_{10}(\theta),\operatorname{id}\bigr),
\] after deduplicating by time index, selected-operator labels, and
rounded parameter values.

\begin{enumerate}
\item \textbf{Projective versus raw residual accumulation.} The clean story distinguishes a projective residual $\bar r$ from a raw Hilbert-space residual $r_{\mathrm{raw}}$. The implemented rollout telemetry is closer to
\end{enumerate}

\[
   \bigl(\Delta J^{\epsilon},\Delta J^{\mathrm{model}},\Delta J^{\mathrm{damp}}\bigr)
   =
   \left(
   \int\epsilon_{\mathrm{step}}^2,
   \int\epsilon_{\mathrm{proj}}^2,
   \int\max\!\bigl(\epsilon_{\mathrm{step}}^2-\epsilon_{\mathrm{proj}}^2,0\bigr)
   \right),
\] rather than a separately stored raw Hilbert-space residual bound.

\begin{enumerate}
\item \textbf{Geometry source.} The clean story is written as if the tangent vectors and $H\psi$ are assembled directly every time. The implementation may instead start from externally supplied local data:
\end{enumerate}

\[
   (\bar G,\bar f,\operatorname{Var}_\psi(H))
   =
   \mathcal O_{\mathrm{geom}}(\psi,H,\mathcal O),
\] after which the host solves the McLachlan system on that returned
geometry.

So the mathematical backbone is the same, but the implemented restricted
surface closes schematically to \[
\begin{aligned}
\text{projective policy geometry}
&+\text{ scalar chart/whitened damping}
+\text{ append-only zero-init growth}\\
&+\text{ checkpointed short-horizon rollout}
+\text{ lexicographic prune key}
+\text{ optional oracle-supplied local geometry}.
\end{aligned}
\] \#\# 17C.20 Discrete spectra and density-sector observables for HH
checkpoint dynamics A controller trajectory row \(r_n\) (index
\(n=0,\dots,N-1\)) carries either \texttt{time} or
\texttt{physical\_time}, and stores raw lattice observables under keys
such as \texttt{site\_occupations}, \texttt{site\_occupations\_exact},
\texttt{staggered}, \texttt{doublon}, and \texttt{doublon\_exact}. The
spectral stack in the repository works directly on these rows and does
not re-simulate dynamics.

\[
t_n := \text{row}_n[\text{time}_\star],\qquad
\text{time}_\star \in \{\texttt{time},\texttt{physical\_time}\}.
\] For each site \(j\in\{0,\dots,L-1\}\), define spin-resolved
occupations \[
n^{\uparrow}_{j}(t_n)=\langle \psi(t_n)|\hat n_{j,\uparrow}|\psi(t_n)\rangle,\quad
n^{\downarrow}_{j}(t_n)=\langle \psi(t_n)|\hat n_{j,\downarrow}|\psi(t_n)\rangle,\quad
n_{j}(t_n)=n^{\uparrow}_{j}(t_n)+n^{\downarrow}_{j}(t_n),
\] with \(\hat n_{j,\sigma}=(I-\hat Z_{j,\sigma})/2\) under this
codebase's convention. Equivalently, in the statevector path used by
\texttt{hh\_realtime\_checkpoint\_controller.py}, these are evaluated as
\[
n^{\sigma}_{j}(t_n)=\sum_{s}p_s\,b^{\sigma}_{j}(s),\qquad
D(t_n)=\sum_{s}p_s\sum_{j=0}^{L-1} b^{\uparrow}_{j}(s)b^{\downarrow}_{j}(s),
\] where \(p_s=|\psi_s|^2\) in the computational basis and
\(b^{\sigma}_{j}(s)\in\{0,1\}\) are basis bits. The repository therefore
records \[
\texttt{site\_occupations}[n,j]=n_j(t_n),\qquad
\texttt{staggered}[n]=\frac1L\sum_{j=0}^{L-1}(-1)^j n_j(t_n),\qquad
\texttt{doublon}[n]=D(t_n),
\] and the corresponding \texttt{\_exact} fields for the driven exact
trajectory. \[
D(t_n)=\sum_{j=0}^{L-1} n^{\uparrow}_{j}(t_n)n^{\downarrow}_{j}(t_n).
\] Here \(j\) denotes the lattice-site index, while \(n\) denotes the
discrete time-sample index.

From \(\{n_j(t_n)\}\) one obtains the secondary diagnostics used in the
spectrum module: \[
\delta n_j(t_n):=n_j(t_n)-\frac1L\sum_{m=0}^{L-1}n_m(t_n),
\qquad
d_{ij}(t_n):=n_i(t_n)-n_j(t_n),
\qquad
m(t_n):=\frac1L\sum_{j=0}^{L-1}(-1)^j n_j(t_n).
\] (For \(L=2\), \(\delta n_0=\frac12\,d_{01}\) and
\(\delta n_1=-\frac12\,d_{01}\).)

For spectral post-processing, let \(x_n:=x(t_n)\) denote any sampled
scalar observable. The implemented V1 pipeline first infers or assumes a
uniform step \(t_{n+1}-t_n=\Delta t\), then detrends and re-centers the
trace, applies either no window or a Hann window, and finally evaluates
the one-sided real FFT. In this notation, \(\tilde x_n\) denotes the
detrended and re-centered signal: for constant detrending it is \(x_n\)
with its sample mean removed, whereas for linear detrending it is the
residual after subtracting the least-squares line in \(t_n\),
re-centered exactly as in the code path. The windowed sequence is
\(y_n:=w_n\tilde x_n\), with \(w_n=1\) for no window and \(w_n\) equal
to the standard Hann taper for Hann windowing.

\[
Y_k = \sum_{n=0}^{N-1} y_n e^{-i2\pi kn/N},\qquad
k=0,\dots,\lfloor N/2\rfloor,
\] which is the one-sided coefficient returned by the real-FFT stage
(implemented with \texttt{np.fft.rfft}) for a real sampled trace. The
associated frequency grid is \[
\omega_k=\frac{2\pi k}{N\Delta t},\qquad
f_k=\frac{k}{N\Delta t},\qquad
\omega_{\mathrm{Ny}}=\frac{\pi}{\Delta t},\qquad
\Delta\omega=\frac{2\pi}{N\Delta t}.
\] The reported one-sided amplitude \(A_k\) is obtained from \(|Y_k|\)
by normalizing with the window-weight sum and doubling
positive-frequency bins away from DC and Nyquist: \[
A_k=
\begin{cases}
\dfrac{|Y_k|}{\sum_{n=0}^{N-1} w_n}, & k=0,\\
\dfrac{2|Y_k|}{\sum_{n=0}^{N-1} w_n}, & 1\le k< N/2,\\
\dfrac{|Y_k|}{\sum_{n=0}^{N-1} w_n}, & k=N/2 \text{ when } N \text{ is even}.
\end{cases}
\] Error-time signals are also treated as observables: \[
x^{\mathrm{err}}_n=x^{\mathrm{ctrl}}_n-x^{\mathrm{exact}}_n
\quad\Rightarrow\quad
\text{apply the same FFT-processing machinery to }x^{\mathrm{err}}_n.
\] When the drive frequency \(\omega_d\) is known, the same centered
trace is also fit to drive-locked harmonics: \[
x(t)\approx c_0+\sum_{\ell=1}^{K_{\max}}\!\bigl(a_\ell\cos(\ell\omega_d t)+b_\ell\sin(\ell\omega_d t)\bigr),
\quad
R_\ell=\sqrt{a_\ell^2+b_\ell^2},\ \phi_\ell=\operatorname{atan2}(b_\ell,a_\ell).
\] Here \(R_\ell\) is the amplitude of the \(\ell\)-th drive-locked
harmonic and \(\phi_\ell\) is its phase.

The spectral stack is therefore a full transform layer on top of
trajectory observables: raw site-resolved occupations, staggered
density, doublon, and pair imbalance are each converted to sampled
series, then detrend-window-FFT processed in the same
one-sided-amplitude form, with controller-versus-exact residuals using
the same operators.

\hypertarget{final-primitive-closed-summary}{%
\section{18. Final Primitive-Closed
Summary}\label{final-primitive-closed-summary}}

The symbolic substitution chain is linear and closed:

\begin{enumerate}
\def\labelenumi{\arabic{enumi}.}
\tightlist
\item
  choose the fermion ordering \(p(i,\sigma)\),
\item
  write the Jordan-Wigner ladder primitives \(\hat c_p^\dagger\) and
  \(\hat c_p\),
\item
  reduce number operators to \((I-Z_p)/2\),
\item
  substitute them into \(\hat H_t\), \(\hat H_U\), and \(\hat H_v\),
\item
  write \(\hat b_i\), \(\hat b_i^\dagger\), \(\hat n_{b,i}\),
  \(\hat x_i\), and \(\hat P_i\),
\item
  substitute them into \(\hat H_{\mathrm{ph}}\) and \(\hat H_g\),
\item
  reduce the drive term to explicit identity-plus-\(Z\) coefficients,
\item
  propagate those operators into statevector primitives,
\item
  build fixed and adaptive variational manifolds from the same
  operators,
\item
  define projective McLachlan tangent geometry on those manifolds,
\item
  define branch growth, probe rollout, and beam pruning on top of the
  local defect geometry,
\item
  define continuation and handoff on top of those manifolds, with
  historical replay paths kept in the appendix,
\item
  define benchmark and noise-validation contracts relative to the same
  exact targets.
\end{enumerate}

This is the point of the symbolic duplicate: the same mathematical
skeleton remains, but implementation labels are stripped away so the
document can serve as a clean theory-facing reference.

\hypertarget{data}{%
\section{19. Data}\label{data}}

This final chapter records empirical evidence for the selector and
pruning claims developed in the main body. The goal is not to replace
the symbolic formulas above, but to show that the live Phase-3 HH
selector surfaces described in §§11 and 17 materially change the
observed outcomes on matched runs.

\hypertarget{matched-novelty-ablation}{%
\subsection{19.1 Matched novelty
ablation}\label{matched-novelty-ablation}}

The freshest clean ablation available is a matched four-run HH study
with \[
L=2,\qquad t=1,\qquad U=4,\qquad \omega_0=1,\qquad g=0.5,\qquad n_{\mathrm{ph,max}}=1,
\] using binary bosons, blocked ordering, open boundaries, the same
seed, the same compiled backend path, the same windowed reoptimization
policy, the same live pruning machinery, and the same
rescue/symmetry/lifetime settings. The only changes are:

\begin{enumerate}
\def\labelenumi{\arabic{enumi}.}
\tightlist
\item
  operator pool: \texttt{pareto\_lean\_l2} versus \texttt{full\_meta},
\item
  novelty exponent: \(\gamma_N=1\) versus \(\gamma_N=0\).
\end{enumerate}

Writing the final filtered-energy error as \[
\lvert \Delta E \rvert = \bigl|E_{\mathrm{final}}-E_{\mathrm{exact,filtered}}\bigr|,
\] the matched runs give:

\begin{longtable}[]{@{}
  >{\raggedright\arraybackslash}p{(\columnwidth - 14\tabcolsep) * \real{0.10}}
  >{\raggedleft\arraybackslash}p{(\columnwidth - 14\tabcolsep) * \real{0.13}}
  >{\raggedleft\arraybackslash}p{(\columnwidth - 14\tabcolsep) * \real{0.13}}
  >{\raggedleft\arraybackslash}p{(\columnwidth - 14\tabcolsep) * \real{0.13}}
  >{\raggedleft\arraybackslash}p{(\columnwidth - 14\tabcolsep) * \real{0.13}}
  >{\raggedleft\arraybackslash}p{(\columnwidth - 14\tabcolsep) * \real{0.13}}
  >{\raggedright\arraybackslash}p{(\columnwidth - 14\tabcolsep) * \real{0.10}}
  >{\raggedleft\arraybackslash}p{(\columnwidth - 14\tabcolsep) * \real{0.13}}@{}}
\toprule
Pool & \(\gamma_N\) & \(E_{\mathrm{final}}\) &
\(\lvert \Delta E \rvert\) & Depth & Parameters & Stop reason & Prune
kept \\
\midrule
\endhead
\texttt{pareto\_lean\_l2} & \texttt{1.0} & \texttt{0.15872408236037028}
& \texttt{5.6178234644\ \textbackslash{}times\ 10\^{}\{-5\}} &
\texttt{14} & \texttt{25} & \texttt{drop\_plateau} & \texttt{4\ /\ 5} \\
\texttt{pareto\_lean\_l2} & \texttt{0.0} & \texttt{0.48696021436469644}
& \texttt{3.2829231024\ \textbackslash{}times\ 10\^{}\{-1\}} &
\texttt{14} & \texttt{23} & \texttt{drop\_plateau} & \texttt{4\ /\ 5} \\
\texttt{full\_meta} & \texttt{1.0} & \texttt{0.15872408236037047} &
\texttt{5.6178234644\ \textbackslash{}times\ 10\^{}\{-5\}} & \texttt{13}
& \texttt{28} & \texttt{drop\_plateau} & \texttt{5\ /\ 5} \\
\texttt{full\_meta} & \texttt{0.0} & \texttt{0.48699090951182056} &
\texttt{3.2832300539\ \textbackslash{}times\ 10\^{}\{-1\}} & \texttt{13}
& \texttt{22} & \texttt{drop\_plateau} & \texttt{5\ /\ 5} \\
\bottomrule
\end{longtable}

If \[
R_{\mathrm{nov}}(\mathcal G)
=
\frac{\lvert \Delta E \rvert_{\gamma_N=0,\mathcal G}}{\lvert \Delta E \rvert_{\gamma_N=1,\mathcal G}},
\] then the two matched pool ratios are \[
R_{\mathrm{nov}}(\texttt{pareto\_lean\_l2})\approx 5.84376\times 10^3,
\qquad
R_{\mathrm{nov}}(\texttt{full\_meta})\approx 5.84431\times 10^3.
\]

So, on this matched HH benchmark, turning novelty weighting off worsens
the final filtered-energy error by roughly a factor of
\(5.8\times 10^3\) in both a lean pool and a heavy pool.

\hypertarget{interpretation-of-the-novelty-ablation}{%
\subsection{19.2 Interpretation of the novelty
ablation}\label{interpretation-of-the-novelty-ablation}}

Several points survive this comparison cleanly.

\begin{enumerate}
\def\labelenumi{\arabic{enumi}.}
\item
  \textbf{Novelty-on recovers the low-error line.} In both pools,
  \(\gamma_N=1\) reaches the same final error scale, \[
  \lvert \Delta E \rvert \approx 5.62\times 10^{-5}.
  \]
\item
  \textbf{Novelty-off collapses to a much worse terminal scaffold.} In
  both pools, \(\gamma_N=0\) ends near \[
  E_{\mathrm{final}}\approx 0.487,
  \qquad
  \lvert \Delta E \rvert \approx 3.28\times 10^{-1},
  \] far above the filtered exact target near \(0.158668\).
\item
  \textbf{This is not merely a depth effect.} The on/off pairs stop at
  the same logical depth within each pool, so the dominant difference is
  the operator-ranking trajectory rather than a simple ``more steps''
  explanation.
\item
  \textbf{This is not merely a pruning-on versus pruning-off effect.}
  Pruning stays active in all four runs. What changes is the upstream
  ranking rule. In particular, novelty telemetry is still assembled when
  \(\gamma_N=0\); what is ablated is the multiplicative use of novelty
  inside the selector score, not the underlying geometric information
  itself.
\end{enumerate}

These runs therefore support the claim that the reduced-tangent-space
novelty factor is doing real work inside the ADAPT ranking rule rather
than serving as passive telemetry.

\hypertarget{pool-family-redundancy}{%
\subsection{19.3 Pool-family redundancy}\label{pool-family-redundancy}}

The same data story also sharpens the pool-reduction claim. For the
studied HH regime \[
L=2,\qquad n_{\mathrm{ph,max}}=1,
\] the heavy \texttt{full\_meta} pool has size \(46\), while the lean
\texttt{pareto\_lean\_l2} pool has size \(11\). Hence the pool-reduction
fraction is \[
\eta_{\mathrm{pool}}
=
1-\frac{11}{46}
=
\frac{35}{46}
\approx 0.7609.
\]

Equivalently, the lean pool removes about \(76\%\) of the heavy-pool
candidates while still retaining the same practical final energy scale,
\[
E_{\mathrm{final}}\approx 0.15872408,
\qquad
\lvert \Delta E \rvert \approx 5.62\times 10^{-5},
\] on the current-code lean rerun.

The generator classes discussed in this L=2 comparison may be
represented symbolically by \[
\begin{aligned}
A^{\mathrm{sing}}_{ia,\sigma}
&=
i\!\left(\hat c^\dagger_{a\sigma}\hat c_{i\sigma}-\hat c^\dagger_{i\sigma}\hat c_{a\sigma}\right),\\
A^{\mathrm{dbl}}_{ijab}
&=
i\!\left(\hat c^\dagger_a\hat c^\dagger_b\hat c_j\hat c_i-\mathrm{h.c.}\right),\\
G_i^{(\mathrm{cloud})}
&=
\tilde n_i\sum_{r\in\mathcal N(i)}\alpha_{ir}P_r,\\
G_{ij}^{(\mathrm{hd})}
&=
T_{ij}^{(+)}(P_i-P_j),\\
G_i^{(\mathrm{dd})}
&=
D_iP_i,\\
G_{ij}^{(\mathrm{od})}
&=
J_{ij}^{(-)}(x_i-x_j),\\
G_{ij}^{(2)}
&=
T_{ij}^{(+)}(x_i-x_j)^2,\\
x_i
&=
\hat b_i+\hat b_i^\dagger,
\qquad
P_i
=
i(\hat b_i^\dagger-\hat b_i).
\end{aligned}
\] On this \(L=2\) benchmark, the retained/useful classes were
\(A^{\mathrm{sing}}\), the \(P\)-type polaronic classes
\(G_i^{(\mathrm{cloud})}\), \(G_{ij}^{(\mathrm{hd})}\), and
\(G_i^{(\mathrm{dd})}\), together with the quadrature seeds. By
contrast, \(A^{\mathrm{dbl}}\), HVA, and the x-type / higher-order
channels \(G_{ij}^{(\mathrm{od})}\) and \(G_{ij}^{(2)}\) were not needed
to preserve the achieved error scale.

Empirically, the following heavy-pool families are absent from the lean
pool yet are not needed to retain that achieved accuracy scale on this
benchmark:

\begin{itemize}
\tightlist
\item
  lifted UCCSD doubles,
\item
  HH unit terms,
\item
  HVA layers,
\item
  \texttt{paop\_cloud\_x},
\item
  \texttt{paop\_disp},
\item
  \texttt{paop\_dbl},
\item
  \texttt{paop\_dbl\_x},
\item
  \texttt{paop\_curdrag},
\item
  \texttt{paop\_hop2}.
\end{itemize}

This should be read as a scoped data claim, not as a universal proof:
for the present HH \(L=2\), \(n_{\mathrm{ph,max}}=1\) target problem,
the heavy pool carries substantial operator-family redundancy relative
to the observed energy objective.

\hypertarget{weak-coupling-l2-class-prune-redo}{%
\subsubsection{\texorpdfstring{19.3.1 Weak-coupling \(L=2\) class-prune
redo}{19.3.1 Weak-coupling L=2 class-prune redo}}\label{weak-coupling-l2-class-prune-redo}}

A second \(L=2\) class-prune line at weaker coupling sharpens the same
conclusion. For \[
L=2,\qquad t=1,\qquad U=0.5,\qquad \omega_0=1,\qquad g=0.2,\qquad n_{\mathrm{ph,max}}=1,
\] using binary bosons, blocked ordering, open boundaries, the compiled
backend path, full-window reoptimization, beam width \(B=3\), and the
direct \texttt{phase3\_v1} HH selector, the heavy parent
\texttt{full\_meta} run reached \[
E_{\mathrm{parent}}=-0.7763497696348323,
\qquad
\lvert\Delta E\rvert_{\mathrm{parent}}=2.7761269927\times 10^{-5},
\qquad
F_{\mathrm{parent}}=0.9999861889450502,
\] at logical depth \(d_{\mathrm{parent}}=7\).

From that parent, a class-first pruning loop removed one operator family
at a time and reran fresh ADAPT after every proposed removal. A removal
was accepted only if the rerun satisfied \[
\lvert\Delta E\rvert \le 5\times 10^{-4},
\qquad
F \ge F_{\mathrm{parent}}-10^{-4}.
\] The accepted removals were

\begin{itemize}
\tightlist
\item
  \texttt{hh\_termwise\_unit},
\item
  \texttt{hva\_layer},
\item
  \texttt{paop\_cloud\_x},
\item
  \texttt{paop\_curdrag},
\item
  \texttt{paop\_dbl},
\item
  \texttt{paop\_dbl\_p},
\item
  \texttt{paop\_dbl\_x},
\item
  \texttt{paop\_disp},
\item
  \texttt{paop\_hop2},
\item
  \texttt{paop\_hopdrag},
\item
  \texttt{hh\_termwise\_quadrature}.
\end{itemize}

Hence the final validated keep-set was simply \[
\mathcal G_{\mathrm{lean,weak}}
=
\{\texttt{uccsd\_sing},\ \texttt{uccsd\_dbl},\ \texttt{paop\_cloud\_p}\}.
\] A fresh rerun using only this class-filtered pool then produced \[
E_{\mathrm{lean}}=-0.7763497696348325,
\qquad
\lvert\Delta E\rvert_{\mathrm{lean}}=2.7761269927\times 10^{-5},
\qquad
F_{\mathrm{lean}}=0.9999861889437017,
\] at logical depth \(d_{\mathrm{lean}}=4\). Thus, to numerical
precision, \[
E_{\mathrm{lean}}\approx E_{\mathrm{parent}},
\qquad
\lvert\Delta E\rvert_{\mathrm{lean}}\approx \lvert\Delta E\rvert_{\mathrm{parent}},
\qquad
F_{\mathrm{lean}}\approx F_{\mathrm{parent}},
\] while the logical depth drops from \(7\) to \(4\).

A separate optimizer check on this same weak-coupling direct Phase-3
beam surface asked a narrower question: whether SPSA, still run
locally/noiselessly with the pre-transpile proxy-cost surface
(\texttt{phase3\_backend\_cost\_mode=proxy}) and still using the same
reduced class-filtered pool with full-window refits, could at least
reach the practical noise-threshold regime even if it did not exactly
recover the Powell basin. The best completed local SPSA artifact on that
surface was \[
E_{\mathrm{SPSA,local}}=-0.7763064877481185,
\qquad
\lvert\Delta E\rvert_{\mathrm{SPSA,local}}=7.1043156641\times 10^{-5},
\qquad
F_{\mathrm{SPSA,local}}=0.9999653825260001,
\] from the direct beam-enabled class-filtered run
\texttt{20260328\_hh\_l2\_u05\_g02\_local\_beam\_b3\_spsa\_case4\_direct\_m3200\_a0p1\_c0p02\_A5}.
This did not match the Powell/class-pruned reference value
\(2.7761269927\times 10^{-5}\), but it did enter the sub-\(10^{-4}\)
regime and therefore satisfied the more practical gate of crossing below
the anticipated noise floor. In that sense the local SPSA question at
weak coupling was answered positively: exact-basin recovery remained
unresolved, but threshold-level local recovery was achieved strongly
enough to justify a direct noisy/backend follow-on on the same
algorithmic surface.

A direct noisy transfer follow-on then kept the same reduced pool
\(\{\texttt{uccsd\_sing},\texttt{uccsd\_dbl},\texttt{paop\_cloud\_p}\}\),
the same beam/full-window \texttt{phase3\_v1} surface, and the same SPSA
recipe \((a,c,A,\texttt{maxiter})=(0.1,0.02,5.0,3200)\), but turned on
\texttt{backend\_scheduled} FakeMarrakesh scoring with
\texttt{expectation\_v1}, \texttt{512} shots, \texttt{1} repeat, and no
mitigation; the exact command is recorded at
\texttt{artifacts/agent\_runs/20260328\_hh\_l2\_u05\_g02\_phase3\_beam\_b3\_fake\_marrakesh\_spsa\_case4\_recipe\_attempt3\_supervised/logs/command.sh}.
During that supervised run, the best observed heartbeat reached \[
E_{\mathrm{SPSA,noisy,best}}=-0.7762934955401648,
\qquad
\lvert\Delta E\rvert_{\mathrm{SPSA,noisy,best}}=8.4035364595\times 10^{-5},
\] at depth \(95\), as preserved in
\texttt{artifacts/agent\_runs/20260328\_hh\_l2\_u05\_g02\_phase3\_beam\_b3\_fake\_marrakesh\_spsa\_case4\_recipe\_attempt3\_supervised/best\_so\_far\_snapshot.json}.
So, before final max-depth termination, the same reduced-pool direct
Phase-3 beam SPSA lane already crossed the practical noisy
sub-\(10^{-4}\) regime under transpiled FakeMarrakesh selection/scoring.

\hypertarget{live-four-run-supervised-status-block}{%
\subsubsection{19.3.2 Live four-run supervised status
block}\label{live-four-run-supervised-status-block}}

\begin{verbatim}
Objective<As of March 28, 2026 at 10:59:12 CDT, all 4 noisy
  reduced-pool backend_scheduled SPSA runs are still genuinely
  running with live child PIDs and no stderr in artifacts/
  agent_runs/20260328_hh_l2_u05_g02_phase3_beam_b3_fake_marrake
  sh_spsa_case4_recipe_attempt4_dropstop/progress.json,
  artifacts/
  agent_runs/20260328_hh_l2_u05_g02_phase3_beam_b3_fake_marrake
  sh_spsa_case4_recipe_attempt5_dropstop_lifetimeoff/
  progress.json, artifacts/
  agent_runs/20260328_hh_l2_u05_g02_phase3_beam_b4_c3_k4_fake_m
  arrakesh_spsa_case4_recipe_attempt6_dropstop_lifetimeoff/
  sh_spsa_case4_recipe_attempt7_dropstop_transpile_single/
  progress.json. Best observed noisy gaps so far are: attempt4
  3.3871153069520155e-05 at depth 76 in artifacts/
  agent_runs/20260328_hh_l2_u05_g02_phase3_beam_b3_fake_marrake
  sh_spsa_case4_recipe_attempt4_dropstop/logs/stdout.log,
  attempt6 3.733262477279009e-05 at depth 49 in artifacts/
  agent_runs/20260328_hh_l2_u05_g02_phase3_beam_b4_c3_k4_fake_m
  arrakesh_spsa_case4_recipe_attempt6_dropstop_lifetimeoff/
  logs/stdout.log, attempt7 3.899366517123859e-05 at depth 67
  in artifacts/
  agent_runs/20260328_hh_l2_u05_g02_phase3_beam_b3_fake_marrake
  sh_spsa_case4_recipe_attempt7_dropstop_transpile_single/logs/
  stdout.log, and attempt5 1.0855648531227224e-04 at depth 67
  in artifacts/
  agent_runs/20260328_hh_l2_u05_g02_phase3_beam_b3_fake_marrake
  sh_spsa_case4_recipe_attempt5_dropstop_lifetimeoff/logs/
  stdout.log.>
  Why/Intent<So the current result ordering is: baseline clean-
  stop attempt4 best, then wider-beam+lifetime-off attempt6,
  then transpile_single_v1 attempt7, with lifetime-off only
  attempt5 trailing. None has written a final JSON yet because
  the clean stop is intentionally gated by --adapt-drop-min-
  depth 96, so even though three runs are already well below
  the practical 1e-4 threshold, the normal drop_plateau exit
  cannot fire before depth 96.>
  Suggested Next step/how this fits into broader picture<I
  recommend we keep these running until the first one crosses
  depth 96 and exits normally, because that will give us the
  first clean final JSON on this surface without interrupting a
  good basin. If the ranking stays similar, attempt4 should
  finish first, while attempt6 and attempt7 already look
  scientifically competitive enough to justify keeping them
  alive as the main comparison axes.>
\end{verbatim}

The rejected removals show which families remained indispensable under
this contract:

\begin{itemize}
\tightlist
\item
  removing \texttt{paop\_cloud\_p} gave
  \(\lvert\Delta E\rvert\approx 1.0813\times 10^{-2}\) and
  \(F\approx 0.996632\),
\item
  removing \texttt{uccsd\_sing} gave
  \(\lvert\Delta E\rvert\approx 2.19934\) and \(F\approx 0.214233\),
\item
  removing \texttt{uccsd\_dbl} gave
  \(\lvert\Delta E\rvert\approx 1.31029\times 10^{-2}\) and
  \(F\approx 0.996691\).
\end{itemize}

So at this weaker-coupling \(L=2\) point the validated lean pool is even
smaller than the earlier heavy/lean \(U=4\), \(g=0.5\) line: the heavy
search may still visit quadrature seeds, but the final class-level
support needed to preserve the achieved energy/fidelity scale is just \[
\texttt{uccsd\_sing}\;\cup\;\texttt{uccsd\_dbl}\;\cup\;\texttt{paop\_cloud\_p}.
\]

For the backend-facing follow-up at this same weak-coupling point, the
Heron/Marrakesh optimization should be done as a \textbf{transpile-only}
sweep rather than through the noisy staged wrapper: run
\texttt{python\ -m\ pipelines.hardcoded.adapt\_circuit\_cost} on the
promoted locked scaffold
\texttt{hh\_l2\_u05\_g02\_full\_meta\_class\_pruned\_lean4\_locked\_scaffold\_v1.json}
over
\texttt{\textbackslash{}texttt\{optimization\textbackslash{}\_level\}\textbackslash{}in\textbackslash{}\{1,2\textbackslash{}\}}
and
\texttt{\textbackslash{}texttt\{seed\textbackslash{}\_transpiler\}\textbackslash{}in\textbackslash{}\{0,\textbackslash{}dots,9\textbackslash{}\}},
then rank by compiled two-qubit count, depth, and size. On
\texttt{\textbackslash{}texttt\{FakeMarrakesh\}} the winning setting was
\texttt{\textbackslash{}texttt\{optimization\textbackslash{}\_level\}=2},
\texttt{\textbackslash{}texttt\{seed\textbackslash{}\_transpiler\}=5},
reducing the same 4-operator, 6-runtime-term scaffold from about
\texttt{25} two-qubit gates, depth \texttt{63}, size \texttt{109} to
\texttt{18} two-qubit gates, depth \texttt{43}, size \texttt{83} without
changing the physics or operator order.

\hypertarget{relation-to-live-phase-3-selectorpruning-implementation}{%
\subsection{19.4 Relation to live Phase-3 selector/pruning
implementation}\label{relation-to-live-phase-3-selectorpruning-implementation}}

The canonical Phase-3 target contract is the one stated in
§§11.4.1--11.4.3. Current live repo code still lags that contract in two
visible ways.

First, the live selector path still applies a hard uncertainty-adjusted
gradient pre-cap before the cheap screen and still uses a raw
remaining-depth / remaining-evaluations proxy inside the lifetime-burden
term. The corrected canonical contract in §11.4.1 instead ranks the full
admissible candidate-position pool with the horizon-smoothed cheap score
\(S_{3,\mathrm{cheap}}(r;t)\), forms a horizon-sized coarse shortlist
\(\mathcal C_{\mathrm{cheap}}\), and uses the useful horizon \[
H_t=\min\!\bigl(D_{\mathrm{left}}(t),\widehat N_{\mathrm{rem}}(t)\bigr)
\] inside the lifetime multiplier. Here \(\widehat N_{\mathrm{rem}}(t)\)
is controller telemetry and therefore a ranking/burden aid, not a hard
admissibility theorem.

Second, the live pruning path still ranks removals primarily by small
amplitude together with weaker prior proxy-benefit tie-breaks and
accepts removals through a standalone post-refit regression threshold.
The corrected canonical contract in §11.4.3 instead separates prune
permissibility from local expendability, ranks permissible removals by
frozen-ablation loss, and uses retained admitted-step gain as the
principal energetic safety condition. Thus frozen ablation is to be read
as a local, solver-relative expendability heuristic, while the
gain-retention inequality is the actual prune safety condition.

So the honest status note is: current code still trails the canonical
Phase-3 mathematics, but the intended direction of alignment is \[
\text{live implementation}\longrightarrow \text{corrected canonical contract of §§11.4.1--11.4.3},
\] not the reverse.

\hypertarget{data-provenance-note}{%
\subsection{19.5 Data provenance note}\label{data-provenance-note}}

The novelty-ablation numbers in this chapter come from the fresh direct
\texttt{adapt\_pipeline} artifacts
\texttt{pi\_ablate\_pareto\_lean\_l2\_novelty\_on.json},
\texttt{pi\_ablate\_pareto\_lean\_l2\_novelty\_off.json},
\texttt{pi\_ablate\_full\_meta\_novelty\_on.json}, and
\texttt{pi\_ablate\_full\_meta\_novelty\_off.json}, all dated
\texttt{2026-03-24}. The heavy-versus-lean pool-family comparison is
anchored by the \texttt{2026-03-21} comparison bundle summarized in
\texttt{pareto\_lean\_l2\_vs\_heavy\_L2\_ecut1\_comparison.md} together
with the corresponding lean rerun artifact. The data are therefore
matched and current enough for PI-facing evidence.

The honesty caveat is the same one stated in the reporting notes: these
matched ablations share the same physics, seed, and selector family, but
they are not claimed to be byte-for-byte reproductions of every earlier
March 21/22 runtime-resolved default. That caveat does not weaken the
on/off comparisons above, because the resolved settings were shared
inside each matched pair.

\hypertarget{machine-agent-redo-note-for-the-successful-l2-pruning-session}{%
\subsection{19.6 Machine-agent redo note for the successful L=2 pruning
session}\label{machine-agent-redo-note-for-the-successful-l2-pruning-session}}

For the validated HH prune line, keep the problem fixed at \[
L=2,\qquad t=1,\qquad U=4,\qquad \omega_0=1,\qquad g=0.5,\qquad n_{\mathrm{ph,max}}=1,
\] with the saved fullhorse parent artifact as the starting point. A
machine agent can reproduce the session in three steps from repo root:

\begin{Shaded}
\begin{Highlighting}[]
\ExtensionTok{python} \AttributeTok{{-}m}\NormalTok{ pipelines.hardcoded.hh\_prune\_nighthawk }\DataTypeTok{\textbackslash{}}
  \AttributeTok{{-}{-}input{-}json}\NormalTok{ artifacts/json/hh\_backend\_adapt\_fullhorse\_powell\_20260322T194155Z\_fakenighthawk.json }\DataTypeTok{\textbackslash{}}
  \AttributeTok{{-}{-}mode}\NormalTok{ both }\DataTypeTok{\textbackslash{}}
  \AttributeTok{{-}{-}prune{-}threshold}\NormalTok{ 1e{-}4 }\DataTypeTok{\textbackslash{}}
  \AttributeTok{{-}{-}output{-}json}\NormalTok{ artifacts/json/hh\_prune\_nighthawk\_}\OperatorTok{\textless{}}\NormalTok{timestamp}\OperatorTok{\textgreater{}}\NormalTok{.json}

\ExtensionTok{python} \AttributeTok{{-}m}\NormalTok{ pipelines.hardcoded.hh\_prune\_nighthawk }\DataTypeTok{\textbackslash{}}
  \AttributeTok{{-}{-}mode}\NormalTok{ export\_fixed\_scaffolds }\DataTypeTok{\textbackslash{}}
  \AttributeTok{{-}{-}source{-}artifact{-}json}\NormalTok{ artifacts/json/hh\_prune\_nighthawk\_aggressive\_5op.json}

\ExtensionTok{python} \AttributeTok{{-}m}\NormalTok{ pipelines.hardcoded.adapt\_circuit\_cost }\DataTypeTok{\textbackslash{}}
  \AttributeTok{{-}{-}artifact{-}json}\NormalTok{ artifacts/json/hh\_prune\_nighthawk\_gate\_pruned\_7term.json }\DataTypeTok{\textbackslash{}}
  \AttributeTok{{-}{-}backend{-}name}\NormalTok{ ibm\_marrakesh }\DataTypeTok{\textbackslash{}}
  \AttributeTok{{-}{-}seed{-}transpiler}\NormalTok{ 7 }\AttributeTok{{-}{-}optimization{-}level}\NormalTok{ 1 }\DataTypeTok{\textbackslash{}}
  \AttributeTok{{-}{-}sweep{-}backends}
\end{Highlighting}
\end{Shaded}

The redo is not complete unless it records \textbf{fidelity as well as
energy}. For each surviving scaffold and exported fixed scaffold, save
at least: \texttt{abs\_delta\_e}, exact-state fidelity
(\texttt{exact\_state\_fidelity},
\texttt{local\_exact\_state\_fidelity}, or
\texttt{expected\_local\_exact\_state\_fidelity}), retained exyz labels,
backend, transpiler seed, optimization level, compiled
\texttt{count\_2q}, compiled depth, and compiled size. The validated
executable anchor is the \texttt{gate\_pruned\_7term} artifact, not
\texttt{circuit\_optimized\_7term}: the trustworthy target is \[
\lvert\Delta E\rvert\approx 4.23\times 10^{-4},\qquad F\approx 0.9998,
\] with Marrakesh compile cost near \texttt{25} two-qubit gates and
depth \texttt{63-75}. A leaner 6-term executable may reach about
\texttt{14} two-qubit gates and depth \texttt{48}, but its exact
regression is materially worse, near \[
\lvert\Delta E\rvert\approx 4.61\times 10^{-3}.
\]

\hypertarget{direct-noisy-backend_scheduled-l2-fastprobe-note}{%
\subsection{\texorpdfstring{19.7 Direct noisy
\texttt{backend\_scheduled} L=2 fastprobe
note}{19.7 Direct noisy backend\_scheduled L=2 fastprobe note}}\label{direct-noisy-backend_scheduled-l2-fastprobe-note}}

For the same HH point \[
L=2,\qquad t=1,\qquad U=4,\qquad \omega_0=1,\qquad g=0.5,\qquad n_{\mathrm{ph,max}}=1,
\] a useful data line is the direct local noisy \texttt{phase3\_v1}
fastprobe on \texttt{FakeNighthawk} with \texttt{backend\_scheduled}
oracle evaluation, \texttt{512} shots, \texttt{2} repeats, and optional
\texttt{mthree} readout mitigation. The March 26 command family is
recorded in
\texttt{artifacts/runstate/hh\_phase3\_v1\_local\_noisy\_backend\_scheduled\_fastprobe\_retry2\_20260326\_cmd.sh}
and the stabilized readout follow-up
\texttt{artifacts/runstate/hh\_phase3\_v1\_local\_noisy\_backend\_scheduled\_fastprobe\_retry3\_20260326\_cmd.sh},
with the no-readout comparator in
\texttt{artifacts/runstate/hh\_phase3\_v1\_local\_noisy\_backend\_scheduled\_noreadout\_retry1\_20260326/}.

These probes were intentionally weak
(\texttt{\textbackslash{}texttt\{adapt\textbackslash{}\_max\textbackslash{}\_depth\}=1})
and were stopped after the first scout evaluation timing split was
captured, so they are debug-quality runtime evidence rather than
convergence-quality VQE runs. The first measured \texttt{plus}
evaluation on the readout+mthree path gave \[
T_{\mathrm{total,readout}}\approx 37.04\,\mathrm{s},
\qquad
T_{\mathrm{term}}\approx 26.39\,\mathrm{s},
\qquad
T_{\mathrm{calib}}\approx 10.64\,\mathrm{s},
\] with readout-apply time negligible at about \texttt{0.01s}. The
no-readout comparator gave \[
T_{\mathrm{total,no\,readout}}\approx 26.84\,\mathrm{s},
\qquad
T_{\mathrm{term}}\approx 26.83\,\mathrm{s}.
\] So on this \(L=2\) HH debug line the dominant cost is
backend-scheduled term execution itself, while readout mitigation
contributes a secondary first-hit warmup rather than the main runtime
burden. The practical lesson is that \texttt{backend\_scheduled} is
honest enough for a very small late shortlist or confirmation pass, but
too expensive to act as the broad pre-shortlist scout surface.

\hypertarget{fixed-scaffold-compile-control-scout-note}{%
\subsection{19.8 Fixed-scaffold compile-control scout
note}\label{fixed-scaffold-compile-control-scout-note}}

For the locked Marrakesh/Heron 6-term HH scaffold
\texttt{artifacts/json/hh\_marrakesh\_fixed\_scaffold\_6term\_drop\_eyezee\_20260323T171528Z.json},
the useful local compile-control test is a transpile-only scout on
\texttt{FakeMarrakesh} with fixed circuit/parameters,
\texttt{shots=4096}, \texttt{oracle\_repeats=8}, readout/mthree
mitigation, and grid
\((\texttt{optimization\_level},\texttt{seed\_transpiler})\in\{1,2\}\times\{0,\dots,4\}\),
ranking by \(\Delta E=E_{\mathrm{noisy}}-E_{\mathrm{ideal}}\) then by
compiled cost; among the saved \texttt{8/10} completed candidates in
\texttt{artifacts/agent\_runs/20260328\_direct\_fixed\_scaffold\_compile\_scout\_fullaccess\_attempt2\_timeout7200/json/20260328\_direct\_fixed\_scaffold\_compile\_scout\_fullaccess\_attempt2\_timeout7200.json},
the best-so-far setting was \texttt{opt1\_seed4} with
\(\Delta E\approx 9.9899\times 10^{-2}\), two-qubit count \texttt{14},
and depth \texttt{48}, whereas the completed lower-depth
\texttt{opt2\_*} candidates kept the same two-qubit count \texttt{14}
but had depth \texttt{38} and worse \(\Delta E\):
\texttt{opt2\_seed0\textbackslash{}approx\ 1.1726\textbackslash{}times\ 10\^{}\{-1\}},
\texttt{opt2\_seed1\textbackslash{}approx\ 1.2245\textbackslash{}times\ 10\^{}\{-1\}},
\texttt{opt2\_seed2\textbackslash{}approx\ 1.2029\textbackslash{}times\ 10\^{}\{-1\}},
so lower compiled depth alone did not improve noisy accuracy and the
current hardware-facing compile preset is
\texttt{(optimization\_level=1,\ seed\_transpiler=4)}.

As an \(L=2\) run-planning prior rather than a theorem, the earlier
local FakeMarrakesh ablations were mixed between readout-only and twirl
while DD usually underperformed, but the completed 12-cell
saved-\(\theta^\star\) mitigation matrix in
\texttt{artifacts/json/20260328\_fixed\_scaffold\_saved\_theta\_mitigation\_matrix\_attempt5\_timeout86400\_delta\_logs\_fixed\_scaffold\_saved\_theta\_mitigation\_matrix.json}
now supersedes that partial prior: the top three cells were
\texttt{opt2\_seed0\_\_zne\_on\_\_twirl\_dd} with
\(\Delta E\approx-6.6431\times 10^{-4}\) and stderr
\(\approx1.2322\times 10^{-2}\) at 14 two-qubit gates and depth 38,
\texttt{opt1\_seed4\_\_zne\_on\_\_twirl\_dd} with
\(\Delta E\approx8.2990\times 10^{-3}\) and stderr
\(\approx2.6104\times 10^{-2}\) at 14 two-qubit gates and depth 48, and
\texttt{opt2\_seed0\_\_zne\_on\_\_twirl} with
\(\Delta E\approx1.4807\times 10^{-2}\) and stderr
\(\approx1.1157\times 10^{-2}\) at 14 two-qubit gates and depth 38. For
a real-QPU first pass on this locked 6-term scaffold, the current best
local starting lane is therefore
\texttt{(optimization\_level=2,\ seed\_transpiler=0)} with
\texttt{ZNE\ on\ +\ twirl\ +\ DD}, with
\texttt{(optimization\_level=1,\ seed\_transpiler=4)} plus the same
mitigation stack as the nearest fallback if one wants to hedge against
ZNE-fit variance while staying on the same two-qubit count.

For the weak-coupling locked scaffold
\texttt{artifacts/json/useful/L2/hh\_l2\_u05\_g02\_full\_meta\_class\_pruned\_lean4\_locked\_scaffold\_v1.json},
the completed 4-cell readout+mthree shortlist in
\texttt{artifacts/json/20260329\_weak\_fixed\_scaffold\_top4\_mitigation\_readout\_opt2seed5\_fixed\_scaffold\_saved\_theta\_mitigation\_matrix.json}
had top two cells \texttt{opt2\_seed5\_\_zne\_on\_\_twirl\_dd} with
\(\Delta E\approx1.0800\times 10^{-2}\) and stderr
\(\approx4.7014\times 10^{-3}\), and
\texttt{opt2\_seed5\_\_zne\_on\_\_dd} with
\(\Delta E\approx1.6984\times 10^{-2}\) and stderr
\(\approx1.9163\times 10^{-3}\), both at 18 two-qubit gates and depth
43.

The later readout-on versus readout-off ablations showed that base
readout mitigation was helping rather than hurting on both
fixed-scaffold winner lines. For the strong-coupling local winner, the
same \texttt{opt2\_seed0\_\_zne\_on\_\_twirl\_dd} lane moved from
\(\Delta E\approx-6.6431\times 10^{-4}\) with readout on in
\texttt{artifacts/json/20260328\_fixed\_scaffold\_saved\_theta\_mitigation\_matrix\_attempt5\_timeout86400\_delta\_logs\_fixed\_scaffold\_saved\_theta\_mitigation\_matrix.json}
to \(\Delta E\approx5.4456\times 10^{-2}\) with readout off in
\texttt{artifacts/json/20260329\_strong\_fixed\_scaffold\_winner\_readout\_off\_ablation\_opt2seed0\_twirl\_dd\_retry2\_fixed\_scaffold\_saved\_theta\_mitigation\_matrix.json},
at the same compiled cost (14 two-qubit gates, depth 38). For the
weak-coupling 18-two-qubit compile \texttt{opt2\_seed5}, the analogous
readout-off follow-up in
\texttt{artifacts/json/20260330\_weak\_fixed\_scaffold\_top2\_readout\_off\_ablation\_opt2seed5\_fixed\_scaffold\_saved\_theta\_mitigation\_matrix.json}
gave \texttt{opt2\_seed5\_\_zne\_on\_\_twirl\_dd} at
\(\Delta E\approx5.6830\times 10^{-2}\) and
\texttt{opt2\_seed5\_\_zne\_on\_\_dd} at
\(\Delta E\approx5.8131\times 10^{-2}\), versus the readout-on values
\(\approx1.0800\times 10^{-2}\) and \(\approx1.6984\times 10^{-2}\)
respectively, so turning readout off worsened the best weak lanes by
about \(4.60\times10^{-2}\) and \(4.11\times10^{-2}\).

The surprising regime comparison is therefore not in the noiseless
scaffold bias but in the noisy response. The strong locked scaffold had
\(E_{\mathrm{ideal}}-E_{\mathrm{exact}}\approx4.6131\times10^{-3}\) from
\texttt{artifacts/json/hh\_marrakesh\_fixed\_scaffold\_6term\_drop\_eyezee\_20260323T171528Z.json},
while the weak locked scaffold had
\(E_{\mathrm{ideal}}-E_{\mathrm{exact}}\approx2.7761\times10^{-5}\) from
\texttt{artifacts/json/useful/L2/hh\_l2\_u05\_g02\_full\_meta\_class\_pruned\_lean4\_locked\_scaffold\_v1.json}.
Thus the weak scaffold was much closer to exact in the noiseless sense,
yet still produced a worse noisy \(\Delta E\) after the current
suppression stack; the deficit is therefore a noise-response /
mitigation-effect issue rather than a larger ideal-state approximation
error.

\hypertarget{heron-transpile-pareto-front-and-noisy-replay-for-fixed-l2-scaffolds}{%
\subsection{19.9 Heron transpile Pareto front and noisy replay for fixed
L=2
scaffolds}\label{heron-transpile-pareto-front-and-noisy-replay-for-fixed-l2-scaffolds}}

All scaffolds share \(L=2\), \(n_{\mathrm{ph,max}}=1\), binary bosons,
blocked ordering, open boundaries. Each was transpiled to
\texttt{FakeMarrakesh} over
\((\texttt{optimization\_level},\texttt{seed\_transpiler})\in\{1,2\}\times\{0,\dots,9\}\);
the table reports the best compiled two-qubit count per scaffold.
Term-pruned variants start from the 7-term locked scaffold and drop
Pauli rotation terms, then reoptimize \(\theta\) via Powell.

\begin{longtable}[]{@{}
  >{\raggedright\arraybackslash}p{(\columnwidth - 12\tabcolsep) * \real{0.12}}
  >{\raggedleft\arraybackslash}p{(\columnwidth - 12\tabcolsep) * \real{0.16}}
  >{\raggedleft\arraybackslash}p{(\columnwidth - 12\tabcolsep) * \real{0.16}}
  >{\raggedleft\arraybackslash}p{(\columnwidth - 12\tabcolsep) * \real{0.16}}
  >{\raggedleft\arraybackslash}p{(\columnwidth - 12\tabcolsep) * \real{0.16}}
  >{\raggedright\arraybackslash}p{(\columnwidth - 12\tabcolsep) * \real{0.12}}
  >{\raggedright\arraybackslash}p{(\columnwidth - 12\tabcolsep) * \real{0.12}}@{}}
\toprule
Scaffold & \(\lvert\Delta E\rvert\) & CZ gates & Depth & \(F\) &
Coupling & Source \\
\midrule
\endhead
\texttt{weak\_lean4\_locked} (4 ops) & \(2.78\times 10^{-5}\) & 18 & 43
& 0.9999 & \(U{=}0.5,g{=}0.2\) & class-prune \\
\texttt{aggressive\_5op} (12 terms) & \(5.62\times 10^{-5}\) & 26 & 72 &
--- & \(U{=}4,g{=}0.5\) & op-prune \\
\texttt{gate\_pruned\_7term} (7 terms) & \(4.23\times 10^{-4}\) & 21 &
51 & 0.9998 & \(U{=}4,g{=}0.5\) & term-prune \\
\texttt{6term\_drop\_eyezee} (6 terms) & \(4.61\times 10^{-3}\) & 14 &
38 & 0.9985 & \(U{=}4,g{=}0.5\) & term-prune \\
\texttt{6term\_drop\_eyeeez} (6 terms) & \(4.61\times 10^{-3}\) & 14 &
56 & 0.9985 & \(U{=}4,g{=}0.5\) & term-prune \\
\texttt{5term\_drop\_both\_disp} (5 terms) & \(6.92\times 10^{-3}\) & 12
& 48 & 0.9986 & \(U{=}4,g{=}0.5\) & term-prune \\
\bottomrule
\end{longtable}

The Pareto front over \((\lvert\Delta E\rvert,\,\text{CZ count})\)
contains two points when both coupling regimes are pooled: \[
\bigl(2.78\times 10^{-5},\;18\;\text{CZ}\bigr)_{\text{weak}},
\qquad
\bigl(4.61\times 10^{-3},\;14\;\text{CZ}\bigr)_{\text{strong}}.
\] Within strong coupling alone, the three-point front is
\(\{(5.62\times 10^{-5},26),\;(4.61\times 10^{-3},14),\;(6.92\times 10^{-3},12)\}\).
The dominated scaffolds (\texttt{pruned\_scaffold\_11op} at 51 CZ,
\texttt{ultra\_lean\_6op}/\texttt{readapt\_6op} at 38 CZ,
\texttt{gate\_pruned\_7term} at 21 CZ) offer no accuracy benefit
relative to cheaper alternatives.

\textbf{Noisy replay on \texttt{FakeMarrakesh}
(\texttt{backend\_scheduled})} for the strong-coupling Pareto-optimal
6-term scaffold (14 CZ, depth 38, \texttt{optimization\_level=2},
\texttt{seed\_transpiler=0}):

\begin{longtable}[]{@{}lrr@{}}
\toprule
Configuration & \(E_{\mathrm{noisy}}\) &
\(E_{\mathrm{noisy}}-E_{\mathrm{exact}}\) \\
\midrule
\endhead
Saved \(\theta^*\) + readout/mthree & \(0.2773\) & \(+0.119\) \\
Saved \(\theta^*\) + readout + gate twirling & \(0.2208\) &
\(+0.062\) \\
SPSA-128 best + readout + gate twirling & \(0.1368\) & \(-0.022\) \\
Saved \(\theta^*\) ideal (noiseless) & \(0.1633\) & \(+0.005\) \\
Exact GS & \(0.1587\) & \(0\) \\
\bottomrule
\end{longtable}

Mitigation stack: readout/mthree plus local 2Q Pauli twirling, symmetry
\texttt{verify\_only}, no ZNE, no DD. SPSA: 128 iterations, 1024 shots,
1 repeat, seed 19. The SPSA-optimized noisy energy undershoots exact by
\(0.022\), indicating noise-induced bias below the variational floor;
the ideal energy at the SPSA-best \(\theta\) is \(0.189\), confirming
parameter drift. At the noiseless-optimal \(\theta^*\), gate twirling
halves the noise bias from \(0.119\) to \(0.062\).

Reproduction:

\begin{Shaded}
\begin{Highlighting}[]
\ExtensionTok{python} \AttributeTok{{-}m}\NormalTok{ pipelines.exact\_bench.heron\_pareto\_sweep }\AttributeTok{{-}{-}backend}\NormalTok{ FakeMarrakesh }\AttributeTok{{-}{-}coupling}\NormalTok{ both}
\ExtensionTok{python} \AttributeTok{{-}m}\NormalTok{ pipelines.hardcoded.hh\_staged\_noise }\DataTypeTok{\textbackslash{}}
  \AttributeTok{{-}{-}L}\NormalTok{ 2 }\AttributeTok{{-}{-}t}\NormalTok{ 1.0 }\AttributeTok{{-}{-}u}\NormalTok{ 4.0 }\AttributeTok{{-}{-}g{-}ep}\NormalTok{ 0.5 }\AttributeTok{{-}{-}omega0}\NormalTok{ 1.0 }\AttributeTok{{-}{-}n{-}ph{-}max}\NormalTok{ 1 }\DataTypeTok{\textbackslash{}}
  \AttributeTok{{-}{-}ordering}\NormalTok{ blocked }\AttributeTok{{-}{-}boundary}\NormalTok{ open }\AttributeTok{{-}{-}boson{-}encoding}\NormalTok{ binary }\DataTypeTok{\textbackslash{}}
  \AttributeTok{{-}{-}include{-}fixed{-}scaffold{-}noisy{-}replay} \DataTypeTok{\textbackslash{}}
  \AttributeTok{{-}{-}fixed{-}final{-}state{-}json}\NormalTok{ artifacts/json/hh\_marrakesh\_fixed\_scaffold\_6term\_drop\_eyezee\_20260323T171528Z.json }\DataTypeTok{\textbackslash{}}
  \AttributeTok{{-}{-}use{-}fake{-}backend} \AttributeTok{{-}{-}backend{-}name}\NormalTok{ FakeMarrakesh }\DataTypeTok{\textbackslash{}}
  \AttributeTok{{-}{-}fixed{-}scaffold{-}runtime{-}transpile{-}optimization{-}level}\NormalTok{ 2 }\DataTypeTok{\textbackslash{}}
  \AttributeTok{{-}{-}fixed{-}scaffold{-}runtime{-}seed{-}transpiler}\NormalTok{ 0 }\DataTypeTok{\textbackslash{}}
  \AttributeTok{{-}{-}final{-}method}\NormalTok{ SPSA }\AttributeTok{{-}{-}final{-}maxiter}\NormalTok{ 128 }\AttributeTok{{-}{-}final{-}seed}\NormalTok{ 19 }\DataTypeTok{\textbackslash{}}
  \AttributeTok{{-}{-}mitigation}\NormalTok{ readout }\AttributeTok{{-}{-}local{-}readout{-}strategy}\NormalTok{ mthree }\AttributeTok{{-}{-}local{-}gate{-}twirling} \DataTypeTok{\textbackslash{}}
  \AttributeTok{{-}{-}symmetry{-}mitigation{-}mode}\NormalTok{ verify\_only }\DataTypeTok{\textbackslash{}}
  \AttributeTok{{-}{-}shots}\NormalTok{ 1024 }\AttributeTok{{-}{-}oracle{-}repeats}\NormalTok{ 1 }\AttributeTok{{-}{-}noise{-}seed}\NormalTok{ 7 }\DataTypeTok{\textbackslash{}}
  \AttributeTok{{-}{-}smoke{-}test{-}intentionally{-}weak} \AttributeTok{{-}{-}skip{-}pdf}
\end{Highlighting}
\end{Shaded}

Artifacts:
\texttt{artifacts/json/heron\_pareto\_front\_20260327T223120Z.json},
\texttt{artifacts/json/hh\_6term\_pareto\_readout\_twirl\_fast\_20260328.json}.

\hypertarget{appendix-a.-drive-exact-propagation-and-propagator-semantics}{%
\section{Appendix A. Drive, Exact Propagation, and Propagator
Semantics}\label{appendix-a.-drive-exact-propagation-and-propagator-semantics}}

\hypertarget{a.1-static-drive-labels-and-time-dependent-coefficients}{%
\subsection{A.1 Static drive labels and time-dependent
coefficients}\label{a.1-static-drive-labels-and-time-dependent-coefficients}}

Let \(\mathcal L_{\mathrm{drive}}\) denote the fixed set of density
operators appearing in the drive. Then the drive always has the form \[
\hat H_{\mathrm{drive}}(t)=\sum_{\lambda\in\mathcal L_{\mathrm{drive}}} c_\lambda(t) P_\lambda,
\] with fixed labels \(P_\lambda\) and time-dependent coefficients
\(c_\lambda(t)\).

\hypertarget{a.2-reference-coefficient-map}{%
\subsection{A.2 Reference coefficient
map}\label{a.2-reference-coefficient-map}}

For a site-density drive, \[
\Delta c[Z_{p(i,\uparrow)}](t)=\Delta c[Z_{p(i,\downarrow)}](t)=-\frac{1}{2}v_i(t),
\] and the identity shift is \[
\Delta c[I](t)=\sum_i v_i(t)
\] if one chooses to retain the scalar offset explicitly.

\hypertarget{a.3-propagator-surfaces}{%
\subsection{A.3 Propagator surfaces}\label{a.3-propagator-surfaces}}

A common propagator menu is:

\begin{itemize}
\tightlist
\item
  second-order Suzuki,
\item
  piecewise exact propagation.
\end{itemize}

\hypertarget{a.4-reference-propagator-versus-reported-trajectory-propagator}{%
\subsection{A.4 Reference propagator versus reported trajectory
propagator}\label{a.4-reference-propagator-versus-reported-trajectory-propagator}}

Write the reported propagator as \(U_{\mathrm{traj}}(t)\) and the
high-accuracy reference propagator as \(U_{\mathrm{ref}}(t)\). These
need not be the same approximation.

\hypertarget{a.4.1-static-reference-branch}{%
\subsubsection{A.4.1 Static reference
branch}\label{a.4.1-static-reference-branch}}

In the static case, one may take \[
U_{\mathrm{ref}}(t)=e^{-it\hat H}.
\]

\hypertarget{a.4.2-drive-enabled-reference-branch}{%
\subsubsection{A.4.2 Drive-enabled reference
branch}\label{a.4.2-drive-enabled-reference-branch}}

For time-dependent \(\hat H(t)\), a refined piecewise-constant reference
is \[
U_{\mathrm{ref}}(t_{n+1},t_n) \approx e^{-i\Delta t\hat H(t_n^*)},
\] with \(t_n^*\) sampled on a refined grid and the full reference
obtained as an ordered product over the refined slices.

\hypertarget{a.5-filtered-exact-manifold-and-subspace-fidelity}{%
\subsection{A.5 Filtered exact manifold and subspace
fidelity}\label{a.5-filtered-exact-manifold-and-subspace-fidelity}}

Let \(\Pi_{\mathrm{sector}}\) denote the projector onto the target
symmetry sector. Then the filtered exact energy is \[
E_{\mathrm{exact,filtered}} = \min_{\psi\in\Pi_{\mathrm{sector}}\mathcal H} \langle \psi|\hat H|\psi\rangle,
\] and a sector-aware fidelity may be defined by \[
F_{\mathrm{sector}}(t) = \frac{|\langle \psi_{\mathrm{exact}}(t)|\Pi_{\mathrm{sector}}|\psi(t)\rangle|^2}{\langle \psi(t)|\Pi_{\mathrm{sector}}|\psi(t)\rangle}.
\]

\hypertarget{a.7-branch-resolved-observable-contracts}{%
\subsection{A.7 Branch-resolved observable
contracts}\label{a.7-branch-resolved-observable-contracts}}

If several trajectory branches are carried simultaneously, each
observable should be indexed by branch label \(b\): \[
\mathcal O_b(t),
\qquad
b\in\{\mathrm{var},\mathrm{traj},\mathrm{ref},\mathrm{noisy},\mathrm{ideal}\}.
\] Typical observables are energy, density, doublon density, and
fidelity.

\hypertarget{a.8-state-import-and-handoff-contract}{%
\subsection{A.8 State import and handoff
contract}\label{a.8-state-import-and-handoff-contract}}

A portable state bundle can be written abstractly as \[
\mathcal B = (\Lambda,\theta,a,E,\mathcal C),
\] where \(\Lambda\) is the physical manifest, \(\theta\) are
variational parameters when available, \(a\) is an amplitude map for the
working state, \(E\) stores energy summaries, and \(\mathcal C\) is
continuation metadata.

\hypertarget{a.9-highest-tier-continuation-surfaces}{%
\subsection{A.9 Highest-tier continuation
surfaces}\label{a.9-highest-tier-continuation-surfaces}}

At the highest continuation tier, \(\mathcal C\) may include split-event
summaries, motif data, symmetry diagnostics, replay policy data, and
limited branch histories. These are naturally viewed as auxiliary
mathematical metadata rather than as part of the operator algebra
itself.

\hypertarget{appendix-b.-continuation-and-handoff-contract}{%
\section{Appendix B. Continuation and Handoff
Contract}\label{appendix-b.-continuation-and-handoff-contract}}

\hypertarget{b.1-handoff-state-bundle}{%
\subsection{B.1 Handoff state bundle}\label{b.1-handoff-state-bundle}}

A handoff bundle may contain:

\begin{itemize}
\tightlist
\item
  the parameter manifest \(\Lambda\),
\item
  the current energy summary,
\item
  the normalized statevector amplitudes,
\item
  an exact comparison target,
\item
  optional continuation metadata.
\end{itemize}

One convenient symbolic amplitude map is \[
\{b_{N_q-1}\cdots b_0 \mapsto (\Re a_b,\Im a_b)\}.
\]

\hypertarget{b.2-continuation-block}{%
\subsection{B.2 Continuation block}\label{b.2-continuation-block}}

The continuation payload can be treated abstractly as \[
\mathcal C = (m,\mathcal S,\mathcal M_{\mathrm{opt}},\Gamma,\Sigma,\mathcal L,\mathcal U,\Xi,\mathcal R,\mathcal P),
\] where:

\begin{itemize}
\tightlist
\item
  \(m\) is the continuation mode,
\item
  \(\mathcal S\) is inherited manifold-structure data,
\item
  \(\mathcal M_{\mathrm{opt}}\) is auxiliary optimizer state,
\item
  \(\Gamma\) is selected-generator metadata,
\item
  \(\Sigma\) is split-event data,
\item
  \(\mathcal L\) is a motif library,
\item
  \(\mathcal U\) is motif usage,
\item
  \(\Xi\) is symmetry information,
\item
  \(\mathcal R\) is rescue or recovery history,
\item
  \(\mathcal P\) is optional historical replay-policy metadata when such
  provenance is still being carried.
\end{itemize}

When controller forecast telemetry is persisted, quantities such as
\(D_{\mathrm{left}}(t)\), \(\widehat N_{\mathrm{rem}}(t)\), and \(H_t\)
belong naturally to this policy/provenance layer. They remain heuristic
controller state rather than operator-level observables.

\hypertarget{b.3-general-handoff-map}{%
\subsection{B.3 General handoff map}\label{b.3-general-handoff-map}}

Let \(\mathcal H_{A\to B}\) be a handoff map from source manifold
\(\mathcal M_A\) to target manifold \(\mathcal M_B\): \[
\mathcal H_{A\to B}: (|\psi_A\rangle,\theta_A,\mathcal C_A) \mapsto (|\psi_B^{(0)}\rangle,\theta_B^{(0)},\mathcal C_B).
\] The map is intentionally abstract here: it records state,
coordinates, and continuation payload transfer between manifolds without
committing the main body to any obsolete replay-specific burn-in policy.

\hypertarget{b.4-symmetry-and-tier-three-payload-note}{%
\subsection{B.4 Symmetry and tier-three payload
note}\label{b.4-symmetry-and-tier-three-payload-note}}

Symmetry information belongs in the handoff if later stages will use it
for validation, postselection, or projector-renormalized estimation.
Mathematically, the carried object is a projector or a list of commuting
symmetry generators.

\hypertarget{b.5-historical-warmrefine-chain-note}{%
\subsection{B.5 Historical warm/refine chain
note}\label{b.5-historical-warmrefine-chain-note}}

The obsolete warm-start \(\to\) refine \(\to\) ADAPT \(\to\) replay path
is kept only in Appendix A.2 and is not part of the active main-body
symbolic contract.

\hypertarget{appendix-c.-cross-checks-and-exact-benchmark-contracts}{%
\section{Appendix C. Cross-Checks and Exact-Benchmark
Contracts}\label{appendix-c.-cross-checks-and-exact-benchmark-contracts}}

\hypertarget{c.1-trial-matrix-by-problem-family}{%
\subsection{C.1 Trial matrix by problem
family}\label{c.1-trial-matrix-by-problem-family}}

For the Hubbard model, a natural benchmark trial set is \[
\mathcal T_{\mathrm{Hub}} = \{\text{layerwise},\;\text{excitation-based},\;\text{adaptive(exc)},\;\text{adaptive(full-H)}\}.
\] For the HH model, the canonical sector-preserving core set is \[
\mathcal T_{\mathrm{HH}} = \{\text{HH-layerwise},\;\text{HH-physical-termwise},\;\text{adaptive(full-H)}\}.
\] A seed-surface extension can enlarge \(\mathcal T_{\mathrm{HH}}\) by
adding polaronic refinements built on the same sector-preserving
fixed-family surface.

\hypertarget{c.2-auto-scaled-parameter-resolution}{%
\subsection{C.2 Auto-scaled parameter
resolution}\label{c.2-auto-scaled-parameter-resolution}}

Let \[
S_{\mathrm{trot}}(L,n_{\mathrm{ph,max}}),
\quad
R_{\mathrm{restart}}(L,n_{\mathrm{ph,max}}),
\quad
I_{\max}(L,n_{\mathrm{ph,max}})
\] be monotone nondecreasing control maps for Trotter steps, restart
count, and optimizer effort. The key contract is monotonicity with
problem size, not a single universal constant setting.

\hypertarget{c.3-exact-target-and-shared-reference-state-construction}{%
\subsection{C.3 Exact target and shared reference-state
construction}\label{c.3-exact-target-and-shared-reference-state-construction}}

Within a benchmark matrix, all trial families should be compared against
the same exact target and the same conserved-sector reference
construction. Thus every trial is judged against the same
\(E_{\mathrm{exact,sector}}\) and the same physical Hamiltonian
instance.

\hypertarget{c.4-per-trial-energy-and-trajectory-semantics}{%
\subsection{C.4 Per-trial energy and trajectory
semantics}\label{c.4-per-trial-energy-and-trajectory-semantics}}

Each trial should at minimum record:

\begin{itemize}
\tightlist
\item
  the variational ground-state energy,
\item
  the exact comparison target,
\item
  trajectory observables on a shared time grid when dynamics are
  studied,
\item
  fidelity or overlap relative to a common reference branch.
\end{itemize}

\hypertarget{c.5-artifact-contracts}{%
\subsection{C.5 Artifact contracts}\label{c.5-artifact-contracts}}

Every report artifact should begin with a concise parameter manifest
listing:

\begin{itemize}
\tightlist
\item
  model family,
\item
  ansatz family,
\item
  drive enabled or disabled,
\item
  \(t\), \(U\), \(\omega_0\), \(g\), and \(n_{\mathrm{ph,max}}\),
\item
  any propagation parameters needed to reproduce the result.
\end{itemize}

\hypertarget{c.6-hh-seed-surface-preset-and-resource-proxies}{%
\subsection{C.6 HH seed-surface preset and resource
proxies}\label{c.6-hh-seed-surface-preset-and-resource-proxies}}

A seed-surface comparison is naturally evaluated not only by energy but
also by simple resource proxies such as \[
C_{\mathrm{2q}},
\qquad
D_{\mathrm{proxy}},
\qquad
C_{\mathrm{1q}}.
\] One can then compare improvement per marginal resource, \[
\frac{\Delta E}{C_{\mathrm{2q}}},
\qquad
\frac{\Delta E}{D_{\mathrm{proxy}}}.
\]

\hypertarget{appendix-d.-noise-validation-symmetry-mitigation-and-parity-contracts}{%
\section{Appendix D. Noise Validation, Symmetry-Mitigation, and Parity
Contracts}\label{appendix-d.-noise-validation-symmetry-mitigation-and-parity-contracts}}

\hypertarget{d.1-filtered-versus-full-exact-targets}{%
\subsection{D.1 Filtered versus full exact
targets}\label{d.1-filtered-versus-full-exact-targets}}

The filtered target is \[
E_{\mathrm{exact,filtered}} = \min_{\psi\in\mathcal H_{\mathrm{sector}}} \langle \psi|\hat H|\psi\rangle,
\] while the unrestricted target is \[
E_{\mathrm{exact,full}} = \min \operatorname{spec}(\hat H).
\] The two coincide only if the physical and computational sectors match
exactly.

\hypertarget{d.2-noisy-minus-ideal-observable-deltas}{%
\subsection{D.2 Noisy-minus-ideal observable
deltas}\label{d.2-noisy-minus-ideal-observable-deltas}}

For any observable \(\mathcal O(t)\), define \[
\Delta \mathcal O(t) = \mathcal O_{\mathrm{noisy}}(t)-\mathcal O_{\mathrm{ideal}}(t).
\] In particular, \[
\Delta E(t)=E_{\mathrm{noisy}}(t)-E_{\mathrm{ideal}}(t),
\qquad
\Delta D(t)=D_{\mathrm{noisy}}(t)-D_{\mathrm{ideal}}(t).
\]

\hypertarget{d.3-anchor-parity-gate}{%
\subsection{D.3 Anchor parity gate}\label{d.3-anchor-parity-gate}}

If a new idealized path is compared against a fixed anchor trajectory, a
strict parity gate takes the form \[
\max_{t\in\mathcal T}\,|\mathcal O_{\mathrm{new}}(t)-\mathcal O_{\mathrm{anchor}}(t)| \le \varepsilon_{\mathrm{parity}}.
\] This is an anchor-based equality test, not merely a qualitative
agreement statement.

\hypertarget{d.4-same-anchor-paired-comparison-semantics}{%
\subsection{D.4 Same-anchor paired comparison
semantics}\label{d.4-same-anchor-paired-comparison-semantics}}

Two noisy methods should be compared against the same ideal anchor, on
the same time grid, with the same initial state, so that differences can
be attributed to the noise or mitigation model rather than to different
starting conditions.

\hypertarget{d.5-symmetry-mitigation-mode-surface}{%
\subsection{D.5 Symmetry-mitigation mode
surface}\label{d.5-symmetry-mitigation-mode-surface}}

Three increasingly strong symmetry surfaces are:

\begin{enumerate}
\def\labelenumi{\arabic{enumi}.}
\tightlist
\item
  verify-only sector diagnostics,
\item
  bitstring postselection,
\item
  projector-renormalized estimators.
\end{enumerate}

If \(\Pi\) is the projector onto the target symmetry sector, the
renormalized estimator is \[
\langle O\rangle_{\Pi} = \frac{\langle \psi|\Pi O \Pi|\psi\rangle}{\langle \psi|\Pi|\psi\rangle}.
\]

\hypertarget{d.6-parameter-manifest-and-provenance-digest}{%
\subsection{D.6 Parameter manifest and provenance
digest}\label{d.6-parameter-manifest-and-provenance-digest}}

Even in a symbolic setting, it is useful to view each artifact as
carrying a manifest \(\mathfrak M\) and a provenance digest
\(h(\mathfrak M)\). The purpose is not algebraic; it is to guarantee
that the reported mathematics can be traced back to a unique physical
parameter set.

\hypertarget{appendix-e.-spec-only-or-not-yet-formalized-surfaces-kept-out-of-the-main-body}{%
\section{Appendix E. Spec-only or not-yet-formalized surfaces kept out
of the main
body}\label{appendix-e.-spec-only-or-not-yet-formalized-surfaces-kept-out-of-the-main-body}}

\hypertarget{e.1-broader-continuation-theory-items}{%
\subsection{E.1 Broader continuation-theory
items}\label{e.1-broader-continuation-theory-items}}

Several topics belong naturally in a later symbolic pass rather than in
the present core manuscript:

\begin{itemize}
\tightlist
\item
  multi-source motif transfer and motif tiling,
\item
  richer split-event formalism,
\item
  more detailed rescue dynamics,
\item
  exact handoff-policy closure beyond the present abstract maps,
\item
  measured-geometry and backend-specific closure beyond the exact
  symbolic statevector surface.
\end{itemize}

\hypertarget{e.2-obsolete-historical-warm-start-to-refine-to-adapt-to-replay-path}{%
\subsection{\texorpdfstring{E.2 Obsolete historical warm-start \(\to\)
refine \(\to\) ADAPT \(\to\) replay
path}{E.2 Obsolete historical warm-start \textbackslash to refine \textbackslash to ADAPT \textbackslash to replay path}}\label{e.2-obsolete-historical-warm-start-to-refine-to-adapt-to-replay-path}}

The historical staged path can be written symbolically as \[
\begin{aligned}
\mathcal M_{\mathrm{warm}}
&=\{\,U_{\mathrm{warm}}(\theta_{\mathrm{w}})|\phi_0\rangle:\theta_{\mathrm{w}}\in\mathbb R^{m_{\mathrm{w}}}\,\},\\
\mathcal M_{\mathrm{refine}}
&=\{\,U_{\mathrm{refine}}(\theta_{\mathrm{r}})U_{\mathrm{warm}}(\theta_{\mathrm{w}}^*)|\phi_0\rangle:\theta_{\mathrm{r}}\in\mathbb R^{m_{\mathrm{r}}}\,\},\\
\mathcal M_{\mathrm{replay}}(\mathcal O_*)
&=\{\,U_{\mathrm{replay}}(\varphi;\mathcal O_*)|\phi_0\rangle:\varphi\in\mathbb R^{m_{\mathrm{rep}}}\,\}.
\end{aligned}
\] The corresponding historical handoff chain is \[
\begin{aligned}
(\theta_{\mathrm{w}}^*,|\psi_{\mathrm{warm}}\rangle)
&\in\mathcal M_{\mathrm{warm}}
\mapsto
(\theta_{\mathrm{r}}^*,|\psi_{\mathrm{refine}}\rangle)
\in\mathcal M_{\mathrm{refine}}\\
&\mapsto
(\mathcal O_*,\theta_*^{\mathrm{adapt}},|\psi_*\rangle)\\
&\mapsto
(\varphi^{(0)},|\psi_{\mathrm{replay}}^{(0)}\rangle)
=
\mathcal H_{\mathrm{replay}}(|\psi_*\rangle,\theta_*^{\mathrm{adapt}},\mathcal O_*,\mathcal C_*)
\in\mathcal M_{\mathrm{replay}}(\mathcal O_*).
\end{aligned}
\] A typical replay-style historical policy is:

\begin{enumerate}
\def\labelenumi{\arabic{enumi}.}
\tightlist
\item
  preserve inherited scaffold parameters,
\item
  initialize residual parameters at zero,
\item
  optionally freeze inherited coordinates for a burn-in stage,
\item
  then unfreeze and refit in a controlled window.
\end{enumerate}

This path is kept for archival provenance only and is not part of the
active main-body symbolic contract.

\hypertarget{e.3-drive-amplitude-comparison-path}{%
\subsection{E.3 Drive-amplitude comparison
path}\label{e.3-drive-amplitude-comparison-path}}

A drive-amplitude comparison path can be expressed symbolically by
comparing several amplitudes \(A_0,A_1,\dots\) at fixed Hamiltonian
parameters and examining \[
\Delta E_A(t),\qquad \Delta D_A(t),\qquad F_A(t)
\] across the amplitude family. The formal structure is clear even if
the exact report layout is deferred.

\hypertarget{e.4-appendix-rule}{%
\subsection{E.4 Appendix rule}\label{e.4-appendix-rule}}

The rule of this symbolic manuscript is therefore:

\begin{enumerate}
\def\labelenumi{\arabic{enumi}.}
\tightlist
\item
  main body = closed mathematical definitions and comparison contracts,
\item
  appendix = promising but not yet fully formalized extensions.
\end{enumerate}

\hypertarget{references}{%
\section{References}\label{references}}

\begin{enumerate}
\def\labelenumi{\arabic{enumi}.}
\item
  Grimsley, Economou, Barnes, and Mayhall, An adaptive variational
  algorithm for exact molecular simulations on a quantum computer,
  Nature Communications 10, 3007 (2019), DOI:
  10.1038/s41467-019-10988-2.
\item
  Provost and Vallée, Riemannian structure on manifolds of quantum
  states, Communications in Mathematical Physics 76, 289-301 (1980),
  DOI: 10.1007/BF02193559.
\item
  Stokes, Izaac, Killoran, and Carleo, Quantum Natural Gradient, Quantum
  4, 269 (2020), DOI: 10.22331/q-2020-05-25-269.
\item
  Nocedal and Wright, Trust-Region Methods, in Numerical Optimization,
  Springer, DOI: 10.1007/0-387-22742-3\_4.
\item
  Ramôa, Santos, Mayhall, Barnes, and Economou, Reducing measurement
  costs by recycling the Hessian in adaptive variational quantum
  algorithms, Quantum Science and Technology 10, 015031 (2025), DOI:
  10.1088/2058-9565/ad904e.
\item
  Ramôa, Anastasiou, Santos, Mayhall, Barnes, and Economou, Reducing the
  resources required by ADAPT-VQE using coupled exchange operators and
  improved subroutines, npj Quantum Information (2025), DOI:
  10.1038/s41534-025-01039-4.
\item
  Verteletskyi, Yen, and Izmaylov, Measurement optimization in the
  variational quantum eigensolver using a minimum clique cover, The
  Journal of Chemical Physics 152, 124114 (2020), DOI:
  10.1063/1.5141458.
\item
  Yen, Verteletskyi, and Izmaylov, Measuring All Compatible Operators in
  One Series of Single-Qubit Measurements Using Unitary Transformations,
  Journal of Chemical Theory and Computation (2020), DOI:
  10.1021/acs.jctc.0c00008.
\item
  Ikhtiarudin, Sunnardianto, Fathurrahman, Agusta, and Dipojono,
  Shot-Efficient ADAPT-VQE via Reused Pauli Measurements and
  Variance-Based Shot Allocation, arXiv:2507.16879 (2025).
\item
  Stadelmann, Übelher, Ramôa, Sambasivam, Barnes, and Economou,
  Strategies for Overcoming Gradient Troughs in the ADAPT-VQE Algorithm,
  arXiv:2512.25004 (2025).
\end{enumerate}

\hypertarget{addendum-2026-04-04-manuscript-aligned-direct-hh-adapt-phase3_v1-sweep-notes}{%
\subsubsection{\texorpdfstring{Addendum (2026-04-04): manuscript-aligned
direct HH ADAPT / \texttt{phase3\_v1} sweep
notes}{Addendum (2026-04-04): manuscript-aligned direct HH ADAPT / phase3\_v1 sweep notes}}\label{addendum-2026-04-04-manuscript-aligned-direct-hh-adapt-phase3_v1-sweep-notes}}

This addendum records the new manuscript-aligned direct HH ADAPT
behavior on the public \texttt{phase3\_v1} surface, with
\texttt{runtime\_split=off} treated as the canonical manuscript/public
setting.

Physics point for the main recovery study. - \texttt{L=2} -
\texttt{t=1.0} - \texttt{U=4.0} - \texttt{g\_ep=0.5} -
\texttt{omega0=1.0} - \texttt{n\_ph\_max=1} - \texttt{problem=hh} -
binary bosons - blocked ordering - open boundary

Authority targets for interpretation. - Manuscript matched-ADAPT target
from §19.1: \texttt{E\ =\ 0.1587240823603703},
\texttt{\textbar{}ΔE\textbar{}\ =\ 5.6178234644e-05},
\texttt{stop\_reason\ =\ drop\_plateau}. - Fixed-scaffold anchor from
§19.6: \texttt{\textbar{}ΔE\textbar{}\ ≈\ 4.23e-04},
\texttt{F\ ≈\ 0.9998}.

\hypertarget{current-status-on-the-new-manuscript-aligned-surface}{%
\paragraph{Current status on the new manuscript-aligned
surface}\label{current-status-on-the-new-manuscript-aligned-surface}}

The best live strong-coupling beam line on the direct public surface is
\texttt{beam\_runtime\_split-off\_\_batching-off\_\_lifetime-off\_\_profile-tight\_\_gamma-1.0}
with - \texttt{E\ =\ 0.15872511715315915} -
\texttt{\textbar{}ΔE\textbar{}\ =\ 5.721302743347256e-05} -
\texttt{ansatz\_depth\ =\ 20} - \texttt{logical\_num\_parameters\ =\ 20}
- \texttt{num\_parameters\ =\ 48} -
\texttt{stop\_reason\ =\ drop\_plateau} -
\texttt{exact\_state\_fidelity\ =\ 0.9999715655261739}

This line clearly beats the fixed-scaffold anchor, but it does not yet
numerically match the manuscript matched-ADAPT target. The residual gap
to the manuscript matched-ADAPT line is \texttt{1.0347927889e-06} in
\texttt{\textbar{}ΔE\textbar{}}.

A historical diagnostic reoptimization run with windowed \texttt{COBYLA}
reached - \texttt{E\ =\ 0.15872468271899687} -
\texttt{\textbar{}ΔE\textbar{}\ =\ 5.6778593271189504e-05} - residual
gap to the manuscript matched-ADAPT line: \texttt{6.0035862640e-07}

This is useful as a numerical diagnostic, but it is not the forward
recipe for new runs. Going forward, the intended optimizer family is
\texttt{SPSA}, not \texttt{COBYLA}.

\hypertarget{executed-run-families-and-outcomes}{%
\paragraph{Executed run families and
outcomes}\label{executed-run-families-and-outcomes}}

\begin{enumerate}
\def\labelenumi{\arabic{enumi}.}
\tightlist
\item
  Strong-coupling beam sweep.
\end{enumerate}

\begin{itemize}
\tightlist
\item
  Artifact root:
  \texttt{artifacts/agent\_runs/20260403\_hh\_l2\_u4\_g05\_beam\_sweep\_stage1/}
\item
  Best case:
  \texttt{beam\_runtime\_split-off\_\_batching-off\_\_lifetime-off\_\_profile-tight\_\_gamma-1.0}
\item
  Result: \texttt{\textbar{}ΔE\textbar{}\ =\ 5.721302743347256e-05}
\item
  Interpretation: beam search nearly recovers the manuscript
  matched-ADAPT quality.
\end{itemize}

\begin{enumerate}
\def\labelenumi{\arabic{enumi}.}
\setcounter{enumi}{1}
\tightlist
\item
  Strong-coupling runtime sweep.
\end{enumerate}

\begin{itemize}
\tightlist
\item
  Artifact root:
  \texttt{artifacts/agent\_runs/20260403\_hh\_l2\_u4\_g05\_runtime\_sweep\_stage1/}
\item
  Best case still bad:
  \texttt{\textbar{}ΔE\textbar{}\ =\ 0.3282923042519191}
\item
  Interpretation: the single-trajectory runtime line collapses into the
  same bad basin seen in the manuscript novelty-off failure scale.
\end{itemize}

\begin{enumerate}
\def\labelenumi{\arabic{enumi}.}
\setcounter{enumi}{2}
\tightlist
\item
  Phase-3 beam-vs-batching diagnostic.
\end{enumerate}

\begin{itemize}
\tightlist
\item
  Artifact root:
  \texttt{artifacts/agent\_runs/20260403\_hh\_l2\_u4\_g05\_phase3\_diag\_beam\_vs\_batch\_v1/}
\item
  Single-trajectory cases stayed at
  \texttt{\textbar{}ΔE\textbar{}\ ≈\ 0.3282923}.
\item
  Beam \texttt{b3,c2,k3} immediately recovered
  \texttt{\textbar{}ΔE\textbar{}\ =\ 5.721302743347256e-05}.
\item
  Interpretation: the decisive lever is frontier diversity from beam
  search, not batching.
\end{itemize}

\begin{enumerate}
\def\labelenumi{\arabic{enumi}.}
\setcounter{enumi}{3}
\tightlist
\item
  Beam-capacity recovery slice.
\end{enumerate}

\begin{itemize}
\tightlist
\item
  Artifact root:
  \texttt{artifacts/agent\_runs/20260403\_hh\_l2\_u4\_g05\_phase3\_beam\_capacity\_recovery\_v1/}
\item
  \texttt{b3,c2,k3} and \texttt{b4,c2,k4} tied at
  \texttt{\textbar{}ΔE\textbar{}\ =\ 5.721302743347256e-05}.
\item
  Interpretation: larger beam width by itself does not improve the
  recovered solution once the beam is already large enough.
\end{itemize}

\begin{enumerate}
\def\labelenumi{\arabic{enumi}.}
\setcounter{enumi}{4}
\tightlist
\item
  Prune/shortlist micro-sweep.
\end{enumerate}

\begin{itemize}
\tightlist
\item
  Artifact root:
  \texttt{artifacts/agent\_runs/20260404\_hh\_l2\_u4\_g05\_phase3\_prune\_shortlist\_micro\_v1/}
\item
  All cases were numerically identical to saved precision.
\item
  Interpretation: \texttt{phase1\_prune\_fraction},
  \texttt{phase1\_prune\_max\_candidates}, and shortlist width are not
  the active driver of the remaining manuscript gap on this beam-fixed
  line.
\end{itemize}

\begin{enumerate}
\def\labelenumi{\arabic{enumi}.}
\setcounter{enumi}{5}
\tightlist
\item
  Reoptimization slice.
\end{enumerate}

\begin{itemize}
\tightlist
\item
  Artifact root:
  \texttt{artifacts/agent\_runs/20260404\_hh\_l2\_u4\_g05\_phase3\_reopt\_slice\_v1/}
\item
  Windowed \texttt{COBYLA} improved the gap to
  \texttt{\textbar{}ΔE\textbar{}\ =\ 5.6778593271189504e-05}.
\item
  \texttt{POWELL} windowed stayed at the original beam winner.
\item
  \texttt{POWELL} periodic/full refits got worse.
\item
  \texttt{append\_only} failed badly at
  \texttt{\textbar{}ΔE\textbar{}\ =\ 0.6244148638396572}.
\item
  Interpretation: optimizer/refit policy can move the line, but the
  \texttt{COBYLA} improvement should be treated as diagnostic only, not
  as the new manuscript-forward recipe.
\end{itemize}

\begin{enumerate}
\def\labelenumi{\arabic{enumi}.}
\setcounter{enumi}{6}
\tightlist
\item
  Score-weight slice.
\end{enumerate}

\begin{itemize}
\tightlist
\item
  Artifact root:
  \texttt{artifacts/agent\_runs/20260404\_hh\_l2\_u4\_g05\_phase3\_score\_weight\_slice\_v1/}
\item
  \texttt{near=0.90} gave a small improvement to
  \texttt{\textbar{}ΔE\textbar{}\ =\ 5.721150236287498e-05} and also
  reduced the scaffold to \texttt{logical\_num\_parameters\ =\ 19},
  \texttt{num\_parameters\ =\ 46}, \texttt{ansatz\_depth\ =\ 19}.
\item
  Other \texttt{rho} and \texttt{lambda\_H} variants were essentially
  flat at the saved precision.
\end{itemize}

\begin{enumerate}
\def\labelenumi{\arabic{enumi}.}
\setcounter{enumi}{7}
\tightlist
\item
  Seed robustness slice.
\end{enumerate}

\begin{itemize}
\tightlist
\item
  Artifact root:
  \texttt{artifacts/agent\_runs/20260404\_hh\_l2\_u4\_g05\_phase3\_seed\_robustness\_slice\_v1/}
\item
  Seeds \texttt{5,\ 7,\ 11,\ 13,\ 17} all reproduced the same
  \texttt{POWELL} baseline result exactly to saved precision.
\item
  Interpretation: the current beam-fixed \texttt{POWELL} plateau is
  deterministic across this seed set.
\end{itemize}

\begin{enumerate}
\def\labelenumi{\arabic{enumi}.}
\setcounter{enumi}{8}
\tightlist
\item
  Weak-coupling diagnostic lanes.
\end{enumerate}

\begin{itemize}
\tightlist
\item
  Canonical weak sweep root:
  \texttt{artifacts/agent\_runs/20260403\_hh\_l2\_u05\_g02\_weak\_logical\_scaffold\_canonical\_sweep\_v1/}
\item
  Weak beam sweep root:
  \texttt{artifacts/agent\_runs/20260403\_hh\_l2\_u05\_g02\_beam\_sweep\_stage1\_parent\_full\_meta/}
\item
  Best weak result in both lanes stayed at
  \texttt{\textbar{}ΔE\textbar{}\ ≈\ 0.0131029358626}.
\item
  Parent-quality target from §19.3.1 is much smaller, so weak coupling
  remains a diagnostic regime rather than a recovered manuscript-quality
  regime.
\end{itemize}

\hypertarget{good-parameters-versus-bad-parameters}{%
\paragraph{Good parameters versus bad
parameters}\label{good-parameters-versus-bad-parameters}}

Parameters/choices that consistently helped or were at least required
for recovery. - Direct HH ADAPT on the public \texttt{phase3\_v1}
surface. - \texttt{runtime\_split=off}. - Beam search on. - Beam at or
above \texttt{live\_branches=3}, \texttt{children\_per\_parent=2},
\texttt{terminated\_keep=3}. - \texttt{gamma\_N=1.0} on the public
strong-coupling line. - \texttt{batching=off} was slightly better than
\texttt{batching=on} in the best recovered beam lane. -
\texttt{lifetime\_cost\_mode=off} was at least as good as
\texttt{phase3\_v1} within the explored resolution. - \texttt{near=0.90}
was the only score-weight variation that gave a measurable improvement
in the explored \texttt{POWELL} lane.

Parameters/choices that were neutral or nearly inert in the explored
region. - Beam width beyond the \texttt{b3,c2,k3} threshold. -
\texttt{profile=tight} versus \texttt{profile=wide}. -
\texttt{phase1\_prune\_fraction} in the explored \texttt{0.15} to
\texttt{0.35} range. - \texttt{phase1\_prune\_max\_candidates} in the
explored \texttt{4} to \texttt{8} range. - shortlist width in the
explored micro-sweep. - Most \texttt{rho} and \texttt{lambda\_H}
score-weight perturbations.

Parameters/choices that consistently hurt. - Single-trajectory runtime
mode with no beam. - Runtime line behavior that reproduces
\texttt{\textbar{}ΔE\textbar{}\ ≈\ 0.3282923}. -
\texttt{runtime\_split=shortlist\_pauli\_children\_v1} relative to
\texttt{off} on this manuscript/public lane. - \texttt{POWELL}
periodic/full refits relative to the windowed baseline. -
\texttt{append\_only} refit policy, which failed catastrophically.

\hypertarget{practical-interpretation}{%
\paragraph{Practical interpretation}\label{practical-interpretation}}

For the manuscript-aligned direct HH ADAPT path at strong coupling, the
main qualitative lesson is now clear: beam-induced frontier diversity is
the key ingredient that recovers the near-manuscript solution, while
pruning pressure, shortlist width, profile choice, and modest
score-weight changes are second-order. The live public-surface recipe
already beats the fixed-scaffold anchor by a large margin, but it still
does not numerically match the manuscript matched-ADAPT target.

\hypertarget{update-public-cli-runtime-weights-and-reduced-pool-sweeps}{%
\paragraph{2026-04-04 update: public-CLI runtime weights and
reduced-pool
sweeps}\label{update-public-cli-runtime-weights-and-reduced-pool-sweeps}}

After exposing the manuscript-facing selector/runtime knobs on the
public direct \texttt{phase3\_v1} CLI, two additional conclusions
emerged for the same strong-coupling \(L=2\) point.

\begin{enumerate}
\def\labelenumi{\arabic{enumi}.}
\tightlist
\item
  Public-CLI \texttt{full\_meta\ +\ SPSA} improved once frontier ratios
  were decoupled from the near-degenerate batch shell.
\end{enumerate}

\begin{itemize}
\tightlist
\item
  Artifact root:
  \texttt{artifacts/agent\_runs/20260404\_hh\_l2\_u4\_g05\_phase3\_public\_weights\_frontier\_spsa\_v1/}
\item
  Best public-CLI \texttt{full\_meta\ +\ SPSA} line: \[
  E=0.15878350593338042,\qquad |\Delta E|=1.1560180765\times 10^{-4},
  \] with \texttt{logical\_num\_parameters=21},
  \texttt{num\_parameters=53}, \texttt{ansatz\_depth=21}, and
  \texttt{fidelity=0.9999480673}.
\item
  The corresponding \texttt{weights\_zero} line was only slightly worse,
  \[
  |\Delta E|=1.1722877711\times 10^{-4},
  \] while the tighter coupled-frontier setting \[
  (\text{phase2 frontier},\text{phase3 frontier})=(0.98,0.98)
  \] reproduced the older worse basin \[
  |\Delta E|=1.4383993770\times 10^{-4}.
  \]
\item
  So the main gain from that patch was not a delicate burden-weight
  choice; it was fixing the implementation coupling between shortlist
  frontiering and the batch-shell threshold so that
  \texttt{phase2\_frontier\_ratio}, \texttt{phase3\_frontier\_ratio},
  and \texttt{eta\_\{nd\}} became genuinely independent controls.
\end{itemize}

\begin{enumerate}
\def\labelenumi{\arabic{enumi}.}
\setcounter{enumi}{1}
\tightlist
\item
  The reduced-pool \texttt{pareto\_lean\ +\ SPSA} lane remained rigid.
\end{enumerate}

\begin{itemize}
\tightlist
\item
  Artifact roots:

  \begin{itemize}
  \tightlist
  \item
    \texttt{artifacts/agent\_runs/20260404\_hh\_l2\_u4\_g05\_phase3\_pool\_family\_spsa\_fixed\_v1/}
  \item
    \texttt{artifacts/agent\_runs/20260404\_hh\_l2\_u4\_g05\_phase3\_pareto\_lean\_gamma\_spsa\_fixed\_v1/}
  \item
    \texttt{artifacts/agent\_runs/20260404\_hh\_l2\_u4\_g05\_phase3\_pareto\_lean\_drop\_policy\_spsa\_fixed\_v1/}
  \item
    \texttt{artifacts/agent\_runs/20260404\_hh\_l2\_u4\_g05\_pareto\_lean\_new\_cli\_weights\_spsa\_v1/}
  \item
    \texttt{artifacts/agent\_runs/20260404\_hh\_l2\_u4\_g05\_pareto\_lean\_compat\_proxy\_spsa\_v1/}
  \item
    \texttt{artifacts/agent\_runs/20260404\_hh\_l2\_u4\_g05\_pareto\_lean\_frontier\_ratio\_spsa\_v1/}
  \end{itemize}
\item
  Baseline reduced-pool result: \[
  E=0.1587895750333393,\qquad |\Delta E|=1.2167090761\times 10^{-4},
  \] with \texttt{logical\_num\_parameters=21},
  \texttt{num\_parameters=52}, \texttt{ansatz\_depth=21}, and
  \texttt{fidelity=0.9999443566}.
\item
  This value persisted across all tested variations of:

  \begin{itemize}
  \tightlist
  \item
    \texttt{gamma\_N} in the explored \texttt{0.75} to \texttt{1.5}
    range,
  \item
    drop-policy parameters,
  \item
    motif bonus,
  \item
    duplicate penalty,
  \item
    leakage penalty \texttt{eta\_L},
  \item
    compatibility-prescreen weights,
  \item
    remaining-evaluations proxy mode,
  \item
    frontier ratios from \texttt{(0.85,0.85)} up through
    \texttt{(0.95,0.95)} and the asymmetric \texttt{(0.90,0.75)},
    \texttt{(0.90,0.95)}, \texttt{(0.95,0.90)} cases.
  \end{itemize}
\item
  Loosening the phase-2 frontier too far did hurt: \texttt{(0.75,0.75)}
  gave \[
  |\Delta E|=1.5234589586\times 10^{-4},
  \] and \texttt{(0.75,0.90)} gave \[
  |\Delta E|=1.5457582580\times 10^{-4}.
  \]
\item
  Hence the reduced \texttt{pareto\_lean} line appears structurally
  plateaued under the currently exposed runtime-law knobs: it is not
  rescued by more local tuning of \texttt{gamma\_N}, batching pressure,
  motif/repetition heuristics, or compatibility/proxy terms.
\end{itemize}

\begin{enumerate}
\def\labelenumi{\arabic{enumi}.}
\setcounter{enumi}{2}
\tightlist
\item
  Cost-aware ranking versus raw-energy ranking now separate.
\end{enumerate}

\begin{itemize}
\tightlist
\item
  On raw strong-coupling energy error alone, the best historical rerank
  observed in these sweeps was the diagnostic \texttt{COBYLA} line \[
  E=0.15872468271899687,\qquad |\Delta E|=5.6778593271\times 10^{-5},
  \] with \texttt{logical\_num\_parameters=20},
  \texttt{num\_parameters=50}, \texttt{ansatz\_depth=20}.
\item
  However, the best currently observed \(|\Delta E|/K\)-style point
  among runs that reported logical scaffold size was the
  \texttt{near=0.90} strong-coupling line from the score-weight slice:
  \[
  E=0.15872511562808855,\qquad |\Delta E|=5.7211502363\times 10^{-5},
  \] with \[
  K_{\text{logical}}=19,\qquad K_{\text{runtime}}=46,\qquad d=19.
  \]
\item
  So the present picture is:

  \begin{itemize}
  \tightlist
  \item
    best raw-\(|\Delta E|\) line: slightly better than the old beam
    winner, but achieved on a non-preferred optimizer branch,
  \item
    best cost-aware line: slightly worse in \(|\Delta E|\) but leaner in
    both logical scaffold size and runtime parameter count,
  \item
    best reduced-pool SPSA line: still stuck near \(1.22\times 10^{-4}\)
    and therefore materially behind the broad-beam public lane.
  \end{itemize}
\end{itemize}

\hypertarget{update-completed-reduced-pool-l2l3-matrix}{%
\paragraph{\texorpdfstring{2026-04-04 update: completed reduced-pool
\texttt{L=2/L=3}
matrix}{2026-04-04 update: completed reduced-pool L=2/L=3 matrix}}\label{update-completed-reduced-pool-l2l3-matrix}}

The overnight reduced-pool matrix
\texttt{artifacts/agent\_runs/20260404\_hh\_phase3\_reduced\_pool\_matrix\_v1/}
completed through the strong confirmation and weak-transfer waves. Its
campaign summary and final strong templates are recorded in

\begin{itemize}
\tightlist
\item
  \texttt{artifacts/agent\_runs/20260404\_hh\_phase3\_reduced\_pool\_matrix\_v1/campaign\_summary.json},
\item
  \texttt{artifacts/agent\_runs/20260404\_hh\_phase3\_reduced\_pool\_matrix\_v1/final\_strong\_templates\_by\_L.json},
\item
  \texttt{artifacts/agent\_runs/20260404\_hh\_phase3\_reduced\_pool\_matrix\_v1/matrix\_leaderboard.json}.
\end{itemize}

Data field: the completed matrix contained \texttt{212} successful
cases, \texttt{36} failed cases, and \texttt{184} unique successful
configurations.

For the strong-coupling \(L=2\) target point, the best matrix result was
\[
E=0.15877612622031317,\qquad |\Delta E|=1.0822209459\times 10^{-4},
\] with \[
K_{\mathrm{logical}}=20,\qquad K_{\mathrm{runtime}}=52,\qquad d=20,\qquad F=0.9999519211,
\] from the confirmed \texttt{pareto\_lean} template with \[
(\text{beam},\text{frontier})=(2,2,2;\,0.85,0.85).
\] This improved over the earlier reduced-pool SPSA plateau near
\(1.2167\times10^{-4}\), but it still did not recover the
manuscript-quality strong-coupling target at
\(5.6178234644\times10^{-5}\) and did not beat the earlier broad-lane
live recovery line near \(5.7213027433\times10^{-5}\).

For the strong-coupling \(L=3\) point, the reduced-pool matrix failed
badly. Its best confirmed case was \[
E=-0.2564741176143438,\qquad |\Delta E|=5.0141481774\times 10^{-1},
\] with fidelity only \[
F\approx 5.81\times 10^{-8}.
\] The best matrix lane there was \texttt{uccsd\_paop\_lf\_full}, not
the more aggressively reduced families. Hence the completed reduced-pool
matrix should be read as negative evidence for the current public
reduced execution surface at \(L=3\), not as evidence that the
historical \(L=3\) HH target itself is unattainable.

For weak coupling, the transferred \(L=2\) reduced-pool line remained
poor: \[
E=-0.7632521982953032,\qquad |\Delta E|=1.3125332609\times10^{-2},
\] so the matrix did not close the weak-coupling parent-quality gap
there. By contrast, the transferred \(L=3\) weak baseline was
numerically much better, \[
E=-1.0439438365685594,\qquad |\Delta E|=4.1313925567\times10^{-4},
\] with fidelity \[
F=0.9998116001.
\] That \(L=3\) weak result is still not a manuscript anchor, but it
does show that the reduced families are not uniformly broken across all
\(L=3\) HH regimes; the main failure is the strong-coupling \(L=3\)
execution surface.

To make that contrast explicit, the surviving historical broad-pool
\(L=3\) reference from
\texttt{artifacts/reports/hh\_heavy\_scaffold\_best\_yet\_20260321.md}
was \[
E=0.24502564220194084,\qquad |\Delta E|=8.4942074026\times10^{-5},
\] while the matched lean comparison report
\texttt{artifacts/reports/pareto\_lean\_vs\_heavy\_L3\_comparison.md}
recorded the lean/class-pruned line at about \[
|\Delta E|\approx 1.15\times10^{-4}.
\] Both are therefore qualitatively incompatible with the present
reconstructed direct-surface reruns, which all collapsed to the same bad
basin.

In particular, the explicit three-way reproduction run
\texttt{artifacts/agent\_runs/20260404\_hh\_l3\_u4\_g05\_historical\_surface\_pool\_compare\_v1/}
reconstructed the documented historical settings \[
\texttt{phase3\_v1},\ \texttt{POWELL},\ \texttt{phase1 shortlist}=256,\ \texttt{phase2 shortlist}=128,\ \texttt{split}=\texttt{shortlist\_pauli\_children\_v1},
\] and compared

\begin{itemize}
\tightlist
\item
  \texttt{full\_meta},
\item
  class-filtered \texttt{full\_meta},
\item
  \texttt{pareto\_lean\_l3}.
\end{itemize}

All three produced the same failed result, \[
E=-0.25838912148429044,\qquad |\Delta E|=5.0332982161\times10^{-1},
\] with negligible fidelity. Hence the current issue is not simply which
reduced family is chosen. Rather, some additional execution-surface
ingredient from the older successful \(L=3\) lane is still missing from
the present direct CLI reconstruction.

\hypertarget{appendix-raw-hf-start-direct-adapt-burdens-versus-historical-post-processed-fixed-scaffold-submissions}{%
\paragraph{Appendix: raw HF-start direct ADAPT burdens versus historical
post-processed / fixed-scaffold
submissions}\label{appendix-raw-hf-start-direct-adapt-burdens-versus-historical-post-processed-fixed-scaffold-submissions}}

To avoid conflating the heavy raw recent ADAPT replay artifacts with the
older locked fixed-manifold scaffold family, the following tables record
the exact transpile-style cost snapshots discussed during the April 7
accounting pass.

These counts are all on the same local hardware proxy surface, namely
\texttt{FakeMarrakesh}, and they should be interpreted as scaffold
burden rather than full campaign burden. In particular, the recent raw
ADAPT ladders below are \emph{not} the old locked 6/7-term artifacts.

Throughout this appendix, \textbf{raw HF-start direct ADAPT} means
\texttt{adapt\_ref\_json\ =\ null} with no later JSON prune/recovery
substitution. The post-processed \texttt{31/37/57}-2Q readapt/prune
artifacts listed later are comparison-only and are not the object a
fresh real-QPU ADAPT run from HF would execute.

\hypertarget{a.-raw-hf-start-direct-adapt-replay-scaffold-burden-snapshots}{%
\subparagraph{A. Raw HF-start direct ADAPT / replay scaffold burden
snapshots}\label{a.-raw-hf-start-direct-adapt-replay-scaffold-burden-snapshots}}

\begin{longtable}[]{@{}
  >{\raggedright\arraybackslash}p{(\columnwidth - 12\tabcolsep) * \real{0.11}}
  >{\raggedleft\arraybackslash}p{(\columnwidth - 12\tabcolsep) * \real{0.15}}
  >{\raggedleft\arraybackslash}p{(\columnwidth - 12\tabcolsep) * \real{0.15}}
  >{\raggedleft\arraybackslash}p{(\columnwidth - 12\tabcolsep) * \real{0.15}}
  >{\raggedleft\arraybackslash}p{(\columnwidth - 12\tabcolsep) * \real{0.15}}
  >{\raggedleft\arraybackslash}p{(\columnwidth - 12\tabcolsep) * \real{0.15}}
  >{\raggedleft\arraybackslash}p{(\columnwidth - 12\tabcolsep) * \real{0.15}}@{}}
\toprule
Case & Logical blocks & Runtime rotations & Raw \texttt{cx} & Transpiled
2Q & Depth & Size \\
\midrule
\endhead
Noisy weak raw run (\texttt{20260407} ladder) & 19 & 47 & 142 & 221 &
631 & 922 \\
Noisy mid raw run (\texttt{20260407} ladder) & 21 & 53 & 154 & 256 & 697
& 1084 \\
Noisy high raw run (\texttt{20260407} ladder) & 23 & 56 & 160 & 262 &
764 & 1113 \\
Noiseless Math front-page anchor
(\texttt{w5\_l2\_strong\_pareto\_lean\_b2c2k2\_f0p85\_0p85\_f1\_seed5})
& 20 & 52 & 152 & 218 & 633 & 939 \\
\bottomrule
\end{longtable}

Interpretation: the validated front-page
\texttt{pareto\_lean\ +\ phase3\_v1} deployment anchor is still a heavy
raw HF-start ADAPT replay scaffold on Marrakesh (\texttt{218} two-qubit
gates), i.e.~it sits near the recent weak raw-noise lane and far above
the post-processed prune/readapt and locked fixed-scaffold families. The
lighter \texttt{81}-2Q direct HF-start winner is the separate scientific
winner listed on the front page.

\hypertarget{b.-historical-post-processed-or-fixed-scaffold-runtime-submissions}{%
\subparagraph{B. Historical post-processed or fixed-scaffold runtime
submissions}\label{b.-historical-post-processed-or-fixed-scaffold-runtime-submissions}}

The next table isolates the \(L=2\) runtime submissions that were
actually sent to the real \texttt{ibm\_marrakesh} backend. The listed
two-qubit counts are obtained by transpiling the exact submitted
scaffold artifacts on \texttt{FakeMarrakesh} with the same transpiler
settings recorded in the runtime-result JSON.

\begin{longtable}[]{@{}
  >{\raggedright\arraybackslash}p{(\columnwidth - 10\tabcolsep) * \real{0.15}}
  >{\raggedright\arraybackslash}p{(\columnwidth - 10\tabcolsep) * \real{0.15}}
  >{\raggedright\arraybackslash}p{(\columnwidth - 10\tabcolsep) * \real{0.15}}
  >{\raggedleft\arraybackslash}p{(\columnwidth - 10\tabcolsep) * \real{0.20}}
  >{\raggedleft\arraybackslash}p{(\columnwidth - 10\tabcolsep) * \real{0.20}}
  >{\raggedright\arraybackslash}p{(\columnwidth - 10\tabcolsep) * \real{0.15}}@{}}
\toprule
Real-QPU submission & Runtime source artifact & Runtime transpile
settings & Transpiled 2Q & Depth & Notes \\
\midrule
\endhead
\texttt{hh\_gatepruned7\_fixedtheta\_runtime\_ibm\_marrakesh\_20260323T145219Z\_direct}
& \texttt{artifacts/json/hh\_prune\_nighthawk\_gate\_pruned\_7term.json}
& \texttt{opt=1}, \texttt{seed=42} & 25 & 63 & Locked 7-term fixed-theta
Marrakesh submission \\
\texttt{hh\_readapt5\_energyonly\_runtime\_ibm\_marrakesh\_20260323T153244Z}
& \texttt{artifacts/json/hh\_prune\_nighthawk\_readapt\_5op.json} &
\texttt{opt=1}, \texttt{seed=7} & 37 & 126 & Post-prune/recovery runtime
candidate; not fresh HF-start ADAPT \\
\texttt{hh\_gatepruned6\_fixedtheta\_runtime\_ibm\_marrakesh\_20260323T171528Z}
&
\texttt{artifacts/json/hh\_marrakesh\_fixed\_scaffold\_6term\_drop\_eyezee\_20260323T171528Z.json}
& \texttt{opt=1}, \texttt{seed=7} & 14 & 48 & Locked 6-term fixed-theta
Marrakesh submission \\
\bottomrule
\end{longtable}

So the historical real-QPU gate counts for the three recorded Marrakesh
submissions were \[
25,\qquad 37,\qquad 14
\] two-qubit gates respectively for the locked 7-term, readapt-5, and
locked 6-term submissions.

\end{document}
